% Self-contained release source. Compile this file to obtain the delivered PDF.
\documentclass[reqno]{amsart}
\usepackage[T1]{fontenc}
\usepackage{lmodern}
\usepackage{amsmath,amssymb,mathtools,mathrsfs}
\usepackage{booktabs,longtable,array}
\usepackage{enumitem}
\usepackage{xcolor}
\usepackage{microtype}
\usepackage[unicode,pdfusetitle,hidelinks]{hyperref}
\usepackage{listings}
\setlength{\emergencystretch}{2em}
\setcounter{tocdepth}{1}
\allowdisplaybreaks[2]
\newtheorem{theorem}{Theorem}[section]
\newtheorem{proposition}[theorem]{Proposition}
\newtheorem{lemma}[theorem]{Lemma}
\newtheorem{corollary}[theorem]{Corollary}
\theoremstyle{definition}
\newtheorem{definition}[theorem]{Definition}
\theoremstyle{remark}
\newtheorem{remark}[theorem]{Remark}
\newtheorem{example}[theorem]{Example}
\newcommand{\C}{\mathbb C}
\newcommand{\R}{\mathbb R}
\newcommand{\N}{\mathbb N}
\newcommand{\D}{\mathbb D}
\newcommand{\HH}{\mathbb H}
\DeclareMathOperator{\Tr}{Tr}
\DeclareMathOperator{\Id}{Id}
\DeclareMathOperator{\supp}{supp}
\DeclareMathOperator{\rank}{rank}
\DeclareMathOperator{\diag}{diag}
\DeclareMathOperator{\osc}{osc}
\DeclareMathOperator{\dist}{dist}
\newcommand{\sectionfile}[1]{\input{sections/#1}}
\title[Stability semigroups and ground-state algorithms]{Stability-preserving semigroups,\\ capacity dynamics, and classical\\ ground-state algorithms}
\author{Dongsheng Wei}
\address{Independent Researcher}
\email{dongshengwei2025@icloud.com}
\date{}
\subjclass[2020]{Primary 47D06, 30C15; Secondary 15A42, 60J27, 81P45, 68Q25}
\keywords{stable polynomial, positive semigroup, Lee--Yang theorem, generator classification, capacity, spectral gap, quantum Hamiltonian, classical algorithm, ground-state sampling}
\begin{document}
\begin{abstract}
We classify the linear generators that preserve all complex
half-plane-stable polynomials on a finite box of degrees. Their canonical
differential order is at most two, and their coefficients admit
certificates of size independent of the capacities. In disk coordinates,
strictly positive coordinate damping makes the actual semigroup kernel
strictly zero-free and gives a sharp, capacity-independent real-part
spectral gap, without self-adjointness or coefficient positivity.
For positive degree-preserving evolutions, the homogeneous-sector
spectral bounds form a concave sequence. This yields sharp
token-graph spectral concavity and resolves the adjacency and
signless-Laplacian monotonicity conjectures of Apte, Parekh, and Sud.
For bounded-degree two-body generators on binary variables, this structure yields
deterministic classical algorithms for the principal eigenvalue and
relative approximations to stable product tests of the principal eigenvectors. The algorithms
combine local perturbation coefficients with a nonvanishing
reciprocal-field domain for the scalar tests. At fixed local norm bound
and positive damping, their bit complexity is polynomial in the system
size, inverse error, and input length. Applied to the full
Suzuki--Fisher family, they compute ground-state energies, relative
amplitudes against product states in a hemisphere, and the corresponding product-event
probabilities, and give classical sampling of adaptive equatorial
measurement records. We also prove an all-temperature Hamiltonian
converse and classify quantum Markov semigroups that completely preserve
the Lee--Yang property at arbitrary finite local spins: interacting generators are excluded,
and every spin above one-half has isotropic noise and axial drift.
A separate positive branch classifies finite, infinite, and
mixed-capacity dynamics, with the generation and domain hypotheses
stated explicitly.
\end{abstract}
\maketitle
\clearpage
\tableofcontents
% BEGIN SECTION introduction
\section{Introduction}\label{intro:section}

A zero-free polynomial can encode a probability distribution, a quantum
state, or a linear operator. In each setting the location of its zeros
imposes a global constraint on the encoded object. For probability
generating functions, real stability is a source of negative dependence.
For quantum systems, the Lee--Yang property controls the coherent tensor
of a state or a Gibbs operator. The same algebraic question appears in both
settings: which continuous evolutions preserve the relevant zero-free
region?

The classification of individual stability-preserving operators by their
symbols gives a starting point for this question
\cite{BB09,BB10,BB10II}. A semigroup imposes an additional condition:
every operator $e^{tA}$, including those at arbitrarily small positive
times, must be a preserver. This infinitesimal requirement is especially
restrictive at collisions of roots. It is also useful. A generator
description can identify all admissible interactions, distinguish
physical positivity from complex stability, and supply a spectral
structure that is not apparent from the entries of the operator.

Infinitesimal stability preservers have also been studied directly.
Purbhoo constructs a family on multiaffine polynomials indexed by
matrices with real diagonal and nonnegative off-diagonal entries
\cite[Proposition~4.5]{Purbhoo18}. His result gives sufficient
generators. The structural theorem below supplies necessary and
sufficient conditions for arbitrary endomorphisms on every finite
coordinate box.

This paper develops these consequences in the following chain:
\[
 \begin{gathered}
 \text{root-collision constraints}
 \ \Longrightarrow\
 \text{generator structure}\\
 \ \Longrightarrow\
 \text{strict complex spectral control}
 \ \Longrightarrow\
 \text{classical computation}.
 \end{gathered}
\]
There are two preservation classes. The first tests all complex stable
polynomials. Its finite-dimensional generator cone is particularly rigid,
and underlies our spectral and quantum results. The second tests
nonnegative real stable polynomials, as required for probability
generating functions. It allows additional dynamics, including
coalescence, and requires a separate capacity and analytic theory.
Neither preservation class is silently substituted for the other.

\subsection{The finite-dimensional structure theorem}

Let $V_\kappa$ be the space of polynomials with coordinatewise degree
bounds $\kappa=(\kappa_1,\ldots,\kappa_n)$. Every linear operator on this
space has a canonical normally ordered differential representation.
The first main theorem describes which such operators generate
stability-preserving semigroups, without assuming positivity, locality,
or a differential order in advance.

In real half-plane coordinates, the answer is second order. If
\[
 A=\sum_{\alpha\leq\kappa}P_\alpha(z)\partial^\alpha
\]
is the canonical expression, preservation forces $P_\alpha=0$ for
$|\alpha|>2$. Each diagonal second-order coefficient is a polynomial
of its own variable and is nonpositive on the real line. Each mixed
coefficient is a polynomial of the corresponding pair of variables and
is nonpositive on the real plane. Together with the condition that
$A$ acts on $V_\kappa$, these conditions are also sufficient.
This is Theorem~\ref{gen:main-real}. The finite-box constraint is
substantial: it couples the highest terms of the lower-order
coefficients to the second-order part.

For complex generators there is one further component.
Every generator is the sum of such a real generator and
\begin{equation}\label{intro:imaginary}
 \mathrm i\left(cI+\sum_i
   \left(q_i(z_i)\partial_i-\frac{\kappa_i}{2}q_i'(z_i)\right)\right),
 \qquad c\in\R,\quad q_i\geq0\ \hbox{on }\R,
\end{equation}
where each $q_i$ is a real polynomial of degree at most two.
Conversely every such sum generates a complex stability-preserving
semigroup (Theorem~\ref{gen:main-complex}).
The operators in \eqref{intro:imaginary} are the infinitesimal
one-variable M\"obius motions pointing into the half-plane.
Thus arbitrary interactions in the imaginary part are excluded.

The proof has two different ingredients. Root-cluster asymptotics force
the order bound and the signs. The local multiple-root mechanism has
an antecedent in the variational analysis of Burke and Overton
\cite[Theorem~1.2 and Lemma~1.4]{BurkeOverton01}.
A finite-dimensional algebraic
decomposition then reduces sufficiency to explicit binary primitives.
The required certificates use Gram matrices of orders at most four,
independently of the numerical capacities. Polarization gives an exact
binary-slot realization, with the binomial factors retained.
These are generator statements; they are not obtained by differentiating
a smooth boundary of the stable set. That boundary is singular, and its
higher-multiplicity strata matter.

The fixed-degree one-variable geometry can nevertheless be expressed
spectrally. If $H(p)$ is the Hermite matrix of a monic real-rooted
polynomial $p$, then
\[
 T_{\mathcal R_n}(p)
   =\{h:DH_p[h]|_{\ker H(p)}\succeq0\}.
\]
Appendix~\ref{herm:section} proves this formula for arbitrary root
multiplicities and identifies exactly when the associated multiplier
image is closed. This supplies a local geometric result, with its
fixed-degree convention stated explicitly; it does not replace the
degree-changing arguments used for the full generator theorem.

\subsection{Strict damping and a complex spectral gap}

Pass to the unit disk by a Cayley transform. Suppose $A_0$ generates
a semigroup preserving all complex disk-stable polynomials in
$V_\kappa$, and set
\begin{equation}\label{intro:damping}
 A=A_0-2\sum_i\mu_iE_i,\qquad E_i=z_i\partial_i,\qquad \mu_i>0.
\end{equation}
The damping has a simple interpretation: it contracts each polynomial
variable towards the center of the disk. The conclusion is stronger
than preservation of the open disk. For every $t>0$ the actual
binomial-weighted kernel of $e^{tA}$ is nonzero on the closed unit
polydisk, and it is stable in a polydisk of some radius $r_t>1$.
For any fixed positive time, a common strict radius works at all later
times (Theorem~\ref{spec:strict}).

Consequently $A$ has an algebraically simple principal eigenvalue $a_*$
and
\begin{equation}\label{intro:gap}
 \Re a_*-\Re a\geq2\min_i\mu_i
 \quad\text{for every other eigenvalue }a.
\end{equation}
The constant is sharp, and the bound is independent of the capacities
and number of variables (Theorem~\ref{spec:gap}). Neither self-adjointness
nor coefficient positivity is assumed. The principal right and left
objects are strictly stable.

The proof uses the finite-dimensional complex-cone theory of Rugh
\cite{Rugh10}, together with a damping argument and contraction of stable
kernels. The sharp radius-to-spectrum mechanism has a Hermitian
Lee--Yang antecedent in \cite{BGLW26}. Our extension concerns arbitrary
complex finite-box stable generators and their actual kernels.
It must not be read as a bound on all transients in a fixed ordinary
operator norm. The appropriate quantity is instead the oscillation
of a quotient $h/p$ over the polydisk. Its contraction, including
time-dependent flows, is developed in Section~\ref{spec:section}.
The projective-metric and oscillation constructions themselves have
antecedents in \cite{Dubois09,GaubertQu15}.

\subsection{Comparing homogeneous sectors}

Coefficient positivity gives another spectral consequence. Suppose a
finite-box generator preserves real stability, has nonnegative
off-diagonal entries, and preserves every homogeneous degree.
The spectral bounds of its degree blocks form a concave sequence
(Theorem~\ref{sector:growth}). This also holds for nonself-adjoint,
reducible blocks. The proof evolves a single polynomial with positive
coefficients, applies Newton's inequalities to its diagonal restriction,
and takes exponential growth rates. The underlying growth-exponent
principle is classical \cite{Branden10Discrete}.

For a nonnegatively weighted graph \(G\), let \(A_k,D_k\) be the
adjacency and degree matrices of its \(k\)-token graph. The generator
classification applies to the whole family
\[
 B_k(a)=A_k+aD_k,\qquad -1\le a\le1.
\]
Its largest eigenvalues satisfy
\[
 2\lambda_{\max}(B_k(a))
 \ge\lambda_{\max}(B_{k-1}(a))+\lambda_{\max}(B_{k+1}(a)),
\]
and are symmetric under \(k\leftrightarrow n-k\).
They are therefore nondecreasing up to half filling. This proves
the adjacency and signless-Laplacian monotonicity assertions of
\cite[Conjectures~5--6]{APS26}, with a stronger weighted concavity
statement and a sharp universal interval for \(a\)
(Theorem~\ref{sector:token}).
The \(a=0\) semigroup is already contained in Purbhoo's construction
\cite{Purbhoo18}; the new conclusion compares its particle-number
sectors. The corresponding EPR interpretation uses the established
token-graph correspondence of \cite{APS26}.

\subsection{A classical algorithm beyond self-adjoint operators}

Here is the algorithmic consequence in a form that specifies its input
model. On $n$ binary variables let $A_0$ be given as a sum of one-site
and two-site matrices with Gaussian-rational entries. Suppose the
interaction graph has degree at most $d$, each local matrix has norm
at most $J$, and $e^{tA_0}$ preserves all complex disk-stable
multiaffine polynomials. The last condition is a promise about the
whole generator; individual displayed summands need not be preservers.
For fixed $d,J$, and a fixed rational $\mu>0$, the principal eigenvalue of
\[
 A_0-2\mu\sum_iE_i
\]
can be approximated to absolute complex error $\delta$ by a deterministic
classical algorithm in time
\[
 \operatorname{poly}_{d,J,\mu}
       \bigl(n,\delta^{-1},\text{input bit length}\bigr).
\]
This is Theorem~\ref{alg:complex}. The matrix has exponentially many
entries, but the input is its sparse local description. There is no
claim of a running time polynomial uniformly in $\mu^{-1}$.

The same local input also gives relative access to the principal
eigenvectors. Normalize the right and transpose-left vectors $v,w$
to have vacuum coefficient one. For any tensor product $B=\bigotimes_iB_i$
of Gaussian-rational $2$-by-$2$ matrices with $(B_i)_{00}=1$ and
disk-stable kernels, Theorem~\ref{alg:product-test} computes
\[
 w^{\mathsf T}Bv
 \quad\hbox{and}\quad
 \Tr(B\Pi)=\frac{w^{\mathsf T}Bv}{w^{\mathsf T}v},
 \qquad \Pi=\frac{vw^{\mathsf T}}{w^{\mathsf T}v},
\]
to relative complex error $\epsilon$ in
$\operatorname{poly}_{d,J,\mu}(n,\epsilon^{-1},L_{\rm in})$ time.
Both quantities are nonzero. This is scalar query access; no
ordinary-norm conditioning bound for $\Pi$ is required.

The algorithm uses two distinct analytic regions. A local expansion
at a product vacuum has a size-independent convergence radius.
The perturbation-coefficient algorithm of Bravyi, DiVincenzo, and
Loss \cite{BDL08} supplies these coefficients. A complexification
argument extends the coefficient bounds from Hermitian perturbations
to arbitrary complex ones. The stability gap then continues the same
principal branch through a right half-plane. An explicit disk map
joining these regions turns logarithmically many local coefficients
into a value estimate at the desired fixed damping. We give the
branch-identification argument, the coefficient construction, and the
rational bit bounds in Section~\ref{alg:section}.

For the state queries, the extra step is to obtain a nonvanishing
scalar test on that same type of domain. The creation-coefficient
bounds in \cite{BDL08} give a uniform vacuum disk for both eigenvectors.
The half-plane continuation and strict disk contraction then show
that $w^{\mathsf T}Bv$ never vanishes. Its logarithm has a one-sided
real-part bound of order $n$ throughout the domain
(Lemma~\ref{alg:test-domain}). Connected coefficients of this logarithm
can be computed from subsystems of logarithmic size, without requiring
the induced subsystems to satisfy the global preservation promise.

This method does not introduce a new perturbation expansion or a new
general interpolation principle; connected-cluster computation combined
with analytic continuation is an established strategy
\cite{PatelRegts17,YYZ22,WildAlhambra23}.
Its point is that complex stability
provides a global analytic domain for a class of sparse operators that
need not be self-adjoint, positive, or Markovian. The computational
statement uses that domain to go beyond the perturbative neighborhood
where the coefficient expansion is initially justified.

\subsection{Quantum consequences and their scope}

For qubits, consider
\begin{equation}\label{intro:sf}
 \begin{gathered}
 H=-\sum_{ij}
 \left(J^z_{ij}Z_iZ_j+
       \sum_{a,b\in\{x,y\}}J^{ab}_{ij}\sigma_i^a\sigma_j^b\right)
   -\sum_i(h_i^xX_i+h_i^yY_i+h_i^zZ_i),
 \\
 J^z_{ij}\geq\|J^\perp_{ij}\|_{\rm op},\quad h_i^z\geq0 ,
 \end{gathered}
\end{equation}
with real coefficients and
$J^\perp_{ij}=(J^{ab}_{ij})_{a,b\in\{x,y\}}$.
This is the phase-rotated Suzuki--Fisher family.
Its Lee--Yang sufficiency, including the full transverse operator-norm
condition, is classical: see \cite[equation~(2.47)]{SuzukiFisher71}
and \cite[Theorem~2]{Dunlop79}. Asano's contemporaneous account already
uses a full matrix generating function and its composition law for the
spin-one-half XXZ subfamily \cite[equation~(1.15), Theorem~1]{AsanoJapanese70}.
The tensor language of \cite{WBG26} is a modern formulation of this
tradition. We do not claim the sufficient family or the use of full
matrix generating functions as new.
The generator theorem gives the converse: up to an additive
scalar, these are precisely the finite-qubit Hamiltonians whose
full Gibbs tensor is Lee--Yang at every temperature
(Theorem~\ref{eq:sf}). In particular no genuine higher-body term is
possible under that full-tensor hypothesis. Ruelle's all-temperature converse
\cite[Theorem 9]{Ruelle10} already forces pair interactions for classical
coefficientwise Boltzmann Lee--Yang polynomials. The object here is the
full matrix tensor of an arbitrary Hermitian Hamiltonian, including
noncommuting interactions.

For rational Pauli coefficients, bounded degree and bounded local
strength, and longitudinal fields bounded below by a fixed $\mu>0$,
Theorem~\ref{alg:sf} gives deterministic classical ground-energy
approximation to arbitrary absolute accuracy.
Theorem~\ref{alg:sf-state} additionally computes nonzero relative
amplitudes against product states
$(|0\rangle+s_i|1\rangle)/\sqrt{1+|s_i|^2}$, $|s_i|\le1$,
and relative probabilities of the corresponding product events.
It gives a classical randomized sampler for adaptive equatorial
measurements, with any prescribed total variation error. Each site
is measured at most once; the strategy's own running time and
the bit lengths of its phases are included explicitly.
At nonnegative fields, adding a fixed field of size proportional
to $\varepsilon$ gives an additive $\varepsilon n$ energy-value
approximation scheme (Theorem~\ref{alg:value}).
This covers bounded-degree EPR models on arbitrary graphs and,
after a local unitary transformation, bounded-degree bipartite
Quantum MaxCut. In the unweighted latter setting the estimate
also gives a relative approximation scheme for the optimum value.

Earlier classical algorithms treat important subclasses of quantum
ferromagnets \cite{BG17,HMS20}. The broader classical computation
question is stated in \cite{BGLW26}; a related interpolation strategy
was announced in \cite{WongTalk23}, with further background in
\cite{WongThesis26}.
Our comparison is between precise hypotheses and error guarantees,
not between names for the methods. Relative \(X\)-basis amplitudes, phases, and
measurement marginals already appear in
\cite[Section~4]{WBG26}, for general Lee--Yang states with a fixed
spatial zero-free radius greater than one and state-dependent
quasipolynomial circuits. Here the local Hamiltonian is the input,
the field is any fixed positive constant, and the complexity is
polynomial; a size-independent spatial radius margin is not assumed.
The new analytic input is the nonvanishing reciprocal-field domain
for the principal-state tests. It is distinct from the thermal
partition-function domains used for product measurements in
\cite[Theorem~15]{YYZ22}.
The zero-field energy-value scheme does not
prepare an approximately optimal quantum state. It does not solve
the unrestricted inverse-polynomial-accuracy EPR problem
\cite{MarwahaSud26}, and it is not a replacement for approximation
algorithms that also construct states \cite{ALMPSS26}.

Positive damping also proves strict zero-freeness of the actual Gibbs
kernel and of its unique ground-state polynomial, addressing the
qualitative radius question in \cite{WBG26}.
The stronger proposed radius $s^{-1/2}$ for a deformed EPR ground
state is false. Appendix~\ref{cx:section} gives an unweighted
six-vertex tree at $s=1/2$, together with an exact rational certificate
of a zero strictly inside that radius. It leaves the qualitative
radius-greater-than-one conclusion intact.

Physical quantum Markov evolution exhibits a different rigidity.
For at least two qubits, ordinary Lee--Yang preservation by a
continuous completely positive trace-preserving semigroup is
equivalent to full-Choi preservation. Its generator is a sum of
independent one-qubit generators satisfying two fixed-size matrix
inequalities (Theorem~\ref{dyn:main}).
Locality is already forced by testing finitely many fixed trine
product states at times tending to zero
(Theorem~\ref{dyn:trine}). This is a statement about physical positive
inputs, for which the full complex preserver theorem cannot simply
be assumed. It contrasts interacting equilibrium with independent
Lee--Yang-preserving Markovian evolution.

Theorem~\ref{dyn:arbitrary-spin} extends the dynamical classification
to arbitrary finite local spins, without assuming a local or linear
Lindblad representation. If at least one site has spin above one-half,
it is enough to test pure Lee--Yang inputs near time zero. All
intersite terms vanish. On every site of spin $j\ge1$, the generator is
\[
 \mathcal L_i(\rho)=-\mathrm i[h_iS_i^z,\rho]
       -\frac{\gamma_i}{2}\sum_{\alpha=x,y,z}
                   [S_i^\alpha,[S_i^\alpha,\rho]],
 \qquad h_i\in\R,\quad\gamma_i\ge0.
\]
Qubit sites retain the two matrix inequalities above. The isotropic
spin semigroup itself is classical; the conclusion here is that
Lee--Yang preservation forces this form out of an arbitrary finite
GKLS generator. A higher-spin site also supplies the pure-state
test that rules out the extra qubit-only boundary possibilities.
As a short application, the spin-one channel identifies axial
isotropic Brownian convolution on $\mathrm{SO}(3)$
(Corollary~\ref{dyn:brownian-recognition}).
The finite-representation recognition mechanism is an instance
of Voit's Gaussian characterization \cite{Voit03}, not a new
general characterization principle for Lie-group processes.

The appendices give the corresponding static invertible
classification, the isometric and pure-basis-output cases,
one-qubit singular strata, and the independently covariant
higher-capacity case. They retain their hypotheses: arbitrary
singular noisy channels and arbitrary noncovariant higher-spin
static channels are not classified here.

\subsection{Positive dynamics and capacities}

For probability generating functions the positivity restriction
changes the answer. The positive branch includes exclusion,
immigration, transfer, and coalescence mechanisms. Its infinitesimal
tests arise from finite root collisions and from Poisson or
Laguerre--P\'olya limits. We prove arbitrary-jump Markov rigidity,
a canonical logarithmic-symbol theorem, the finite quadratic-contact
criterion, and the finite/infinite/mixed-capacity classification
in Section~\ref{pos:section}. The connection with earlier
stability-preserving particle systems is made explicitly
\cite{BBL09,LVR10}.

Infinite capacity requires more than an algebraic generator condition.
The weighted spaces, polynomial-domain result, killed
finite-capacity approximations, and generation criteria are proved
separately in Appendix~\ref{ana:section}.
In particular, dilation by a factor greater than one is not treated
as a bounded similarity on unweighted $\ell^1$.
We also distinguish polynomial-valued semigroups from entire-valued
extensions; the latter need not preserve every entire zero-free
function.

This organization separates three kinds of input: inherited symbol,
complex-cone, and perturbation theory; the polynomial-specific
generator and strictification arguments; and the analytic or physical
conditions needed for each application. It permits the positive
capacity theory and the complex spectral theory to share their
root-collision mechanism without asserting that their preservation
cones coincide.

\subsection{Organization}
Section~\ref{pre:section} fixes all normalizations.
Section~\ref{gen:section} proves the finite-box generator theorem,
Section~\ref{spec:section} develops strictification, spectral
contraction, and sector concavity; Section~\ref{alg:section} proves the eigenvalue,
principal-state query, and measurement-sampling algorithms.
Sections~\ref{eq:section} and \ref{dyn:section} treat equilibrium and
quantum Markov dynamics. Section~\ref{pos:section} gives the positive
capacity theory. The appendices contain the local Hermite geometry,
analytic generation, static and higher-capacity classifications,
quantitative diagonal rigidity, and the exact radius counterexample.

% END SECTION introduction

% BEGIN SECTION preliminaries
\section{Conventions and zero-free kernels}\label{pre:section}

\subsection{Regions, degrees, and preservation}
Write $\HH=\{z\in\C:\Im z>0\}$ and $\D_r=\{z\in\C:|z|<r\}$,
with $\D=\D_1$.
\begin{definition}[Stability and preservation]\label{pre:stability}
A nonzero polynomial is \emph{$\Omega$-stable} if it
has no zero when all its variables lie in $\Omega$.
For a product of different regions the corresponding meaning is used.
A polynomial is \emph{real stable} if its coefficients are real and it is
$\HH$-stable. In one variable this is equivalent to real-rootedness.
A \emph{positive stable polynomial} has nonnegative real coefficients and
is real stable. Positivity throughout refers to coefficients, unless an
operator on a Hilbert space is explicitly said to be positive semidefinite.
An operator \emph{preserves} one of these classes when it sends every
member to another member or to zero.
\end{definition}

Let $\N=\{0,1,2,\ldots\}$. For $\kappa\in\N^n$ put
\[
 V_\kappa=\{f\in\C[z_1,\ldots,z_n]:\deg_{z_i}f\leq\kappa_i\},
 \qquad
 \binom{\kappa}{\alpha}=\prod_i\binom{\kappa_i}{\alpha_i}.
\]
Variables of capacity zero can be deleted. A degree bound is not an
assertion that the polynomial has that degree. All finite-box statements
include degree drops and constants. For a real operator, its action on
complex polynomials always means its complex-linear extension.
Every operator on $V_\kappa$ has a unique canonical expression
\begin{equation}\label{pre:normal-order}
 A=\sum_{\alpha\leq\kappa}P_\alpha(z)\partial^\alpha,
\end{equation}
as an operator on $V_\kappa$. Uniqueness follows by applying the expression
successively to monomials, starting with $1$: the equation for $z^\alpha$
determines $\alpha!P_\alpha$ after all lower terms are known. The phrase
\emph{differential order} refers to this expression.

A finite-dimensional generator $A$ preserves a class when $e^{tA}$
preserves it for all $t\geq0$. For an infinite-dimensional space, the
existence and topology of a strongly continuous semigroup are additional
requirements, treated in Appendix~\ref{ana:section}. A formal
differential expression alone does not assert generation.

\begin{definition}[Complete stability preservation]\label{pre:complete}
An operator is \emph{completely stable} if its tensor product with the
identity in arbitrarily many extra polynomial variables preserves the
same complex stable class, for every finite choice of the extra degree
bounds.
\end{definition}
This is a statement about polynomial
stability, distinct from complete positivity of a quantum channel.
Positive stability with nonnegative coefficients and complete complex
stability must also be distinguished.

\subsection{Polarization and the disk kernel}
The normalized polarization of $z_i^a$ into $\kappa_i$ binary slots is
\[
 z_i^a\longmapsto
 \binom{\kappa_i}{a}^{-1}
 e_a(z_{i1},\ldots,z_{i\kappa_i}).
\]
It is symmetric in the slots and specializes to $z_i^a$ on the diagonal.
Polarization preserves stability in the convex circular regions used
here, including open disks and half-planes; diagonalization preserves
it directly. We use the Grace--Walsh--Szeg\H{o} coincidence
theorem in this form; see \cite{BB09,BB10II,BB09Circular}.
The binomial normalization is part of the definition, not an optional
choice of basis.\label{pre:polarization}

For $T:V_\kappa\to V_\lambda$ define its disk kernel by
\begin{equation}\label{pre:kernel}
 K_T(z,w)=T_z\prod_i(1+z_iw_i)^{\kappa_i}
 =\sum_{\alpha\leq\kappa}\binom{\kappa}{\alpha}
       (Tz^\alpha)(z)w^\alpha.
\end{equation}
The notation $z_i$ in the first expression denotes the input variables
on which $T$ acts; its result is a polynomial in the output variables.
The kernel has output degree bound $\lambda$ and input degree bound
$\kappa$. In particular $K_{\Id}(z,w)=\prod_i(1+z_iw_i)^{\kappa_i}$.

For two polynomials of degree at most $\kappa$ in paired variables define
their bilinear contraction by
\begin{equation}\label{pre:contraction}
 [f(u),g(v)]_\kappa
   =\sum_{\alpha\leq\kappa}
      \binom{\kappa}{\alpha}^{-1} f_\alpha g_\alpha,
 \qquad
 f(u)=\sum_\alpha f_\alpha u^\alpha,\quad
 g(v)=\sum_\alpha g_\alpha v^\alpha .
\end{equation}
There is no complex conjugation in this formula. Matrix composition is
exactly contraction of the corresponding kernels:
\[
 K_{ST}(z,w)=[K_S(z,u),K_T(v,w)]_\kappa,
\]
where $\kappa$ is the intermediate capacity.

\begin{lemma}[Disk contraction]\label{pre:grace}
Suppose $f(z,u)$ and $g(v,w)$ are stable in products of disks and
$\deg_{u_i}f,\deg_{v_i}g\leq\kappa_i$. If their paired radii
$r_i,s_i$ satisfy $r_is_i\geq1$, then their contraction
\eqref{pre:contraction} is stable in the remaining disks or is identically
zero. If all $r_is_i>1$, contraction at any point in the remaining
disks is nonzero.
\end{lemma}
\begin{proof}
First consider a multiaffine polynomial
$F(u,v)=a+bu+cv+duv$ that is nonzero on $\D_r\times\D_s$.
Then $a\ne0$ and $F(rt,st)$ is nonzero for $|t|<1$.
If $d\ne0$, the product of its two roots has modulus at least one,
so $|d|rs\le|a|$; if $d=0$, this inequality is immediate.
Thus $a+d\ne0$ when $rs>1$.
The same argument applies after fixing any uncontracted variables
in their respective disks. Consequently the operation
$a+bu+cv+duv\mapsto a+d$ preserves stability under strict
separation of the paired radii. It may be iterated on a general
multiaffine polynomial, without a factorization assumption.

Polarize the paired variables of $f$ and $g$ into $\kappa_i$
slots for coordinate $i$, and apply this operation to each pair
of corresponding slots in their product.
For a fixed multi-index $\alpha$, the matching slot subsets have
cardinality $\binom{\kappa}{\alpha}$, while the two normalized
polarizations contribute the factor
$\binom{\kappa}{\alpha}^{-2}$.
The result is therefore exactly \eqref{pre:contraction}.
This proves strict nonvanishing, including degree deficiencies.
It is the weighted form of circular-domain contraction;
see \cite{AsanoJapanese70,BB09,BB10II}.
For equality of paired radii, replace every paired variable in each input
polynomial by $(1-\varepsilon)$ times that variable, apply the strict
statement, and let $\varepsilon\downarrow0$. Hurwitz's theorem gives
the stable-or-zero conclusion. The argument is valid after fixing
any values of the unpaired variables in their respective disks.
\end{proof}

The finite-box symbol theorem \cite{BB09,BB10II} identifies stability
preservers of rank at least two by the appropriate stable kernel;
the low-rank alternatives are retained whenever they arise.
In particular an invertible disk-stability preserver has a stable
kernel \eqref{pre:kernel}. We do not infer a kernel criterion for
arbitrary positive-coefficient preservers from this theorem.

The Cayley isomorphism is
\begin{equation}\label{pre:cayley}
 (\mathcal C_\kappa f)(x)
   =\prod_i(x_i+\mathrm i)^{\kappa_i}
      f\left(\frac{x_1-\mathrm i}{x_1+\mathrm i},\ldots,
              \frac{x_n-\mathrm i}{x_n+\mathrm i}\right).
\end{equation}
It is a linear isomorphism between the degree boxes, and $f$ is
$\D$-stable if and only if $\mathcal C_\kappa f$ is $\HH$-stable.
Thus the generator classifications in half-plane coordinates and
the spectral results in disk coordinates concern conjugate
finite-dimensional semigroups.

\subsection{Coherent coordinates and physical tensors}
Equip $V_\kappa$ with the inner product for which
$\binom{\kappa}{\alpha}^{1/2}z^\alpha$ is an orthonormal basis.
Equivalently, identify the standard basis $|\alpha\rangle$ of
$\bigotimes_i\C^{\kappa_i+1}$ with
$\binom{\kappa}{\alpha}^{1/2}z^\alpha$. The polynomial of a vector
$\psi$ is
\[
 f_\psi(z)=\sum_{\alpha\leq\kappa}
          \binom{\kappa}{\alpha}^{1/2}\psi_\alpha z^\alpha.
\]
An operator $M$ on this Hilbert space has coherent tensor
\begin{equation}\label{pre:coherent}
 F_M(z,w)=\sum_{\alpha,\beta\leq\kappa}
 \binom{\kappa}{\alpha}^{1/2}\binom{\kappa}{\beta}^{1/2}
 M_{\alpha\beta}z^\alpha w^\beta .
\end{equation}
This is precisely \eqref{pre:kernel} for its action in coherent
coordinates. Bra variables are independent complex variables;
they are not conjugated when testing zero-freeness.
For a positive semidefinite density matrix $\rho$, the assertion
that $F_\rho$ is $\D$-stable is called its \emph{Lee--Yang property}.

For a linear map $\Phi$ between matrix algebras the full coherent
Choi tensor is defined by replacing $M_{\alpha\beta}$ in
\eqref{pre:coherent} with the coefficients
$\Phi(|\gamma\rangle\langle\delta|)_{\alpha\beta}$ and adding
input monomials $u^\gamma v^\delta$ with the corresponding square-root
binomial weights. Its stability is called \emph{full-Choi
Lee--Yang preservation}. For qubits all capacities are one, so
these weights are one. They cannot be omitted for higher spin.
Ordinary physical preservation tests only positive semidefinite
Lee--Yang inputs; it is a weaker condition for isolated maps.

We use the Pauli matrices
\[
 X=\begin{pmatrix}0&1\\1&0\end{pmatrix},\qquad
 Y=\begin{pmatrix}0&-\mathrm i\\\mathrm i&0\end{pmatrix},\qquad
 Z=\begin{pmatrix}1&0\\0&-1\end{pmatrix}.
\]
Thus $E_i=z_i\partial_i$ represents $|1\rangle\langle1|_i$
and $Z_i=I-2E_i$ on a qubit.
For general capacity $E_i=z_i\partial_i$ has eigenvalues
$0,\ldots,\kappa_i$. A \emph{principal eigenvalue} below means an
algebraically simple eigenvalue whose real part is strictly greater
than the real parts of all other eigenvalues; for a discrete operator
it means the corresponding strict modulus dominance.

% END SECTION preliminaries

% BEGIN SECTION generators
\section{Generators on a finite polynomial box}
\label{gen:section}

We first classify semigroups acting on the whole stable class, with no
assumption on the signs of coefficients of the polynomials or of the
operators.  The distinction from preservation of the coefficient-positive
stable class is essential: the two problems have different generator cones.

Let \(\kappa=(\kappa_1,\ldots,\kappa_d)\) have finite positive integer
coordinates, and write
\[
 V_\kappa^{\mathbb F}
 =\mathbb F[z_1,\ldots,z_d]_{\deg_{z_i}\leq\kappa_i},
 \qquad \mathbb F\in\{\R,\C\}.
\]
Coordinates of capacity zero are omitted.  If no coordinate remains, the
space consists of constants and every scalar generator has the requisite
preservation property.  Stability and complete preservation have the meanings
of Definitions~\ref{pre:stability} and~\ref{pre:complete}; in this section the
domain of stability is \(\HH^d\).
A real operator acts on \(V_\kappa^\C\) by complexification.

For a polynomial \(q\) in one variable of degree at most two, define
\begin{equation}
 J_{i,\kappa_i}(q)
 =q(z_i)\partial_i-\frac{\kappa_i}{2}q'(z_i).
 \label{gen:eq-J}
\end{equation}
These are endomorphisms of the box.  The three choices \(q=1,z_i,z_i^2\)
give
\[
 J_i^-=\partial_i,\qquad
 J_i^0=z_i\partial_i-\frac{\kappa_i}{2},\qquad
 J_i^+=z_i^2\partial_i-\kappa_i z_i.
\]
Their span is the sitewise \(\mathfrak{sl}_2\) algebra.

Every endomorphism has a unique canonical normal representation
\begin{equation}
 A=\sum_{\alpha\leq\kappa}P_\alpha(z)\partial^\alpha.
 \label{gen:eq-canonical}
\end{equation}
Here the coefficient polynomials are not assumed to belong to the original
box.  The order of \(A\) means the largest \(|\alpha|\) for which its
canonical coefficient \(P_\alpha\) is nonzero.  The uniqueness of this
representation, proved below, is important at small capacities: in particular
\(\partial_i^2\) is absent when \(\kappa_i=1\).

\begin{theorem}[Real generators]
\label{gen:main-real}
Let \(A\in\operatorname{End}_{\R}(V_\kappa^\R)\).  The following conditions
are equivalent.
\begin{enumerate}
\item For every \(t\geq0\), \(e^{tA}\) preserves real stability.
\item For every \(t\geq0\), its complexification preserves complex stability,
also after adjoining arbitrary inert polynomial variables.
\item The canonical representation of \(A\) has order at most two, and its
second-order coefficients satisfy
\begin{align}
 P_{2e_i}(s)&\leq0
 &&\text{for all }s\in\R,\quad \kappa_i\geq2,
 \label{gen:eq-real-unary}\\
 P_{e_i+e_j}(s,t)&\leq0
 &&\text{for all }(s,t)\in\R^2,\quad i<j.
 \label{gen:eq-real-pair}
\end{align}
\end{enumerate}
In the third condition, box preservation forces \(P_{2e_i}\) to depend only
on \(z_i\), with degree at most four, and \(P_{e_i+e_j}\) to depend only on
\(z_i,z_j\), with bidegree at most \((2,2)\).  There are no further
inequalities on the zeroth- and first-order coefficients beyond the
identities already forced by \(A(V_\kappa^\R)\subseteq V_\kappa^\R\).
\end{theorem}

The coefficient \(P_{e_i+e_j}\) in this statement is the full coefficient
of \(\partial_i\partial_j\), with \(i<j\).  Thus a convention writing the
second-order part as \(\sum_{i,j}Q_{ij}\partial_i\partial_j\), with \(Q\)
symmetric, has \(P_{e_i+e_j}=2Q_{ij}\).
We use the full-coefficient convention throughout this section.

\begin{theorem}[Complex generators]
\label{gen:main-complex}
Let \(A\in\operatorname{End}_{\C}(V_\kappa^\C)\).  The following conditions
are equivalent.
\begin{enumerate}
\item For every \(t\geq0\), \(e^{tA}\) preserves complex stability.
\item For every \(t\geq0\), \(e^{tA}\) preserves complex stability after
adjoining arbitrary inert polynomial variables.
\item There are a real generator \(R\) satisfying
Theorem~\ref{gen:main-real}, a real scalar \(c\), and real polynomials
\(q_i\) of degree at most two, nonnegative on \(\R\), such that
\begin{equation}
 A=R+i\left(c\Id+\sum_{i=1}^dJ_{i,\kappa_i}(q_i)\right).
 \label{gen:eq-complex-decomposition}
\end{equation}
\end{enumerate}
Equivalently, the canonical representation has order at most two, all its
second-order coefficients are real and satisfy
\eqref{gen:eq-real-unary}--\eqref{gen:eq-real-pair}, and every
\(\Im P_{e_i}\) is nonnegative on \(\R^d\).  In this formulation box
preservation forces
\(\Im P_{e_i}(z)=q_i(z_i)\) with \(\deg q_i\leq2\), and forces the
corresponding imaginary zeroth-order term in
\eqref{gen:eq-complex-decomposition}.
\end{theorem}

The two theorems identify ordinary and complete preservation for
\emph{semigroups starting at the identity}.  We will prove this by constructing
the generators.  It is not an assertion that an arbitrary isolated real
stability preserver preserves complex stability or admits every tensor
extension.

\subsection{Canonical differential operators and box preservation}
\label{gen:subsec-box}

We record the algebraic restrictions separately from the stability
inequalities.  All statements in this subsection hold over either \(\R\)
or \(\C\).

\begin{lemma}[Canonical normal representation]
\label{gen:canonical}
Every linear endomorphism of \(V_\kappa^\mathbb F\) has a unique
representation \eqref{gen:eq-canonical}.
\end{lemma}

\begin{proof}
Apply the proposed representation successively to monomials in increasing
total degree.  The constant monomial determines \(P_0=A1\).
Once \(P_\beta\) is known for all \(|\beta|<|\alpha|\), the equation on
\(z^\alpha\) determines
\[
 P_\alpha
 =\frac1{\alpha!}
 \left(Az^\alpha-
 \sum_{\substack{\beta\leq\alpha\\\beta\neq\alpha}}
 P_\beta(z)(\alpha)_\beta z^{\alpha-\beta}\right).
\]
Here \((\alpha)_\beta=\prod_i\alpha_i!/(\alpha_i-\beta_i)!\).
No derivative whose index is not bounded by \(\alpha\) acts nontrivially on
\(z^\alpha\).  The recursion therefore proves existence and uniqueness
on the monomial basis.
\end{proof}

Descriptions of differential operators on finite polynomial modules by
$\mathfrak{sl}_2$ enveloping algebras are classical in quasi-exact
solvability \cite{Turbiner92,Turbiner94}; we include the following
decomposition to make the box conventions and small-capacity degeneracies
explicit.

\begin{proposition}[Operators of order at most two on a box]
\label{gen:box-decomposition}
Every box endomorphism of canonical order at most two is a linear
combination of \(\Id\), the operators \(J_i^-,J_i^0,J_i^+\), and their
products of length two, including products at the same site.
Every endomorphism of order at most one has a unique expression
\begin{equation}
 c\Id+\sum_{i=1}^dJ_{i,\kappa_i}(q_i),
 \qquad c\in\mathbb F,\quad q_i\in\mathbb F[z_i]_{\leq2}.
 \label{gen:eq-first-order}
\end{equation}
Consequently the second-order coefficients have the variable dependence
and degree bounds stated in Theorem~\ref{gen:main-real}.  The vector space
of endomorphisms of order at most two has dimension
\begin{equation}
 1+3d+5\#\{i:\kappa_i\geq2\}+9\binom d2.
 \label{gen:eq-order-two-dimension}
\end{equation}
\end{proposition}

\begin{proof}
Write \(p_{\alpha,\beta}\) for the coefficient of \(z^\beta\) in
\(P_\alpha\), and group terms according to the integer monomial shift
\(\nu=\beta-\alpha\).  For input exponent
\(x\in\prod_i\{0,\ldots,\kappa_i\}\), the coefficient of \(z^{x+\nu}\) is
the restriction to this grid of
\begin{equation}
 F_\nu(x)=
 \sum_{\substack{|\alpha|\leq2,\ \alpha\leq\kappa\\
                  \alpha+\nu\geq0}}
 p_{\alpha,\alpha+\nu}(x)_\alpha.
 \label{gen:eq-shift-polynomial}
\end{equation}
This polynomial has total degree at most two and degree at most
\(\kappa_i\) in each \(x_i\).  We use the polynomial
\eqref{gen:eq-shift-polynomial} also when \(x+\nu\) is outside the exponent
grid.

Suppose first that \(\nu_i<0\).  Every summand has
\(\alpha_i\geq-\nu_i\), so \(F_\nu\) is divisible by
\((x_i)_{-\nu_i}\).  If \(-\nu_i>\kappa_i\), there are no summands and
\(F_\nu=0\).  In particular, whenever a grid point would give a negative
output exponent, this falling factorial makes \(F_\nu\) vanish there.

If \(\nu_i>0\), box preservation gives vanishing on every overflow face
\[
 x_i\in\{\kappa_i-\nu_i+1,\ldots,\kappa_i\}
          \cap\{0,\ldots,\kappa_i\}.
\]
To justify the divisibility as a polynomial identity, fix one such value
of \(x_i\).  On the remaining tensor grid the value of \(F_\nu\) is zero:
either the output has a negative coordinate, in which case the preceding
falling-factorial argument applies, or it is a genuine monomial outside
the box, whose coefficient must vanish.  Tensor-product interpolation,
using the coordinate degree bounds in
\eqref{gen:eq-shift-polynomial}, shows that the restriction to this face
is identically zero as a polynomial in all remaining variables.  Thus
the corresponding linear factor in \(x_i\) divides \(F_\nu\).
If \(\nu_i>\kappa_i\), the same argument gives vanishing on all
\(\kappa_i+1\) coordinate values and hence \(F_\nu=0\).

For every remaining nonzero shift polynomial, these factors give
\begin{equation}
 F_\nu=D_\nu G_\nu,\qquad
 D_\nu(x)=\prod_i(x_i)_{\nu_i^-}
                       (\kappa_i-x_i)_{\nu_i^+},
 \qquad
 \deg G_\nu\leq2-|\nu|_1.
 \label{gen:eq-shift-factor}
\end{equation}
The factors belonging to distinct coordinate values are coprime.
Thus \(F_\nu\neq0\) implies \(|\nu|_1\leq2\).

Put \(E_i=z_i\partial_i\).  For \(\nu=0\), every possible \(F_\nu\)
is realized by a combination of
\(\Id,E_i,E_i^2,E_iE_j\).
For \(\nu=\pm e_i\), the required constant or affine factor \(G_\nu\)
is realized by \(J_i^\pm\) and \(J_i^\pm E_j\).
For \(\nu=\pm2e_i\) use \((J_i^\pm)^2\); these shifts can be nonzero
only if \(\kappa_i\geq2\).
For \(\nu=\pm e_i\pm e_j\), with \(i\neq j\), use the corresponding
products \(J_i^\pm J_j^\pm\).
The sign of the raising coefficient
\(x_i-\kappa_i=-(\kappa_i-x_i)\) is absorbed in the scalar multiplier.
This accounts for every possibility in \eqref{gen:eq-shift-factor}
on the input grid, and hence for every matrix entry of \(A\).

This argument also covers binary sites.  For example, if
\(\kappa_i=1\), the canonical coordinate degree of \(F_\nu\) in
\(x_i\) is at most one; any inadmissible quadratic dependence has
already been excluded before the realization step.
Some displayed products may then be linearly dependent or zero, which
does not affect the spanning assertion.

Repeating the same argument with total degree at most one leaves only
a scalar and the three sitewise operators at each coordinate.
They are linearly independent: the off-diagonal shifts separate the
raising and lowering terms, while an affine polynomial in the diagonal
exponents vanishing on the full tensor grid has every coefficient zero.
This proves \eqref{gen:eq-first-order}, including uniqueness.

The first-order coefficients of \(J_i^-,J_i^0,J_i^+\) span
\(1,z_i,z_i^2\).  Products at one site therefore have principal
coefficients spanning all polynomials of degree at most four.
Products at two different sites span all biforms of bidegree at most
\((2,2)\).  At a site with \(\kappa_i\geq2\), these are canonical
second-order coefficients; at a binary site that derivative is absent.
The kernel of this principal-coefficient map is exactly the
order-at-most-one space.  Its dimension is \(1+3d\), and the respective
principal-coefficient spaces have dimensions five and nine.
Surjectivity gives \eqref{gen:eq-order-two-dimension}.
\end{proof}

The proposition says that the operators of canonical order at most two
form the image of the degree-at-most-two part of the enveloping algebra
of the sitewise \(\mathfrak{sl}_2\)'s.  This is an algebraic statement
about the whole box.  Stability will impose precisely the signs of
its principal coefficients and, in the complex case, an inward sign
on the imaginary first-order fields.

\subsection{Root clusters and infinitesimal necessity}
\label{gen:subsec-roots}

The relevant localization uses an isolated root cluster and is valid
even when the ambient polynomial degree changes.  No full-degree
tangent formula is applied to a lower-degree polynomial.
The underlying multiple-root contact principle is classical:
after rotation of the half-plane, it is contained in the polynomial
abscissa tangent theorem of Burke and Overton
\cite[Theorem~1.2 and Lemma~1.4]{BurkeOverton01}.
We give a direct moment proof to record the signs and the local
factorization needed for the operator argument.

\begin{lemma}[Lower-half-plane root clusters]
\label{gen:cluster-moments}
Let \(p_t(s)\), \(t\geq0\), be a curve of complex polynomials of uniformly
bounded degree whose coefficients are right differentiable at zero.
Suppose
\[
 p_0(s)=s^m g(s),\qquad m\geq1,\qquad g(0)\neq0,
\]
and that \(p_t\) has no zero in \(\HH\) for all sufficiently small \(t\).
Let \(C_t\) be the monic factor consisting of the \(m\) roots converging
to zero.  Its coefficients are right differentiable, and
\begin{equation}
 \left.\frac{d}{dt}C_t(s)\right|_{t=0}
 =c_1s^{m-1}+c_2s^{m-2},
 \qquad
 \Im c_1\geq0,\quad c_2\in\R,\quad c_2\leq0.
 \label{gen:eq-cluster-variation}
\end{equation}
When \(m=1\), the \(c_2\) term is absent.  If every \(p_t\) is real
rooted, then \(c_1\) is real as well.
\end{lemma}

\begin{proof}
Choose a small circle about zero containing no other zero of \(p_0\).
Coefficient continuity and Rouch\'e's theorem isolate exactly \(m\)
roots inside it for small \(t\).  If these roots are
\(\lambda_1(t),\ldots,\lambda_m(t)\), their power sums are
\[
 M_k(t)=\sum_{j=1}^m\lambda_j(t)^k
       =\frac1{2\pi i}\int s^k\frac{p_t'(s)}{p_t(s)}\,ds,
       \qquad k\geq1,
\]
where the contour is the chosen circle.  The integrands are right
differentiable with a denominator bounded away from zero.
Thus each \(M_k\) is right differentiable and \(M_k(t)=O(t)\).
Newton's identities give the same differentiability for the coefficients
of \(C_t\).

Write \(\lambda_j=a_j+ib_j\).  The roots lie in the closed lower
half-plane, so \(b_j\leq0\), and
\[
 \sum_j|b_j|=-\Im M_1(t)=O(t),\qquad
 \sum_jb_j^2=O(t^2).
\]
Moreover,
\[
 \sum_ja_j^2=\Re M_2(t)+\sum_jb_j^2=O(t).
\]
Consequently
\[
 \Im M_2(t)=2\sum_ja_jb_j=O(t^{3/2})=o(t),
\]
and, for \(k\geq3\),
\[
 |M_k(t)|
 \leq\sum_j|\lambda_j(t)|^k
 \leq\left(\sum_j|\lambda_j(t)|^2\right)^{k/2}
 =O(t^{k/2})=o(t).
\]
The second inequality uses the ordinary finite counting sum over
the roots.

Write
\(C_t=s^m+\gamma_1(t)s^{m-1}+\cdots+\gamma_m(t)\).
At \(t=0\), all \(\gamma_j\) and \(M_j\) vanish.
Differentiating Newton's identities gives
\(\gamma_j'(0)=-M_j'(0)/j\).
Hence every derivative below the coefficient of \(s^{m-2}\)
vanishes.  The first coefficient has nonnegative imaginary part
because \(\Im M_1(t)\leq0\).
For the second coefficient, \(M_2'(0)\) is real, and
\[
 \Re M_2'(0)=\lim_{t\downarrow0}
       \frac{\sum_ja_j(t)^2-\sum_jb_j(t)^2}{t}\geq0.
\]
This proves \eqref{gen:eq-cluster-variation}.  Real roots make
\(M_1\) real.
\end{proof}

\begin{proposition}[Order and sign from product contacts]
\label{gen:contact-necessity}
Suppose \(A\in\operatorname{End}_{\C}(V_\kappa^\C)\) generates a
complex stability-preserving semigroup.  Then \(A\) has canonical
order at most two, every second-order canonical coefficient is real
and nonpositive on \(\R^d\), and
\(\Im P_{e_i}\geq0\) on \(\R^d\).
If \(A\) is real and its semigroup is only assumed to preserve real
stability, the order bound and the nonpositivity of its
second-order coefficients still follow.
\end{proposition}

\begin{proof}
For \(0\neq\alpha\leq\kappa\) and \(y\in\R^d\), use the stable
product
\[
 f_{\alpha,y}(z)=\prod_i(z_i-y_i)^{\alpha_i}.
\]
In the canonical representation, every derivative at \(z=y\)
vanishes except for \(\partial^\alpha\), and therefore
\begin{equation}
 (Af_{\alpha,y})(y)=\alpha!P_\alpha(y).
 \label{gen:eq-contact-isolation}
\end{equation}
Choose \(v\in\R_{>0}^d\) and restrict the polynomial curve to
\(z=y+sv\).  Its initial value is
\[
 \left(\prod_iv_i^{\alpha_i}\right)s^{|\alpha|}.
\]
Its coefficients are differentiable and its degree is bounded by
\(|\kappa|\).  It has no zero for \(s\in\HH\), because such an \(s\)
puts every active coordinate in \(\HH\).
The restricted polynomial cannot be identically zero, since that
would contradict multivariate stability.  Invertibility of
\(e^{tA}\) also excludes the zero multivariate output.

Apply Lemma~\ref{gen:cluster-moments}.  Factoring the restriction
as \(C_t(s)G_t(s)\), its constant first variation is
\[
 G_0(0)\left.\frac{d}{dt}C_t(0)\right|_{t=0}.
\]
The term involving the derivative of \(G_t\) vanishes because \(C_0(0)=0\).
Here \(G_0(0)=\prod_iv_i^{\alpha_i}>0\).
For \(|\alpha|\geq3\), the constant variation is zero, so
\eqref{gen:eq-contact-isolation} gives \(P_\alpha(y)=0\).
For \(|\alpha|=2\), it is real and nonpositive, giving the
corresponding assertion for \(P_\alpha(y)\).
For \(\alpha=e_i\), its imaginary part is nonnegative.

Since \(y\) is arbitrary, polynomial identity proves the vanishing
of all coefficients of order greater than two.  A complex polynomial
real-valued on all of \(\R^d\) has real coefficients; this proves
reality of the principal coefficients.  For a real semigroup,
the same inputs and restrictions are real rooted, so the real
case follows from the same lemma.  Proposition~\ref{gen:box-decomposition}
then supplies the asserted variable dependence.
\end{proof}

The moment calculation identifies the obstruction to higher differential
order: one-sided control of the imaginary root depths, together with
the second root moment, eliminates every higher first-order cluster
moment.  Positivity of coefficients of the input polynomial is not
used; every real contact point is available.

\subsection{Binary apolar generators and sums of squares}
\label{gen:subsec-binary}

The sufficiency proof starts in two binary variables.
For
\(f=f_{00}+f_{10}s+f_{01}t+f_{11}st\in V_{(1,1)}^\R\), put
\[
 \Delta(f)=f_{10}f_{01}-f_{00}f_{11}.
\]
A nonzero real \(f\) is stable if and only if \(\Delta(f)\geq0\).
Indeed, if \(f=a+bs+ct+dst\) and \(b+dt\neq0\), its zero in \(s\)
is \(-(a+ct)/(b+dt)\), whose imaginary part equals
\(-\Delta(f)\Im t/|b+dt|^2\).
The cases in which \(f\) is independent of \(s\) follow directly.

\begin{lemma}[Low-degree nonnegative polynomials]
\label{gen:sos}
Every nonnegative real univariate polynomial of degree at most four is a sum
of squares of real polynomials of degree at most two.
Every nonnegative real polynomial of bidegree at most \((2,2)\)
is a sum of squares of real bilinear polynomials.
\end{lemma}

\begin{proof}
For a nonzero nonnegative univariate polynomial, every real root has
even multiplicity and its nonreal roots occur in conjugate pairs.
A positive real quadratic is a sum of two squares of linear
polynomials.  Repeated use of
\[
 (a^2+b^2)(c^2+d^2)=(ac-bd)^2+(ad+bc)^2
\]
gives the first assertion, with square degrees at most two.
The zero polynomial presents no exception.

For the second assertion let \(R(s,t)\geq0\) and homogenize it to
\[
 \widetilde R(x,y,u)=u^4R(x/u,y/u).
\]
The bidegree bound makes this a ternary quartic.  It is nonnegative
when \(u\neq0\), and hence everywhere by continuity.
Hilbert's theorem on nonnegative ternary quartics writes it as a sum
of squares of real quadratic forms; we use the version including
all degenerate nonnegative forms \cite{PRSS04}.
The coefficients of \(x^4\) and \(y^4\) in \(\widetilde R\) are zero.
They are sums of the squares of the \(x^2\) and \(y^2\) coefficients
of the quadratic forms, respectively.  Thus every form is a linear
combination of \(xy,xu,yu,u^2\).  Setting \(u=1\) yields bilinear
squares.
\end{proof}

\begin{lemma}[Apolar rank-one primitive]
\label{gen:binary-primitive}
For a real bilinear polynomial \(F=a+bs+ct+dst\), define
\[
 \ell_F(f)=d f_{00}-c f_{10}-b f_{01}+a f_{11},
 \qquad
 B_Ff=-F\ell_F(f).
\]
The full mixed canonical coefficient of \(B_F\) is \(-F^2\).
Moreover
\begin{align}
 B_F^2&=2\Delta(F)B_F,
 \label{gen:eq-apolar-square}\\
 \Delta(f+kB_Ff)
 &=\Delta(f)+k\bigl(1+k\Delta(F)\bigr)\ell_F(f)^2
 \qquad(k\in\R).
 \label{gen:eq-apolar-discriminant}
\end{align}
For all \(t\geq0\), \(e^{tB_F}\) preserves real stability.
\end{lemma}

\begin{proof}
For any endomorphism \(B\) of the binary box, its mixed coefficient is
\begin{equation}
 P_B(s,t)=B(st)-sB(t)-tB(s)+stB(1).
 \label{gen:eq-binary-coefficient}
\end{equation}
The four evaluations of \(B_F\) give \(P_{B_F}=-F^2\).
Also
\(\ell_F(F)=2ad-2bc=-2\Delta(F)\), which gives
\eqref{gen:eq-apolar-square}.
The polarization of the quadratic form \(\Delta\) satisfies
\(D\Delta_f[F]=-\ell_F(f)\).  Expanding
\(\Delta(f-kF\ell_F(f))\) proves
\eqref{gen:eq-apolar-discriminant}.

If \(D_F=\Delta(F)\), equation~\eqref{gen:eq-apolar-square} gives
\[
 e^{tB_F}=\Id+k_tB_F,\qquad
 k_t=
 \begin{cases}
 (e^{2tD_F}-1)/(2D_F),&D_F\neq0,\\
 t,&D_F=0.
 \end{cases}
\]
For \(t\geq0\), \(k_t\geq0\) and
\[
 1+k_tD_F=\frac{1+e^{2tD_F}}2>0.
\]
Thus \eqref{gen:eq-apolar-discriminant} preserves the inequality
\(\Delta(f)\geq0\).
The exponential is invertible, so a nonzero input does not become
the zero polynomial.  These formulas also include \(F=0\) and
\(D_F=0\).
\end{proof}

For use below, we make explicit the passage from real to complex
preservation in this binary semigroup.

\begin{lemma}[Orientation at the identity]
\label{gen:binary-completeness}
Let \(B\) be a real endomorphism of \(V_{(1,1)}^\R\).
Then \(e^{tB}\) preserves real stability for all \(t\geq0\) if and
only if its mixed coefficient \(P_B\) is nonpositive on \(\R^2\).
In this case every \(e^{tB}\) is a complete complex stability
preserver.
\end{lemma}

\begin{proof}
For necessity, take \(f=(s-x)(t-y)\).  It has \(\Delta(f)=0\).
Differentiating \(\Delta(e^{\tau B}f)\) at \(\tau=0\), and using
\eqref{gen:eq-binary-coefficient}, gives \(-P_B(x,y)\).
Real preservation implies that this derivative is nonnegative.

Conversely, Lemma~\ref{gen:sos} gives
\(-P_B=\sum_\ell F_\ell^2\).
The operator \(B-\sum_\ell B_{F_\ell}\) has zero mixed coefficient
and hence has canonical order at most one.
Proposition~\ref{gen:box-decomposition} writes it as a real scalar
plus two sitewise operators \(J_{i,1}(q_i)\), with real quadratic
\(q_i\), without a sign restriction on these two quadratics.

For clarity, the real sitewise flows preserve complex stability in
both time directions.  If \(q(z)=q_0+q_1z+q_2z^2\), let
\[
 M_\tau=\exp\left[
 \tau\begin{pmatrix}q_1/2&q_0\\-q_2&-q_1/2\end{pmatrix}
 \right]
 =\begin{pmatrix}a_\tau&b_\tau\\c_\tau&d_\tau\end{pmatrix}.
\]
This is a real matrix of determinant one.  The polynomial action
\begin{equation}
 U_\tau f(z)
 =(c_\tau z+d_\tau)^\kappa
 f\left(\frac{a_\tau z+b_\tau}{c_\tau z+d_\tau}\right)
 \label{gen:eq-real-mobius}
\end{equation}
has infinitesimal generator \(J_\kappa(q)\).
The quotient maps \(\HH\) to itself, since its imaginary part is
\(\Im z/|c_\tau z+d_\tau|^2\), and the denominator has no zero there.
Formula~\eqref{gen:eq-real-mobius} is polynomial after cancellation
term by term and holds for every real \(\tau\).
It also works with inert variables.

Lemma~\ref{gen:binary-primitive} and the finite-dimensional Lie
product formula now show that \(e^{\tau B}\) preserves real
stability.  Products of preserving factors converge coefficientwise;
Hurwitz's theorem and invertibility of the limiting exponential
exclude a zero output.

It remains to justify complex preservation of the apolar factors
and of the resulting binary semigroup.  We use the finite-box
classification of Borcea and Br\"and\'en
\cite[Theorems~1.1 and~1.2]{BB09}.
A real stability preserver of rank greater than two has a stable
plus symbol or a stable reflected-minus symbol.  For the binary
box these are
\[
 G_T^+(s,t;u,v)=T\bigl((s+u)(t+v)\bigr),\qquad
 G_T^-(s,t;u,v)=T\bigl((s-u)(t-v)\bigr),
\]
where \(T\) acts only in \(s,t\).
Every \(e^{\tau B}\) has rank four.
At \(\tau=0\) the minus symbol is \((s-u)(t-v)\), which is not
stable: it vanishes when \(s=u\in\HH\).
The cone consisting of stable polynomials and zero is closed
in this fixed coefficient space.  Continuity in \(\tau\) therefore
excludes the minus alternative for all sufficiently small
nonnegative \(\tau\).  The plus symbol is then stable, and the
complex classification gives complex preservation.

More explicitly, adjoining an identity operator on an auxiliary
box of capacity \(\lambda\) multiplies the plus symbol by
\(\prod_j(\xi_j+\eta_j)^{\lambda_j}\).  This remains stable.
The same complex classification proves preservation on every
enlarged box.  Finally, divide an arbitrary time into sufficiently
small equal pieces and compose.  This proves complete preservation
for every time and completes the lemma.
\end{proof}

This use of the symbol theorem is limited and precise.
The real rank-two exceptional case is not invoked: the binary
pair operator has rank four at every time.
A single binary site is handled directly by
\eqref{gen:eq-real-mobius}.

\subsection{Polarization and the proof of the real classification}
\label{gen:subsec-real-proof}

Let \(\mathcal P_\kappa\) be normalized polarization into
\(N=|\kappa|\) binary variables, grouped into \(\kappa_i\) slots at
site \(i\):
\[
 \mathcal P_\kappa(z_i^r)
 =\frac{e_r(z_{i,1},\ldots,z_{i,\kappa_i})}
             {\binom{\kappa_i}{r}}.
\]
Let \(\mathcal D_\kappa\) identify all slots belonging to each site.
The two maps are inverse between \(V_\kappa^\C\) and the subspace
symmetric within each group of slots.
Normalized polarization preserves upper-half-plane stability, and
diagonalization preserves it in the reverse direction
\cite[Proposition~2.4]{BB09}, with the normalization fixed in the preliminaries.

We will need the exact derivative factors
\begin{align}
 \mathcal D_\kappa\partial_{i,a}\mathcal P_\kappa f
 &=\frac{\partial_i f}{\kappa_i},
 \label{gen:eq-polar-one}\\
 \mathcal D_\kappa\partial_{i,a}\partial_{j,b}
                          \mathcal P_\kappa f
 &=\frac{\partial_i\partial_j f}{\kappa_i\kappa_j},
 &&i\neq j,
 \label{gen:eq-polar-cross}\\
 \mathcal D_\kappa\partial_{i,a}\partial_{i,b}
                          \mathcal P_\kappa f
 &=\frac{\partial_i^2f}{\kappa_i(\kappa_i-1)},
 &&a\neq b.
 \label{gen:eq-polar-same}
\end{align}
The last formula is used only for \(\kappa_i\geq2\).
They follow by differentiating the elementary symmetric polynomial
on a monomial and then diagonalizing.

\begin{proof}[Proof of Theorem~\ref{gen:main-real}]
The implication from complete complex preservation to real
preservation is immediate, since the operator is real.
Necessity of the order and sign conditions is
Proposition~\ref{gen:contact-necessity}.
We prove sufficiency constructively.

Write \(a_i=P_{2e_i}\) and \(b_{ij}=P_{e_i+e_j}\).
For each \(i<j\), apply Lemma~\ref{gen:sos} to obtain
\[
 -b_{ij}(s,t)=\sum_\ell F_{ij,\ell}(s,t)^2
\]
with real bilinear \(F_{ij,\ell}\).
The binary operator
\(B_{ij}=\sum_\ell B_{F_{ij,\ell}}\) has mixed coefficient
\(b_{ij}\) and generates a complete preserving semigroup.
On the full slot space, place a copy of \(B_{ij}\) on each pair
consisting of one slot at site \(i\) and one at site \(j\), and
sum these \(\kappa_i\kappa_j\) copies.
Every summand is a complete generator, so the sum is one by the
Lie product formula.
It commutes with permutations of slots at either site and therefore
preserves the symmetric subspace.

Descend this operator by \(\mathcal D_\kappa\) and
\(\mathcal P_\kappa\).
Equation~\eqref{gen:eq-polar-cross} shows that its full mixed
coefficient is exactly \(b_{ij}(z_i,z_j)\).
No other second-order coefficient is created.
Indeed, after writing each binary summand in normal order,
zeroth- and first-order terms remain of those respective orders
after diagonalization, and the only second-order derivative
involves its two selected slots.

For a site with \(\kappa_i\geq2\), write
\[
 -a_i(z)=\sum_\ell p_{i,\ell}(z)^2,\qquad
 p_{i,\ell}(z)=a_\ell z^2+b_\ell z+c_\ell,
\]
using the univariate part of Lemma~\ref{gen:sos}.
Set
\[
 F_{i,\ell}(s,t)=a_\ell st+\frac{b_\ell}{2}(s+t)+c_\ell.
\]
These forms are symmetric in their two variables and satisfy
\(F_{i,\ell}(z,z)=p_{i,\ell}(z)\).
The symmetric binary operator
\(2\sum_\ell B_{F_{i,\ell}}\) has mixed coefficient
\[
 P_i(s,t)=-2\sum_\ell F_{i,\ell}(s,t)^2,
 \qquad P_i(z,z)=2a_i(z).
\]
Symmetry of the operator follows directly from symmetry of both
\(F_{i,\ell}\) and its apolar functional.
Sum copies of this operator over all unordered pairs of distinct
slots at site \(i\).
It is again a complete generator and commutes with all slot
permutations at that site.
Its descended same-site principal coefficient, by
\eqref{gen:eq-polar-same}, is
\[
 \binom{\kappa_i}{2}
 \frac{P_i(z_i,z_i)}{\kappa_i(\kappa_i-1)}
 =a_i(z_i).
\]
This is the reason for the factor two in the binary operator.
There is no same-site construction when \(\kappa_i=1\).

Let \(A_{\mathrm{pr}}\) be the sum of all descended generators.
For each summand, permutation invariance implies that its
exponential on the symmetric subspace is intertwined by
polarization with the exponential of the descended operator.
Polarization, complete preservation on the full slot space,
and diagonalization therefore show that \(A_{\mathrm{pr}}\)
generates a complete preserving semigroup on the original box.
It has exactly the same second-order canonical coefficients as
\(A\).

The difference \(A-A_{\mathrm{pr}}\) has order at most one.
By Proposition~\ref{gen:box-decomposition}, it is a real scalar
plus a sum of real sitewise \(J_{i,\kappa_i}(q_i)\).
Formula~\eqref{gen:eq-real-mobius} proves complete preservation
for each of their flows, with no sign restriction on \(q_i\).
A final Lie product formula proves complete preservation for
\(e^{tA}\).
The argument takes place on any fixed enlarged finite box,
and the limiting exponential is invertible there, so Hurwitz's
zero alternative never occurs for a nonzero stable input.
This proves the theorem.
\end{proof}

\subsection{Imaginary inward fields and the complex classification}
\label{gen:subsec-complex-proof}

The remaining complex generators have a particularly explicit
one-site realization.

\begin{lemma}[Imaginary parabolic primitive]
\label{gen:imaginary-primitive}
Let \(a,b\in\R\) and \(\kappa\geq1\).
The generator \(iJ_\kappa((az+b)^2)\) has the semigroup
\begin{equation}
 (U_tf)(z)=d_t(z)^\kappa f(\phi_t(z)),\qquad
 d_t(z)=1-iat(az+b),\qquad
 \phi_t(z)=\frac{z+ibt(az+b)}{1-iat(az+b)}.
 \label{gen:eq-imaginary-flow}
\end{equation}
For every \(t\geq0\), this is a complete complex stability preserver.
\end{lemma}

\begin{proof}
The corresponding fractional-linear matrix is
\[
 M_t=\Id+it
 \begin{pmatrix}ab&b^2\\-a^2&-ab\end{pmatrix}.
\]
The displayed nonidentity matrix has square zero, so
\(M_{t+s}=M_tM_s\) and \(\det M_t=1\).
Denominators in \eqref{gen:eq-imaginary-flow} cancel on each
monomial of degree at most \(\kappa\), giving a polynomial
operator and an exact semigroup.
Differentiating at zero gives
\[
 \left.\frac{d}{dt}U_tf\right|_{t=0}
 =i\bigl((az+b)^2f'-\kappa a(az+b)f\bigr)
 =iJ_\kappa((az+b)^2)f.
\]

For \(z\in\HH\) and \(t\geq0\), direct multiplication by the
conjugate denominator gives
\begin{equation}
 \Im\phi_t(z)
 =\frac{\Im z+t|az+b|^2}{|1-iat(az+b)|^2}>0.
 \label{gen:eq-imaginary-height}
\end{equation}
When \(a\neq0\) and \(t>0\), the denominator has its only zero at
\(-b/a-i/(a^2t)\), below the real axis.
When \(a=0\), it is identically one and
\(\phi_t(z)=z+ib^2t\).
Thus the substitution maps \(\HH\) to itself and the prefactor
does not vanish there.
The same argument with all other variables held inert proves
complete preservation.  The cases \(t=0\) and \(a=b=0\) give
the identity.
\end{proof}

\begin{proof}[Proof of Theorem~\ref{gen:main-complex}]
Assume ordinary complex preservation and decompose \(A=R+iS\)
by entrywise real and imaginary parts in the ordinary monomial
basis.  Both \(R\) and \(S\) are real box endomorphisms.
The canonical recursion in Lemma~\ref{gen:canonical} is real
linear, so it commutes with this decomposition.

By Proposition~\ref{gen:contact-necessity}, \(A\) has order
at most two and its second-order coefficients are real and
nonpositive.  The same coefficients are therefore the
principal coefficients of \(R\), which satisfies
Theorem~\ref{gen:main-real}.  Meanwhile \(S\) has order at most
one.  Proposition~\ref{gen:box-decomposition} gives
\[
 S=c\Id+\sum_iJ_{i,\kappa_i}(q_i),
 \qquad c\in\R,\quad q_i\in\R[z_i]_{\leq2}.
\]
Its first-order coefficients are exactly
\(q_i=\Im P_{e_i}\); the product-contact test makes them
nonnegative on \(\R\).  This proves the necessity of
\eqref{gen:eq-complex-decomposition}.

Conversely, every real quadratic nonnegative on \(\R\) is a
sum of squares of real linear polynomials.
For \(q(z)=az^2+bz+c\) with \(a>0\), complete the square:
\[
 q(z)=\left(\sqrt a\,z+\frac b{2\sqrt a}\right)^2
            +c-\frac{b^2}{4a}.
\]
The last constant is nonnegative.
If \(a=0\), nonnegativity forces \(b=0\) and \(c\geq0\).
Lemma~\ref{gen:imaginary-primitive} and the Lie product formula
therefore give complete preservation for each \(iJ(q_i)\).
The generator \(R\) has the same property by the real theorem,
and \(ic\Id\) contributes a nonzero scalar phase.
One further product formula proves complete preservation for
their sum.  Invertibility of the finite-dimensional exponential
again excludes the zero limit.
Complete preservation implies ordinary preservation, finishing
the equivalence.
\end{proof}

In particular, every complex second-order coefficient in an
admissible generator is forced to be real.  The imaginary part
has no pair term: it consists entirely of sitewise inward
quadratic fields and a scalar.
The sign in \eqref{gen:eq-imaginary-flow} is fixed by the
upper-half-plane convention: \(i\partial\) shifts the argument
upward and moves polynomial zeros downward.

\subsection{Finite certificates, quotient cones, and exact lifts}
\label{gen:subsec-certificates}

The construction records both a finite membership certificate
and a representation by local binary generators.

\begin{corollary}[Gram certificates and lineality]
\label{gen:gram-cones}
In Theorem~\ref{gen:main-real}, the principal inequalities are
equivalent to the existence of real positive semidefinite matrices
\(G_i\) of size three and \(G_{ij}\) of size four such that
\begin{align}
 -P_{2e_i}(s)
  &=(1,s,s^2)G_i(1,s,s^2)^{\top},
  &&\kappa_i\geq2,
 \label{gen:eq-gram-unary}\\
 -P_{e_i+e_j}(s,t)
  &=(1,s,t,st)G_{ij}(1,s,t,st)^{\top},
  &&i<j.
 \label{gen:eq-gram-pair}
\end{align}
The extra inequalities in Theorem~\ref{gen:main-complex} have
size-two certificates
\begin{equation}
 q_i(s)=(1,s)K_i(1,s)^{\top},\qquad K_i\succeq0.
 \label{gen:eq-gram-imaginary}
\end{equation}
These matrix sizes do not depend on the numerical capacities.

Let \(\mathcal N_4,\mathcal N_{2,2},\mathcal N_2\) denote,
respectively, the cones of nonnegative real univariate quartics
of degree at most four, nonnegative real biforms of bidegree
at most \((2,2)\), and nonnegative real univariate quadratics.
The real generator cone has lineality
\[
 \mathfrak g_\kappa
 =\R\Id+\sum_iJ_{i,\kappa_i}(\R[z_i]_{\leq2}),
\]
and its quotient by this lineality is naturally isomorphic to
\begin{equation}
 \prod_{i:\kappa_i\geq2}\mathcal N_4
 \ \times\!
 \prod_{i<j}\mathcal N_{2,2}.
 \label{gen:eq-real-quotient}
\end{equation}
The complex generator cone, regarded as a real cone, has
lineality
\[
 \C\Id+\sum_iJ_{i,\kappa_i}(\R[z_i]_{\leq2}),
\]
and its quotient is the product in
\eqref{gen:eq-real-quotient} times \(\prod_i\mathcal N_2\).
\end{corollary}

\begin{proof}
The square decompositions in Lemma~\ref{gen:sos} give
\eqref{gen:eq-gram-unary} and \eqref{gen:eq-gram-pair} by
taking sums of outer products of the coefficient vectors.
Conversely a real positive semidefinite matrix is a sum
of rank-one positive semidefinite matrices, yielding the
required square decomposition.
The quadratic case is the same argument with the vector
\((1,s)\).

The negative-principal-coefficient map has kernel precisely
the order-at-most-one space, by
Proposition~\ref{gen:box-decomposition}.
The polarization construction realizes every tuple of
nonnegative principal polynomials, so it is onto the product
in \eqref{gen:eq-real-quotient}.
A real \(A\) and its negative are both admissible if and only
if all their second-order coefficients vanish.
This gives exactly the stated real lineality and quotient.

For complex generators also record the imaginary first-order
polynomials \(q_i\).
Their nonnegativity is independent of the real principal
coefficients, as
\eqref{gen:eq-complex-decomposition} shows.
Both signs of a generator can be admissible only if all these
\(q_i\) and all second-order coefficients vanish.
The remaining imaginary scalar is unrestricted.
This proves the complex assertion.
\end{proof}

\begin{corollary}[Exact binary-slot lift]
\label{gen:slot-lift}
For every generator \(A\) in Theorem~\ref{gen:main-complex},
there is a complete stability-preserving generator
\(\widetilde A\) on the full multiaffine space with
\(|\kappa|\) variables such that
\begin{equation}
 \widetilde A\mathcal P_\kappa=\mathcal P_\kappa A,
 \qquad
 e^{t\widetilde A}\mathcal P_\kappa
          =\mathcal P_\kappa e^{tA}\quad(t\geq0).
 \label{gen:eq-exact-lift}
\end{equation}
It commutes with every permutation of slots within a site.
It can be chosen as a sum of a scalar, one-slot real
M\"obius generators, one-slot imaginary parabolic generators,
and nonnegative multiples of the rank-one apolar generators
\(B_F\) on pairs of slots.
\end{corollary}

\begin{proof}
For the real second-order part use exactly the slot sums
constructed in the proof of
Theorem~\ref{gen:main-real}.
They commute with within-site permutations and descend to the
chosen principal-coefficient realization.
Lift each remaining sitewise drift by
\[
 J_{i,\kappa_i}(q)
 \longmapsto
 \sum_{a=1}^{\kappa_i}
 \left(q(z_{i,a})\partial_{i,a}
                    -\frac12q'(z_{i,a})\right).
\]
Equation~\eqref{gen:eq-polar-one} shows that its diagonalization
on a polarized input is precisely \(J_{i,\kappa_i}(q)\).
The lifted sum is permutation invariant, so this equality
upgrades to intertwining with \(\mathcal P_\kappa\).
The same lift applies to the imaginary quadratics, which are
sums of the parabolic primitives of
Lemma~\ref{gen:imaginary-primitive}.
The scalar lifts as the same scalar identity.

All summands generate complete preserving flows on the full
slot space, hence so does their sum.
This proves the first relation in
\eqref{gen:eq-exact-lift}; exponentiation gives the second.
\end{proof}

The lift is a generator intertwining, not merely a static
polarization followed by an unrelated preserving map.
It is not asserted to have a nonnegative coefficient matrix.
Consequently it imposes no coefficient-positive lifting
principle on a different preservation problem.
Likewise, the signed generation statements proved above are
finite dimensional; they contain no infinite-capacity domain assertion.

\subsection{A degree-two ancillary test for positive semigroups}

The preceding classifications concern preservation of all stable inputs.
For positive semigroups, preservation of stable inputs with nonnegative
coefficients is a weaker condition.  The exact gap can be detected with
only two additional polynomial degrees.  This assertion also has a
mixed-capacity version, using the positive classification proved in
Section~\ref{pos:section}.

Let \(\kappa\in(\N_{>0}\cup\{\infty\})^d\), and let \(T(t)\) be an existing
coefficientwise-positive \(C_0\)-semigroup on
\(X_\kappa=\ell^1(\Omega_\kappa)\), with generator \(A\).
For a finite list of ancillary capacities \(\lambda\), the inert extension
\(T(t)\otimes\Id_\lambda\) acts coefficientwise in the original variables.
In the next theorem preservation means preservation of the corresponding
positive Laguerre--P\'olya class; equivalently, it suffices to test the
allowed stable polynomial inputs with nonnegative coefficients.
No assertion of semigroup generation is inferred from a formal operator.

\begin{theorem}[The sharp ancillary test]\label{gen:ancilla}
The following conditions are equivalent.
\begin{enumerate}
\item \(T(t)\otimes\Id_{(1,1)}\) preserves positive stability for every
      \(t\geq0\).
\item \(T(t)\otimes\Id_{(2)}\) preserves positive stability for every
      \(t\geq0\).
\item Every inert extension by finitely many finite- or infinite-capacity
      coordinates preserves positive stability for every \(t\geq0\).
\item All allowed polynomials belong to \(D(A)\), the generator has the
      mixed-capacity rate form of Theorem~\ref{pos:mixed}, and its symmetric
      second-order matrix in \eqref{pos:principal-matrix} satisfies
      \begin{equation}\label{gen:eq-entrywise-cone}
       Q_{ij}(y)\leq0\qquad(y\in\R^d,\ 1\leq i,j\leq d).
      \end{equation}
\item The same domain and rate form hold, and the rate parameters satisfy
      \begin{equation}\label{gen:eq-positive-complete-rates}
       \alpha_i=\chi_i=0,\quad \tau_i\leq0,\qquad
       a_{ij}=g_{ij}=0,\qquad
       \theta_{ij}^2\leq4c_{ij}c_{ji}.
      \end{equation}
      The capacity conventions of Theorem~\ref{pos:mixed} remain in force;
      in particular \(\tau_i=0\) for binary sites and \(c_{ij}=0\) when
      the destination \(j\) has infinite capacity.
\end{enumerate}
The two ancillary degrees are necessary uniformly over capacity vectors:
one inert binary coordinate cannot replace the two in condition (1).
For finite \(\kappa\), these conditions are also equivalent to complete
preservation of arbitrary complex stable polynomials by \(T(t)\).
\end{theorem}

\begin{proof}
First note that
\(X_{\kappa\sqcup\lambda}=\ell^1(\Omega_\lambda;X_\kappa)\).
The inert extension is therefore a positive \(C_0\)-semigroup of norm
\(\|T(t)\|\).  If polynomials belong to \(D(A)\), each finite polynomial
in the enlarged space belongs to its generator domain, and its generator
is \(A\otimes\Id\) there.  This verifies the analytic hypotheses whenever
Theorem~\ref{pos:mixed} is applied below.  On a preserving side that theorem
itself supplies the polynomial domain.  Clearly (3) implies (1) and (2).

Suppose (1) holds.  Inputs independent of the ancillary variables give
ordinary positive preservation.  Theorem~\ref{pos:mixed} supplies the rate
form, and its application on the enlarged space gives the contact test
\eqref{pos:contact-test} with \(\widetilde Q=Q\oplus0\).
Let \(u,v\) denote the ancillary coordinates.  For distinct original
indices \(i,j\), fix any \(y\in\R^d\), and set
\[
 a=e_i+e_u,\qquad b=e_j+e_v,\qquad
 H=\frac{ab^\top+ba^\top}{8},\qquad
 \widetilde y=(y,-y_i,-y_j).
\]
Here \(e_i\) are coordinate vectors and \(e\) below is the all-ones vector.
Direct calculation gives
\begin{equation}\label{gen:eq-ancillary-contact}
 \begin{gathered}
 H_{ab}\geq0,\qquad e^\top He=1,\qquad H\widetilde y=0,\\
 (He)(He)^\top-H=\frac{(a-b)(a-b)^\top}{16}\succeq0.
 \end{gathered}
\end{equation}
All diagonals required to vanish at binary sites do vanish.  Thus this is
an admissible contact in \eqref{pos:contact-domain}, and
\[
 0\geq\operatorname{tr}(\widetilde Q(\widetilde y)H)=Q_{ij}(y)/4.
\]
For \(i=j\) use \(a=e_i+e_u\), \(b=e_i+e_v\).
The same identities hold; the only nonzero original diagonal of \(H\)
is at site \(i\), so the construction is allowed when \(\kappa_i\geq2\).
It again gives \(Q_{ii}(y)\leq0\).  For binary sites that coefficient is
zero by canonical convention.  Hence (1) implies (4).

For (2), again obtain the original rate form first.  Fix \(y\) and the
indices \(i,j\), permitting \(i=j\) only at a nonbinary site.  Choose
\(R\geq\max(y_i,y_j)\), and put \(c=R-y_i\), \(d=R-y_j\).
With the single capacity-two ancillary coordinate \(w\), set
\[
 a=e_i+e_w,\quad b=e_j+e_w,\quad
 H=(ab^\top+ba^\top)/8,\quad \widetilde y=(y,-R).
\]
The normalization and semidefinite identity in
\eqref{gen:eq-ancillary-contact} still hold, while
\[
 H\widetilde y=-(da+cb)/8\leq0,\qquad
 \widetilde y^\top H\widetilde y=cd/4\geq0.
\]
The diagonal in coordinate \(w\) is now allowed.  The positive stable
quadratic is a positive scalar multiple of
\((z_i+w+c)(z_j+w+d)\).  The contact inequality once more isolates
\(Q_{ij}(y)/4\), proving (2)$\Rightarrow$(4).

A univariate polynomial \(\alpha s+\tau s^2+\chi s^3\) is nonpositive
on \(\R\) exactly when \(\alpha=\chi=0\) and \(\tau\leq0\).
For the bivariate entry in \eqref{pos:principal-matrix}, simultaneous
scaling of its two variables to both signs of infinity kills the cubic
homogeneous part; scaling towards zero kills the linear part.  The remaining
quadratic is nonpositive exactly when
\[
 \begin{pmatrix}c_{ji}&-\theta_{ij}/2\\
                 -\theta_{ij}/2&c_{ij}\end{pmatrix}\succeq0.
\]
The nonnegativity of the off-diagonal rate parameters makes this condition
precisely \eqref{gen:eq-positive-complete-rates}.  This proves (4)$\Leftrightarrow$(5).

Assume (4) and adjoin any finite list of ancillary capacities.  The
extended rate form is unchanged, with all new rates zero.  Every contact
matrix \(H\) in \eqref{pos:contact-domain} is entrywise nonnegative, so
\[
 \operatorname{tr}((Q(y)\oplus0)H)\leq0.
\]
Theorem~\ref{pos:mixed}, with the tensor-domain observation above, gives
preservation.  This proves (4)$\Rightarrow$(3).
In finite dimensions, (4) is exactly the second-order sign condition of
Theorem~\ref{gen:main-real}: the full mixed differential coefficient is
\(2Q_{ij}\), and the diagonal coefficient is \(Q_{ii}\).
That theorem proves the final complete-preservation assertion and its converse.

For sharpness take an active capacity-two coordinate \(z\) and
\(A=(z-z^2)\partial_z^2\).  With one inert binary coordinate its contact
matrix has the form \(H=\left(\begin{smallmatrix}h&a\\a&0\end{smallmatrix}\right)\),
where \(h,a\geq0\).  If \(h>0\), the condition \(Hy\leq0\) forces
\(y_z\leq0\), using the second row when \(a>0\), and the first when
\(a=0\).  Consequently
\(\operatorname{tr}(Q(y)H)=h(y_z-y_z^2)\leq0\); for \(h=0\) it is zero.
The finite positive criterion proves one-binary preservation.
With two inert binary variables the stable input \((z+u)(z+v)\) evolves to
\[
 e^{-2t}z^2+z(u+v)+uv+(1-e^{-2t})z.
\]
Identifying \(u=v=w\) and specializing \(z=r\in(0,1)\) gives a real
quadratic in \(w\) of discriminant
\(-4(1-e^{-2t})r(1-r)<0\) for \(t>0\).  Stability is preserved by both
operations, so the evolved polynomial is not stable.  This proves the
uniform sharpness, without asserting that every fixed capacity vector
requires two ancillary degrees.
\end{proof}

The complete positive cone has a particularly small pairwise representation.
Apart from scalar and linear Euler terms, its one-site and directed
first-order terms are deaths, finite births, infinite immigration,
transfers to infinite sites, and the damping operators
\(-\eta_iE_i(E_i-1)\), where \(E_i=z_i\partial_i\) and \(\eta_i\geq0\).
A pair with one finite and one infinite site additionally permits the
saturated transfer into the finite site,
\(c_{ij}z_j\partial_i(\kappa_j-E_j)\) for \(i\in U\), \(j\in F\),
with \(c_{ij}\geq0\); its pair diagonal parameter is necessarily zero.
For two finite sites the remaining terms are exactly
\begin{equation}\label{gen:eq-positive-pair}
 c_{ij}z_j\partial_i(\kappa_j-E_j)
 +c_{ji}z_i\partial_j(\kappa_i-E_i)+\theta_{ij}E_iE_j,
 \qquad \theta_{ij}^2\leq4c_{ij}c_{ji}.
\end{equation}
Indeed expansion recovers \eqref{pos:principal-matrix} and the prescribed
first-order rates; all remaining diagonal coefficients are scalar and
linear Euler terms.  Each pair cone is linearly isomorphic to
\(\mathbb S_+^2\).  Spectral decomposition therefore expresses every pair
as the sum of at most two rank-one terms, whose parameters have the form
\(c_{ij}=u^2\), \(c_{ji}=v^2\), \(\theta_{ij}=\pm2uv\).
This representation concerns positive stability and is distinct from the
signed apolar slot lift in Corollary~\ref{gen:slot-lift}.

% END SECTION generators

% BEGIN SECTION spectral
\section{Spectral contraction and sector growth}\label{spec:section}

We first study arbitrary complex stability-preserving
semigroups on a finite polynomial space.  No self-adjointness, positivity of
coefficients, or differential representation of the generator is assumed.
Positive damping of polynomial degree has two consequences: it improves the
zero-free domain of the actual evolution, and it gives a sharp contraction
estimate in a norm determined by stable polynomials.
The last subsection returns to real half-plane stability and assumes
coefficient positivity and preservation of homogeneous degree.
It compares spectral bounds between sectors without damping.

\subsection{Weighted kernels and the stable cone}

Zero-capacity coordinates may be omitted.  If no coordinate remains,
the kernel and eigenvector assertions reduce to nonzero constants.
For the gap statements assume that at least one coordinate remains.
Fix \(d\ge1\) and \(\kappa\in\N_{>0}^d\), write
\[
 V_\kappa=\{p\in\C[z_1,\ldots,z_d]:\deg_{z_i}p\le\kappa_i\},
 \qquad b_\alpha=\binom\kappa\alpha,
\]
and let \(\D_r^d=\{z:|z_i|<r\text{ for all }i\}\).
Let \(\mathcal C_\kappa\) consist of zero and all polynomials in
\(V_\kappa\) that do not vanish on \(\D_1^d\).  We call a polynomial
\emph{strictly stable} if it does not vanish on \(\overline{\D_1^d}\).
For \(T\in\operatorname{End}_{\C}(V_\kappa)\), use the kernel
of \eqref{pre:kernel}, namely
\begin{equation}\label{spec:kernel}
 K_T(z,w)=T_z\prod_{i=1}^d(1+z_iw_i)^{\kappa_i}
 =\sum_{\alpha\le\kappa}b_\alpha(Tz^\alpha)(z)w^\alpha.
\end{equation}
The corresponding coefficient pairing is bilinear:
\begin{equation}\label{spec:pairing}
 [f,g]_\kappa=\sum_{\alpha\le\kappa}b_\alpha^{-1}f_\alpha g_\alpha.
\end{equation}
In particular, \([K_T(z,\cdot),p]_\kappa=Tp(z)\).  There is no
complex conjugation in this identity.  In the orthonormal coherent basis
\(\sqrt{b_\alpha}z^\alpha\), the matrix kernel has one square-root
binomial weight on each of its two legs.

We shall use the classical disk contraction theorem, Lemma~\ref{pre:grace},
in the following form.
If two nonzero stable coefficient tensors have matched variables of the
same capacities, their contraction by \eqref{spec:pairing} is stable or
zero whenever each pair of contracted radii has product at least one.
If every such product is greater than one, the contraction is nonzero.
The radii of all uncontracted variables are retained.  This is the
Grace--Asano contraction principle, with the finite-capacity formulation
obtained by polarization; see \cite{BB09,BB10II,WBG26}.  Indeed, a
coefficient of multidegree \(\alpha\) is divided by \(b_\alpha\)
under normalized polarization, and summation over its \(b_\alpha\)
equal subsets leaves exactly the reciprocal weight in
\eqref{spec:pairing}.  Applied to two kernels, this pairing gives the
kernel of their operator product.

\begin{lemma}\label{spec:kernel-criterion}
An invertible operator \(T\) preserves \(\mathcal C_\kappa\) if and
only if \(K_T\) is zero-free on \(\D_1^d\times\D_1^d\).
If \(K_T\) is zero-free on \(\D_r^d\times\D_r^d\), where \(r>1\),
then \(Tp\) is nonzero and zero-free on \(\D_r^d\) for every
\(p\in\mathcal C_\kappa\setminus\{0\}\), without an invertibility
assumption on \(T\) in this last assertion.
\end{lemma}
\begin{proof}
For fixed \(w\in\D_1^d\), the polynomial
\(\prod_i(1+z_iw_i)^{\kappa_i}\) is stable and nonzero.  Its image
under an invertible stability preserver is stable and nonzero.  Evaluating
at any \(z\in\D_1^d\) proves the forward implication.  Conversely,
contract the input variables of \(K_T\) against a stable polynomial.
The contraction theorem gives a stable polynomial or zero; invertibility
excludes zero on a nonzero input.  For the final assertion the contracted
radii have product \(r>1\), so the strict contraction theorem itself
excludes zero and retains output radius \(r\).
\end{proof}

\begin{lemma}\label{spec:cone-interior}
The set \(\mathcal C_\kappa\) is a closed, complex-scalar-invariant
cone containing no complex two-dimensional subspace.  Its interior consists
exactly of the strictly stable polynomials.  Every invertible operator
preserving this cone maps its interior into its interior.
\end{lemma}
\begin{proof}
Closedness follows from Hurwitz's theorem on the connected polydisk.
Every complex two-dimensional subspace contains a nonzero polynomial
annihilated by evaluation at zero, and such a polynomial is not stable.
A polynomial nonzero on the closed polydisk remains so under sufficiently
small coefficient perturbations, by compactness.  Conversely, a boundary
zero \(\zeta\) of a stable polynomial moves into the open polydisk for
the arbitrarily small perturbation \(p((1+\varepsilon)z)\), at
\(z=\zeta/(1+\varepsilon)\).  Zero is not an interior point, since
arbitrarily small nonzero multiples of \(z_1\) do not belong to the
cone.  Finally, an invertible linear operator is open, so its image of
the cone interior is an open subset of the cone.
\end{proof}

\subsection{Strictification by degree damping}

Put \(E_i=z_i\partial_i\), \(E=\sum_iE_i\), and \(P_0p=p(0)\).
The following theorem improves the zero-free domain of the exponential
itself, rather than that of an auxiliary splitting operator.

\begin{theorem}[Degree-damping strictification]\label{spec:strict}
Suppose \(A_0\in\operatorname{End}_{\C}(V_\kappa)\) satisfies
\(e^{tA_0}\mathcal C_\kappa\subseteq\mathcal C_\kappa\) for
\(t\ge0\).  Let \(\mu_i>0\) and
\[
 A=A_0-2\sum_{i=1}^d\mu_i E_i.
\]
For every \(t>0\), the kernel \(K_{e^{tA}}\) is nonzero on the
closed unit polydisk in all \(2d\) variables.  Consequently it has a
zero-free radius \(r_t>1\).  The same radius works for the image of
every nonzero unit-disk-stable input polynomial.

Moreover, for each \(t_0>0\), one radius \(r>1\) works simultaneously
for the kernels \(K_{e^{tA}}\), \(t\ge t_0\), and for their action
on all nonzero unit-disk-stable inputs.
\end{theorem}

\begin{proof}
First consider uniform damping \(A=A_0-2\mu E\).  Fix \(t>0\)
and allow \(\mu\) to vary in the complex half-plane \(\Re\mu>0\).
For \(s\ge0\),
\[
 e^{-2s\mu E}p(z)=p(e^{-2s\mu}z)
\]
preserves disk stability.  Thus every Lie product
\[
 \left(e^{tA_0/m}e^{-2t\mu E/m}\right)^m
\]
preserves stability and is invertible.  The finite-dimensional product
formula, followed by Hurwitz and invertibility of the limiting
exponential, proves preservation by
\(T_t(\mu)=e^{t(A_0-2\mu E)}\).  Lemma
\ref{spec:kernel-criterion} gives
\begin{equation}\label{spec:open-complex-field}
 K_{T_t(\mu)}(z,w)\ne0
 \quad\text{if }z,w\in\D_1^d,\quad\Re\mu>0.
\end{equation}
Each kernel coefficient is an entire function of \(\mu\).

We next establish a real large-damping limit, valid without normality.
Let \(a_{00}=(A_01)(0)\).  In the Euclidean norm on monomial
coefficients, \(E\) is diagonal, nonnegative, and has the constants
as its zero eigenspace.  Hence, for real \(\mu>0\),
\[
 \|e^{-2s\mu E}\|\le1,
 \qquad e^{-2s\mu E}\longrightarrow P_0\quad(s>0).
\]
Expand \(T_t(\mu)\) in its Dyson series about \(-2\mu E\).
The term of order \(k\ge1\) is the integral over
\(0<s_1<\cdots<s_k<t\) of
\[
 e^{-2\mu(t-s_k)E}A_0e^{-2\mu(s_k-s_{k-1})E}
 \cdots A_0e^{-2\mu s_1E}.
\]
Every time interval is positive almost everywhere.  The integrand
therefore converges to \(a_{00}^kP_0\), and its norm is bounded by
\(\|A_0\|^k\), independently of \(\mu\).  Dominated convergence
on the simplex gives the limit \(t^ka_{00}^kP_0/k!\).  These integrated
terms are uniformly bounded by the summable sequence
\(t^k\|A_0\|^k/k!\).  Including the order-zero term and passing
the limit through the sum proves
\begin{equation}\label{spec:large-damping}
 T_t(\mu)\longrightarrow e^{ta_{00}}P_0
 \quad\text{in operator norm as real }\mu\longrightarrow+\infty.
\end{equation}
This is a finite-dimensional instance of strong-coupling limits such as
\cite[Theorem~1]{Burgarth19}; the displayed argument fixes its
normalization directly.  Since \(K_{P_0}=1\), the limiting kernel is
the nonzero constant \(e^{ta_{00}}\).

Fix any \((z,w)\in\overline{\D_1^d}\times\overline{\D_1^d}\).
For \(0<\rho<1\), the functions
\[
 f_\rho(\mu)=K_{T_t(\mu)}(\rho z,\rho w)
\]
are holomorphic and nonvanishing on \(\Re\mu>0\).  Their limit as
\(\rho\uparrow1\) is locally uniform in that half-plane, because
the kernel has finitely many entire coefficient functions.  Hurwitz
therefore says that \(K_{T_t(\mu)}(z,w)\), as a function of
\(\mu\), is identically zero or nowhere zero.  The limit
\eqref{spec:large-damping} excludes the first alternative.  This proves
closed-polydisk nonvanishing at the desired real \(\mu>0\).
Compactness then supplies a radius \(r_t>1\).  Lemma
\ref{spec:kernel-criterion} gives the common input-output assertion.

For nonuniform rates, let \(\mu=\min_i\mu_i>0\).  The operator
\(A_0-2\sum_i(\mu_i-\mu)E_i\) generates a stable semigroup by the
same product-formula argument.  Apply the uniform result to this
operator and the remaining damping \(-2\mu E\).

Finally fix \(t_0>0\), and choose \(r>1\) for
\(B=e^{t_0A/2}\).  For every \(t\ge t_0\),
\[
 e^{tA}=B e^{(t-t_0)A}B.
\]
The middle kernel has radius one, including at \(t=t_0\), and the
outer kernels have radius \(r\).  The two kernel contractions pair
radii whose products are \(r>1\), retaining radius \(r\) on both
external legs.  A further contraction with any stable input proves the
uniform input assertion.
\end{proof}

\begin{remark}\label{spec:strict-scope}
The radius can depend on the capacity, generator, and initial positive
time.  The theorem gives neither a dimension-independent lower bound for
\(r-1\) nor an improvement on inert ancillary variables.  Complex
scalar shifts of the generator do not affect its conclusions.  No
normalization by a trace is used.
\end{remark}

\subsection{Dominant eigenvectors and their dual certificates}

We use the following finite-dimensional complex-cone theorem of Rugh
\cite[Theorem~8.4]{Rugh10}: if a nontrivial closed complex-scalar-invariant cone
contains no complex plane and a linear map sends its nonzero part into
its interior, then the map has an algebraically simple nonzero
eigenvalue of strictly largest modulus.  Convexity of the cone and
self-adjointness of the map are not hypotheses.

\begin{proposition}\label{spec:perron}
Under the hypotheses of Theorem~\ref{spec:strict}, \(A\) has an
algebraically simple eigenvalue \(a_\star\) uniquely largest in real
part.  Choose its right eigenvector \(v\) and left eigenfunctional
\(\ell\) so that \(\ell(v)=1\).  Both \(v\) and
\begin{equation}\label{spec:dual-polynomial}
 Q_\ell(w)=\sum_{\alpha\le\kappa}b_\alpha\ell(z^\alpha)w^\alpha
\end{equation}
are zero-free on a common polydisk of radius greater than one.
Furthermore \(\ell(p)\ne0\) for every nonzero stable \(p\), and
\(e^{-ta_\star}e^{tA}p\to\ell(p)v\).
\end{proposition}
\begin{proof}
For each fixed \(t>0\), Theorem~\ref{spec:strict} and Lemma
\ref{spec:cone-interior} verify all the hypotheses of Rugh's theorem
for \(T=e^{tA}\); the cone contains all constant polynomials
and is therefore nontrivial. Its simple dominant eigenspace is invariant under
\(A\), because \(A\) commutes with \(T\).  Thus it is an
\(A\)-eigenline, say for \(a_\star\).  Spectral mapping gives
strict real-part dominance: another generator eigenvalue with equal
real part would give the same modulus for the exponential.  If two
distinct generator eigenvalues exponentiated to the same value, that
value would have algebraic multiplicity at least two.  A nontrivial
Jordan block at \(a_\star\) would also give a nontrivial exponential
Jordan block, since \(t e^{ta_\star}\ne0\).  Both alternatives
contradict algebraic simplicity for \(T\).

Fix a kernel radius \(r>1\) for \(T\), and put
\(\Lambda=e^{ta_\star}\).  Finite-dimensional spectral theory gives
\[
 \Lambda^{-m}T^m\longrightarrow P=v\otimes\ell,
 \qquad \ell(v)=1.
\]
Every power has kernel radius \(r\), because its internal contraction
radii have product \(r^2>1\).  The kernel transform is injective, so
\(K_P\) is not the zero polynomial.  Hurwitz yields
\[
 K_P(z,w)=v(z)Q_\ell(w)\ne0\qquad(z,w\in\D_r^d).
\]
Both factors consequently have radius \(r\).  The strict weighted
contraction of \(Q_\ell\) with a nonzero unit-stable polynomial is
\(\ell(p)\), so it is nonzero.  The strict real-part separation and
the finite-dimensional spectral expansion then give the asserted
convergence, including when the remaining spectrum has Jordan blocks.
\end{proof}

\subsection{Sharp spectral bounds}

We first record the scalar estimate used in both the spectral and
dynamical arguments.  It is the diameter form of the Schwarz lemma
used in \cite[Theorem~2]{BGLW26}.

\begin{lemma}\label{spec:schwarz}
If \(g\) is holomorphic on \(\D_1^d\), continuous on its closure,
and \(0\le q\le1\), then
\[
 \operatorname{diam}g(q\overline{\D_1^d})
 \le q\operatorname{diam}g(\overline{\D_1^d}).
\]
The same conclusion holds for a coordinatewise dilation all of whose
moduli are at most \(q\).
\end{lemma}
\begin{proof}
The cases \(q=1\), zero diameter, and \(q=0\) are immediate.
For \(x,y\in q\overline{\D_1^d}\), put
\(\rho=\max(\|x\|_\infty,\|y\|_\infty)\).  If \(\rho>0\),
\[
 f(\zeta)=
 \frac{g(\zeta x/\rho)-g(\zeta y/\rho)}
 {\operatorname{diam}g(\overline{\D_1^d})}
\]
is holomorphic on the unit disk, bounded in modulus by one, and
vanishes at zero.  Schwarz at \(\zeta=\rho<1\) bounds
\(|g(x)-g(y)|\) by \(\rho\) times the denominator.  Take the
supremum over \(x,y\).  A coordinatewise dilation has its image in
\(q\overline{\D_1^d}\).
\end{proof}

\begin{theorem}[Complex radius bound]\label{spec:radius-gap}
Suppose \(r>1\) and \(K_T\) is zero-free on
\(\D_r^d\times\D_r^d\).  Then \(T\) has an algebraically simple
nonzero eigenvalue \(\Lambda_\star\) of strictly largest modulus,
with an eigenvector of radius \(r\).  Every other eigenvalue satisfies
\begin{equation}\label{spec:static-ratio}
 |\Lambda/\Lambda_\star|\le r^{-2}.
\end{equation}
No invertibility, Hermiticity, or diagonalizability is assumed.
\end{theorem}
\begin{proof}
The kernel maps the nonzero stable cone into its interior by Lemma
\ref{spec:kernel-criterion}.  Rugh's theorem therefore supplies the
simple dominant eigenvalue.  Its rank-one spectral projection is
nonzero on some interior cone vector, since a proper linear hyperplane
cannot contain an open set.  Normalized iterates of that vector converge
to a nonzero eigenvector \(\psi\) in the closed cone.  Applying \(T\)
once improves \(\psi\) to radius \(r\).

Let \(\phi\) be an eigenvector of a different nonzero eigenvalue,
and put \(g=\phi/\psi\), \(\alpha=\Lambda/\Lambda_\star\).
The function \(g\) is nonconstant and holomorphic on \(\D_r^d\).
For \(1/r<s<r\), we claim
\[
 \alpha g(\D_r^d)\subseteq g(\D_s^d).
\]
Otherwise, for some \(v\in\D_r^d\), the nonzero polynomial
\(f=\phi-\alpha g(v)\psi\) would be stable on \(\D_s^d\), but
\(Tf=\Lambda(\phi-g(v)\psi)\) would vanish at \(v\).
This contradicts strict contraction, since \(rs>1\).

For \(1/r<s<t<r\), let
\(\Delta_u=\operatorname{diam}g(\overline{\D_u^d})\).
These diameters are finite and positive.  The inclusion gives
\(|\alpha|\Delta_t\le\Delta_s\).  Applying Lemma
\ref{spec:schwarz} after rescaling radius \(t\) gives
\(\Delta_s\le(s/t)\Delta_t\).  Letting \(s\downarrow1/r\)
and \(t\uparrow r\) proves \eqref{spec:static-ratio}.  Zero
eigenvalues satisfy the inequality trivially.
\end{proof}

The quantitative mechanism in this proof is the ratio-squeezing argument
of \cite{BGLW26}; the complex-cone theorem supplies the eigenvector
without a positive-semidefinite hypothesis.

\begin{theorem}[Sharp degree-damping gap]\label{spec:gap}
Under the hypotheses of Theorem~\ref{spec:strict}, set
\(\mu=\min_i\mu_i\).  The dominant generator eigenvalue satisfies
\begin{equation}\label{spec:real-gap}
 \Re a_\star-\Re a\ge2\mu
 \qquad(a\in\operatorname{spec}(A),\ a\ne a_\star).
\end{equation}
The constant is sharp for every positive finite capacity vector.
\end{theorem}
\begin{proof}
For small \(h>0\), put
\[
 D_h=e^{-h\sum_i\mu_iE_i},\qquad
 T_h=D_he^{hA_0}D_h.
\]
The kernel of \(e^{hA_0}\) has radius one by Lemma
\ref{spec:kernel-criterion}.  If \(q_i=e^{-h\mu_i}\), direct
coefficient calculation gives
\[
 K_{T_h}(z,w)=K_{e^{hA_0}}(q_1z_1,\ldots,q_dz_d,
                              q_1w_1,\ldots,q_dw_d).
\]
It therefore has radius \(e^{h\mu}\).  Theorem
\ref{spec:radius-gap} gives a simple dominant eigenvalue
\(\Lambda_\star(h)\) and bounds every other eigenvalue by
\(e^{-2h\mu}|\Lambda_\star(h)|\).

Now \(B_h=(T_h-I)/h\to A\) in norm.  Its eigenvalues are exactly
\(b(h)=(\Lambda(h)-1)/h\), with algebraic multiplicities.  These
multisets are bounded and converge to the spectrum of \(A\).  Uniformly
for bounded complex \(b\),
\[
 h^{-1}\log|1+hb|=\Re b+O(h).
\]
The modulus inequality consequently implies
\(\Re b_\star(h)-\Re b(h)\ge2\mu+O(h)\).
Choose a sequence along which \(b_\star(h)\) converges.  After removing
this one eigenvalue from each multiset, the remaining eigenvalues
converge, with multiplicity, to the remaining copies of the spectrum
of \(A\).  The limiting inequality separates all those copies by
\(2\mu\).  In particular, the distinguished limit is algebraically
simple and is the unique maximal-real-part eigenvalue, hence is
\(a_\star\).  This proves \eqref{spec:real-gap} without
differentiable eigenvalue labels or a normal-matrix perturbation bound.
For \(A_0=0\), the eigenvalues are \(-2\sum_i\mu_i\alpha_i\);
a degree-one monomial at a site of minimum rate attains equality.
\end{proof}

\subsection{Contraction along the actual evolution}

For strictly stable \(p\) and arbitrary \(h\in V_\kappa\), define
\begin{equation}\label{spec:osc-definition}
 \operatorname{Osc}_p(h)
 =\operatorname{diam}\{h(z)/p(z):z\in\overline{\D_1^d}\}.
\end{equation}
It is a complex seminorm with kernel \(\C p\): homogeneity and the
triangle inequality follow from the diameter definition, and zero
diameter means that the rational function is constant.

\begin{lemma}\label{spec:range-exclusion}
If \(T\) is an invertible stability preserver, \(p\) is strictly
stable, and \(h\in V_\kappa\), then \(Tp\) is strictly stable and
\[
 (Th/Tp)(\D_1^d)\subseteq(h/p)(\D_1^d).
\]
In particular \(\operatorname{Osc}_{Tp}(Th)\le
\operatorname{Osc}_p(h)\).
\end{lemma}
\begin{proof}
Strictness follows from Lemma~\ref{spec:cone-interior}.  If \(w\)
does not belong to the input ratio image, then \(h-wp\) is a nonzero
stable polynomial.  Its image \(Th-wTp\) is nonzero and stable,
excluding \(w\) from the output ratio image.  Taking diameters proves
the assertion; the open and closed polydisks give the same diameter by
continuity and density.
\end{proof}

\begin{theorem}[Moving-denominator contraction]\label{spec:oscillation}
Let \(A_0\) generate a stability-preserving semigroup, and let
\(A=A_0-2\sum_i\mu_iE_i\), with \(\mu_i\ge\mu\ge0\).
For every strictly stable \(p\), every \(h\in V_\kappa\), and
\(t\ge0\),
\begin{equation}\label{spec:osc-bound}
 \operatorname{Osc}_{e^{tA}p}(e^{tA}h)
 \le e^{-2\mu t}\operatorname{Osc}_p(h).
\end{equation}
The constant is sharp.
\end{theorem}
\begin{proof}
The damping operator \(D_\tau=e^{-2\tau\sum_i\mu_iE_i}\) acts
by a coordinatewise dilation of maximum modulus \(e^{-2\mu\tau}\).
Lemma~\ref{spec:schwarz}, applied to \(h/p\), gives the corresponding
oscillation contraction.  Each \(e^{\tau A_0}\) is nonexpansive by
Lemma~\ref{spec:range-exclusion}.  Iterating along
\((e^{tA_0/m}D_{t/m})^m\) therefore gives exactly the factor
\(e^{-2\mu t}\), with all intermediate denominators strictly stable.

The products converge to \(e^{tA}\), which is invertible and preserves
stability by the product formula and Hurwitz.  Its image of \(p\)
is strictly stable.  The minimum modulus of that limiting polynomial
on the closed polydisk is positive, so coefficient convergence gives
uniform convergence of the ratios there.  Diameters are continuous
under uniform convergence, proving \eqref{spec:osc-bound}.  Equality
holds for pure damping with \(p=1\) and \(h=z_j\) at a minimum-rate
site.
\end{proof}

\begin{corollary}[Nonautonomous evolution]\label{spec:nonautonomous}
On each compact time interval, suppose \(A_0(t)\) and \(\mu_i(t)\)
are bounded and piecewise continuous, with finitely many continuity
pieces, each \(A_0(t)\) a stability generator and \(\mu_i(t)\ge0\).
Let \(U(t,s)\) solve the linear evolution equation with generator
\(A_0(t)-2\sum_i\mu_i(t)E_i\).  Then, for strictly stable \(p\),
\[
 \operatorname{Osc}_{U(t,s)p}(U(t,s)h)
 \le \exp\!\left(-2\int_s^t\min_i\mu_i(u)\,du\right)
       \operatorname{Osc}_p(h).
\]
\end{corollary}
\begin{proof}
Approximate on each continuity piece by step functions taking actual
values of the coefficients.  They converge in \(L^1\).  Their frozen
evolutions satisfy Theorem~\ref{spec:oscillation}, including when the
minimum rate is zero.  Multiplying the bounds gives the exponential
of the integral of the stepwise minimum rates.  These integrals converge
to the stated integral, since the minimum of finitely many real
coordinates is Lipschitz.  The integral-equation estimate for finite
matrix ODEs gives uniform convergence of the propagators.  The limiting
fundamental matrix is invertible and preserves stability by Hurwitz;
Lemma~\ref{spec:cone-interior} gives a strictly stable limiting
denominator.  Uniform convergence of the ratios passes the estimate
to the limit, as in the autonomous proof.
\end{proof}

\begin{corollary}[Adapted norm and resolvent]\label{spec:resolvent}
Assume positive damping, and normalize \(v,\ell\) as in
Proposition~\ref{spec:perron}.  On \(W=\ker\ell\), set
\(\|h\|_v=\operatorname{Osc}_v(h)\) and
\(B=(A-a_\star I)|_W\).  Then \(\|\cdot\|_v\) is a norm and
\[
 \|e^{tB}\|_{v\to v}\le e^{-2\mu t},\qquad
 \|(\zeta I-B)^{-1}\|_{v\to v}
 \le\frac1{\Re\zeta+2\mu}\quad(\Re\zeta>-2\mu).
\]
Every eigenvalue of \(B\) with real part \(-2\mu\) is semisimple.
\end{corollary}
\begin{proof}
The only possible null vectors of the seminorm are multiples of \(v\),
and \(\ell(v)=1\).  Thus it is a norm on \(W\).  The left
eigenfunctional identity makes \(W\) invariant.  Since
\(e^{tA}v=e^{ta_\star}v\), Theorem~\ref{spec:oscillation} gives
the first estimate, including the complex phase of the scalar.
For \(\Re\zeta>-2\mu\), the absolutely convergent integral
\(\int_0^\infty e^{-t\zeta}e^{tB}\,dt\) is the inverse of
\(\zeta I-B\), by integration of its derivative.  Integrating the
norm estimate gives the stated bound.  A nontrivial Jordan block at a
boundary eigenvalue would produce unbounded polynomial growth after
multiplication by \(e^{2\mu t}\), contradicting the first estimate
and norm equivalence on its finite-dimensional generalized eigenspace.
\end{proof}

\begin{example}[No uniform strict radius]\label{spec:no-uniform-radius}
Fix \(\mu,t>0\).  On \(V_1\), in the coefficient basis \((1,z)\),
let \(X\) and \(Z\) be the usual Pauli matrices and put
\[
 A_{0,R}=RX,\qquad
 A_R=A_{0,R}-2\mu E=
 \begin{pmatrix}0&R\\R&-2\mu\end{pmatrix},\qquad R>0.
\]
The undamped flow preserves disk stability, also with arbitrary inert
variables.  Indeed, writing \(c=\cosh(Rs)\), \(b=\sinh(Rs)\),
its action is the explicit weighted substitution
\[
 (e^{sA_{0,R}}p)(z)=(c+bz)
 p\!\left(\frac{b+cz}{c+bz}\right).
\]
The prefactor is nonzero on the disk and the fraction is a disk
automorphism, since \(c^2-b^2=1\).
Put \(\omega_R=\sqrt{R^2+\mu^2}\).  The identity
\((A_R+\mu I)^2=\omega_R^2I\) gives
\[
 e^{tA_R}=e^{-\mu t}\left[
 \cosh(\omega_Rt)I+
 \frac{\sinh(\omega_Rt)}{\omega_R}(A_R+\mu I)\right].
\]
Thus the image of the fixed input \(1\) has its unique zero at
\[
 z_R=-\frac{\omega_R\coth(\omega_Rt)+\mu}{R}.
\]
For every finite \(R\), \(|z_R|>1\), but \(|z_R|\to1\) as
\(R\to\infty\).  Since \(K_{e^{tA_R}}(z,0)=e^{tA_R}1(z)\),
the same zero obstructs a common kernel radius depending only on
\(\mu,t\).  This failure already occurs for Hermitian generators
on one binary variable.  The dominant eigenvector has its zero at
modulus \((\omega_R+\mu)/R\), which also tends to one.
\end{example}

\begin{example}[Unbounded mixing remainder at a fixed gap]\label{spec:nonnormal}
Fix \(\mu>0\), and on \(V_1\) put
\[
 G_R=R(X+iZ),\qquad
 B_R=G_R-2\mu E=
 \begin{pmatrix}iR&R\\R&-iR-2\mu\end{pmatrix},\qquad R>0.
\]
Here \(G_R^2=0\) and \(G_R^*Z+ZG_R=0\).  Its stability-preserving
action is explicit: with \(a=1+iRs\) and \(b=Rs\),
\[
 (e^{sG_R}p)(z)=(a+bz)
 p\!\left(\frac{b+\overline a z}{a+bz}\right),\qquad
 |a+bz|^2-|b+\overline a z|^2=1-|z|^2.
\]
The denominator is nonzero in the disk, and the fraction maps the
disk to itself.  The formula also proves preservation with inert
variables.

Choose the square root \(\omega_R=\sqrt{\mu^2+2i\mu R}\) with
\(x_R=\Re\omega_R>0\).  The two eigenvalues of \(B_R\) are
\(a_\pm=-\mu\pm\omega_R\), and
\[
 x_R=\sqrt{\frac{\mu\sqrt{\mu^2+4R^2}+\mu^2}{2}}\ge\mu.
\]
The dominant spectral projection and normalized exponential are
\[
 P_R=\frac12\left(I+\frac{B_R+\mu I}{\omega_R}\right),\qquad
 S_R(t)=e^{-ta_+}e^{tB_R}
       =P_R+e^{-2\omega_Rt}(I-P_R).
\]
Their off-diagonal entries give
\[
 \|P_R\|_2\ge\frac{R}{2|\omega_R|},\qquad
 |[S_R(t)]_{10}|
 =\frac{R}{2|\omega_R|}|1-e^{-2\omega_Rt}|
 \ge\frac{R}{2|\omega_R|}(1-e^{-2x_Rt}).
\]
Since \(|\omega_R|=(\mu^4+4\mu^2R^2)^{1/4}\), both quantities
diverge as \(R\to\infty\), the second for every fixed \(t>0\).

To keep the actual spectral gap fixed, pass to \(V_{(1,1)}\) and set
\[
 \mathcal A_R=G_R\otimes I-2\mu(E_1+E_2)
             =B_R\otimes I+I\otimes(-2\mu E).
\]
The undamped part is still a joint stability generator by the displayed
weighted substitution.  Its four eigenvalues are
\(a_+,a_-,a_+-2\mu,a_--2\mu\).  Hence its dominant eigenvalue
is \(a_+\), its real-part gap is exactly \(2\mu\), and its
dominant projector is \(\mathcal P_R=P_R\otimes P_0\).
For the fixed strictly stable input \(p(z_1,z_2)=1+z_2/2\), define
\(\mathcal R_R(t)=e^{-ta_+}e^{t\mathcal A_R}-\mathcal P_R\).
The coefficient of \(z_1z_2\) is exactly
\[
 [z_1z_2]\mathcal R_R(t)p
 =\frac12e^{-2\mu t}\frac{R}{2\omega_R}
   (1-e^{-2\omega_Rt}).
\]
The projector term contains no \(z_2\), which proves this identity
directly from the tensor exponential.  Its modulus tends to infinity
for every fixed \(\mu,t>0\).  Thus even after subtraction of the
dominant projector there is no coefficient-norm mixing bound depending
only on the damping, time, and exact spectral gap, in a fixed two-variable
box and on one fixed strictly stable input.  Equivalence of fixed
finite-dimensional norms gives the same obstruction for every fixed
coefficient norm.  This is consistent with Corollary~\ref{spec:resolvent},
whose norm depends on the dominant polynomial.  The example uses linear
normalization by \(e^{ta_+}\); it makes no claim about nonlinear
normalization by an evolving coefficient.
\end{example}

\begin{remark}\label{spec:projective-attribution}
For the stable cone, the forbidden-pencil projective distance of
\cite{Dubois09} is
\(\operatorname{osc}_{\D_1^d}\log|q/p|\).  Its real infinitesimal
form is \(\operatorname{osc}\Re(h/p)\); maximizing over a scalar
phase of \(h\) gives \eqref{spec:osc-definition}.  Thus the metric
and the general complex-cone contraction principle are inherited.
The estimates above specify the exact degree-damping rate for arbitrary
tangent vectors along the actual, possibly nonautonomous evolution.
\end{remark}

% END SECTION spectral

% BEGIN SECTION sector_spectra
\subsection{Sector growth and token-graph spectral concavity}
\label{sector:section}

We now return to real upper-half-plane stability.
There is a second spectral consequence when coefficient positivity and
total degree are preserved. It does not require degree damping or
self-adjointness. A single stable evolution compares the exponential
growth rates of all homogeneous sectors.

\begin{theorem}[Concavity of sector growth]\label{sector:growth}
Let \(V_\kappa\) be a finite real coordinate box and
\(N=\sum_i\kappa_i\). Suppose that \(A\) has nonnegative off-diagonal
entries in the monomial basis, preserves every homogeneous subspace
\(V_{\kappa,k}\), and that \(e^{tA}\) preserves nonnegative real
stability for every \(t\ge0\). Write
\[
 s_k=\max\{\Re\lambda:\lambda\in
             \operatorname{spec}(A|_{V_{\kappa,k}})\},
 \qquad 0\le k\le N.
\]
Then
\[
 2s_k\ge s_{k-1}+s_{k+1}\qquad(1\le k<N).
\]
Neither irreducibility nor symmetry of \(A\) is assumed.
\end{theorem}

\begin{proof}
Start at the stable polynomial
\(f_0(z)=\prod_i(1+z_i)^{\kappa_i}\).
Every allowed monomial has a strictly positive coefficient.
Choose \(\alpha\) such that \(A+\alpha I\) is entrywise nonnegative.
Then \(e^{tA}=e^{-\alpha t}e^{t(A+\alpha I)}\) is entrywise
nonnegative and has a strictly positive diagonal. Consequently
\(f_t=e^{tA}f_0\) still has every allowed coefficient strictly positive.
Its diagonal restriction is a degree-\(N\) polynomial
\[
 P_t(x)=f_t(x,\ldots,x)=\sum_{k=0}^N c_k(t)x^k,
 \qquad c_k(t)>0,
\]
whose roots are all strictly negative.

Newton's inequalities give
\begin{equation}\label{sector:newton}
 \left(\frac{c_k(t)}{\binom Nk}\right)^2
 \ge
 \frac{c_{k-1}(t)}{\binom N{k-1}}
 \frac{c_{k+1}(t)}{\binom N{k+1}}.
\end{equation}
Here is the particular form needed. The derivative
\(R=P_t^{(k-1)}\) has \(m=N-k+1\) negative roots, so
\(R(x)=R(0)\prod_{j=1}^m(1+u_jx)\) with \(u_j>0\).
Cauchy's inequality gives
\[
 \sum_{i<j}u_iu_j
 \le\frac{m-1}{2m}\left(\sum_i u_i\right)^2.
\]
The first three coefficients of \(R\) are
\((k-1)!c_{k-1}\), \(k!c_k\), and
\((k+1)!c_{k+1}/2\). Substitution gives
\eqref{sector:newton}, including its binomial normalization.

Let \(A_k=A|_{V_{\kappa,k}}\), and let \(v_k\) be the coefficient
vector of the degree-\(k\) part of \(f_0\). Degree preservation gives
\[
 c_k(t)=\mathbf1^{\mathsf T}e^{tA_k}v_k.
\]
Since every entry of \(v_k\) is positive, this quantity is bounded
above and below by fixed positive multiples of the sum of all
entries of \(e^{tA_k}\). For an entrywise nonnegative \(m_k\)-by-\(m_k\)
matrix \(M\), that sum lies between \(\|M\|_1\) and
\(m_k\|M\|_1\), where \(\|\cdot\|_1\) is the induced maximum
column-sum norm. Spectral mapping gives
\(\|e^{tA_k}\|_1\ge e^{ts_k}\), and Jordan normal form gives
\[
 \|e^{tA_k}\|_1
 \le C_k(1+t^{m_k-1})e^{ts_k}\qquad(t\ge0)
\]
for a fixed constant \(C_k\). Therefore
\[
 \lim_{t\to\infty}t^{-1}\log c_k(t)=s_k.
\]
Take logarithms in \eqref{sector:newton}, divide by \(t\), and
let \(t\to\infty\). The fixed binomial factors disappear, proving
the assertion. For \(N\le1\) there is no inequality to check.
\end{proof}

The proof uses coefficient sums, not
\(\operatorname{Tr}(e^{tA_k})\). No real-rooted generating function
for the latter traces is asserted. The passage from coefficient
log-concavity to concavity of growth exponents is the one-variable
tropical-stability principle
\cite[Theorem~5 and the discussion following it]{Branden10Discrete}.
The proof above establishes the needed exponential version directly,
including nonsymmetric blocks and their Jordan factors.
The preservation theorem supplies the input needed to apply it.

Let \(G\) be a finite simple graph on \([n]\), with nonnegative edge
weights \(w_{ij}\). Its \(k\)-token graph has the \(k\)-subsets as
vertices. The edge from \(S\) to \(S-\{i\}+\{j\}\) has weight
\(w_{ij}\) whenever \(i\in S\), \(j\notin S\), and \(ij\) is an
edge of \(G\). Denote its adjacency and weighted degree matrices by
\(A_k(G)\) and \(D_k(G)\). The sectors \(k=0,n\) are singletons
with zero matrices.

\begin{theorem}[Token-graph spectral concavity]\label{sector:token}
For every \(a\in[-1,1]\), put
\[
 \Lambda_k(a)=\lambda_{\max}(A_k(G)+aD_k(G)).
\]
For every nonnegatively weighted \(G\),
\[
 2\Lambda_k(a)\ge\Lambda_{k-1}(a)+\Lambda_{k+1}(a)
 \qquad(1\le k<n),\qquad
 \Lambda_k(a)=\Lambda_{n-k}(a).
\]
In particular the sequence is nondecreasing up to half filling.
The interval \([-1,1]\) is the exact interval of real parameters
for which this concavity holds for every such graph.
\end{theorem}

\begin{proof}
On the binary polynomial box define, for every unordered edge,
\[
 \mathcal B_{ij}=w_{ij}\bigl[
 (z_j+az_i)\partial_i+(z_i+az_j)\partial_j
 -(z_i^2+2az_iz_j+z_j^2)\partial_i\partial_j\bigr].
\]
It annihilates \(1,z_iz_j\) and sends
\(z_i\) to \(w_{ij}(az_i+z_j)\) and
\(z_j\) to \(w_{ij}(z_i+az_j)\).
Thus \(\mathcal B=\sum_{ij}\mathcal B_{ij}\) preserves total
degree, is a real symmetric Metzler matrix, and restricts on the
degree-\(k\) sector to \(A_k(G)+aD_k(G)\).
Its mixed second-order coefficient is
\[
 -w_{ij}\bigl[(s+at)^2+(1-a^2)t^2\bigr]\le0.
\]
Theorem~\ref{gen:main-real} therefore gives stability preservation.
Theorem~\ref{sector:growth} yields concavity, since for a real
symmetric matrix its spectral bound is its largest eigenvalue.

Complementation of subsets gives a weight-preserving isomorphism
between the \(k\)- and \((n-k)\)-token graphs, proving symmetry.
Set \(\delta_k=\Lambda_k-\Lambda_{k-1}\). Concavity says that
the \(\delta_k\) are nonincreasing, while symmetry gives
\(\delta_{n-k+1}=-\delta_k\). If \(k\le(n+1)/2\), it follows that
\(\delta_k\ge-\delta_k\), proving monotonicity.

For sharpness, take the unweighted star \(K_{1,3}\). Its two-token
graph is a six-cycle. On the equal-leaf subspace, the one-token
matrix is
\(\left(\begin{smallmatrix}3a&\sqrt3\\\sqrt3&a\end{smallmatrix}\right)\);
its orthogonal complement has eigenvalue \(a\). Thus
\[
 \Lambda_1(a)=2a+\sqrt{a^2+3},\qquad
 \Lambda_2(a)=2a+2.
\]
For \(a>1\), \(\Lambda_1(a)>\Lambda_2(a)\); symmetry then
contradicts concavity at \(k=2\). For \(a<-1\), a single edge
has \(\Lambda_0=\Lambda_2=0\) and \(\Lambda_1=a+1<0\), again
violating concavity. At \(a=1\), the star has
\(\Lambda_1=\Lambda_2=4\), so connectedness does not imply
strict monotonicity.
\end{proof}

For completeness, the local preservation used here has a short
verification independent of the generator classification. For
\(\tau\ge0\), set \(b=e^{a\tau}\sinh\tau\) and
\(c=e^{a\tau}\cosh\tau\). The unweighted edge exponential fixes
\(1,xy\) and sends \(x,y\) to \(cx+by,bx+cy\).
Its algebraic symbol is
\[
 xy+uv+b(xu+yv)+c(xv+yu).
\]
Writing this symbol as \(\zeta^{\mathsf T}M\zeta/2\) in the
order \(\zeta=(x,y,u,v)^{\mathsf T}\), one has
\[
 M=\begin{pmatrix}
 0&1&b&c\\
 1&0&c&b\\
 b&c&0&1\\
 c&b&1&0
 \end{pmatrix}.
\]
Its eigenvalues are
\[
 1+e^{(a+1)\tau},\quad 1-e^{(a+1)\tau},\quad
 -1-e^{(a-1)\tau},\quad -1+e^{(a-1)\tau}.
\]
Exactly one is positive when \(-1\le a\le1\). The quadratic
form is positive on every strictly positive real vector \(Y\).
A zero at \(X+\mathrm iY\) would imply
\(X^{\mathsf T}MY=0\) and
\(X^{\mathsf T}MX=Y^{\mathsf T}MY>0\), contradicting the
one-positive-eigenvalue property. Hence the symbol is stable.
The classical symbol theorem \cite{BB09} gives complete
preservation for each edge. The finite-dimensional product formula
and coefficientwise closure give preservation for their sum;
invertibility of the limiting exponential excludes zero outputs.

Taking \(a=1\) and \(a=0\) gives, respectively, monotonicity of
the largest eigenvalues of the signless Laplacian and the adjacency
matrix, as conjectured in \cite[Conjectures~5--6]{APS26}.
Theorem~\ref{sector:token} proves the stronger weighted concavity
throughout the displayed interpolation interval.
The adjacency-preserving generator at \(a=0\) is already contained
in Purbhoo's construction \cite[Proposition~4.5]{Purbhoo18}.
The contribution here is the cross-sector spectral conclusion.
The bounds in Conjectures~1--4 of \cite{APS26} were settled by
\cite{BBKL26} and are not claimed here.

In the occupation basis, the same block matrices arise from
\[
 \mathsf B_a(G)=\frac12\sum_{ij}w_{ij}
 \bigl[X_iX_j+Y_iY_j+a(I-Z_iZ_j)\bigr].
\]
Thus its maximum energy over fixed particle number is concave in
that number. The negative Hamiltonian has the corresponding convex
sector ground-state energies. At \(a=1\), the common one-qubit
rotation \(R=e^{-\mathrm i\pi X/4}\), with
\(RYR^*=Z\) and \(RZR^*=-Y\), identifies this operator with the
EPR Hamiltonian in \cite{APS26}. There is a
maximum-energy state in a half-filled sector in the rotated basis.
This locates a maximizing sector; it does not give an efficient
algorithm within that exponentially large sector. For \(a<0\)
the quantity in Theorem~\ref{sector:token} is the largest
eigenvalue, not in general the spectral radius.

% END SECTION sector_spectra

% BEGIN SECTION algorithms
\section{Classical computation of principal eigenvalues and states}
\label{alg:section}

The spectral theorem of the preceding section has an algorithmic consequence
for operators with a bounded-degree local description.  A positive damping
rate selects a holomorphic eigenvalue throughout a complex half-plane.
This domain can be joined to a disk on which perturbation coefficients are
efficiently computable.  An explicit change of variable then reduces
evaluation at any fixed positive damping rate to logarithmically many
coefficients.

The coefficient algorithm is due to Bravyi, DiVincenzo, and Loss
\cite{BDL08}.  Our use of it includes a complexification argument, so that
every invocation of their algorithm remains within its Hermitian
hypotheses.  Neither an extension of Hermitian eigenvalue ordering to
nonnormal matrices nor numerical continuation of eigenvectors is needed.
We first estimate the principal eigenvalue.  Starting in
Section~\ref{alg:states-subsection}, we then compute stable product
tests of its eigenvectors, including relative ground-state amplitudes,
product-event probabilities, and adaptive equatorial measurement
records.  These scalar queries do not require a bound on the condition
number of the spectral projection.

\subsection{Local input and the approximation theorem}
\label{alg:input-subsection}

Identify the coefficient space
\[
 \mathcal P_n
 =\{p\in\C[z_1,\ldots,z_n]:\deg_{z_i}p\le1\}
 \cong(\C^2)^{\otimes n},
 \qquad z^\alpha\longleftrightarrow|\alpha\rangle.
\]
We use the Euclidean norm in this coefficient basis and the induced
operator norm.  Let \(\mathcal S_n\) be the set of nonzero polynomials in
\(\mathcal P_n\) that do not vanish in \(\D^n\), where
\(\D=\{z\in\C:|z|<1\}\).  The degree operator is
\begin{equation}
 E=\sum_{i=1}^n |1\rangle\langle1|_i
   =\sum_{i=1}^n z_i\partial_{z_i}.
 \label{alg:degree}
\end{equation}
All local matrices below are supplied in the displayed tensor-product
basis.  A Gaussian rational means an element of \(\mathbb Q(i)\).

\begin{theorem}[Principal-eigenvalue approximation]
\label{alg:complex}
Fix \(d\in\mathbb N\cup\{0\}\) and positive rational numbers \(J,\mu\).
For \(n\ge1\), suppose that the input specifies
\begin{equation}
 A_0=\sum_{i=1}^n K_i+
       \sum_{\{i,j\}\in\mathcal E}K_{ij},
 \label{alg:local-input}
\end{equation}
where the graph \((\{1,\ldots,n\},\mathcal E)\) has maximum degree at
most \(d\), every local matrix has Gaussian-rational entries and operator
norm at most \(J\), and
\begin{equation}
 e^{tA_0}\mathcal S_n\subseteq\mathcal S_n
 \qquad(t\ge0).
 \label{alg:promise}
\end{equation}
Let \(a_\star\) be the eigenvalue of
\(A=A_0-2\mu E\) with largest real part.  This eigenvalue is unique
and algebraically simple.  Given a positive rational \(\delta\), a
deterministic classical algorithm returns
\(\widehat a\in\mathbb Q(i)\) such that
\[
 |\widehat a-a_\star|\le\delta.
\]
For fixed \(d,J,\mu\), its bit complexity is polynomial in \(n\),
the input bit length, and \(\delta^{-1}\).
\end{theorem}

Only the sum \(A_0\) is assumed to satisfy
\eqref{alg:promise}; the displayed local terms need not individually
generate preserving semigroups.  The theorem uses the supplied local
representation and does not include an algorithm for finding such a
representation or verifying the preservation promise.
No positivity, Hermiticity, stochastic normalization, or diagonalizability
is assumed.

We will give explicit choices of all continuation parameters.  In
particular, put
\begin{equation}
 N=2n,\qquad D=d+1,\qquad
 \rho=\frac{2^{-18}}{DJ},\qquad c=\frac1\mu.
 \label{alg:parameters}
\end{equation}
The number \(p\) of perturbation coefficients required by the proof
satisfies
\begin{equation}
 p=O\!\left(
   \left(1+\frac{c}{\rho}\right)^2
   \log\!\left(2+\frac{C_{d,J,\mu}n}{\delta}\right)
 \right),
 \label{alg:order-bound}
\end{equation}
for a constant \(C_{d,J,\mu}\) independent of \(n,\delta\).
The arithmetic operation count is
\(N\exp(O_D(p))\), with a polynomial bit overhead.
The large constants in \eqref{alg:parameters} are inherited from
the perturbation estimate.  They are sufficient for a complexity
statement, not a claim of practical efficiency.  In particular,
Theorem~\ref{alg:complex} does not assert polynomial dependence
on \(\mu^{-1}\).

\subsection{Perturbation coefficients and complexification}
\label{alg:perturbation-subsection}

We first state precisely the published perturbation input.
For a graph on \(N\) qubits with maximum degree at most \(D\ge1\), set
\begin{equation}
 H_0=2\sum_{u=1}^N |1\rangle\langle1|_u,
 \qquad \Omega=|0\rangle^{\otimes N}.
 \label{alg:unperturbed}
\end{equation}
For any matrix \(V\), the eigenvalue of \(H_0+xV\) issuing from the
simple eigenvalue \(0\) has a convergent germ at \(x=0\), which we write
\begin{equation}
 F_V(x)=\sum_{k\ge1}e_k(V)x^k.
 \label{alg:energy-germ}
\end{equation}

\begin{proposition}[Bravyi--DiVincenzo--Loss]
\label{alg:bdl}
Let \(W=\sum_{\{u,v\}}W_{uv}\) be Hermitian, with each
\(W_{uv}\) Hermitian and of norm at most \(L>0\).
For \(H_0\) as in \eqref{alg:unperturbed}, put
\[
 R_L=\frac{2^{-16}}{DL}.
 \]
Then
\begin{equation}
 |e_k(W)|\le2^{-15}N R_L^{-k}\qquad(k\ge1).
 \label{alg:bdl-bound}
\end{equation}
For every \(p\ge1\), the coefficients
\(e_1(W),\ldots,e_p(W)\) can be computed in
\(N\exp(O_D(p))\) arithmetic operations.
\end{proposition}

These are Lemma~6 and Theorem~3 of \cite{BDL08}, specialized to
unperturbed gap \(2\).  Their Theorem~2 also gives convergence of the
energy expansion throughout the corresponding high-field disk.
The coefficient algorithm is formal: computing the first \(p\)
coefficients does not require the eventual evaluation point to lie in
that disk.  We use the coefficient bound and algorithm exactly in this
Hermitian form.  The required rational-input bit implementation is
specified below.

\begin{lemma}[Complex perturbation coefficients]
\label{alg:complexification}
Let \(V=\sum_{\{u,v\}}V_{uv}\), where the local matrices are arbitrary
complex matrices with \(\|V_{uv}\|\le J\).
The germ \eqref{alg:energy-germ} extends to a holomorphic eigenvalue
function on
\[
 |x|<\rho,\qquad \rho=\frac{2^{-18}}{DJ},
 \]
and its coefficients satisfy
\begin{equation}
 |e_k(V)|\le2^{-15}N\rho^{-k}\qquad(k\ge1).
 \label{alg:complex-bound}
\end{equation}
If the entries of the local matrices are Gaussian rational, for each
integer \(p\ge1\) its first \(p\) coefficients can be computed exactly using
\(N\exp(O_D(p))\) arithmetic operations on Gaussian rationals.
\end{lemma}

\begin{proof}
Let \(Q=\Id-|\Omega\rangle\langle\Omega|\), and let \(S\) be the
inverse of \(H_0\) on \(Q(\C^2)^{\otimes N}\), extended by zero on
\(\Omega\).  Normalize the eigenvector germ by
\[
 v(x)=\Omega+\sum_{k\ge1}v_kx^k,\qquad
 \langle\Omega,v_k\rangle=0.
 \]
Writing \(v_0=\Omega\), coefficient comparison in
\((H_0+xV)v(x)=F_V(x)v(x)\) gives
\begin{align}
 e_k(V)&=\langle\Omega,Vv_{k-1}\rangle,
 \label{alg:formal-energy}\\
 v_k&=-S Q\left(
        Vv_{k-1}-\sum_{j=1}^{k-1}e_j(V)v_{k-j}
       \right).
 \label{alg:formal-vector}
\end{align}
The vacuum coordinate functional in these formulas is fixed; it does
not conjugate any variable matrix entry of \(V\).
Induction shows that \(v_k\) and \(e_k(V)\) are homogeneous
polynomials of degree \(k\) in the entries of \(V\), with no conjugate
variables.  The same recursion also proves uniqueness of the formal
normalized eigenpair.  Thus the coefficients returned by
Proposition~\ref{alg:bdl} are the restrictions of these polynomials
to Hermitian inputs.

Decompose the local terms as
\[
 B_{uv}=\frac{V_{uv}+V_{uv}^*}{2},\qquad
 C_{uv}=\frac{V_{uv}-V_{uv}^*}{2i},
 \]
and write \(V=B+iC\) for the resulting sums.
Both \(B_{uv}\) and \(C_{uv}\) are Hermitian of norm at most \(J\).
For real \(\theta\), the local perturbation
\[
 W_\theta=B\cos\theta+C\sin\theta
 =\frac12 e^{-i\theta}V+\frac12 e^{i\theta}V^*
 \]
is Hermitian, with each local norm at most \(2J\).
Since \(e_k\) is homogeneous of degree \(k\), the Fourier coefficient
of frequency \(-k\) in \(e_k(W_\theta)\) is
\(2^{-k}e_k(V)\).  Indeed, a monomial of degree \(k\) contributes
that frequency only when its \(V\)-term is selected in all \(k\)
factors.  Consequently
\begin{equation}
 e_k(V)=\frac{2^k}{2\pi}
       \int_0^{2\pi} e^{ik\theta}e_k(W_\theta)\,d\theta.
 \label{alg:fourier}
\end{equation}
Apply \eqref{alg:bdl-bound} with \(L=2J\).  The corresponding
radius is \(2^{-17}/(DJ)\), so \eqref{alg:fourier} gives
\[
 |e_k(V)|
 \le 2^k\,2^{-15}N
          \left(\frac{2^{-17}}{DJ}\right)^{-k}
 =2^{-15}N\rho^{-k}.
 \]
This proves normal convergence on compact subsets of \(|x|<\rho\).
Near zero the sum is the actual eigenvalue germ, by
\eqref{alg:formal-energy}--\eqref{alg:formal-vector} and the
simple eigenvalue of \(H_0\).  The identity theorem applied to
\[
 \det\bigl(H_0+xV-F_V(x)\Id\bigr)
 \]
therefore proves the eigenvalue identity on the entire disk.
This step does not assert simplicity there or make any ordering
claim about the other eigenvalues.

For the coefficient algorithm, take the distinct rational nodes
\[
 s_\ell=-1+\frac{2\ell}{p},\qquad 0\le\ell\le p.
 \]
At each node, \(B+s_\ell C\) is a Hermitian local perturbation
with norm bound \(2J\).  Apply
Proposition~\ref{alg:bdl} at each of these \(p+1\) inputs.
For \(k\le p\), the function
\[
 s\longmapsto e_k(B+sC)
 \]
is a polynomial of degree at most \(k\).
Its values at the nodes determine its value at \(s=i\), namely
\(e_k(V)\), by exact rational interpolation.
The factor \(p+1\) in the number of calls and the polynomial
interpolation cost are absorbed in \(N\exp(O_D(p))\).
Thus the integral in \eqref{alg:fourier} is used only for a bound;
the algorithm uses no numerical quadrature.
\end{proof}

\begin{remark}
\label{alg:complexification-scope}
The factor four between \(R_J\) and \(\rho\) has two sources:
the Hermitian family has local norm bound \(2J\), and extracting
the extremal Fourier coefficient contributes \(2^k\).
The argument transfers the coefficient estimate, not a Hermitian
spectral-ordering or spectral-simplicity theorem.  This distinction
is necessary for arbitrary nonnormal inputs.
\end{remark}

One-site perturbations can be incorporated without altering this
published input.  Attach a private auxiliary qubit \(i'\) to each
physical site \(i\), and represent a one-site term \(V_i\) by
\(V_i\otimes\Id_{i'}\) on the edge \(\{i,i'\}\).
Give every auxiliary qubit the unperturbed term
\(2|1\rangle\langle1|_{i'}\).  The resulting operator is exactly
\begin{equation}
 \widetilde H(x)
  =(2E+xV)\otimes\Id_{\mathrm{aux}}
   +\Id_{\mathrm{phys}}\otimes
        2\sum_{i'}|1\rangle\langle1|_{i'}.
 \label{alg:auxiliary}
\end{equation}
Its vacuum eigenvalue germ and all its coefficients agree with
those of the physical operator.
The number of qubits is \(N=2n\), the maximum degree is at most
\(D=d+1\), and the local norm bound is unchanged.
The identity on the auxiliary factor is important: the auxiliary
qubit is not coupled to the physical perturbation.
This construction is only a device for using the coefficient
algorithm; it imposes no stability assumption on the auxiliary
description.

\subsection{A global eigenvalue branch}
\label{alg:branch-subsection}

\begin{lemma}[Half-plane continuation]
\label{alg:half-plane}
Under \eqref{alg:promise}, the matrix
\[
 A(h)=A_0-2hE
 \]
has an algebraically simple eigenvalue \(a(h)\), uniquely maximizing
real part, for every \(h\) with \(\Re h>0\).
The function \(a\) is holomorphic on that half-plane.
For \(V=-A_0\), the function
\begin{equation}
 F(x)=-x\,a(1/x),\qquad \Re x>0,
 \label{alg:reciprocal}
\end{equation}
is a holomorphic eigenvalue of \(2E+xV\).
It agrees with the vacuum germ of
Lemma~\ref{alg:complexification} near the positive real origin.
\end{lemma}

\begin{proof}
Write \(h=u+iv\), where \(u>0\).
The imaginary Euler flow acts on polynomials by
\[
 e^{-2itvE}p(z_1,\ldots,z_n)
   =p(e^{-2itv}z_1,\ldots,e^{-2itv}z_n),
 \]
and therefore preserves \(\mathcal S_n\).
The finite-dimensional product formula gives
\[
 e^{t(A_0-2ivE)}
   =\lim_{m\to\infty}
       \bigl(e^{tA_0/m}e^{-2itvE/m}\bigr)^m.
 \]
For each stable input, every approximating output is stable.
The limit is stable or zero by Hurwitz's theorem, and it cannot
be zero on a nonzero input because the limiting exponential is
invertible.  Hence \(A_0-2ivE\) also generates a preserving
semigroup.

Theorem~\ref{spec:gap}, applied with damping \(u\), supplies an
algebraically simple eigenvalue of
\((A_0-2ivE)-2uE=A(h)\) whose real part exceeds those of all the
others by at least \(2u\).  At each point, the implicit function
theorem for the characteristic polynomial makes this eigenvalue
locally holomorphic.  Spectral continuity preserves its strict
real-part inequality in a sufficiently small neighborhood.
Thus the locally defined functions agree wherever their
neighborhoods overlap: they select the same uniquely distinguished
eigenvalue.  This proves the global holomorphic assertion.

The reciprocal map takes the right half-plane into itself, so
\eqref{alg:reciprocal} is holomorphic there.  The matrix identity
\[
 -x A(1/x)=2E+xV
 \]
proves the eigenvalue assertion.  For positive real \(x\), the
scalar \(-x\) reverses real-part order, so \(F(x)\) is the unique
eigenvalue of \(2E+xV\) with smallest real part.  We do not make
this ordering assertion for nonreal \(x\).

It remains to identify the high-field germ without assuming
Hermiticity of \(V\).  Choose positive \(x\) so small that
\(x\|V\|<1/4\).
Normality of \(2E\) implies that every eigenvalue of
\(2E+xV\) is within distance \(1/4\) of
\(\{0,2,\ldots,2n\}\).  To see the inclusion, outside these disks
the resolvent of \(2E\) has norm at most \(4\), so its
product with \(xV\) has norm less than \(1\).
The Neumann series therefore makes the perturbed resolvent
invertible, including when \(V=0\).
The Riesz projection inside \(|z|=1\) has rank one:
the contour stays in the resolvent set throughout
\(2E+\tau xV\), \(0\le\tau\le1\), and its rank at \(\tau=0\)
is one.  Hence the eigenvalue near zero is precisely the vacuum
branch.  All other eigenvalues have real part greater than one.
The vacuum branch is therefore the uniquely minimal real-part
eigenvalue, proving its agreement with \(F(x)\) on a positive
real interval.
\end{proof}

Apply the identity theorem on the connected overlap
\(\{\Re x>0\}\cap\{|x|<\rho\}\).
Lemma~\ref{alg:half-plane} and
Lemma~\ref{alg:complexification}, with the auxiliary construction
when needed, yield one holomorphic function on
\begin{equation}
 \Omega_\rho=\{x\in\C:\Re x>0\}\cup\{x\in\C:|x|<\rho\}.
 \label{alg:glued-domain}
\end{equation}
It remains an eigenvalue of \(\widetilde H(x)\) throughout this
domain, by \eqref{alg:auxiliary} and the characteristic-polynomial
identity.  In particular,
\begin{equation}
 |F(x)|\le\|\widetilde H(x)\|
          \le2N+|x|W,\qquad W=NDJ.
 \label{alg:norm-bound}
\end{equation}
The displayed \(W\) is a convenient rational upper bound on the
sum of all local norms.  The scalar estimate
\(|\lambda|\le\|T\|\) for an eigenvalue of a matrix \(T\)
does not require normality.
Finally,
\begin{equation}
 a_\star=-\mu F(1/\mu).
 \label{alg:target}
\end{equation}

\subsection{Effective continuation and exact arithmetic}
\label{alg:continuation-subsection}

The next lemma gives the continuation step independently of its
operator-theoretic application.

\begin{lemma}[An explicit continuation map]
\label{alg:map}
Let \(\rho,c>0\), and suppose that \(F\) is holomorphic on
\(\Omega_\rho\).
Define
\begin{equation}
 \eta=\frac{\rho^2}{4c+2\rho},\qquad
 R=c+\eta,\qquad q=\frac cR,\qquad
 b=\frac{R^2-c^2}{R},\qquad
 \phi(w)=\frac{bw}{1-qw}.
 \label{alg:map-parameters}
\end{equation}
Then \(0<q<1\), \(\phi(\D)\subset\Omega_\rho\),
\(\phi(0)=0\), and \(\phi(q)=c\).
If \(|F(x)|\le L_0+L_1|x|\) on \(\Omega_\rho\), the function
\(G=F\circ\phi\) satisfies
\[
 |G(w)|\le M:=L_0+L_1(c+R)\qquad(w\in\D).
 \]
Writing \(F(x)=\sum_{j\ge0}f_jx^j\) at zero and
\(G(w)=\sum_{k\ge0}g_kw^k\), one has \(g_0=f_0\) and
\begin{equation}
 g_k=\sum_{j=1}^k
        f_j b^j q^{k-j}\binom{k-1}{j-1}
       \qquad(k\ge1).
 \label{alg:coefficient-transform}
\end{equation}
For every \(p\ge0\),
\begin{equation}
 \left|F(c)-\sum_{k=0}^p g_kq^k\right|
    \le\frac{Mq^{p+1}}{1-q}.
 \label{alg:tail}
\end{equation}
\end{lemma}

\begin{proof}
The definitions give \(q\in(0,1)\), \(\eta\le\rho/2\), and
\[
 0<R^2-c^2=2c\eta+\eta^2<\rho^2.
 \]
If \(|x-c|<R\) and \(\Re x\le0\), then
\[
 |x|^2=|x-c|^2+2c\Re x-c^2<R^2-c^2<\rho^2.
 \]
Thus the disk of center \(c\) and radius \(R\) lies in
\(\Omega_\rho\).
The identity
\[
 \phi(w)=c+R\frac{w-q}{1-qw}
 \]
shows that \(\phi\) maps \(\D\) biholomorphically onto this
disk.  Direct substitution gives its values at \(0\) and \(q\).
Since \(|\phi(w)|<c+R\), the bound on \(G\) follows.

For \(j\ge1\), expansion at zero gives
\[
 \phi(w)^j=b^j w^j(1-qw)^{-j}
  =b^j\sum_{k\ge j}
       \binom{k-1}{j-1}q^{k-j}w^k.
 \]
Only \(j\le k\) contributes to the coefficient of \(w^k\),
which proves \eqref{alg:coefficient-transform}.
Applying Cauchy's estimate on circles with radii increasing to
one gives \(|g_k|\le M\).
The geometric tail at \(w=q\) now proves \eqref{alg:tail}.
No holomorphic extension to the boundary of \(\D\) is used.
\end{proof}

\begin{proof}[Proof of Theorem~\ref{alg:complex}]
The existence, uniqueness, and simplicity of \(a_\star\) follow
from Theorem~\ref{spec:gap}.
Use \(V=-A_0\), the auxiliary construction
\eqref{alg:auxiliary}, and the parameters
\eqref{alg:parameters}.
Lemmas~\ref{alg:complexification} and
\ref{alg:half-plane} give the function \(F\) on
\(\Omega_\rho\), with \eqref{alg:norm-bound}.
Apply Lemma~\ref{alg:map} with
\[
 L_0=2N,\qquad L_1=W,\qquad
 M=2N+(c+R)W.
 \]
All the parameters in \eqref{alg:map-parameters} are rational.

Choose the smallest integer \(p\ge0\) for which
\begin{equation}
 \mu M q^{p+1}\le\delta(1-q).
 \label{alg:stopping}
\end{equation}
The algorithm finds this integer by rational multiplication and
comparison; evaluation of logarithms is unnecessary.
If \(p>0\), compute \(e_1(V),\ldots,e_p(V)\) by
Lemma~\ref{alg:complexification}, and then compute
\(g_0,\ldots,g_p\) by
\eqref{alg:coefficient-transform}.
Here \(g_0=F(0)=0\).
Return the Gaussian rational
\begin{equation}
 \widehat a=-\mu\sum_{k=0}^p g_kq^k.
 \label{alg:output}
\end{equation}
For \(p=0\), no coefficient call is made and the same formula
returns zero.
Equations \eqref{alg:target}, \eqref{alg:tail}, and
\eqref{alg:stopping} give
\(|\widehat a-a_\star|\le\delta\).

We first bound the number of arithmetic operations.
A useful exact identity is
\begin{equation}
 \frac1{1-q}=\frac R\eta
     =1+\frac{4c^2}{\rho^2}+\frac{2c}{\rho}.
 \label{alg:contraction-parameter}
\end{equation}
Since \(-\log q\ge1-q\), the stopping rule implies
\[
 p=O\!\left(
   \frac1{1-q}
   \log\!\left(2+\frac{\mu M}{\delta(1-q)}\right)
 \right).
 \]
The bound \eqref{alg:order-bound} follows, because
\(M=O_{d,J,\mu}(n)\).
For fixed \(d,J,\mu\), this is
\(p=O_{d,J,\mu}(\log(2+n/\delta))\).
The coefficient algorithm uses
\(N\exp(O_D(p))\) arithmetic operations.
Interpolation, the triangular transform, and evaluation in
\eqref{alg:output} add polynomially many operations in \(p\).

We finish by verifying that exact arithmetic has only polynomial
bit overhead.  This also specifies a finite-input algorithm rather
than a real-arithmetic oracle.
Let \(B_0\) bound the bit length of each numerator and denominator
in the original local matrices.
For \(p>0\), the interpolation nodes
\(s_\ell=-1+2\ell/p\) have \(O(\log(p+1))\) bits.
Taking local Hermitian parts and forming \(B+s_\ell C\)
therefore produces Gaussian-rational inputs with bit lengths
polynomial in \(B_0\) and \(\log(p+1)\).

Consider one Hermitian coefficient computation at these nodes.
Let \(Q_0\) be the product of the positive denominators of the
real and imaginary parts of its local entries.  There are
\(O_D(N)\) such entries, so \(\log Q_0\) is polynomial in the
input length.  Set
\begin{equation}
 K_p=48(p+1)!,\qquad
 D_j=(Q_0K_p)^j\quad(0\le j\le p).
 \label{alg:common-denominator}
\end{equation}
We recall the properties of the coefficient recursion used in
Sections~3--5 of \cite{BDL08}.  An order-\(j\) creation
coefficient \(C_j(M)\) is zero unless \(M\) is contained in a
connected set of at most \(j+1\) vertices.  Each nonzero term in
its recursion contains one local matrix entry, previous creation
coefficients of total order \(j-1\), and a divisor
\[
 k!\,E_0(M),\qquad k\le4,\qquad E_0(M)=2|M|.
 \]
The basis creation matrices
\(\prod_{u\in M}|1\rangle\langle0|_u\) have entries zero or one;
expanding their nested commutators therefore produces signed sums
of local perturbation entries.
The corresponding energy recursion has the same form without
the outer divisor \(E_0(M)\).
These are the explicit recursions used to implement the published
coefficient algorithm, rather than an assumption about arbitrary
symbolic perturbation expressions.

For \(j\le p\), every product \(k!E_0(M)\) just displayed
divides \(K_p\).  Induction in \(j\) consequently shows that
\begin{equation}
 D_j C_j(M)\in\mathbb Z[i],
 \qquad D_j e_j(B+s_\ell C)\in\mathbb Z[i].
 \label{alg:denominator-induction}
\end{equation}
Indeed, the previous coefficients contribute denominator
dividing \((Q_0K_p)^{j-1}\), the new local entry contributes
a factor dividing \(Q_0\), and the remaining divisor divides
\(K_p\).  All terms at a fixed order may be accumulated over
this common denominator.  Its bit length is
\[
 O\bigl(p\log Q_0+p^2\log(p+1)\bigr).
 \]

Intermediate numerators have polynomial bit length as well.
Here a deliberately loose bound suffices.
Fully unfold a term of perturbation order at most \(p\).
It has at most \(p\) local insertions and bounded-arity recursion
nodes.  Choose its local edges in at most \((ND)^{O(p)}\)
ways.  Their union has \(O(p)\) vertices, and all intermediate
creation sets lie in this union.  Choosing these sets at
\(O(p)\) nodes, together with the tree shape, order partitions,
and bounded-depth commutator terms, gives at most
\[
 \exp\bigl(O_D(p^2+p\log(N+1))\bigr)
 \]
fully unfolded summands.  Each summand contains at most \(p\)
bounded local entries; the positive integer energy and factorial
divisors cannot increase its absolute value.
Its magnitude is therefore at most
\(\max(1,2J)^{O(p)}\) times a fixed exponential factor.
After introducing \eqref{alg:common-denominator}, this bounds
the logarithmic magnitude of every intermediate numerator by
a polynomial in \(p,N,B_0\).
This argument concerns bit sizes only.  The algorithm does not
unfold these trees: its sharper operation count is the inherited
\(N\exp(O_D(p))\) count.

Exact interpolation also has polynomial bit cost.
For the nodes \(s_\ell=-1+2\ell/p\), its weights at \(i\)
can be written
\begin{equation}
 L_\ell(i)=
 \frac{\displaystyle
       \prod_{\substack{0\le r\le p\\r\ne\ell}}
           (pi+p-2r)}
      {2^p(-1)^{p-\ell}\ell!(p-\ell)!}.
 \label{alg:interpolation-weights}
\end{equation}
Both numerator and denominator have
\(O(p\log(p+1))\) bits.
The coefficients \(e_k(V)\) are obtained by multiplying these
weights by the computed node values and adding \(p+1\) terms.
The rational powers and binomial coefficients in
\eqref{alg:coefficient-transform}, and the powers in
\eqref{alg:stopping}--\eqref{alg:output}, increase bit length
only polynomially.
Thus a safe total bit bound is
\[
 N\exp(O_D(p))\operatorname{poly}(N,B_0,p,L_{\mathrm{in}}),
 \]
where \(L_{\mathrm{in}}\) is the total input bit length,
including the requested accuracy.
Together with \eqref{alg:order-bound}, this proves the theorem.
\end{proof}

\subsection{Suzuki--Fisher ground energies}
\label{alg:sf-subsection}

The preceding result applies to the Suzuki--Fisher class of
Theorem~\ref{eq:sf}.  We keep its sign convention:
the longitudinal term is
\(-\sum_i\mu_i^z Z_i\), with \(\mu_i^z\ge0\).
The Lee--Yang theorem gives disk stability preservation by the
Gibbs semigroup \(e^{-tH}\); see
\cite{SuzukiFisher71,WBG26,BGLW26}.

\begin{theorem}[Fixed positive longitudinal field]
\label{alg:sf}
Fix a maximum degree \(d\), a positive rational bound \(J\) on
the absolute values of the Pauli coefficients, and a positive
rational \(\mu\).
Let \(H\) be a Suzuki--Fisher Hamiltonian with rational Pauli
coefficients on an \(n\)-vertex graph of maximum degree at most
\(d\), and suppose
\[
 \mu_i^z\ge\mu\qquad(1\le i\le n).
 \]
For every positive rational \(\delta\), its ground energy
\(E_0(H)\) can be approximated to absolute error at most
\(\delta\) by a deterministic classical algorithm polynomial in
\(n\), the input bit length, and \(\delta^{-1}\).
The polynomial may depend on \(d,J,\mu\).
\end{theorem}

\begin{proof}
Define
\[
 V=H+\mu\sum_i Z_i.
 \]
The longitudinal fields of \(V\) are \(\mu_i^z-\mu\ge0\);
all its interaction inequalities are unchanged.
Thus \(V\) is again Suzuki--Fisher, and
\(A_0=-V\) satisfies \eqref{alg:promise}.
Its local matrices have Gaussian-rational entries.
Combining the finitely many Pauli terms on each edge or site
gives a local norm bound depending only on \(J,\mu\), and leaves
the degree bound unchanged.
For example, sums of the absolute Pauli coefficients give
explicit rational norm upper bounds, so the algorithm need not
compute exact matrix norms.

Since \(2E=n\Id-\sum_i Z_i\),
\begin{equation}
 A_0-2\mu E=-H-\mu n\Id.
 \label{alg:sf-shift}
\end{equation}
The matrix on the right is Hermitian.  Its eigenvalue with
largest real part is therefore
\(-E_0(H)-\mu n\).
Apply Theorem~\ref{alg:complex} with error \(\delta\).
If it returns \(\widehat a\), return the real rational
\[
 \widehat E=-\Re\widehat a-\mu n.
 \]
Equation \eqref{alg:sf-shift} gives the required error bound
and the asserted complexity.
\end{proof}

\begin{remark}[Hermitian coefficient bound]
\label{alg:hermitian-radius}
For this application, the complexification step can be omitted:
the perturbation \(V\) is Hermitian.  Applying
Proposition~\ref{alg:bdl} directly gives the larger guaranteed
disk \(2^{-16}/(DJ')\), where \(J'\) is the local norm bound
after the auxiliary construction.  The same continuation proof
then applies to the ground-energy branch of \(2E+xV\).
The factor-four improvement in this disk changes constants,
not the fixed-parameter complexity conclusion.
\end{remark}

Theorem~\ref{alg:sf} includes nonstoquastic interactions and
does not require conservation of total degree.  Its computational
output and parameter dependence differ from several earlier
results.  Bravyi and Gosset \cite{BG17} give randomized
algorithms for a two-axis stoquastic family on arbitrary graphs.
Harrow, Mehraban, and Soleimanifar \cite[Section~7 of the full version]{HMS20}
give an all-temperature quasipolynomial partition-function
algorithm for a number-conserving XXZ subclass.  Their
fugacity-polynomial reduction uses the commutation of the
interaction with total magnetization.
Bravyi, Gosset, Liu, and Wong \cite{BGLW26} give a polynomial-time
quantum ground-energy algorithm for the full Suzuki--Fisher
family, including zero fields and inverse-polynomial absolute
accuracy.  Their sharp Hermitian gap and quantum algorithm are
not consequences newly claimed here.
Theorem~\ref{alg:sf} gives a deterministic classical conclusion
on bounded-degree inputs at every fixed positive field.
It does not replace the unrestricted quantum theorem.

\subsection{Energy-value approximation at zero field}
\label{alg:value-subsection}

A fixed positive field can also be used as a controlled
perturbation when the requested error is proportional to the
number of sites.

\begin{theorem}[Additive energy-value approximation]
\label{alg:value}
Fix a maximum degree \(d\), a bound \(J\) on the absolute
Pauli coefficients, and a positive rational \(\varepsilon\).
For every rational Suzuki--Fisher Hamiltonian on an
\(n\)-vertex graph satisfying these bounds, a deterministic
classical algorithm polynomial in \(n\) and the input bit length
returns \(\widehat E\in\mathbb Q\) such that
\[
 |\widehat E-E_0(H)|\le\varepsilon n.
 \]
The longitudinal fields may vanish.
The polynomial may depend on \(d,J,\varepsilon\).
The output is an energy value; no approximating state is
asserted.
\end{theorem}

\begin{proof}
It suffices to consider \(0<\varepsilon\le1\), since a smaller
accuracy parameter also meets a larger requested tolerance.
Put \(\nu=\varepsilon/4\) and
\[
 H_\nu=H-\nu\sum_i Z_i.
 \]
This is a Suzuki--Fisher Hamiltonian whose longitudinal fields
are all at least the fixed positive number \(\nu\).
Its local coefficients are bounded in terms of \(J,\varepsilon\).
The variational characterization of the lowest eigenvalue gives
\[
 |E_0(H_\nu)-E_0(H)|
 \le\|H_\nu-H\|\le\nu n.
 \]
Use Theorem~\ref{alg:sf} to estimate \(E_0(H_\nu)\) with
absolute error at most \(\varepsilon n/2\).
The returned value differs from \(E_0(H)\) by at most
\[
 \frac{\varepsilon n}{4}+\frac{\varepsilon n}{2}
   =\frac{3\varepsilon n}{4}
   \le\varepsilon n.
 \]
For fixed \(d,J,\varepsilon\), the field \(\nu\) and all local
bounds are fixed.  The runtime in Theorem~\ref{alg:sf} is
therefore polynomial in the input size and \(n\).
\end{proof}

We spell out the consequence for Quantum MaxCut and EPR
Hamiltonians, keeping value approximation separate from state
construction.
For a graph \(G=(V,\mathcal E)\) with nonnegative rational weights, let
\[
 Q_G=\sum_{\{i,j\}\in\mathcal E}
      w_{ij}|\Psi^-\rangle\langle\Psi^-|_{ij},
 \qquad
 |\Psi^-\rangle=\frac{|01\rangle-|10\rangle}{\sqrt2},
 \]
and write \(\operatorname{OPT}(G)=\lambda_{\max}(Q_G)\).
Similarly, define the EPR objective
\[
 P_G=\sum_{\{i,j\}\in\mathcal E}
      w_{ij}|\Phi^+\rangle\langle\Phi^+|_{ij},
 \qquad
 |\Phi^+\rangle=\frac{|00\rangle+|11\rangle}{\sqrt2}.
 \]

\begin{corollary}[EPR and bipartite Quantum MaxCut values]
\label{alg:qmc}
For fixed degree and upper weight bounds and fixed rational
\(\varepsilon>0\), the optimal values of \(P_G\) on arbitrary
graphs and of \(Q_G\) on bipartite graphs admit deterministic
classical additive-\(\varepsilon |V|\) approximation in
polynomial time.
For unweighted bipartite Quantum MaxCut, this gives a relative
energy-value approximation scheme after isolated vertices are
deleted.  In each case the algorithm estimates the optimum
value and does not produce a state attaining the estimate.
\end{corollary}

\begin{proof}
The negative EPR objective is Suzuki--Fisher up to a scalar:
\begin{equation}
 -P_G
  =-\frac14\sum_{\{i,j\}\in\mathcal E}w_{ij}
       \bigl(\Id+X_iX_j-Y_iY_j+Z_iZ_j\bigr).
 \label{alg:epr-expansion}
\end{equation}
On each edge the transverse matrix has operator norm
\(w_{ij}/4\), equal to its longitudinal interaction
coefficient.  Subtract the known scalar
\(-\frac14\sum_{\{i,j\}}w_{ij}\) when applying
Theorem~\ref{alg:value}, and restore it afterward.
This proves the assertion for \(P_G\).

If \(G\) is bipartite with parts \(A,B\), let
\(U=\prod_{i\in A}Y_i\).
On every edge, \(U\) maps the singlet vector to the EPR vector
up to a scalar of modulus one.  Hence
\[
 UQ_GU^*=P_G.
 \]
The spectra agree, proving the additive assertion for
\(Q_G\).  This is the same local change of basis used in
\cite[Corollary~2]{BGLW26}.

For the unweighted bipartite Quantum MaxCut statement, remove isolated
vertices and let \(N\) and \(m\) be the remaining vertex and
edge counts.  If \(m=0\), the value is zero and is returned
exactly.  Otherwise \(N\le2m\).
The product vector which is \(0\) on
one part and \(1\) on the other has expectation \(m/2\).
Thus \(\operatorname{OPT}(G)\ge m/2\).
Apply the additive algorithm with error
\(\varepsilon N/4\).  Its error is bounded by
\[
 \frac{\varepsilon N}{4}
 \le\frac{\varepsilon m}{2}
 \le\varepsilon\,\operatorname{OPT}(G).
 \]
For each fixed \(\varepsilon\), the runtime is polynomial,
which proves the relative value approximation assertion.
\end{proof}

The additive statement for bounded weights does not automatically
give a relative guarantee when positive weights may be arbitrarily
small.  Nor is the value scheme a state-producing approximation
algorithm of the type studied, for example, in \cite{ALMPSS26}.
The inverse-polynomial absolute-accuracy conjectures for general
EPR problems in \cite{MarwahaSud26} remain outside its scope.
Indeed, introducing a field of order \(\delta/n\) would make
the factor \(c/\rho\) in \eqref{alg:order-bound} grow with \(n\).
The resulting coefficient count need no longer be logarithmic,
and the algorithm proved here need no longer run in polynomial
time.

% END SECTION algorithms

% BEGIN SECTION principal_states
\subsection{Stable tests of principal eigenvectors}
\label{alg:states-subsection}

We now pass from one eigenvalue to queries about its principal
eigenvectors.  The input remains \eqref{alg:local-input}--\eqref{alg:promise}.
For the binary coefficient space, the transpose exchanges the two
groups of variables in a matrix kernel.  Proposition~\ref{spec:perron}
therefore gives strictly disk-stable right eigenvectors of both
\(A=A_0-2\mu E\) and \(A^{\mathsf T}\).  Normalize them by
\[
 v_\varnothing=w_\varnothing=1.
\]
The transpose, rather than the adjoint, makes the dependence on a
complex parameter holomorphic.  A \emph{stable product test} is
\(B=\bigotimes_{i=1}^n B_i\), where \(B_i\in M_2(\C)\),
\((B_i)_{00}=1\), and
\[
 K_{B_i}(z,w)=
 (B_i)_{00}+(B_i)_{10}z+(B_i)_{01}w+(B_i)_{11}zw
\]
does not vanish in \(\D^2\).  The matrices \(B_i\) need not be
Hermitian, positive, or invertible.

\begin{theorem}[Classical product tests of principal states]
\label{alg:product-test}
Fix \(d,J,\mu\) as in Theorem~\ref{alg:complex}, and suppose that
\(A_0\) satisfies that theorem's input assumptions.  Given a stable
product test with Gaussian-rational entries, both
\[
 G_B=w^{\mathsf T}Bv,\qquad
 \Tr(B\Pi)=\frac{w^{\mathsf T}Bv}{w^{\mathsf T}v},
 \qquad \Pi=\frac{vw^{\mathsf T}}{w^{\mathsf T}v},
\]
are nonzero.  For every rational \(0<\epsilon<1\), a deterministic
classical algorithm returns Gaussian-rational estimates of these
two quantities with relative complex error at most \(\epsilon\).
Its bit complexity is
\[
 \operatorname{poly}_{d,J,\mu}
       (n,\epsilon^{-1},L_{\mathrm{in}}),
\]
where \(L_{\mathrm{in}}\) includes the local generator, the test
matrices, and the accuracy.  Relative complex error means
\(|\widehat z/z-1|\le\epsilon\).
\end{theorem}

This is query access to specified scalar tests, not a listing of an
exponentially long vector.  No polynomial bound on \(\|\Pi\|\) is
assumed or obtained.  The proof instead constructs a nonvanishing
scalar function with a controlled logarithm in the reciprocal-field
parameter.

\subsection{A uniform domain for normalized eigenvectors}

Write \(V=-A_0\), \(H(x)=2E+xV\), and set
\begin{equation}\label{alg:state-radii}
 r_0=\frac{2^{-18}}{(d+1)J},\qquad \rho_s=\frac{r_0}{8}.
\end{equation}
The following use of the Kirkwood--Thomas expansion concerns its
creation coefficients, not just the energy estimate of
Proposition~\ref{alg:bdl}.
For a nonempty set \(M\subseteq[n]\), put
\[
 a_M^\dagger=\prod_{i\in M}|1\rangle\langle0|_i .
\]
These operators commute, and products with overlapping supports vanish.
The vacuum-normalized eigenvector germ has a unique expression
\[
 v(x)=\exp\left(-\sum_{\varnothing\ne M\subseteq[n]}
                   C(M;x)a_M^\dagger\right)\Omega,
 \qquad C(M;x)=\sum_{k\ge1}C_k(M)x^k.
\]
Uniqueness follows by taking the finite logarithm in the nilpotent
commutative algebra generated by the \(a_i^\dagger\).

\begin{lemma}[Creation coefficients and a zero-free vacuum disk]
\label{alg:creation-disk}
For every local complex \(V\) with the bounds above,
\begin{equation}\label{alg:creation-bound}
 \max_i\sum_{M\ni i}|C_k(M)|\le2^{-15}r_0^{-k},
 \qquad C_k(M)=0\quad\hbox{if }|M|>k+1.
\end{equation}
The normalized right eigenvector germs of \(H(x)\) and
\(H(x)^{\mathsf T}\) extend holomorphically to \(|x|<\rho_s\).
Their coefficient polynomials are nonzero on the closed polydisk
\(\{|z_i|\le2\}\) throughout this parameter disk.
\end{lemma}

\begin{proof}
For Hermitian perturbations with local norm bound \(L\) on a graph
of maximum degree \(D\), Lemma~5, equation~(32), and Lemma~7,
equation~(47), of \cite{BDL08}, with unperturbed gap \(2\), give
\[
 \max_i\sum_{M\ni i}|C_k(M)|\le2^{-15}R_L^{-k},
 \qquad R_L=2^{-16}/(DL),
\]
and say that a nonzero \(C_k(M)\) is supported inside a connected
set of at most \(k+1\) vertices.  In particular \(|M|\le k+1\).
Use the private auxiliary qubits in \eqref{alg:auxiliary} for
one-site terms, with \(D=d+1\).  The normalized eigenvector germ
of the augmented operator is exactly the physical germ tensored
with the auxiliary vacuum.  Its creation coefficients involving
an auxiliary vertex vanish; all other coefficients agree with
the physical ones.

The recursion \eqref{alg:formal-energy}--\eqref{alg:formal-vector}
and the finite logarithm show that \(C_k(M)\) is a homogeneous
polynomial of degree \(k\) in the local entries of \(V\), without
their conjugates.  With \(W_\theta\) from the proof of
Lemma~\ref{alg:complexification},
\[
 C_k(M;V)=\frac{2^k}{2\pi}\int_0^{2\pi}
              e^{ik\theta}C_k(M;W_\theta)\,d\theta.
\]
Taking absolute values, summing over \(M\ni i\), and applying the
Hermitian estimate with \(L=2J\) proves
\eqref{alg:creation-bound}.  The support conclusion is inherited
by the same coefficient identity.  Transposition gives the
identical estimates for the left germ.
For \(|x|<r_0/8\), it follows that
\begin{equation}\label{alg:weighted-creation}
 \max_i\sum_{M\ni i}4^{|M|}|C(M;x)|
 \le 2^{-15}\sum_{k\ge1}4^{k+1}(|x|/r_0)^k
 <2^{-13}.
\end{equation}
All these series converge locally uniformly.

We give the elementary polymer argument that turns this estimate
into zero-freeness.  For weights \(u_M\) on nonempty subsets of a
finite set, define
\[
 Z(S)=\sum_{\substack{\mathcal A\text{ a collection of}\\
              \text{pairwise disjoint nonempty subsets of }S}}
                       \prod_{M\in\mathcal A}u_M .
\]
Suppose that
\(\sup_i\sum_{M\ni i}|u_M|2^{|M|-1}\le1/2\).
Induction on \(|S|\) proves
\begin{equation}\label{alg:polymer-ratio}
 Z(S)\ne0,\qquad
 \left|\frac{Z(S)}{Z(S\setminus\{i\})}-1\right|\le\frac12
 \quad(i\in S).
\end{equation}
Indeed,
\[
 Z(S)=Z(S\setminus\{i\})+
           \sum_{\substack{M\subseteq S\\i\in M}}u_M Z(S\setminus M).
\]
By the induction hypothesis, deleting the \(|M|-1\) additional
vertices gives
\(\left|Z(S\setminus M)/Z(S\setminus\{i\})\right|
 \le2^{|M|-1}\).
The claimed ratio estimate follows, and its disk does not contain
zero.  The empty set starts the induction.

The polynomial of \(v(x)\) is this polymer sum with
\(u_M=-C(M;x)z^M\).  The factorials in the exponential cancel
the orderings of pairwise disjoint sets.  For \(|z_i|\le2\),
the left side of the polymer condition is at most half the
left side of \eqref{alg:weighted-creation}, hence is less than
\(1/2\).  This proves the asserted zero-freeness.

The convergent creation series and its finite exponential give
holomorphic vectors on \(|x|<\rho_s\).  Their formal eigenvalue
is the convergent energy series of
Lemma~\ref{alg:complexification}.  The eigenvector equations hold
near zero and then throughout the disk by the identity theorem.
No assertion of simplicity on the entire disk is needed.
\end{proof}

\begin{lemma}[Nonvanishing principal tests]\label{alg:test-domain}
The vectors in Lemma~\ref{alg:creation-disk} and the normalized
principal vectors of \(A(1/x)\) glue to holomorphic vectors
\(v(x),w(x)\) on
\[
 \Omega_{\rho_s}=\{\Re x>0\}\cup\{|x|<\rho_s\}.
\]
For every stable product test,
\begin{equation}\label{alg:test-function}
 G_B(x)=w(x)^{\mathsf T}Bv(x)
\end{equation}
is nonvanishing on this domain, \(G_B(0)=1\), and
\begin{equation}\label{alg:test-log-bound}
 |G_B(x)|\le4^n,\qquad
 \Re F_B(x)\le n\log4,\quad F_B=\log G_B,\quad F_B(0)=0 .
\end{equation}
\end{lemma}

\begin{proof}
The half-plane argument in Lemma~\ref{alg:half-plane} applies to
the principal eigenvectors as well.  Local analytic eigenvectors
exist for a simple eigenvalue; their vacuum coordinates are nonzero
because their coefficient polynomials are strictly stable.
Vacuum normalization makes these local choices agree on overlaps.
Use \(A(h)^{\mathsf T}\) for the left vectors.
On a sufficiently small positive real interval they coincide
with the vacuum vectors, by the same Riesz-projection argument
as in Lemma~\ref{alg:half-plane}.  That interval may depend on \(n\);
it is used only to identify the branches, not to choose
\(\rho_s\).  The identity theorem on the connected overlap
gives the required gluing.

In the half-plane, \(v(x)\) and \(w(x)\) each have a strict
polydisk radius greater than one.  In the small disk they have
radius at least two by Lemma~\ref{alg:creation-disk}.
Contract their legs against the two legs of every \(K_{B_i}\).
The strict version of the disk contraction rule, equivalently
polarized apolarity, applies because each contracted pair has
product of radii greater than one; see
Lemma~\ref{pre:grace}.
The resulting scalar is \eqref{alg:test-function}, hence is
nonzero.  This argument allows singular tests and needs no
uniform lower bound for the strict radius.

For completeness, a multiaffine stable polynomial
\(f(z)=\sum_S f_Sz^S\) with \(f_\varnothing=1\) satisfies
\(|f_S|\le1\).  Set variables outside \(S\) to zero and all
remaining variables to \(t\).  If \(f_S\ne0\), the resulting
polynomial has degree \(|S|\), constant term one, leading
coefficient \(f_S\), and all roots of modulus at least one.
Their product proves the bound.  The case \(f_S=0\) is immediate.
Apply this observation to \(v,w\) and each \(K_{B_i}\).
The expansion of \(w^{\mathsf T}Bv\) has \(4^n\) terms of modulus
at most one.  Also \(G_B(0)=\prod_i(B_i)_{00}=1\).
The domain is star-shaped about zero, so its nonvanishing
holomorphic function has a unique logarithm vanishing there.
Taking its real part gives \eqref{alg:test-log-bound}.
\end{proof}

\begin{lemma}[Continuation with a one-sided logarithmic bound]
\label{alg:log-continuation}
Use the parameters \eqref{alg:map-parameters} with
\(\rho=\rho_s\) and \(c=1/\mu\).  Write
\(F_B(x)=\sum_{j\ge1} f_jx^j\) and
\(F_B(\phi(u))=\sum_{k\ge1}g_ku^k\).
The transform \eqref{alg:coefficient-transform} remains valid,
and
\begin{equation}\label{alg:log-tail}
 |g_k|\le4n,\qquad
 \left|F_B(c)-\sum_{k=1}^p g_kq^k\right|
       \le \frac{4nq^{p+1}}{1-q}.
\end{equation}
In particular \(|F_B(c)|\le Kn\), where \(K=4q/(1-q)\).
\end{lemma}

\begin{proof}
The map lemma puts \(\phi(\D)\) in \(\Omega_{\rho_s}\).
The holomorphic function
\(U(u)=2n-F_B(\phi(u))\) has nonnegative real part and \(U(0)=2n\).
On a circle of radius \(r<1\), the positive function
\(\Re U(re^{it})\) has average \(2n\).
Its \(k\)-th Fourier coefficient, for \(k\ge1\), equals
\(-g_kr^k/2\), so \(|g_k|r^k\le4n\).
Let \(r\uparrow1\) and sum the geometric tail.
The coefficient transform is a formal identity already proved
in Lemma~\ref{alg:map}.
\end{proof}

\subsection{Connected coefficients and finite bit complexity}

We compute the coefficients of \(F_B\) without expanding the
full eigenvectors.  This is an instance of the connected-logarithm
method for multiplicative analytic graph functions; compare
\cite{PatelRegts17,YYZ22}.  The needed multiplicativity is formal
at the vacuum and does not require a preservation promise on
induced subgraphs.

For \(U\subseteq[n]\), keep just the local terms supported inside
\(U\), obtaining \(V_U\), and let \(B_U=\bigotimes_{i\in U}B_i\).
Define the vacuum-normalized formal vectors \(v_U,w_U\) for
\(2E_U+xV_U\) and its transpose.  Put
\[
 F_{B,U}(x)=\log(w_U(x)^{\mathsf T}B_Uv_U(x)),\qquad
 \Phi_{B,U}(x)=\sum_{W\subseteq U}(-1)^{|U|-|W|}F_{B,W}(x).
\]
For the empty subsystem all logarithms are zero.

\begin{lemma}[Local computation of test logarithms]
\label{alg:connected-tests}
For \(k\ge1\),
\[
 [x^k]F_B(x)=
 \sum_{\substack{U\subseteq[n]\ \mathrm{connected}\\ |U|\le k+1}}
                       [x^k]\Phi_{B,U}(x).
\]
For Gaussian-rational input, the first \(p\) coefficients can
be computed exactly with
\(n\exp(O_d(p))\operatorname{poly}(L_{\mathrm{in}},n,p)\)
bit operations.
\end{lemma}

\begin{proof}
If the induced graph on \(U\) splits into components, the
normalized formal vectors tensor across components and the
test does likewise.  Thus \(F_{B,U}\) is the sum of the
component logarithms.  Subset inversion shows that
\(\Phi_{B,U}=0\) when \(U\) is disconnected.

To obtain the order bound as well, give each local perturbation
term its own parameter.  An order-\(k\) coefficient is a
homogeneous polynomial of degree \(k\) in these parameters,
by the unique formal eigenvector recursion.  If a parameter
monomial has disconnected union of supports, set all unused
parameters to zero.  Tensor factorization makes the logarithm
additive across those components, so that monomial has
coefficient zero.  A connected union of at most \(k\) one- or
two-site supports has at most \(k+1\) vertices.
Finally the coefficient of a monomial with support union \(T\)
is unchanged in every induced subsystem containing \(T\).
The alternating sum over \(W\subseteq U\) retains it only if
\(T=U\).  This proves the formula, including both vanishings.

Here is a direct finite-input implementation.  In each subsystem
of size \(m\le p+1\), use the exact recursions
\eqref{alg:formal-energy}--\eqref{alg:formal-vector}, and their
transposes, through order \(p\).
If
\[
 a_k=\sum_{i+j=k}w_{U,i}^{\mathsf T}B_Uv_{U,j},
 \qquad \log\left(1+\sum_{k\ge1}a_kx^k\right)
                  =\sum_{k\ge1}f_kx^k,
\]
then coefficient comparison after logarithmic differentiation gives
\begin{equation}\label{alg:formal-test-log}
 f_k=a_k-\frac1k\sum_{j=1}^{k-1}j f_j a_{k-j}.
\end{equation}
Every vector has dimension at most \(2^{p+1}\).
Even dense matrix arithmetic and enumeration of every
\(W\subseteq U\) cost only \(\exp(O(p))\operatorname{poly}(p)\)
operations per \(U\).  Connected subsets of size at most \(p+1\)
can be enumerated by traversals of their spanning trees, with
deduplication, in \(n\exp(O_d(p))\) time; a traversal has length
at most \(2p\).  The case \(d=0\) consists just of singletons.

We record denominator and magnitude bounds to justify the bit
claim.  Let \(Q_0\) be a common positive denominator for all local
entries of \(V\), let \(Q_B\) be the product of the denominators
of the entries of the \(B_i\), and put
\[
 D_p=2(p+1)!,\qquad T=Q_0D_p .
 \]
In a subsystem of size at most \(p+1\), each nonzero diagonal
entry of the reduced inverse \(S_U\) is \(1/(2r)\),
\(1\le r\le p+1\).  Induction in the vector recursion gives
denominators dividing \(T^{2k}\) for \(v_{U,k},w_{U,k}\)
and \(T^{2k-1}\) for the energy coefficient.
The denominator of \(a_k\) divides \(Q_BT^{2k}\).
Equation~\eqref{alg:formal-test-log} then bounds the denominator
of \(f_k\) by \(k!Q_B^kT^{2k}\).
At each order, accumulate the sums over these common
denominators; their logarithms are polynomial in
\(p\) and the total input length.

Intermediate numerators also have polynomial bit length.
For a subsystem \(U\), \(\|V_U\|\le m(d+1)J\).
On \(|x|\le r_U=[16(1+m(d+1)J)]^{-1}\), the resolvent
Neumann series on the circle \(|z|=1\) bounds the spectral
projection at the vacuum and its difference from the
unperturbed projection.  In particular, the projection has
norm at most \(16/15\), its vacuum entry differs from one by
at most \(1/15\), and both vacuum-normalized eigenvectors have
norm less than two.  Cauchy's estimate gives
\(\|v_{U,k}\|,\|w_{U,k}\|\le2r_U^{-k}\).
As \(\|B_U\|\le2^m\),
\[
 |a_k|\le4(k+1)2^m r_U^{-k}.
\]
For \(k\le p\), choose \(A_p=4(p+1)2^{p+1}\).
The finite formula
\(\log(1+a)=\sum_{\ell=1}^p(-1)^{\ell+1}a^\ell/\ell
 \pmod{x^{p+1}}\)
bounds \(|f_k|\) by \((2A_p)^kr_U^{-k}\).
Matrix products, vector sums and the partial sums in
\eqref{alg:formal-test-log} have the same type of bound, up to
factors \(2^{O(p)}\operatorname{poly}(p)\).
Subset sums add at most \(2^{p+1}\) terms, and the final
connected sum has at most \(n\exp(O_d(p))\) terms.
Combining these magnitude bounds with the common denominators
controls every stored numerator, including before cancellation.
This proves the bit complexity.
\end{proof}

\begin{proof}[Proof of Theorem~\ref{alg:product-test}]
Use \(c=1/\mu\), \(\rho_s\) from \eqref{alg:state-radii}, and the
rational map parameters of Lemma~\ref{alg:map}.
Given \(0<\tau<1\), choose the least \(p\ge0\) for which
\[
 4nq^{p+1}\le\tau(1-q).
\]
Rational multiplication and comparison find \(p\);
\eqref{alg:contraction-parameter} gives
\(p=O_{d,J,\mu}(\log(2+n/\tau))\).
Use Lemma~\ref{alg:connected-tests}, the triangular transform,
and \eqref{alg:log-tail} to compute a Gaussian rational \(S_B\)
with \(|S_B-F_B(c)|\le\tau\).  Also
\(|S_B|\le Kn+1\).  The preceding lemma and the rational
transform give polynomial bit cost.

We make explicit that exponentiation does not introduce an
unbounded-precision oracle.  Let \(T_0\) be an integer upper
bound on \(Kn+1\), and choose an integer \(\ell\) with
\(2^{-\ell}\le\tau\).
Truncate the Taylor series of \(e^{S_B}\) at degree
\[
 M=\left\lceil10(2T_0+\ell+1)\right\rceil
\]
by taking \(\sum_{j=0}^{M}S_B^j/j!\).  The remainder is at most
\(e^{T_0}T_0^{M+1}/(M+1)!\).
Using \(r!\ge(r/3)^r\), it is less than
\(\tau e^{-T_0}\).
Since \(|e^{S_B}|\ge e^{-T_0}\), this is relative error
less than \(\tau\).  All terms are Gaussian rational,
\(M=O_{d,J,\mu}(n+\log\tau^{-1})\), and the numerators and
denominators have polynomial bit length.
Furthermore
\(|e^{S_B}/G_B(c)-1|\le e^\tau-1\).
Choosing \(\tau=\epsilon/32\) supplies the required margin.
Repeat for \(B=I\), and divide the two nonzero rational
estimates to obtain the estimate of \(G_B/G_I\).
For errors at most \(\epsilon/8\) in each factor, the ratio
error is at most
\((\epsilon/4)/(1-\epsilon/8)<\epsilon\).

The nonvanishing assertions follow from
Lemma~\ref{alg:test-domain}; \(I\) is a stable product test
because its local kernel is \(1+zw\).
For a simple eigenvalue the spectral projection is
\(vw^{\mathsf T}/(w^{\mathsf T}v)\), so this quotient is
indeed \(\Tr(B\Pi)\).
\end{proof}

\subsection{Product-state queries and equatorial measurement sampling}
\label{alg:measurements-subsection}

For \(s\in\C\) with \(|s|\le1\), define
\[
 |s\rangle=\frac{|0\rangle+s|1\rangle}{\sqrt{1+|s|^2}} .
\]
These are the pure states in the closed northern hemisphere
of the Bloch sphere in our fixed basis.  The amplitude and
probability queries extend to this whole hemisphere.  Complete
two-outcome projective measurement sampling is stated on the
equator, where the two states \(|s\rangle,|-s\rangle\) are
orthogonal and both satisfy the condition.

\begin{theorem}[Classical access to Suzuki--Fisher ground states]
\label{alg:sf-state}
Fix \(d,J,\mu\) as in Theorem~\ref{alg:sf}.
Let \(H\) be a rational Suzuki--Fisher Hamiltonian on \(n\ge1\)
qubits satisfying that theorem's hypotheses, and let
\(0<\epsilon<1\) be rational.
Its unique unit ground vector \(\psi\) can be normalized by
\(\langle0^n|\psi\rangle>0\).  The query algorithms in (1)--(2)
have bit complexity
\(\operatorname{poly}_{d,J,\mu}(n,\epsilon^{-1},L_{\mathrm{in}})\).
\begin{enumerate}
\item Given \(s_1,\ldots,s_n\in\mathbb Q(i)\) with
      \(|s_i|\le1\), approximate the nonzero complex amplitude
      \(\langle s_1\cdots s_n|\psi\rangle\) with relative
      complex error at most \(\epsilon\).
\item Given \(S\subseteq[n]\) and \(s_i\in\mathbb Q(i)\),
      \(|s_i|\le1\) for \(i\in S\), approximate the positive
      probability
      \[
      \left\langle\psi\left|
        \bigotimes_{i\in S}|s_i\rangle\langle s_i|
                \otimes I_{S^c}\right|\psi\right\rangle
      \]
      by a positive rational with relative error at most
      \(\epsilon\).
\item Given a polynomial-time adaptive strategy that measures
      each site at most once, in a basis
      \(\{|s\rangle,|-s\rangle\}\) with
      \(s\in\mathbb Q(i)\), \(|s|=1\), sample its outcome
      record by a classical randomized algorithm with total
      variation error at most \(\epsilon\).
\end{enumerate}
The strategy in (3) has a finite description, takes previous
outcomes as input, and outputs the next unmeasured site and a
phase.  More precisely, if its total evaluation time along
any record is at most \(T_{\rm strategy}\), and all generated
phases have bit length at most \(L_{\rm phase}\), the sampler
uses at most
\[
 T_{\rm strategy}+
 \operatorname{poly}_{d,J,\mu}
       (n,\epsilon^{-1},L_{\rm in}+L_{\rm phase})
\]
bit operations.  In particular the sampler is polynomial-time
for each fixed polynomial-time strategy.
\end{theorem}

\begin{proof}
Use \(A_0=-(H+\mu\sum_iZ_i)\) as in
\eqref{alg:sf-shift}.  Its local norm bound is a fixed rational
\(J'\) depending on the Pauli bound \(J\) and on \(\mu\);
all radii in the preceding lemmas are evaluated with this
bound, not with an unadjusted \(J\).
At positive real \(x\), \(H(x)\) is Hermitian, so its
vacuum-normalized transpose eigenvector is \(w=\bar v\).
Consequently \(G_I=\|v\|^2>0\) and
\(\psi=v/\sqrt{G_I}\).  Along the positive real path the
normalized logarithm \(F_I\) is the real logarithm, because
it starts at zero and \(G_I\) stays positive.

For amplitudes take
\[
 B_i^{\mathrm{amp}}=
 \begin{pmatrix}1&\bar s_i\\0&0\end{pmatrix}.
\]
Its kernel \(1+\bar s_iw\) is disk-stable, and
\[
 \langle s_1\cdots s_n|\psi\rangle
 =\left(\prod_i(1+|s_i|^2)^{-1/2}\right)
       \exp(F_{B^{\mathrm{amp}}}(c)-F_I(c)/2).
\]
For the event in (2), take \(B_i=I\) off \(S\) and
\[
 B_i^{\mathrm{prob}}=
 \begin{pmatrix}1&\bar s_i\\s_i&|s_i|^2\end{pmatrix}
 =(1+|s_i|^2)|s_i\rangle\langle s_i|
 \quad(i\in S).
\]
Its kernel factors as \((1+s_iz)(1+\bar s_iw)\).
Thus
\[
 P_S=\left(\prod_{i\in S}(1+|s_i|^2)^{-1}\right)
                    \exp(F_B(c)-F_I(c)).
\]
Nonvanishing follows from Lemma~\ref{alg:test-domain};
positivity of the projector expectation then gives \(P_S>0\).

Compute the logarithms with a fixed fraction of the requested
error, using the constructive proof of
Theorem~\ref{alg:product-test}.  For (2), take the real part
of the logarithmic difference before exponentiating; its error
does not increase because the exact difference is real on
the positive path.  For either logarithmic difference use
\(T_0=\lceil2Kn+2\rceil\) in that finite Taylor procedure.
In case (2), the resulting rational estimate is real and
positive when its relative error is less than one.
The probability normalization is rational and exact.
For (1), the product \(R_s=\prod_i(1+|s_i|^2)\) is a positive
rational between \(1\) and \(2^n\).  Rational bisection for
\(\sqrt{R_s}\), followed by inversion, computes its reciprocal
square root to any required relative error with
\(O(n+\log\epsilon^{-1}+L_{\mathrm{in}})\) bits and
polynomial arithmetic cost.  Splitting the error budget
between the two logarithms, exponentiation and this factor
proves (1) and (2), including the phase assertion.

For (3), fix an already generated history.  Its two extensions
have true probabilities \(p_+,p_->0\).  Part (2) computes
positive estimates \(\widehat p_\pm\) with relative error at
most \(\eta=\epsilon/(16n)\).
Normalizing their sum gives a conditional probability with
error at most \(2\eta/(1-\eta)<4\eta\), regardless of the
probability of the parent history.  Round it to a dyadic
probability within \(\epsilon/(4n)\) and draw the corresponding
Bernoulli variable using
\(O(\log(n/\epsilon))\) unbiased random bits.
On a fixed history, projectors on distinct sites commute, so
its probability is exactly the product event in (2), even
though the bases were chosen adaptively.
Coupling the two processes at each step bounds their total
variation distance by the sum of the at most \(n\) conditional
errors, which is less than \(\epsilon\).
At most \(2n\) queries are made.  Their accuracy and phase bit
lengths give the displayed polynomial overhead; evaluating the
strategy contributes \(T_{\rm strategy}\).
\end{proof}

The novelty of this conclusion is its uniform local input and its
fixed-parameter polynomial complexity on the full positive-field
Suzuki--Fisher family.  For general Lee--Yang states with a fixed
spatial radius \(r>1\), \cite[Section~4, Theorems~5--6]{WBG26}
already provides state-dependent quasipolynomial circuits for
relative \(X\)-basis amplitudes, with marginal probabilities and phases in
the state-preparation proof; see also \cite{WongThesis26}.
Our proof does not assume a size-independent spatial radius
margin.  It uses the common reciprocal-field domain instead.
Connected coefficient computation and zero-free interpolation
are established methods
\cite{PatelRegts17,YYZ22,WildAlhambra23}; the new input to those
methods is Lemma~\ref{alg:test-domain} for principal-state tests.
In particular \cite[Theorem~15]{YYZ22} already treats tensorized
measurements of thermal partition functions.  Its thermal
zero-free hypothesis does not itself supply the ground-state
parameter domain used here.

The queries above do not include arbitrary product bases:
the projector onto \(|1\rangle\) has no nonzero vacuum entry
and is not an allowed normalized test.
Nor does the fixed-field theorem extend the zero-field
energy-value scheme to state approximation or sampling.
The broad quantum algorithms in \cite{BGLW26,RT26} have
different parameter ranges and outputs; the present classical
result does not subsume them.

% END SECTION principal_states

% BEGIN SECTION equilibrium
\section{All-temperature Lee--Yang Hamiltonians}\label{eq:section}

We now specialize the complex generator classification to Hermitian
operators on qubits.  The resulting inequalities are precisely the
Suzuki--Fisher conditions, including their full transverse operator-norm
form.  Their sufficiency is classical; the generator theorem supplies
necessity and excludes additional interactions of arbitrary support.
Positive longitudinal fields then improve the zero-free domain of the
actual Gibbs tensor, at every positive temperature.

\subsection{From matrix tensors to polynomial generators}

Fix the computational basis \(\{|a\rangle:a\in\{0,1\}^n\}\),
with \(Z|0\rangle=|0\rangle\).  Associate to a matrix \(M\) the
polynomial
\begin{equation}\label{eq:matrix-tensor}
 F_M(z,w)=\sum_{a,b\in\{0,1\}^n}M_{ab}z^aw^b.
\end{equation}
We say that \(M\) has the tensor Lee--Yang property at radius \(r\)
if \(F_M\) is nonzero on \(\D_r^n\times\D_r^n\).  The two
groups of variables in this definition are independent complex variables.
In particular, \eqref{eq:matrix-tensor} is not the sesquilinear form
\(\langle z|M|z\rangle\).

For a Hermitian Hamiltonian \(H\), the Gibbs operator is
\(M_\beta=e^{-\beta H}\), and its normalized density matrix is
\(\rho_\beta=M_\beta/\Tr M_\beta\).  Since
\(\Tr e^{-\beta H}>0\) for real \(\beta\), the two matrices have
exactly the same polynomial zeros.

The following bridge fixes both the tensor and operator conventions.
It is the binary instance, with a unitary scalar normalization, of the
Cayley correspondence \eqref{pre:cayley}.

\begin{lemma}\label{eq:cayley-bridge}
Let
\[
 C_1=\frac1{\sqrt2}\begin{pmatrix}i&-i\\1&1\end{pmatrix},
 \qquad C=C_1^{\otimes n},\qquad D_1=\diag(-1,1).
\]
For an invertible matrix \(M\), its tensor has the unit-disk
Lee--Yang property if and only if \(CMC^{-1}\), acting on
multiaffine coefficient vectors, preserves upper-half-plane stability.
Consequently, the Gibbs family corresponds to the generator
\begin{equation}\label{eq:generator}
 A=-CHC^*.
\end{equation}
Here operator conjugation is a similarity, whereas transforming both
groups of tensor variables gives the congruence \(CMC^T\).
\end{lemma}
\begin{proof}
The coefficient action of \(C\) is
\[
 (Cf)(s)=2^{-n/2}\prod_i(s_i+i)
 f\!\left(\frac{s_1-i}{s_1+i},\ldots,
           \frac{s_n-i}{s_n+i}\right).
\]
The fractions map the upper half-plane onto the unit disk, and the
prefactors are nonzero there.  Thus \(C\) identifies the two
stability classes.  By the binary form of the contraction theorem,
or Lemma~\ref{spec:kernel-criterion}, the disk-stable tensor of an
invertible \(M\) is equivalent to preservation of disk-stable
coefficient vectors.  This proves the equivalence by similarity.

For completeness, the reverse tensor conversion can also be checked
directly.  The unitary matrix \(C_1\) satisfies \(C_1C_1^T=D_1\),
so the coefficient matrix of
\[
 \prod_i(s_it_i-1)
\]
is \(D_1^{\otimes n}=CC^T\).  Each factor is stable in the two
upper-half-plane variables: if \(s_it_i=1\), then
\(t_i=1/s_i\) has negative imaginary part.  Acting on its first
group of variables by \(N=CMC^{-1}\) gives coefficient matrix
\(NCC^T=CMC^T\).  For fixed \(t\in\HH^n\), invertibility
excludes a zero output, so stability preservation gives nonvanishing
jointly in \((s,t)\in\HH^{2n}\).  Finally,
\[
 F_{CMC^T}(s,t)
 =2^{-n}\prod_i(s_i+i)(t_i+i)
 F_M\!\left(\frac{s-i}{s+i},\frac{t-i}{t+i}\right).
\]
This verifies the congruence and the sign of the half-plane kernel.
As \(C^{-1}=C^*\), exponential similarity yields
\(Ce^{-\beta H}C^{-1}=e^{\beta A}\).
\end{proof}

\subsection{The exact Hamiltonian classification}

Write \(X,Y,Z\) for the Pauli matrices, with
\(Y=\left(\begin{smallmatrix}0&-i\\ i&0\end{smallmatrix}\right)\).
For a real \(2\times2\) matrix, \(\|\cdot\|_{\mathrm{op}}\)
denotes its largest singular value.

\begin{theorem}[All-temperature qubit classification]\label{eq:sf}
For a Hermitian operator \(H\) on \((\C^2)^{\otimes n}\), the
following conditions are equivalent:
\begin{enumerate}
\item \(e^{-\beta H}\) has the tensor Lee--Yang property at radius
one for every \(\beta\ge0\).
\item The same property holds for every \(\beta\) in some interval
\([0,\varepsilon)\), where \(\varepsilon>0\).
\item There are real coefficients such that
\begin{align}
 H={}&h_0 I+\sum_i(h_i^xX_i+h_i^yY_i+h_i^zZ_i)\notag\\
 &+\sum_{i<j}\bigl(a_{ij}X_iX_j+b_{ij}X_iY_j
           +c_{ij}Y_iX_j+d_{ij}Y_iY_j+e_{ij}Z_iZ_j\bigr),
 \label{eq:sf-form}
\end{align}
where
\begin{equation}\label{eq:sf-inequality}
 h_i^z\le0,\qquad
 e_{ij}\le-
 \left\|\begin{pmatrix}a_{ij}&b_{ij}\\c_{ij}&d_{ij}\end{pmatrix}
 \right\|_{\mathrm{op}}.
\end{equation}
There is no restriction on \(h_0,h_i^x,h_i^y\).
\end{enumerate}
Equivalently, in the ferromagnetic convention
\begin{equation}\label{eq:ferromagnetic-convention}
 H=h_0I-\sum_i\sum_{a=x,y,z}\mu_i^a\sigma_i^a
 -\sum_{i<j}\left(J_{ij}^{zz}Z_iZ_j+
       \sum_{a,b=x,y}J_{ij}^{ab}\sigma_i^a\sigma_j^b\right),
\end{equation}
the conditions are
\(\mu_i^z\ge0\) and
\(J_{ij}^{zz}\ge\|J_{ij}^{\perp}\|_{\mathrm{op}}\), where the
rows and columns of \(J_{ij}^{\perp}\) are ordered \(x,y\).
\end{theorem}

\begin{proof}
The first condition implies the second.  By
Lemma~\ref{eq:cayley-bridge}, the second says that \(e^{tA}\)
preserves upper-half-plane stability for small \(t\ge0\).
The semigroup identity extends preservation to all \(t\ge0\):
for any \(t\), choose an integer \(m\) with \(t/m<\varepsilon\)
and compose \(e^{(t/m)A}\) \(m\) times.  We may therefore apply
Theorem~\ref{gen:main-complex} on the binary capacity box.

We first spell out how differential order becomes physical support.
On one binary site a basis of endomorphisms is
\[
 I,\qquad J^- =\partial_s,\qquad
 J^0=s\partial_s-\tfrac12,\qquad
 J^+=s^2\partial_s-s.
\]
The last three operators are traceless; their first-order coefficients
are the linearly independent polynomials \(1,s,s^2\).
Expand \(A\) in tensor products of this basis and choose a support
set \(S\) of maximal cardinality occurring with a nonzero coefficient.
The canonical coefficient of \(\partial_S=\prod_{i\in S}\partial_i\)
is then the sum of its exact-support terms, because no strictly larger
support occurs.  Those terms have distinct principal monomials
\(\prod_{i\in S}s_i^{r_i}\), with \(r_i\in\{0,1,2\}\).
The order-two conclusion of Theorem~\ref{gen:main-complex} forces
this coefficient to vanish if \(|S|\ge3\).  Linear independence
then eliminates every term of that support.  Descending in support
eliminates all interactions on three or more sites.  Local similarity
preserves the identity and traceless subspaces at each site, so the
same support conclusion holds for \(H\).

Next write \(A=R+iS\) with entrywise real matrices \(R,S\).
The imaginary-part conclusion of Theorem~\ref{gen:main-complex}
has the form
\begin{equation}\label{eq:imaginary-generator}
 S=c_0I+\sum_iJ_i(q_i),\qquad
 q_i(s)=\alpha_i+\eta_i s+\gamma_i s^2\ge0\quad(s\in\R),
\end{equation}
where \(J(q)=q\partial-q'/2\).  On coefficients \((1,s)\),
\[
 J(q)=\begin{pmatrix}-\eta/2&\alpha\\-\gamma&\eta/2\end{pmatrix}.
\]
Since \(A\) is Hermitian, \(R^T=R\) and \(S^T=-S\).
Taking the trace in \eqref{eq:imaginary-generator} first gives
\(c_0=0\).  Independence of the remaining one-site traceless
components shows that each \(J(q_i)\) is skew-symmetric.  Thus
\(\eta_i=0\) and \(\gamma_i=\alpha_i\); nonnegativity gives
\(\alpha_i\ge0\).  In particular,
\begin{equation}\label{eq:imaginary-local}
 q_i(s)=\alpha_i(1+s^2),\qquad
 iS=-\sum_i\alpha_iY_i.
\end{equation}

The local Pauli conjugations are
\begin{equation}\label{eq:pauli-cayley}
 C_1XC_1^*=-Z,\qquad C_1YC_1^*=X,\qquad C_1ZC_1^*=-Y.
\end{equation}
Therefore the physical field
\(h_i^xX_i+h_i^yY_i+h_i^zZ_i\) contributes
\(h_i^xZ_i-h_i^yX_i+h_i^zY_i\) to \(A\).
Comparison with \eqref{eq:imaginary-local} gives
\(h_i^z=-\alpha_i\le0\).  A two-site physical Pauli term
containing exactly one \(Z\) becomes a term containing exactly one
\(Y\), and is entrywise imaginary.  The four such terms on each
pair are linearly independent.  Since \(S\) has no two-site
component, all physical \(XZ,YZ,ZX,ZY\) coefficients vanish.
This leaves exactly the form \eqref{eq:sf-form}.

It remains to determine the pair inequality.  Suppress the pair indices.
A physical term \(aXX+bXY+cYX+dYY+eZZ\) contributes
\[
 -aZZ+bZX+cXZ-dXX-eYY
\]
to \(A\).  The first-order coefficients of the coefficient operators
\(X,Z,Y\), respectively, are
\[
 p_X(s)=1-s^2,\qquad p_Z(s)=-2s,\qquad
 p_Y(s)=-i(1+s^2).
\]
Hence the full coefficient of \(\partial_i\partial_j\) is
\begin{align}
 q_{ij}(s,t)={}&e(1+s^2)(1+t^2)-4ast
       -2bs(1-t^2)\notag\\
       &-2ct(1-s^2)-d(1-s^2)(1-t^2).
 \label{eq:mixed-coefficient}
\end{align}
There is no factor of two in this expression: it is the full mixed
coefficient in Theorem~\ref{gen:main-complex}.
Define
\[
 u(s)=\left(\frac{2s}{1+s^2},\frac{1-s^2}{1+s^2}\right),
 \qquad J=\begin{pmatrix}a&b\\c&d\end{pmatrix}.
\]
Then
\begin{equation}\label{eq:opnorm-identity}
 \frac{q_{ij}(s,t)}{(1+s^2)(1+t^2)}=e-u(s)^TJu(t).
\end{equation}
The set \(u(\R)\) is dense in the unit circle.  Consequently
the required inequality \(q_{ij}(s,t)\le0\) for all real \(s,t\)
is equivalent to
\[
 e\le\min_{\|u\|_2=\|v\|_2=1}u^TJv=-\|J\|_{\mathrm{op}}.
\]
This proves necessity of \eqref{eq:sf-inequality}.

Conversely, suppose \eqref{eq:sf-form}--\eqref{eq:sf-inequality}
hold.  The displayed conjugations show that \(A\) has canonical
differential order at most two.  Its mixed coefficients are real and
nonpositive by \eqref{eq:opnorm-identity}, and its imaginary part is
\eqref{eq:imaginary-generator} with
\(q_i(s)=-h_i^z(1+s^2)\ge0\).  Its real part therefore satisfies
the real criterion in Theorem~\ref{gen:main-complex}, and its
imaginary part satisfies that theorem's one-site condition.  The
theorem proves preservation for all \(\beta\ge0\); the bridge
then gives the tensor Lee--Yang property of \(e^{-\beta H}\).
Finally, changing the signs of the coefficients converts
\eqref{eq:sf-inequality} into \eqref{eq:ferromagnetic-convention}.
\end{proof}

\begin{remark}\label{eq:sf-attribution}
The sufficient Suzuki--Fisher family originates in
\cite{SuzukiFisher71,AsanoJapanese70,AsanoPRL70}. The full transverse
operator-norm condition is already contained, in equivalent algebraic
form, in \cite[equation~(2.47)]{SuzukiFisher71} and
\cite[Theorem~2, equation~(8)]{Dunlop79}; see also the modern notation
in \cite[Definition~1 and Fact~3]{BGLW26}. Indeed, for a real matrix
$J=\left(\begin{smallmatrix}a&b\\c&d\end{smallmatrix}\right)$,
\[
 2\|J\|_{\rm op}
 =\sqrt{(a-d)^2+(b+c)^2}+\sqrt{(a+d)^2+(b-c)^2}.
\]
To verify the identity, the squared radicals are
$\|J\|_F^2-2\det J$ and $\|J\|_F^2+2\det J$.
If the singular values are $s_1\geq s_2\geq0$, the two radicals
are therefore $s_1+s_2$ and $s_1-s_2$, in some order.
Full matrix generating functions and their
composition are also explicit in Asano's contemporaneous account for
the XXZ subfamily \cite[equations~(1.15), (3.7)--(3.10)]{AsanoJapanese70}.
Thus a distinction between classical traces and modern full tensors
alone would not identify the contribution. Theorem~\ref{eq:sf}
adds the converse among all finite Hermitian qubit Hamiltonians,
including the exclusion of higher-support interactions.  The proof
also shows that testing only a sufficiently small interval of inverse
temperatures is enough.
\end{remark}

\begin{corollary}\label{eq:complex-support}
Let \(H\) be an arbitrary complex qubit matrix.  If
\(e^{-\beta H}\) has the unit-disk tensor Lee--Yang property for
all sufficiently small real \(\beta\ge0\), then its Pauli
expansion contains no term supported on three or more sites.
\end{corollary}
\begin{proof}
The Cayley bridge and the semigroup identity still apply.  The complex
generator theorem bounds the canonical order by two.  The maximal-support
basis argument in the proof of Theorem~\ref{eq:sf} does not use
Hermiticity and therefore gives the assertion.
\end{proof}

\begin{remark}\label{eq:trace-caveat}
For a complex Hamiltonian, tensor stability does not justify trace
normalization.  For example, \(H=\diag(0,i\pi)\) has Gibbs
polynomial \(1+e^{-i\pi\beta}zw\), which is nonzero on the open
unit bidisk for every real \(\beta\), but
\(\Tr e^{-H}=0\).  The support conclusion in
Corollary~\ref{eq:complex-support} applies to this complex setting;
the normalized Gibbs-state statements use Hermiticity.
\end{remark}

\subsection{Strict positive fields and the actual Gibbs kernel}

\begin{theorem}[Strict-field Gibbs radius]\label{eq:strict-field}
Let \(H\) have the form \eqref{eq:ferromagnetic-convention}, with
\(J_{ij}^{zz}\ge\|J_{ij}^{\perp}\|_{\mathrm{op}}\) and
\(\mu_i^z>0\) at every site.  Then, for every \(\beta>0\),
\(F_{e^{-\beta H}}\) is nonzero on the closed unit polydisk in
its full \(2n\) variables.  In particular, it has some radius
\(r_\beta>1\).

The ground eigenspace of \(H\) is one-dimensional.  Its nonzero
coefficient vector \(\psi\) defines a polynomial
\(f_\psi(z)=\sum_a\psi_a z^a\) of radius greater than one.
For every \(\beta_0>0\), one radius \(r>1\) works for all
Gibbs kernels at \(\beta\ge\beta_0\), and also for the image
under each \(e^{-\beta H}\) of every nonzero unit-disk-stable
coefficient vector.
\end{theorem}
\begin{proof}
Let \(\mu=\min_i\mu_i^z>0\), and write
\(H=H_0-\mu\sum_iZ_i\).  The Hamiltonian \(H_0\) satisfies
the nonstrict Suzuki--Fisher conditions.  By Theorem~\ref{eq:sf}
and the disk version of Lemma~\ref{eq:cayley-bridge}, its
coefficient operator \(A_0=-H_0\) generates a disk-stability
preserver.  In the multiaffine coefficient basis,
\(Z_i=I-2E_i\).  Thus
\[
 -H=A_0-2\mu E+\mu nI.
\]
The harmless scalar factor \(e^{\beta\mu n}\) reduces every
kernel and evolution assertion to Theorem~\ref{spec:strict}.

Proposition~\ref{spec:perron} supplies an algebraically simple
dominant eigenvalue of \(-H\).  Since \(H\) is Hermitian,
this is exactly the negative of its ground energy, so the ground
space is one-dimensional.  The same proposition gives the strict
radius of its right eigenvector.  Equivalently, one may fix a
Gibbs kernel radius \(r>1\) and normalize its powers by their
largest eigenvalue.  These powers retain radius \(r\) by
contraction, and converge to \(|\psi\rangle\langle\psi|\).
Their nonzero limiting polynomial is
\[
 f_\psi(z)\overline{f_\psi(\bar w)}.
\]
Hurwitz gives nonvanishing on \(\D_r^{2n}\), and therefore
nonvanishing of each factor on \(\D_r^n\).
\end{proof}

\begin{remark}\label{eq:gap-attribution}
The spectral conclusion
\[
 E_1(H)-E_0(H)\ge2\min_i\mu_i^z
\]
is the Hermitian specialization of Theorem~\ref{spec:gap} and
was already proved for this family in
\cite[Theorem~1]{BGLW26}.  Here the additional conclusion is
strict nonvanishing of the actual Gibbs tensor at every positive
inverse temperature, and its transfer to the ground-state polynomial.
No numerical lower bound for \(r_\beta-1\), uniform in system
size, is asserted.  This distinction also separates the actual
kernel from the strict kernels used in a product-formula approximation.
\end{remark}

\subsection{A qualitative radius for deformed EPR ground states}

Consider a finite graph without isolated vertices, with positive edge
weights \(w_e\).  On an edge \(e=\{i,j\}\), fix
\(0\le s_e<1\) and \(\theta_e\in\R\), and define
\[
 |\eta_e\rangle=|00\rangle+s_e e^{i\theta_e}|11\rangle,
 \qquad H_{\mathrm{EPR}}=-\sum_e w_e
 |\eta_e\rangle\langle\eta_e|_{ij}.
\]
The edge vectors here are unnormalized; replacing them by normalized
vectors simply rescales the positive weights.

\begin{corollary}\label{eq:epr}
The ground space of \(H_{\mathrm{EPR}}\) is one-dimensional, and
its ground-state polynomial is zero-free on a polydisk of some radius
\(r>1\).
\end{corollary}
\begin{proof}
On one edge, direct expansion in Pauli matrices gives
\begin{align*}
 |\eta_e\rangle\langle\eta_e|
 ={}&\frac{1+s_e^2}{4}(I+Z_iZ_j)
    +\frac{1-s_e^2}{4}(Z_i+Z_j)\\
 &+\frac{s_e}{2}\bigl[
      \cos\theta_e(X_iX_j-Y_iY_j)
      +\sin\theta_e(X_iY_j+Y_iX_j)\bigr].
\end{align*}
Its ferromagnetic coefficients are therefore
\[
 J_e^{zz}=\frac{w_e(1+s_e^2)}4,\qquad
 J_e^\perp=\frac{w_es_e}{2}
 \begin{pmatrix}\cos\theta_e&\sin\theta_e\\
                 \sin\theta_e&-\cos\theta_e\end{pmatrix}.
\]
The displayed real matrix is orthogonal, so
\(\|J_e^\perp\|_{\mathrm{op}}=w_es_e/2\) and
\(J_e^{zz}-\|J_e^\perp\|_{\mathrm{op}}
 =w_e(1-s_e)^2/4\ge0\).  The longitudinal field at site \(i\) is
\[
 \mu_i^z=\sum_{e\ni i}\frac{w_e(1-s_e^2)}4>0.
\]
All hypotheses of Theorem~\ref{eq:strict-field} follow, including
strictness because every vertex has an incident edge.  Apply that
theorem to obtain the stated uniqueness and radius.
\end{proof}

For a common deformation parameter \(s<1\), the corollary gives
the qualitative strict-radius conclusion asked for in
\cite[Conjecture~1]{WBG26}.  It does not assert a prescribed sharp
radius as a function of \(s\).  Edge phases may be chosen
independently; no simultaneous gauge removal is used in the proof.

% END SECTION equilibrium

% BEGIN SECTION quantum_dynamics
\section{Physical quantum Markov dynamics}\label{dyn:section}

The zero-freeness of a Gibbs operator and the preservation of zero-freeness by
a trace-preserving physical evolution impose different restrictions.  We now
classify the latter.  The physical tensor decomposition and the computational
basis are fixed throughout this section.

For a matrix $M$ on $n$ qubits, put
\begin{equation}\label{dyn:polynomial}
 F_M(z,w)=\sum_{a,b\in\{0,1\}^n}M_{ab}z^a w^b.
\end{equation}
A density matrix is called \emph{Lee--Yang} if this polynomial does not vanish
on $\D^{2n}$.  Thus the row and column variables are independent.  For a pure
state, $F_{|\psi\rangle\langle\psi|}(z,w)
=f_\psi(z)\overline{f_\psi(\bar w)}$, where
$f_\psi(z)=\sum_a\psi_a z^a$.  A channel is \emph{ordinarily preserving} if it
preserves these density matrices on its given system, and is \emph{completely
preserving} if it does so after adjoining any number of inert qubits.  Complete
positivity of a channel and complete preservation in this sense are distinct
conditions.

Write $\sigma_x=X$, $\sigma_y=Y$, and $\sigma_z=Z$, with
$Z=\diag(1,-1)$.  The Pauli matrices in all generator formulas below are
unnormalized.  We use the standard finite-dimensional GKLS representation
\cite{GKS76,Lindblad76}.  If $\mathcal P_n$ denotes the nonidentity Hermitian
Pauli strings, every generator of a continuous CPTP semigroup has a
representation
\begin{equation}\label{dyn:gkls}
 \mathcal L(\rho)=-i[H,\rho]
 +\sum_{P,Q\in\mathcal P_n}K_{PQ}
 \left(P\rho Q-\tfrac12\{QP,\rho\}\right),
 \qquad H=H^*,\quad K\succeq0.
\end{equation}
Identity components in the noise have been absorbed into the Hamiltonian.

\begin{theorem}[Ordinary and complete preservation]\label{dyn:main}
Let $n\ge2$, and let $\Phi_t=e^{t\mathcal L}$ be a continuous CPTP semigroup
on $n$ qubits.  The following conditions are equivalent.
\begin{enumerate}
\item $\Phi_t$ ordinarily preserves Lee--Yang density matrices for every $t\ge0$.
\item The same preservation holds for all $t$ in some interval $[0,\varepsilon)$.
\item $\Phi_t$ completely preserves Lee--Yang density matrices for every $t\ge0$.
\item The full Choi tensor of $\Phi_t$ is Lee--Yang for every $t\ge0$.
\item The generator is a sum of one-qubit generators,
\begin{equation}\label{dyn:local-form}
 \mathcal L=\sum_{j=1}^n\mathcal L_j\otimes\Id_{\widehat j},\qquad
 \mathcal L_j(\rho)=-i[h_jZ,\rho]
 +\sum_{\alpha,\beta=x,y,z}(K_j)_{\alpha\beta}
 \left(\sigma_\alpha\rho\sigma_\beta-
 \tfrac12\{\sigma_\beta\sigma_\alpha,\rho\}\right),
\end{equation}
where $h_j\in\R$ and, for a real symmetric $2\times2$ matrix $B_j$, a vector
$d_j\in\R^2$, and $c_j\in\R$,
\begin{equation}\label{dyn:two-lmis}
 K_j=\begin{pmatrix}B_j&i d_j\\-i d_j^{\mathsf T}&c_j\end{pmatrix}\succeq0,
 \qquad
 \begin{pmatrix}
 c_jI_2-\operatorname{adj}B_j&d_j\\
 d_j^{\mathsf T}&\Tr B_j
 \end{pmatrix}\succeq0.
\end{equation}
\end{enumerate}
In this case $\Phi_t=\bigotimes_j e^{t\mathcal L_j}$.  The equivalence between
ordinary and complete preservation is false on one qubit.
\end{theorem}

We prove the theorem in four steps.  A fixed finite family of product inputs
first excludes every interaction.  Two-site pure states then fix the reality
of the remaining local coefficients.  Finally, positive definite separable
states recover the individual local inequalities even though both sites
evolve.  The resulting inequalities are reduced to
\eqref{dyn:two-lmis} without dividing by a possibly vanishing trace.

\subsection{A finite family detecting every interaction}

\begin{lemma}[A one-sided root collision]\label{dyn:collision}
Let $g_t$ be a differentiable curve in a finite-dimensional polynomial space,
with $g_0\not\equiv0$.  Suppose $0$ is a root of $g_0$ of multiplicity $m$,
and $g_t$ is upper-half-plane stable along a sequence $t\downarrow0$ of
positive times.  If $m\ge3$, then $\dot g_0(0)=0$.  If $m=2$ and
$g_0(x)=c x^2+O(x^3)$, then $\dot g_0(0)/c$ is real and nonpositive.
\end{lemma}

\begin{proof}
Choose a small circle about zero containing no other root of $g_0$.
The nearby cluster has a monic factor
$b_t(x)=x^m+a_1(t)x^{m-1}+\cdots+a_m(t)$ with $b_0=x^m$.
Its coefficients are differentiable: the contour integrals of
$x^k g_t'(x)/g_t(x)$ give the cluster power sums, and Newton identities give
its elementary symmetric functions.  This construction also applies when
the ambient degree changes.  The complementary factor is nonzero at zero.

At each of the specified stable times, write the cluster roots as
$\lambda_j=x_j+i y_j$, where $y_j\le0$.  Differentiability of the first
power sum gives $\sum_j|y_j|=O(t)$.  The second power sum is also $O(t)$,
so its real part gives $\sum_jx_j^2=O(t)$, since
$\sum_j y_j^2=O(t^2)$.  Thus every $\lambda_j=O(\sqrt t)$ and
$\sum_j\lambda_j^k=o(t)$ for $k\ge3$.  Newton identities imply
\[
 a_k(t)=o(t)\quad(k\ge3),\qquad
 a_2(t)=-\tfrac12\sum_jx_j^2+o(t).
\]
For the second equality, the imaginary part of the second power sum is
$O(t^{3/2})$ and $a_1(t)^2=O(t^2)$.  Consequently $a_2'(0)$ is real and
nonpositive, while $a_k'(0)=0$ for $k\ge3$.  Multiplication by the
complementary factor proves the assertions.  All derivatives are fixed by
the differentiable curve, so estimates along the given sequence suffice.
\end{proof}

\begin{theorem}[Finite equatorial tests]\label{dyn:trine}
Let $\Phi_t=e^{t\mathcal L}$ be a continuous CPTP semigroup on $n$ qubits.
Set $\mathcal E=\{0,\sqrt3,-\sqrt3\}$ and
\[
 \psi_a=\frac{(1+ia)|0\rangle+(1-ia)|1\rangle}
 {\sqrt{2(1+a^2)}},\qquad
 \rho_{\boldsymbol a}=\bigotimes_{j=1}^n
 |\psi_{a_j}\rangle\langle\psi_{a_j}|,\quad
 \boldsymbol a\in\mathcal E^n.
\]
Suppose that, for each of these $3^n$ densities, there is a sequence of
positive times tending to zero along which its image is Lee--Yang.
Then $\mathcal L$ is a sum of one-qubit GKLS generators.
The sequences may depend on the input.
\end{theorem}

\begin{proof}
The coefficient Cayley transform is
\begin{equation}\label{dyn:cayley}
 C=\frac1{\sqrt2}\begin{pmatrix}i&-i\\1&1\end{pmatrix},\qquad
 f(z)\longmapsto\frac{s+i}{\sqrt2}
 f\left(\frac{s-i}{s+i}\right).
\end{equation}
It takes disk stability to upper-half-plane stability.  Applied to all row
and column legs of a density tensor, it is a coefficient transformation,
not a unitary conjugation of that density matrix.  Direct multiplication
gives $C\psi_a=(-a,1)^{\mathsf T}/\sqrt{1+a^2}$ and
$C\overline{\psi_a}=(a,1)^{\mathsf T}/\sqrt{1+a^2}$.  The transformed
input is therefore a nonzero scalar multiple of
\begin{equation}\label{dyn:paired-product}
 p_{\boldsymbol a}(z,w)=\prod_j(z_j-a_j)(w_j+a_j).
\end{equation}

For a binary coefficient matrix $M=(M_{ab})$, its polynomial action has
the first-order form
\[
 M_{00}+M_{10}s+
 [M_{01}+(M_{11}-M_{00})s-M_{10}s^2]\partial_s.
\]
The identities $CXC^*=-Z$, $CYC^*=X$, $CZC^*=-Y$ consequently give
the derivative coefficients
\begin{equation}\label{dyn:pauli-vectors}
 v(s)=(2s,1-s^2,i(1+s^2))
\end{equation}
for left multiplication.  Right multiplication uses transposed Paulis;
its derivative vector at a general column coordinate $t$ is
$\widetilde v(t)=(2t,-(1-t^2),i(1+t^2))$.  In particular
$\widetilde v(-a)=-\overline{v(a)}$.

Use \eqref{dyn:gkls}.  A zero diagonal entry of a positive semidefinite
matrix has a zero row and column.  Suppose that $k\ge2$ is the greatest
weight of a Pauli string with $K_{PP}>0$, and choose such a string with
support $S$, $|S|=k$.  All noise rows of larger weight vanish.  Put both
coordinates of each site in $S$ at their roots in
\eqref{dyn:paired-product}, keeping every other coordinate away from its
root, and restrict to a line with every coordinate slope positive.
The initial polynomial has a root of multiplicity $2k\ge4$.
Lemma~\ref{dyn:collision} makes its first variation zero at this point.

Hamiltonian and anticommutator terms act only on row legs or only on column
legs, so leave selected roots untouched.  A sandwich $P\rho Q$ can survive
only if both supports contain $S$.  Maximality makes both supports exactly
$S$.  Dividing out the nonzero complementary product and the input scalar
therefore gives
\begin{equation}\label{dyn:psd-contact}
 0=(-1)^k\sum_{\supp P=\supp Q=S}
 K_{PQ}V_P(\boldsymbol a)\overline{V_Q(\boldsymbol a)},
 \qquad V_P(\boldsymbol a)=\prod_{j\in S}v_{P_j}(a_j).
\end{equation}
The sum, without its sign, is a positive semidefinite quadratic form.
Thus the exact-support block $K_{SS}$ kills
$\overline{V(\boldsymbol a)}$ for every $\boldsymbol a\in\mathcal E^S$.
The three components in \eqref{dyn:pauli-vectors} form a basis of the
quadratic polynomials.  Evaluation at the three distinct points of
$\mathcal E$ is invertible.  The resulting tensor products span
$\C^{3^k}$, forcing $K_{SS}=0$, a contradiction.  Every surviving noise
index consequently has weight one.  This still permits cross-site noise
blocks; we eliminate them below.

The anticommutator operator $\sum K_{PQ}QP$ now has support at most two.
If a Hamiltonian term has maximal support $S$ with $|S|=k\ge3$, put just
the row coordinates of $S$ at their roots.  All column coordinates and
remaining row coordinates stay away from their roots.  The $k$-fold
contact again has zero first variation.  No noise term, anticommutator,
or column Hamiltonian removes all these row roots.  The surviving row
Hamiltonian equals a nonzero common factor times
\[
 -i\sum_{\supp P=S}h_P\prod_{j\in S}v_{P_j}(a_j).
\]
The same finite tensor evaluation basis forces every displayed $h_P$ to
vanish.  Descending in the maximal support leaves only Hamiltonian terms
of weight at most two.  The transformed superoperator now acts on at most
two individual row or column legs at a time.

Fix distinct physical sites $i,j$.  A row--row double contact with coordinates
$s,t\in\mathcal E$ uses input parameters $a_i=s,a_j=t$.  A row--column
contact uses $a_i=s,a_j=-t$, which is available because $\mathcal E=-\mathcal E$.
Keep all other coordinates off their roots.  Only the corresponding mixed
principal coefficient survives.  Dividing by the initial quadratic
coefficient divides that coefficient by the positive product of the two
selected slopes.  Lemma~\ref{dyn:collision} therefore makes both mixed
coefficients real and nonpositive on this nine-point grid.

Let $K^{ij}$ be the $3\times3$ cross-site noise block, and write the
Hamiltonian pair as $\sum_{\alpha,\beta}h_{\alpha\beta}
\sigma_{i,\alpha}\sigma_{j,\beta}$ with $h$ real.  The row--column
coefficient matrix is $K^{ij}$.  Since distinct-site Paulis commute,
the two conjugate noise blocks in the anticommutator combine to give
the row--row coefficient matrix
\begin{equation}\label{dyn:cross-matrices}
 -\Re K^{ij}-ih.
\end{equation}
Here the real part is taken entrywise; there is no transpose in this formula.

Put
\[
 u(s)=\frac{(2s,1-s^2)}{1+s^2},\qquad R=\diag(1,-1).
\]
The normalized row and column derivative vectors are $(u(s),i)$ and
$(Ru(t),i)$.  Uniform averaging over $\mathcal E$ gives
\begin{equation}\label{dyn:trine-moments}
 \mathbb E u=0,\qquad \mathbb E uu^{\mathsf T}=\tfrac12I_2.
\end{equation}
Initially write $K^{ij}=\left(\begin{smallmatrix}A&B\\D&e\end{smallmatrix}\right)$
with complex blocks.  The imaginary part of its normalized row--column
coefficient is the imaginary part of
$u^{\mathsf T}ARv+i u^{\mathsf T}B+iDRv-e$.
Its nine values vanish.  Averaging these values, and then their products
with $u$, $Rv$, and $u(Rv)^{\mathsf T}$, proves
$\Im e=0$, $\Re B=0$, $\Re D=0$, and $\Im A=0$.
Hence $B=ib$, $D=id^{\mathsf T}$ with $A,b,d,e$ real, and
\begin{equation}\label{dyn:lr-inequality}
 u^{\mathsf T}ARv-b^{\mathsf T}u-d^{\mathsf T}Rv-e\le0.
\end{equation}
The imaginary part obtained from \eqref{dyn:cross-matrices} is
$-u^{\mathsf T}h_{TT}v+h_{zz}$, where $T=\{x,y\}$.
The same moments give $h_{TT}=0=h_{zz}$.  Write
$k=h_{Tz}$ and $\ell^{\mathsf T}=h_{zT}$.  The real row--row inequality is
\begin{equation}\label{dyn:ll-inequality}
 e-u^{\mathsf T}Av+k^{\mathsf T}u+\ell^{\mathsf T}v\le0.
\end{equation}
Both inequalities have already been divided by the positive factor
$(1+s^2)(1+t^2)$.  Their unweighted averages are respectively $-e$ and $e$,
so $e=0$.  Each of the nine nonpositive values in each inequality has
zero average, and thus each value is zero.  Applying
\eqref{dyn:trine-moments} once more gives $A=b=d=k=\ell=0$.
Consequently $K^{ij}=0$ and the pair Hamiltonian vanishes.

Every remaining block $K^{ii}$ is positive semidefinite.  The remaining
terms of \eqref{dyn:gkls} are therefore genuine one-qubit GKLS generators.
They commute on distinct tensor factors.  The proof compared only
derivatives at zero, so the different stable time sequences present no
compatibility issue.
\end{proof}

\begin{corollary}\label{dyn:immediate-failure}
If a CPTP generator is not a sum of one-qubit generators, one of the fixed
inputs in Theorem~\ref{dyn:trine} has a non-Lee--Yang output at every
sufficiently small positive time.
\end{corollary}

\begin{proof}
Otherwise each input would have Lee--Yang output at positive times
arbitrarily close to zero, providing the sequences in that theorem.
\end{proof}

The theorem makes no assertion that its number of tests is minimal or that
testing the output polynomial is computationally efficient.  Nor can a
single isolated positive time replace the sequences: the interaction
$H=Z_1Z_2$ generates the identity channel at $t=\pi$.

\subsection{The local reality forced by a second physical site}

We first record the elementary one-qubit state geometry.

\begin{lemma}\label{dyn:hemisphere}
A qubit density matrix $\rho=(I+r\cdot\sigma)/2$ is Lee--Yang if and only
if $r_z\ge0$.
\end{lemma}

\begin{proof}
Write $\rho=\left(\begin{smallmatrix}a&b\\\bar b&d\end{smallmatrix}\right)$,
so $a+d=1$, $|b|^2\le ad$, and
$F_\rho(z,w)=g(w)+z h(w)$ with $g=a+bw$, $h=\bar b+dw$.
On $|w|=1$,
\begin{equation}\label{dyn:bilinear-boundary}
 |g(w)|^2-|h(w)|^2
 =(a-d)[a+d+2\Re(bw)].
\end{equation}
If $a>d$, then $a>|b|$, so $g$ is zero-free on the closed disk.
The right side is nonnegative.  The maximum principle applied to $h/g$
proves $|h/g|\le1$ on the disk and hence the desired zero-freeness.
If $a=d$, add a positive multiple of $|0\rangle\langle0|$ and take its
nonzero coefficient limit; Hurwitz's theorem gives the assertion.
If $a<d$, positivity gives $d>|b|$, and
\eqref{dyn:bilinear-boundary} is strictly negative on the circle.
Thus a nearby interior point has $|h|>|g|$ and $h\ne0$;
choosing $z=-g/h$ gives an interior zero.  This proves both directions.
\end{proof}

\begin{proposition}\label{dyn:local-reality}
Suppose $n\ge2$ and a CPTP semigroup ordinarily preserves Lee--Yang
densities on a neighborhood of time zero.  In the local representation
given by Theorem~\ref{dyn:trine}, each noise matrix and Hamiltonian have
the form
\begin{equation}\label{dyn:real-local-shape}
 K=\begin{pmatrix}B&i d\\-i d^{\mathsf T}&c\end{pmatrix},
 \quad B=B^{\mathsf T}\in\R^{2\times2},\quad d\in\R^2,
 \qquad H=hZ.
\end{equation}
Writing $T=\Tr B$ and $J=\left(\begin{smallmatrix}0&-1\\1&0\end{smallmatrix}\right)$,
the local Bloch equation is
\begin{equation}\label{dyn:bloch}
 \dot r_\perp=M r_\perp+k,\qquad \dot r_z=-2T r_z,
 \quad M=2[B-(T+c)I_2]+2hJ,
 \quad k=4Jd.
\end{equation}
In particular the finite-time map has the form
\begin{equation}\label{dyn:affine}
 (r_\perp,r_z)\longmapsto(D_t r_\perp+\tau_t,\lambda_t r_z),
 \quad D_t\in GL_2(\R),\quad \lambda_t=e^{-2Tt}>0.
\end{equation}
\end{proposition}

\begin{proof}
After locality is known, let $q_j(s,w)$ denote the coefficient of
$\partial_s\partial_w$ in the Cayley-transformed generator at site $j$.
For a general Hermitian noise matrix write
\[
 K=\begin{pmatrix}
 \alpha&\delta+i\eta&r+i\mu\\
 \delta-i\eta&\beta&t+i\nu\\
 r-i\mu&t-i\nu&\gamma
 \end{pmatrix}.
\]
Multiplying the row and column derivative vectors in
\eqref{dyn:pauli-vectors} gives
\begin{equation}\label{dyn:imaginary-q}
 \Im q(s,w)=2(s+w)[-\eta(1-sw)+r(1+sw)+t(w-s)].
\end{equation}
The Hamiltonian and anticommutator act on single legs and do not enter
this coefficient.

Fix two sites.  For real $a,b$ and $c_0>0$, the polynomial
$f(z_1,z_2)=(z_1-a)(z_2-b)-c_0$ is upper-half-plane stable: solving for
$z_2$ when $z_1\in\HH$ gives a point in the lower half-plane.
Its inverse Cayley transform defines a pure physical Lee--Yang input.
The transformed conjugate factor is
$g(w_1,w_2)=(w_1+a)(w_2+b)-c_0$, since
$\overline C=\diag(-1,1)C$ and there are two sites.
For arbitrary nonzero real $u,v$, choose the real contact
\[
 z_1=a+u,\quad z_2=b+c_0/u,\qquad
 w_1=-a+v,\quad w_2=-b+c_0/v.
\]
Here $f=g=0$, $\nabla f=(c_0/u,u)$ and
$\nabla g=(c_0/v,v)$.  Their directional derivatives along any strictly
positive direction are nonzero.  The initial product has a double root
with nonzero real leading coefficient.  Lower-order generator terms
vanish at the contact, and its first variation is
\begin{equation}\label{dyn:hyperbola-contact}
 h=\frac{c_0^2}{uv}q_1(a+u,-a+v)
   +uv\,q_2(b+c_0/u,-b+c_0/v).
\end{equation}
Lemma~\ref{dyn:collision} makes $h$ real.  Additional sites may be padded
by equatorial factors kept away from their roots; their generators
contribute zero at this contact.

Formula \eqref{dyn:imaginary-q} gives $\Im q_2(b,-b)=0$.
The coefficient of $c_0$ in the imaginary part of
\eqref{dyn:hyperbola-contact} is therefore
\[
 v\,\partial_s\Im q_2(b,-b)+u\,\partial_w\Im q_2(b,-b)=0.
\]
Independent variation of $u,v$ makes both derivatives vanish.
By \eqref{dyn:imaginary-q}, this says
$-\eta(1+b^2)+r(1-b^2)-2tb=0$ for every real $b$.
Its three coefficients give $\eta=r=t=0$.  Interchanging the sites proves
the same conclusion at the other site, and every site has a partner.
This proves the stated shape of $K$.

At an equatorial Bloch vector the dissipator with this shape has zero
longitudinal velocity.  A Hamiltonian $h_xX+h_yY+h_zZ$ contributes
$\dot r_z=2(h_xr_y-h_yr_x)$.  Product inputs and
Lemma~\ref{dyn:hemisphere} require this quantity to be nonnegative at
every equatorial point.  Replacing that point by its negative gives
$h_x=h_y=0$.  Finally, multiplication of the Pauli matrices in the
generator gives \eqref{dyn:bloch}; in particular
$\mathcal L(I)=-4d_yX+4d_xY$, fixing the displacement sign.
The linear differential equation gives $D_t=e^{tM}$ and
$\lambda_t=e^{-2Tt}$, as claimed.
\end{proof}

\subsection{Separable conditional blocks prevent compensation}

\begin{lemma}[Conditional-block cancellation]\label{dyn:cancellation}
Let a qubit channel $\Psi$ have the affine form
$(r_\perp,r_z)\mapsto(D r_\perp+\tau,\lambda r_z)$, with $\lambda>0$.
If $\alpha,\beta\in\C$ satisfy $|\beta|>|\alpha|$, the polynomial of
$\alpha\Psi(|0\rangle\langle0|)+\beta\Psi(|1\rangle\langle1|)$ has a zero
in $\D^2$.
\end{lemma}

\begin{proof}
The scalar $\alpha+\beta$ is nonzero.  Dividing the matrix by
$(\alpha+\beta)/2$ gives $I+\tau\cdot\sigma_\perp+qZ$, where
\[
 q=\lambda\frac{\alpha-\beta}{\alpha+\beta},\qquad
 \Re q=\lambda\frac{|\alpha|^2-|\beta|^2}{|\alpha+\beta|^2}<0.
\]
Its polynomial is
$(1+q)+(\tau_x+i\tau_y)z+(\tau_x-i\tau_y)w+(1-q)zw$.
Choose a phase $\theta$ making
$(\tau_x+i\tau_y)e^{i\theta}=|\tau|$ and substitute
$z=e^{i\theta}u$, $w=e^{-i\theta}u$; if $\tau=0$ any phase works.
The quadratic has nonzero leading coefficient $1-q$ and root product
$(1+q)/(1-q)$ of modulus less than one.  At least one root has modulus
less than one, giving the asserted pair of disk coordinates.
\end{proof}

The lemma applies to the actual evolving channel at a second site.
It does not require that this channel fix the reference state.  In
particular, if positive matrices $A,B$ give a Lee--Yang input proportional
to $A\otimes|0\rangle\langle0|+B\otimes|1\rangle\langle1|$, ordinary
preservation by $\Phi\otimes\Psi$ forces
\begin{equation}\label{dyn:dominance-necessary}
 |F_{\Phi(B)}(z,w)|\le |F_{\Phi(A)}(z,w)|\qquad (z,w\in\D).
\end{equation}
Indeed a strict reverse inequality supplies $\alpha,\beta$ for
Lemma~\ref{dyn:cancellation}.  This also covers $\alpha=0$.

\begin{proposition}[Recovery of each local inequality]\label{dyn:local-inequality}
Under the assumptions of Proposition~\ref{dyn:local-reality}, each local
matrix in \eqref{dyn:real-local-shape} satisfies
\begin{equation}\label{dyn:circle-criterion}
 r^{\mathsf T}Bv+d\cdot(v-r)\le c
 \qquad (r,v\in\R^2,\ |r|=|v|=1).
\end{equation}
Necessity follows from positive definite separable inputs classical on
one of two selected sites.
\end{proposition}

\begin{proof}
Fix an active site and a distinct reference site.  For $0<\kappa<1/2$
and $e\in\R^2$ with $|e|=1$, put
\[
 U=I+\tfrac12Z,\qquad
 V=\tfrac18I+\tfrac{\sqrt{\kappa(1-\kappa)}}4\sigma_e
        +\tfrac\kappa4Z,\qquad \sigma_e=e_xX+e_yY,
\]
and define
\begin{equation}\label{dyn:separable-family}
 A=\tfrac12(U+V),\quad B_0=\tfrac12(U-V),\qquad
 \rho_{\kappa,e}=\tfrac12
 [A\otimes|0\rangle\langle0|+B_0\otimes|1\rangle\langle1|].
\end{equation}
Both blocks are positive definite, because
$\lambda_{\min}(U)=1/2$ and
$\|V\|=1/8+\sqrt\kappa/4<1/2$.
Their traces sum to $2$, so this is a normalized positive definite
separable density matrix.

For a Hermitian matrix $Q=q_0I+q_\perp\cdot\sigma_\perp+q_zZ$, parameterize
the distinguished boundary by
$z=e^{i(\gamma+\delta)}$, $w=e^{i(\gamma-\delta)}$, and set
$x=\cos\gamma$, $y=\sin\gamma$, $n=(\cos\delta,-\sin\delta)$.
Expansion of \eqref{dyn:polynomial} gives
\begin{equation}\label{dyn:torus-evaluation}
 \tfrac12e^{-i\gamma}F_Q(z,w)=q_0x+q_\perp\cdot n-iq_zy.
\end{equation}
Write $u=1/2$, $v_0=1/8$,
$w_0=2v_0\sqrt{\kappa(1-\kappa)}$ and
$\zeta=\kappa v_0/u$, so $U=I+uZ$ and
$V=v_0I+w_0\sigma_e+\zeta Z$.  Formula
\eqref{dyn:torus-evaluation} yields
\begin{align}\label{dyn:input-dominance}
 \mathscr D_0(x,n)
 &:=\tfrac14(|F_A|^2-|F_{B_0}|^2)\notag\\
 &=v_0[(1-\kappa)x^2
 +2\sqrt{\kappa(1-\kappa)}x(e\cdot n)+\kappa]\notag\\
 &\ge v_0(\sqrt{1-\kappa}|x|-\sqrt\kappa)^2\ge0.
\end{align}
The matrix $A$ is positive definite with strictly positive $Z$ coefficient.
Consequently its diagonal entries obey $a>d>0$, and
\eqref{dyn:bilinear-boundary} is strictly positive on the circle.
The proof of Lemma~\ref{dyn:hemisphere} then gives zero-freeness of $F_A$
on the closed bidisk.  Applying the maximum principle separately in
$z,w$ to $F_{B_0}/F_A$ extends \eqref{dyn:input-dominance} to that bidisk.
Hence $F_A+z_2w_2F_{B_0}$ is nonzero on $\D^4$, proving that
\eqref{dyn:separable-family} is Lee--Yang.  Other sites may be padded by
maximally mixed states.

For the active map write $D_t,\tau_t,\lambda_t$ as in
\eqref{dyn:affine}.  At time zero their derivatives are $M,k,-2T$ in
\eqref{dyn:bloch}.  The reference map has $\lambda_t>0$ at every finite
time.  By \eqref{dyn:dominance-necessary}, ordinary preservation forces
the active output block difference to be nonnegative on the disk and,
by continuity, on its boundary.  Its exact normalized expression is
\begin{equation}\label{dyn:output-dominance}
 \mathscr D_t(x,n)=
 (x+\tau_t\cdot n)
 [v_0(x+\tau_t\cdot n)+w_0n^{\mathsf T}D_te]
 +\lambda_t^2u\zeta(1-x^2).
\end{equation}
Choose $n=-e$ and $x=\rho:=\sqrt{\kappa/(1-\kappa)}\in(0,1)$.
The input difference is zero.  Its derivative must therefore be
nonnegative.  Differentiating \eqref{dyn:output-dominance} gives
\begin{align}\label{dyn:derivative-dominance}
 \frac{\mathscr D'_0}{2v_0\kappa}
 &=\rho k\cdot n-n^{\mathsf T}Mn
      -2T\frac{1-2\kappa}{1-\kappa}\notag\\
 &=2(c-n^{\mathsf T}Bn)+2T\rho^2+\rho k\cdot n\ge0.
\end{align}
For explicit verification, the three derivative contributions before
division are $(k\cdot n)(2v_0\rho-w_0)$,
$-\rho w_0n^{\mathsf T}Mn$, and
$-4Tu\zeta(1-\rho^2)$.  The Hamiltonian term drops out because
$n^{\mathsf T}Jn=0$.

Set $n=Jm$, $|m|=1$.  Since $k=4Jd$ and
$(Jm)^{\mathsf T}B(Jm)=T-m^{\mathsf T}Bm$, we obtain
\begin{equation}\label{dyn:radial-inequality}
 c-T+m^{\mathsf T}Bm+T\rho^2+2\rho\,m\cdot d\ge0
 \qquad (|m|=1,\ 0\le\rho\le1).
\end{equation}
The endpoints follow by continuity.
For arbitrary unit $r,v$, set $(v-r)/2=\rho m$.
Their mean is perpendicular to $m$ and has squared length $1-\rho^2$;
when $\rho=0$, choose either perpendicular direction for $m$.
In two dimensions this gives
\[
 r^{\mathsf T}Bv+d\cdot(v-r)
 =T-m^{\mathsf T}Bm-T\rho^2+2\rho\,m\cdot d.
\]
Apply \eqref{dyn:radial-inequality} with $-m$ to prove
\eqref{dyn:circle-criterion}.
\end{proof}

\subsection{The exact local semidefinite cone}

\begin{lemma}[Reduction of the circle inequality]\label{dyn:lmi-reduction}
Suppose $B=B^{\mathsf T}\in\R^{2\times2}$ and
$\left(\begin{smallmatrix}B&i d\\-i d^{\mathsf T}&c\end{smallmatrix}\right)\succeq0$.
Then \eqref{dyn:circle-criterion} is equivalent to
\begin{equation}\label{dyn:second-lmi}
 \begin{pmatrix}cI_2-\operatorname{adj}B&d\\d^{\mathsf T}&\Tr B\end{pmatrix}
 \succeq0.
\end{equation}
If $T=\Tr B>0$, the equivalent scalar condition is
$c\ge\lambda_{\max}(\operatorname{adj}B+dd^{\mathsf T}/T)$.
If $T=0$, positivity forces $B=d=0$, and every $c\ge0$ is allowed.
\end{lemma}

\begin{proof}
Unitary congruence, followed if necessary by changing the sign of the last
coordinate, gives
$\left(\begin{smallmatrix}B&d\\d^{\mathsf T}&c\end{smallmatrix}\right)\succeq0$.
In particular $B\succeq0$ and $dd^{\mathsf T}\preceq cB$.
The assertion for $T=0$ follows immediately.  Suppose $T>0$.
Writing $(v-r)/2=\rho m$ as above, the expression to maximize is
\[
 T-m^{\mathsf T}Bm-T\rho^2+2\rho\,d\cdot m,
 \qquad |m|=1,\quad0\le\rho\le1.
\]
If $c<T$, then $|d|^2\le c\lambda_{\max}(B)<T^2$.
Choose the sign of $m$ so $d\cdot m\ge0$.
The unconstrained maximum in $\rho$ is attained at
$\rho=|d\cdot m|/T<1$.  The maximum is consequently
\[
 \lambda_{\max}\left(TI_2-B+\frac{dd^{\mathsf T}}T\right)
 =\lambda_{\max}\left(\operatorname{adj}B+\frac{dd^{\mathsf T}}T\right).
\]
This proves equivalence in this case.

If $c\ge T$, put $b_m=m^{\mathsf T}Bm$ and use
$|d\cdot m|\le\sqrt{cb_m}$.  The difference between $c$ and the expression
above is bounded below by
\[
 c-T+b_m+T\rho^2-2\sqrt{cb_m}\rho
 =(\sqrt{b_m}-\sqrt c\rho)^2+(c-T)(1-\rho^2)\ge0.
\]
The scalar matrix condition also holds automatically, because
\[
 dd^{\mathsf T}\preceq cB\preceq T[B+(c-T)I_2],
\]
the second inequality being equivalent to
$(c-T)(TI_2-B)\succeq0$.
Finally, the Schur complement with positive last entry $T$ proves
equivalence of the scalar condition and \eqref{dyn:second-lmi}.
The separate $c\ge T$ argument avoids an invalid unconstrained radial
maximization when $|d|>T$.
\end{proof}

\begin{lemma}[Choi tensors and complete preservation]\label{dyn:choi}
For a CPTP map $\Phi$ on $n$ qubits, the following are equivalent:
its full Choi tensor is Lee--Yang; it completely preserves physical
Lee--Yang states; and its extension by $n$ inert qubits takes the maximally
entangled state to a Lee--Yang state.
For any invertible complex-linear map, without a positivity or trace
assumption, its full Choi tensor is Lee--Yang if and only if its
coefficient action preserves all complex disk-stable polynomials in the
$2n$ row and column variables.
\end{lemma}

\begin{proof}
With $|\Omega\rangle=\sum_a|a\rangle|a\rangle$, the Choi entries are
$J_{ai,bj}=[\Phi(|i\rangle\langle j|)]_{ab}$.
The normalized maximally entangled input is Lee--Yang, since its polynomial
is proportional to $\prod_k(1+z_kz'_k)(1+w_kw'_k)$.
Its output is the normalized Choi matrix.
Conversely, for a joint input $\rho$, the output coefficients are
\[
 [ (\Phi\otimes\Id)(\rho)]_{ar,bs}
 =\sum_{i,j}J_{ai,bj}\rho_{ir,js}.
\]
This is bilinear tensor contraction with no conjugation of entries.
Lemma~\ref{pre:grace} gives a Lee--Yang polynomial or zero;
trace preservation excludes zero.
This proves the physical equivalences.

Let $M$ be the superoperator matrix in coefficient coordinates.  Its full
tensor differs from the Choi tensor only by a permutation of indices.
Contraction with an arbitrary stable coefficient vector gives stability
or zero, and invertibility excludes zero.  In the other direction, if
$M$ preserves all stable polynomials, apply it for each fixed
$w\in\D^{2n}$ to $\prod_j(1+z_jw_j)$.
The resulting joint polynomial is exactly the full tensor of $M$ and
is nonzero for every $(z,w)\in\D^{4n}$.
\end{proof}

\begin{proposition}[One-qubit complete generators]\label{dyn:local-complete}
A one-qubit CPTP semigroup is completely preserving if and only if its
generator has the form \eqref{dyn:real-local-shape} and satisfies the two
positive semidefinite conditions in \eqref{dyn:two-lmis}.
\end{proposition}

\begin{proof}
For necessity, extend the semigroup by one inert qubit.
Propositions~\ref{dyn:local-reality} and \ref{dyn:local-inequality} give the
real local shape and \eqref{dyn:circle-criterion} for the active site.
Lemma~\ref{dyn:lmi-reduction} gives the two matrix conditions.

For sufficiency, transform both coefficient legs by \eqref{dyn:cayley}.
Write $B=\left(\begin{smallmatrix}a&u\\u&b\end{smallmatrix}\right)$ and
$d=(x,y)^{\mathsf T}$.  The transformed generator is real.  This follows
directly from the Pauli identities: the mixed terms pair the purely
imaginary transverse--longitudinal noise entries with the purely imaginary
$Z$ derivative; the anticommutator is
$-(a+b+c)I+y(Z_1+Z_2)+x(X_1-X_2)$; and the longitudinal Hamiltonian is real
after the transform.  Its full coefficient of $\partial_s\partial_t$ is
\begin{align}\label{dyn:local-q}
 q(s,t)={}&4ast-b(1-s^2)(1-t^2)-c(1+s^2)(1+t^2)\notag\\
 &+2u[t(1-s^2)-s(1-t^2)]\notag\\
 &+2x[t(1+s^2)-s(1+t^2)]
 -y[(1-s^2)(1+t^2)+(1+s^2)(1-t^2)].
\end{align}
There is no extra factor of two in this convention.  With
\[
 r(s)=\frac{(2s,1-s^2)}{1+s^2},\qquad
 v(t)=\frac{(2t,t^2-1)}{1+t^2},
\]
we have
\begin{equation}\label{dyn:q-circle}
 \frac{q(s,t)}{(1+s^2)(1+t^2)}
 =r(s)^{\mathsf T}Bv(t)+d\cdot(v(t)-r(s))-c\le0.
\end{equation}
The real two-variable specialization of the finite-box generator theorem
(Theorem~\ref{gen:main-complex}) says that a real endomorphism of the binary
polynomial box generates upper-half-plane stability preservers precisely
when this mixed coefficient is nonpositive on $\R^2$.
It applies here because the transformed operator is an endomorphism of
that box.  Thus its exponential preserves complex stability.
Undoing the Cayley similarity and applying Lemma~\ref{dyn:choi}, with
invertibility of the exponential, proves complete preservation.
\end{proof}

\begin{proof}[Proof of Theorem~\ref{dyn:main}]
Assume preservation on a neighborhood of zero.  Theorem~\ref{dyn:trine}
gives locality.  Propositions~\ref{dyn:local-reality},
\ref{dyn:local-inequality}, and Lemma~\ref{dyn:lmi-reduction} give
\eqref{dyn:local-form}--\eqref{dyn:two-lmis}.  Conversely those conditions
give complete preservation of every local semigroup by
Proposition~\ref{dyn:local-complete}.  Their tensor product completely
preserves by successive application to the individual sites.
Lemma~\ref{dyn:choi} equates this with the full Choi condition.
Complete preservation implies ordinary preservation, which implies its
small-time version.  The one-qubit exception is proved next.
\end{proof}

\subsection{The Hermiticity-preserving branch without positivity}

For the full Choi condition, locality does not require complete positivity.
The local semidefinite reduction does require it, and must be replaced by
the circle inequality when positivity is absent.

\begin{theorem}[Hermiticity- and trace-preserving semigroups]\label{dyn:hp}
Let $T_t=e^{t\mathcal L}$ be a continuous semigroup of complex-linear,
Hermiticity-preserving, trace-preserving maps on $n$ qubits.  Its full
Choi tensor is Lee--Yang for every $t\ge0$ if and only if
$\mathcal L=\sum_j\mathcal L_j\otimes\Id_{\widehat j}$, where each
$\mathcal L_j$ has the algebraic form \eqref{dyn:local-form}, with
\[
 H_j=h_jZ,\qquad
 K_j=\begin{pmatrix}B_j&i d_j\\-i d_j^{\mathsf T}&c_j\end{pmatrix},
 \quad B_j=B_j^{\mathsf T}\in\R^{2\times2},\quad
 d_j\in\R^2,\quad c_j,h_j\in\R,
\]
and satisfies
\begin{equation}\label{dyn:hp-circle}
 c_j\ge\max_{r,v\in S^1}
 [r^{\mathsf T}B_jv+d_j\cdot(v-r)].
\end{equation}
Here no positive semidefiniteness of $K_j$ is imposed.  Requiring the
Choi property only on a neighborhood of zero is equivalent.  If complete
positivity is also imposed, the local conditions reduce exactly to
\eqref{dyn:two-lmis}.
\end{theorem}

\begin{proof}
We first justify the algebraic representation used in the proof.
A Hermiticity-preserving, trace-annihilating generator has a representation
\eqref{dyn:gkls} with $H=H^*$ and $K=K^*$, without the condition
$K\succeq0$.  Indeed, its Choi matrix is Hermitian.  In the orthogonal
vectorized Pauli basis, its compression to the orthogonal complement
of $|I\rangle\!\rangle$ uniquely specifies the Hermitian coefficients
$K_{PQ}$ of the sandwich terms.  Subtract the corresponding dissipator,
which is Hermiticity-preserving and trace-annihilating.  The remaining
Hermitian Choi matrix has zero compression to that complement, hence
is of the form
$|B\rangle\!\rangle\langle\!\langle I|
 +|I\rangle\!\rangle\langle\!\langle B|$ for some matrix $B$.
It represents $\rho\mapsto B\rho+\rho B^*$.
Trace annihilation gives $B+B^*=0$, so $B=-iH$ with $H$ Hermitian.
This proves the claimed representation.

The last assertion of Lemma~\ref{dyn:choi} and the Cayley transform show
that the full Choi hypothesis places the coefficient generator in
Theorem~\ref{gen:main-complex} on $2n$ binary legs.
It consequently has differential order at most two.
This implies tensor support at most two on those individual legs:
expand in the local endomorphism basis
$I,\partial_s,s\partial_s-1/2,s^2\partial_s-s$.
At a nonzero support of maximal size, its full derivative coefficient
is a linear combination of independent products of $1,s,s^2$.
If the support size exceeds two, the order bound forces all its
coefficients to vanish.  Descending eliminates every larger support.
Local Cayley similarity preserves the identity and traceless
subspaces, so this conclusion also holds before the transform.

For nonidentity Pauli strings $P,Q$, the tensor coefficient supported
on their respective row and column legs is exactly the sandwich
coefficient $K_{PQ}$, up to the harmless sign of a transposed $Y$.
Neither a Hamiltonian nor an anticommutator term can cancel a
coefficient involving both row and column legs.  Therefore
\[
 \operatorname{wt}(P)+\operatorname{wt}(Q)>2
 \quad\Longrightarrow\quad K_{PQ}=0.
\]
In particular every row indexed by a string of weight at least two
vanishes, directly and without a PSD argument.  The anticommutator
now has support at most two, so the same support bound excludes every
Hamiltonian term of higher weight.

For distinct physical sites, the cross-site matrices and their normalized
mixed coefficients are precisely \eqref{dyn:cross-matrices}--
\eqref{dyn:ll-inequality}.  The general generator theorem makes these
coefficients real and nonpositive at every real pair.  Restricting to
the trine grid, the discrete moment calculation
\eqref{dyn:trine-moments} eliminates every cross-site block and pair
Hamiltonian exactly as in Theorem~\ref{dyn:trine}.  That calculation
uses Hermiticity and trace annihilation, but never PSD of $K$.
Thus the generator is a sum of local Hermiticity-preserving,
trace-annihilating generators.

At one site, the nine products of the two Pauli derivative triples are
independent.  Reality of their mixed coefficient forces $K_{xy}$ to
be real and $K_{xz},K_{yz}$ to be purely imaginary, giving the displayed
local shape.  Its dissipator is then real after the Cayley transform.
For a Hamiltonian $h_xX+h_yY+h_zZ$, the remaining imaginary first-order
coefficients are
\[
 h_y s^2-2h_xs-h_y,\qquad h_y t^2+2h_xt-h_y.
\]
The complex generator criterion requires each quadratic to be
nonnegative on the real line.  Comparing leading and constant terms
forces $h_y=0$, and then the linear term forces $h_x=0$.
The real mixed coefficient is \eqref{dyn:local-q}, whose normalization
\eqref{dyn:q-circle} is nonpositive exactly under \eqref{dyn:hp-circle}.
This proves necessity without any positivity assumption.

Conversely, the displayed local algebraic generators are
Hermiticity-preserving and trace-annihilating.  Their Cayley transforms
are real binary-box endomorphisms with nonpositive mixed coefficient.
Theorem~\ref{gen:main-complex} and Lemma~\ref{dyn:choi} prove full Choi
stability of their exponentials and their tensor product.  Small-time
preservation extends to all times by the semigroup identity and
Lemma~\ref{pre:grace}, with invertibility excluding zero.
Finally, complete positivity imposes $K_j\succeq0$ in the canonical
representation, and Lemma~\ref{dyn:lmi-reduction} gives exactly the
second matrix condition in \eqref{dyn:two-lmis}.
\end{proof}

\subsection{Cancellation at a fixed positive time}

The mechanism used above has a finite-time version.  It applies to different
local channels and does not assume that they belong to a semigroup.

\begin{theorem}[Static tensor cancellation]\label{dyn:static-cancellation}
Let $n\ge2$.  For each $j$, let $\Psi_j$ be an invertible CPTP qubit map
with Bloch action
\[
 (r_\perp,r_z)\longmapsto(D_jr_\perp+\tau_j,\lambda_jr_z),
 \qquad D_j\in GL_2(\R),\quad\lambda_j>0.
\]
Then $\bigotimes_j\Psi_j$ ordinarily preserves all physical Lee--Yang
densities if and only if every $\Psi_j$ has a Lee--Yang Choi tensor.
If one factor fails this condition, a positive definite separable input,
classical on a second selected site and maximally mixed on the other sites,
detects the failure, whatever the parameters of the second factor.
\end{theorem}

\begin{proof}
The outputs of the two computational basis states give
$|\tau|^2+\lambda^2\le1$, so $|\tau|<1$.
The static one-qubit criterion in Theorem~\ref{loc:classification} says that
the Choi condition for the active factor $\Psi$ is equivalent to
\begin{equation}\label{dyn:static-criterion}
 \|D^{\mathsf T}n\|^2
 \le(x-\tau\cdot n)^2+\lambda^2(1-x^2)
 \qquad(n\in S^1,\ -1\le x\le1).
\end{equation}
Changing $x$ to $-x$ gives the alternative sign convention for its
transverse displacement.
Suppose the condition fails.  By strictness and continuity choose
$x\in(-1,1)$ and $n\in S^1$ such that, with $v=D^{\mathsf T}n$,
$V=|v|$ and $y=x-\tau\cdot n$,
\begin{equation}\label{dyn:static-failure}
 V^2>y^2+\lambda^2(1-x^2).
\end{equation}
In particular $V>0$.  Set $u=yv/V^2$, so $|u|<1$ and $u\cdot v=y$.
The matrix $W=I_2-uu^{\mathsf T}$ is positive definite and
$v^{\mathsf T}Wv=V^2-y^2>\lambda^2(1-x^2)$.
Choose $0<\theta<1$ sufficiently close to one and put
\[
 d_0=\theta\frac{Wv}{\sqrt{v^{\mathsf T}Wv}}.
\]
Then $W-d_0d_0^{\mathsf T}\succ0$, as follows by congruence with
$W^{-1/2}$, and
\begin{equation}\label{dyn:static-separation}
 d_0\cdot v>\lambda\sqrt{1-x^2}.
\end{equation}
Choose $\beta>0$ with $\beta(1-|u|)>|d_0|$ and define
\[
 A_R=R(I+u\cdot\sigma_\perp)+Z,
 \qquad B_0=\beta I+(\beta u+d_0)\cdot\sigma_\perp.
\]
The second matrix is positive definite.  The first is positive definite
once $R^2(1-|u|^2)>1$.

We claim that, for all sufficiently large finite $R$,
$F_{A_R}$ is zero-free on the closed bidisk and
\begin{equation}\label{dyn:static-input-domination}
 |F_{B_0}(s,w)|<|F_{A_R}(s,w)|\qquad(|s|,|w|\le1).
\end{equation}
On the distinguished boundary use \eqref{dyn:torus-evaluation} with
$X=\cos\gamma$, $m\in S^1$, and $\ell=X+u\cdot m$.
The squared-modulus difference divided by four is
\[
 G_R(X,m)=(R^2-\beta^2)\ell^2
 -2\beta\ell(d_0\cdot m)+1-X^2-(d_0\cdot m)^2.
\]
On $\ell=0$, its last two terms are at least
$\delta:=\lambda_{\min}(I_2-uu^{\mathsf T}-d_0d_0^{\mathsf T})>0$.
If $|\ell|\le\epsilon$, then
$|X^2-(u\cdot m)^2|\le2\epsilon$.
Choose $\epsilon>0$ so
$2\epsilon+2\beta\epsilon|d_0|<\delta/2$.
For $R\ge\beta$, this gives $G_R>0$ on that set.
On $|\ell|\ge\epsilon$, using $|\ell|\le2$ gives
\[
 G_R\ge(R^2-\beta^2)\epsilon^2-4\beta|d_0|-|d_0|^2>0
\]
for sufficiently large $R$.  This proves strict torus domination.

The diagonal entries of $A_R$ are $R+1,R-1$ and its off-diagonal
magnitude is $R|u|$.  The constant coefficient in the representation
$F_{A_R}=g(w)+s h(w)$ is nonzero on the closed disk because
$R+1>R|u|$.  On the circle,
\[
 |g(w)|^2-|h(w)|^2
 =4R+4\Re((A_R)_{01}w)>0.
\]
The maximum principle gives closed-bidisk zero-freeness of $F_{A_R}$.
Apply it separately in both coordinates to $F_{B_0}/F_{A_R}$ to obtain
\eqref{dyn:static-input-domination}.
Consequently the density proportional to
$A_R\otimes|0\rangle\langle0|+B_0\otimes|1\rangle\langle1|$
is positive definite, separable, and Lee--Yang.

At the output torus point in \eqref{dyn:torus-evaluation} with
$X=-x$, $m=n$ and $Y=\sqrt{1-x^2}$, the affine action gives
\[
 \tfrac12e^{-i\gamma}F_{\Psi(A_R)}
 =R[-x+\tau\cdot n+u\cdot v]-i\lambda Y=-i\lambda Y,
\]
and
\[
 \tfrac12e^{-i\gamma}F_{\Psi(B_0)}
 =\beta[-x+\tau\cdot n+u\cdot v]+d_0\cdot v=d_0\cdot v.
\]
The first value is nonzero, while \eqref{dyn:static-separation} makes
the second strictly larger in modulus.  The strict inequality persists
when the active coordinates move slightly into the disk.
Lemma~\ref{dyn:cancellation} then gives a zero of the two-site output
after any allowed partner channel.  Padding by maximally mixed states
proves necessity for every $n\ge2$.
Conversely, if all factors have Lee--Yang Choi tensors,
Lemma~\ref{dyn:choi} gives complete preservation at each site and therefore
ordinary preservation of their tensor product.
\end{proof}

The theorem assumes the displayed local affine form.  It is not a
factorization theorem for arbitrary isolated channels.  In particular,
the continuous-time locality theorem does not by itself imply such a
static conclusion.

\subsection{The one-qubit exception and dynamical consequences}

\begin{proposition}[The ordinary one-qubit cone]\label{dyn:one-qubit}
Let a one-qubit CPTP semigroup have Bloch equation
$\dot r=Ar+b$, with real $A,b$.  It ordinarily preserves Lee--Yang states
if and only if
\begin{equation}\label{dyn:hemisphere-generator}
 b_z\ge\sqrt{A_{zx}^2+A_{zy}^2}.
\end{equation}
\end{proposition}

\begin{proof}
By Lemma~\ref{dyn:hemisphere}, the additional condition is invariance of
the northern Bloch hemisphere.  At $r_z=0$, its minimum outward-coordinate
velocity is $b_z-\sqrt{A_{zx}^2+A_{zy}^2}$, since the transverse vector
ranges over the closed unit disk.  A negative value gives immediate
failure.  Conversely, complete positivity keeps $|r(t)|\le1$.
Under \eqref{dyn:hemisphere-generator},
$\dot r_z\ge A_{zz}r_z$ at every time.  Multiplication by
$e^{-A_{zz}t}$ and integration give
$r_z(t)\ge e^{A_{zz}t}r_z(0)\ge0$.
\end{proof}

\begin{example}[Amplitude damping]\label{dyn:amplitude-damping}
The amplitude-damping channel toward $|0\rangle$ has Kraus operators
$K_0=\diag(1,\sqrt\eta)$ and
$K_1=\sqrt{1-\eta}|0\rangle\langle1|$, where $0\le\eta\le1$.
The semigroup $\eta=e^{-\gamma t}$ satisfies
$\dot r_z=\gamma(1-r_z)$ and therefore ordinarily preserves one-qubit
Lee--Yang states.  It is not completely preserving for $0<\eta<1$.
Indeed, its action on one half of a Bell state gives, up to a positive
scalar, the polynomial
\[
 (1+\sqrt\eta\,z_1z_2)(1+\sqrt\eta\,w_1w_2)
 +(1-\eta)z_2w_2.
\]
At $z_1=z_2=w_1=w_2=ir$ its value is
$1-(2\sqrt\eta+1-\eta)r^2+\eta r^4$.
It is positive at $r=0$ and equals $2\eta-2\sqrt\eta<0$ at $r=1$.
An interior root $0<r<1$ proves the claim.
\end{example}

\begin{remark}\label{dyn:rates}
For diagonal $K=\diag(a,b,c)$ the complete criterion is
$c\ge\max(a,b)$, recovering the Pauli-semigroup portion of the known
sufficient channel family \cite[Theorem~3]{WBG26}. The full cone also permits
nonunital noise: $\sqrt\gamma(X+iZ)/2$ is an allowed jump and damps
toward an equatorial pure state. Within the completely preserving cone,
only rotations around the fixed $Z$
axis are allowed as coherent motion.  More generally, ordinary
one-qubit dynamics with Hamiltonian
$(\omega_xX+\omega_yY+\omega_zZ)/2$, downward and upward jump rates
$\gamma_\downarrow,\gamma_\uparrow$, and arbitrary $Z$ dephasing
satisfies \eqref{dyn:hemisphere-generator} exactly when
$\gamma_\downarrow-\gamma_\uparrow\ge
\sqrt{\omega_x^2+\omega_y^2}$.
\end{remark}

\begin{corollary}\label{dyn:no-correlated-preparation}
An ordinarily preserving CPTP semigroup on $n\ge2$ qubits cannot take a
product input to a nonproduct state, either at a finite time or in a limit
as $t\to\infty$.  In particular, a semigroup converging from every input
to a fixed nonproduct density matrix cannot ordinarily preserve the full
Lee--Yang class.
\end{corollary}

\begin{proof}
Theorem~\ref{dyn:main} keeps every product input a tensor product.
The set of product densities is closed, being the continuous image of
a compact product of one-site state spaces.
\end{proof}

This conclusion concerns universal preservation and preparation from
product inputs.  It does not exclude a correlated stationary state of a
nonergodic product semigroup, nor preservation along one specially chosen
trajectory.  The identity semigroup illustrates the first distinction.

\subsection{Explicit boundaries of restricted tests}

The following examples distinguish preservation of selected inputs or
classical populations from preservation of the whole physical class.
Their zeros provide direct certificates independent of the classification.

\begin{lemma}[Pair defects of real hyperbola contacts]\label{dyn:pair-defect}
Let $\mathcal L$ be a sum of local GKLS generators of the real form
\eqref{dyn:real-local-shape}, on $n\ge2$ sites, and let $q_j$ be its
local mixed coefficient \eqref{dyn:local-q}.  Define
\[
 \Delta_j=\sup_{s,w\in\R,\ s+w\ne0}
             \frac{q_j(s,w)}{(s+w)^2}.
\]
All the real double-contact inequalities from the two-site pure
hyperbola family \eqref{dyn:hyperbola-contact} hold if and only if
every $\Delta_j$ is finite and
\begin{equation}\label{dyn:pair-defect-condition}
 \Delta_i+\Delta_j\le0\qquad(i\ne j).
\end{equation}
In particular, these are necessary conditions for ordinary preservation.
They are not asserted to be sufficient for ordinary preservation.
\end{lemma}

\begin{proof}
Fix a pair of sites and number them one and two.
The double-root leading coefficient of the product of the two hyperbola
factors in the proof of Proposition~\ref{dyn:local-reality} has sign $uv$.
Thus its contact inequality is
\[
 \frac{c_0^2}{u^2v^2}q_1(a_0+u,-a_0+v)
 +q_2(b_0+c_0/u,-b_0+c_0/v)\le0.
\]
Writing $S=u+v$ and $R=c_0S/(uv)$, when $S\ne0$ this is the sum of the
two ratios defining $\Delta_1,\Delta_2$.  Every pair of real points with
nonzero sums $S,R$ is attained: choose a nonzero real $P$ with the sign
of $S/R$, taking $0<P<S^2/4$ when it is positive; let $u,v$ be the roots
of $X^2-SX+P$, put $c_0=PR/S>0$, and choose $a_0,b_0$ to fit the first
coordinate in each pair.  Thus the contact inequalities with nonzero
sums say that every value of the first ratio plus every value of the
second is nonpositive.  Fixing a finite value of either ratio bounds
the other ratio above.  Each ratio also has at least one finite value,
so both suprema are real and their sum is nonpositive.  The converse
follows directly from the definition of the suprema.

When $S=0$, also $R=0$.  The paired Pauli derivative vectors obey
$\widetilde v(-s)=-\overline{v(s)}$ in \eqref{dyn:pauli-vectors}, so positivity of the local noise matrix
gives
\[
 q_j(s,-s)=-v(s)^{\mathsf T}K_j\overline{v(s)}\le0.
\]
These contacts therefore impose no additional condition.  For $n>2$,
pad the input by pure local factors that are nonzero at the selected
contact.  Local terms on the other sites then vanish to second order
at that contact, leaving exactly the displayed two-site inequality.
This proves the claim for every pair and the necessity statement.
\end{proof}

\begin{proposition}[Real hyperbola contacts are insufficient]\label{dyn:compensation}
For $a>0$ and $c\ge0$, the two-qubit generator
\[
 \mathcal L(\rho)=a(X_1\rho X_1-\rho)+c(Z_2\rho Z_2-\rho)
\]
does not ordinarily preserve Lee--Yang densities.  Nevertheless, when
$c\ge a$, it satisfies every real double-contact inequality obtained
from the pure input family used in \eqref{dyn:hyperbola-contact}.
\end{proposition}

\begin{proof}
Here $q_1(s,w)=4asw$ and
$q_2(s,w)=-c(1+s^2)(1+w^2)$.  The identities
$4sw\le(s+w)^2$ and
$(1+s^2)(1+w^2)-(s+w)^2=(sw-1)^2$ give
$\Delta_1=a$, $\Delta_2=-c$, with equality attained in each supremum.
Lemma~\ref{dyn:pair-defect} proves the assertion about all real
hyperbola contacts.

For actual failure, use the positive input
\[
 \rho=\frac27\left[
 \begin{pmatrix}2&0\\0&1\end{pmatrix}\otimes|0\rangle\langle0|
 +\frac14\begin{pmatrix}1&1\\1&1\end{pmatrix}\otimes|1\rangle\langle1|
 \right].
\]
It has trace one.  Up to its positive scalar its polynomial is
$2+sw+\tfrac14(1+s)(1+w)zv$, which is zero-free on $\D^4$ because
$|2+sw|>1$ and $|\tfrac14(1+s)(1+w)zv|<1$ there.
The second-site dephasing vanishes on this input and its entire
first-site trajectory.  At $t=(\log2)/a$, the first diagonal block is
$\diag(13/8,11/8)$, whereas the second block is fixed.  At $s=w=19i/20$
the output polynomial becomes
\[
 \frac{1229}{3200}+\frac{39+760i}{1600}zv.
\]
Choose $zv=-1229/[2(39+760i)]$.  Its modulus is $1229/1522<1$, since
$39^2+760^2=761^2$.  Taking both $z,v$ to be the same square root puts
all variables in the open disk and makes the output zero.
\end{proof}

\begin{example}[Random interchange and incoherent exclusion]\label{dyn:exchange-examples}
Let $S$ interchange two qubits.  For $0<p<1$, the channel
$\rho\mapsto p\rho+(1-p)S\rho S$ fails ordinary preservation.
The Lee--Yang input $|+0\rangle\langle+0|$ gives an output polynomial
proportional to
$p(1+z_1)(1+w_1)+(1-p)(1+z_2)(1+w_2)$.
It vanishes at
\[
 z_1=w_1=-1+\sqrt{1-p}\,e^{i\pi/4},\qquad
 z_2=w_2=-1+\sqrt p\,e^{-i\pi/4}.
\]
Every coordinate is in the open disk: for $0<x<1$,
$|-1+xe^{\pm i\pi/4}|^2=1+x^2-\sqrt2x<1$.
In particular the random-interchange semigroup fails at each positive
time, since its mixture parameter is $(1+e^{-2\gamma t})/2$.

The two jumps $\sqrt\gamma|01\rangle\langle10|$ and
$\sqrt\gamma|10\rangle\langle01|$, with $\gamma>0$, also fail at all
sufficiently small positive times.  On the same input, let
$u=e^{-\gamma t/2}$.  The $00$--$10$ coherence decays by $u$, and the two
excited populations exchange at rate $2\gamma$, giving
\[
 2F_t=1+u(z_1+w_1)+\tfrac{1+u^4}{2}z_1w_1
                         +\tfrac{1-u^4}{2}z_2w_2.
\]
At $z_1=w_1=-r$, $z_2=w_2=ir$, this is $1-2ur+u^4r^2$.
Its value at $r=0$ is positive.  At $r=1$ it is $u^4-2u+1$, negative
for $u<1$ sufficiently close to one, because this function vanishes at
one and has derivative two there.  The intermediate value theorem
supplies an interior zero.  Thus this natural quantum extension of
population exchange does not preserve the entire physical class.
\end{example}

\begin{example}[Equal local damping does not restore preservation]\label{dyn:equal-damping}
Apply the same amplitude-damping channel to both halves of a Bell state.
Its output polynomial, up to the positive factor $1/2$, is
\[
 1+(1-\eta)^2+\eta(1-\eta)(z_1w_1+z_2w_2)
 +\eta^2z_1z_2w_1w_2+\eta(z_1z_2+w_1w_2).
\]
At all four variables $ir$, the polynomial is
$1+(1-\eta)^2-2\eta(2-\eta)r^2+\eta^2r^4$.
It is positive at zero and equals $2(1-\eta)(1-2\eta)<0$ at one for
$1/2<\eta<1$.  Hence equal local damping fails ordinary joint
preservation at every sufficiently small positive time.  For instance,
$\eta=3/4$ has the exact interior root
$r^2=(5-2\sqrt2)/3$.  Vacuum padding extends the example to any larger
system.
\end{example}

% END SECTION quantum_dynamics

% BEGIN SECTION high_spin_dynamics
\subsection{Physical dynamics at arbitrary finite spins}
\label{dynspin:subsection}

The physical classification becomes more restrictive when a site has spin
at least one.  We keep the physical basis fixed and use the coherent
normalization of Appendix~\ref{spin:section}.  For positive integers
$\kappa_i=2j_i$, put
\[
 \mathcal H=\bigotimes_{i=1}^n\C^{\kappa_i+1},\qquad
 b_a=\prod_i\binom{\kappa_i}{a_i},\qquad
 f_\psi(z)=\sum_{a\leq\kappa}\sqrt{b_a}\psi_a z^a,
\]
and define
\begin{equation}\label{dynspin:kernel}
 F_\rho(z,w)=\sum_{a,b\leq\kappa}\sqrt{b_ab_b}\rho_{ab}z^aw^b.
\end{equation}
A density matrix is coherently Lee--Yang when this polynomial is nonzero
on $\D^{2n}$, with independent row and column variables.  Ordinary
preservation refers to this class on $\mathcal H$.  Complete preservation
means preservation after adjoining any finite collection of inert finite
spins, with the same encoding.  A pure input means a rank-one density
matrix with disk-stable vector polynomial; the vector need not be a
spin-coherent state.  Adjoints are always taken in the physical
orthonormal basis, equivalently in the polynomial inner product
$\langle z^a,z^b\rangle=\delta_{ab}/b_a$.

\begin{theorem}[Finite-spin physical generators]\label{dyn:arbitrary-spin}
Suppose that $\kappa_i\geq2$ for at least one site, and let
$\Phi_t=e^{t\mathcal L}$ be a continuous CPTP semigroup on $\mathcal H$.
The following conditions are equivalent.
\begin{enumerate}
\item $\Phi_t$ ordinarily preserves coherently Lee--Yang densities for
      every $t\geq0$.
\item For some $\varepsilon>0$, $\Phi_t$ preserves every pure coherently
      Lee--Yang input for $0\leq t<\varepsilon$.
\item $\Phi_t$ completely preserves coherently Lee--Yang densities for
      every $t\geq0$.
\item The generator is a sum of local generators,
\begin{equation}\label{dynspin:locality}
 \mathcal L=\sum_i\mathcal L_i\otimes\Id_{\widehat i}.
\end{equation}
At a site with $j_i\geq1$, its local term is
\begin{equation}\label{dynspin:isotropic}
 \mathcal L_i(\rho)=-i[h_iS_i^z,\rho]
 -\frac{\gamma_i}{2}\sum_{\alpha=x,y,z}
       [S_i^\alpha,[S_i^\alpha,\rho]],
 \qquad h_i\in\R,\quad\gamma_i\geq0.
\end{equation}
At a qubit site, its local term is
\[
 \mathcal L_i(\rho)=-i[h_iZ,\rho]
 +\sum_{\alpha,\beta=x,y,z}(K_i)_{\alpha\beta}
 \left(\sigma_\alpha\rho\sigma_\beta
       -\tfrac12\{\sigma_\beta\sigma_\alpha,\rho\}\right),
 \qquad h_i\in\R,
\]
where $B_i=B_i^{\mathsf T}\in\R^{2\times2}$, $d_i\in\R^2$, and
\begin{equation}\label{dynspin:qubit-cone}
 K_i=\begin{pmatrix}B_i&i d_i\\-i d_i^{\mathsf T}&c_i\end{pmatrix}
 \succeq0,\qquad
 \begin{pmatrix}
 c_iI_2-\operatorname{adj}B_i&d_i\\
 d_i^{\mathsf T}&\Tr B_i
 \end{pmatrix}\succeq0.
\end{equation}
\end{enumerate}
In particular, $\Phi_t=\bigotimes_i e^{t\mathcal L_i}$.
\end{theorem}

The dissipator in \eqref{dynspin:isotropic} is the familiar isotropic
spin diffusion, with Casimir multipole rates
\cite{DenisMartin22,FilippovKuzhamuratova19}.  The assertion here is its
necessity from physical zero-freeness, without an assumption of covariance,
linearity of the noise, locality, or unitality.  The axial Hamiltonian
commutes with the isotropic dissipator; a nonzero axial drift need not
make the full semigroup rotation covariant.

We first establish the reduction from arbitrary GKLS operators using
only physical pure inputs.  The spin realization is
\begin{equation}\label{dynspin:spin-action}
 S_i^x=\tfrac12[(1-z_i^2)\partial_i+\kappa_i z_i],\quad
 S_i^y=\tfrac1{2i}[(1+z_i^2)\partial_i-\kappa_i z_i],\quad
 S_i^z=\kappa_i/2-z_i\partial_i.
\end{equation}
The capacity-dependent Cayley transform is
\begin{equation}\label{dynspin:cayley}
 (Cf)(s)=2^{-|\kappa|/2}\prod_i(s_i+i)^{\kappa_i}
 f\left(\frac{s_1-i}{s_1+i},\ldots,
         \frac{s_n-i}{s_n+i}\right).
\end{equation}
It is the symmetric-power realization of \eqref{dyn:cayley}, and is
unitary in the weighted inner product.  Applied separately to both
coefficient legs of a pure density, it gives
\begin{equation}\label{dynspin:reflection}
 p(z)\mathscr Rp(w),\qquad
 \mathscr Rp(w)=(-1)^{|\kappa|}\overline{p(-\bar w)},\qquad p=Cf_\psi.
\end{equation}
Indeed this follows in capacity one from
$\overline C=\diag(-1,1)C$, and then by symmetric powers and tensor
products.  This two-leg transformation is not unitary conjugation of
the density matrix.  Every stable nonzero $p$ in the capacity box is
obtained from a physical pure input after normalization.  Its degree
may be strictly below the capacity; the inverse of
\eqref{dynspin:cayley} still belongs to the prescribed box.

\begin{lemma}[Physical order reduction]\label{dynspin:order}
Under condition (2) of Theorem~\ref{dyn:arbitrary-spin}, every jump in any
finite GKLS representation has, after removal of its scalar part, the form
\[
 V_\nu\in\operatorname{span}_{\C}\{S_i^x,S_i^y,S_i^z:1\leq i\leq n\}.
\]
The Hamiltonian has spin degree at most two: it is a sum of a scalar,
linear spin terms, on-site symmetric quadratics, and two-site bilinears.
At qubit sites the on-site quadratics reduce to a scalar and linear terms.
\end{lemma}

\begin{proof}
Write a finite-dimensional GKLS representation
\[
 \mathcal L(\rho)=-i[H,\rho]
 +\sum_\nu\left(V_\nu\rho V_\nu^*
                  -\tfrac12\{V_\nu^*V_\nu,\rho\}\right).
\]
In the vector-polynomial coordinates, let
\[
 Q_\nu=CV_\nu C^{-1}
       =\sum_{\alpha\leq\kappa}P_{\nu,\alpha}(z)\partial^\alpha
\]
be its canonical representation.  Let $R$ be the largest order among
all these operators.  Suppose $R\geq2$, and fix an index
$|\alpha|=R$.  For arbitrary $a\in\R^n$ take the stable pure input
$p_a(z)=\prod_i(z_i-a_i)^{\kappa_i}$.  In local coordinates
$u=z-a$, $v=w+a$, its doubled polynomial is exactly $u^\kappa v^\kappa$,
up to a positive normalization constant.  Put $M=2|\kappa|$.

Restrict the transformed output to $(u,v)=\tau\xi$ for any
$\xi\in\R_{>0}^{2n}$.  The initial polynomial is a nonzero multiple
of $\tau^M$, and the output remains upper-half-plane stable near time
zero.  Lemma~\ref{gen:cluster-moments} makes every coefficient of its
first variation below degree $M-2$ vanish.  Varying $\xi$ over the
positive orthant shows, by polynomial identity, that each multivariate
homogeneous part below that degree vanishes as well.

Consider the coefficient of $u^{\kappa-\alpha}v^{\kappa-\alpha}$.
Hamiltonian and anticommutator terms act on one leg only, leaving the
opposite full root factor untouched, and do not contribute.  The
sandwich term from $V_\nu$ is
$(Q_\nu p_a)(z)\mathscr R(Q_\nu p_a)(w)$.  To obtain
$u^{\kappa-\alpha}$ from a canonical term with derivative index $\beta$,
even allowing Taylor terms of its coefficient polynomial at $a$, one
must have $\beta\geq\alpha$.  Since $|\beta|\leq R=|\alpha|$, this
requires $\beta=\alpha$ and the constant Taylor coefficient.  Reflection
therefore gives the total coefficient
\begin{equation}\label{dynspin:squares}
 (-1)^R(\kappa)_\alpha^2
       \sum_\nu|P_{\nu,\alpha}(a)|^2,
 \qquad
 (\kappa)_\alpha=\prod_i\frac{\kappa_i!}{(\kappa_i-\alpha_i)!}.
\end{equation}
Its degree is $M-2R<M-2$, so it is zero.  Every square vanishes for
every real $a$.  Thus every coefficient of maximal order vanishes
identically, contradicting the definition of $R$.

All jumps have order at most one.  Proposition~\ref{gen:box-decomposition}
identifies this space as a scalar plus the local $\mathfrak{sl}_2$
actions.  Formula~\eqref{dynspin:spin-action} and the Cayley rotation
identify it with the stated spin span.  Absorb scalar jump parts into
the Hamiltonian.  The dissipative superoperator now has differential
order at most two on the doubled box.

If $CHC^{-1}$ had maximal canonical order $R\geq3$, use the same input
and choose $|\alpha|=R$.  The coefficient of
$u^{\kappa-\alpha}v^\kappa$ in the first variation can come only from
the row Hamiltonian and its constant maximal-order coefficient at $a$.
The column Hamiltonian retains $u^\kappa$, and the dissipator cannot
lower total degree by $R$.  Its degree $M-R<M-2$ forces that Hamiltonian
coefficient to vanish for every $a$, again a contradiction.
Proposition~\ref{gen:box-decomposition} identifies the resulting
order-two space with products of at most two local spin operators.
Hermiticity gives real coefficients in the Hermitian basis consisting
of scalars, local spins, local symmetric quadrupoles, and products on
distinct sites.  The antisymmetric local products are spin commutators,
and the scalar symmetric component is the Casimir.  At capacity one
the remaining quadrupoles vanish.  This proves the assertion.
\end{proof}

\begin{lemma}[Locality at finite spins]\label{dynspin:local}
Under the same hypothesis, the generator is a sum of local GKLS
generators as in \eqref{dynspin:locality}.
\end{lemma}

\begin{proof}
Put $T_i^\alpha=2S_i^\alpha$.  Lemma~\ref{dynspin:order} gives a PSD
Gram matrix $K$ in the basis of these local spin operators.  Its
$(i,j)$ block is denoted $K^{ij}$.  The transformed left and right
actions have respective derivative vectors
\begin{equation}\label{dynspin:vectors}
 \mathbf v(s)=(2s,1-s^2,i(1+s^2)),\qquad
 \widetilde{\mathbf v}(t)=(2t,t^2-1,i(1+t^2)).
\end{equation}
These vectors are independent of the capacity.  Only their zeroth-order
coefficients depend on $\kappa_i$.

Fix distinct sites $i,j$.  The physical pure product with transformed
vector $\prod_k(z_k-a_k)$ allows independent row--row contacts at
$z_i=s,z_j=t$ and row--column contacts at $z_i=s,w_j=t$, by choosing
$a_i=s$ and respectively $a_j=t$ or $a_j=-t$.  Keep all other legs
off their roots and restrict to a line with all slopes positive.
The initial root has multiplicity two.  Only the chosen mixed
principal coefficient survives at the contact; after division by the
initial leading coefficient its multiplier is the reciprocal of a
positive product of slopes.  Lemma~\ref{gen:cluster-moments} makes it
real and nonpositive for every real $s,t$.  On-site quadrupoles and
other terms leave at least one selected root untouched.

Write the pair Hamiltonian as
$\sum_{\alpha,\beta}h_{\alpha\beta}T_i^\alpha T_j^\beta$, with $h$
real.  The row--column coefficient is
$\mathbf v(s)^{\mathsf T}K^{ij}\widetilde{\mathbf v}(t)$.
The row--row coefficient is
$\mathbf v(s)^{\mathsf T}(-\Re K^{ij}-ih)\mathbf v(t)$: the conjugate Gram blocks
combine in the anticommutator, since the two sites commute.
Set
\[
 u(s)=\frac{(2s,1-s^2)}{1+s^2},\qquad R_0=\diag(1,-1).
\]
Its range is dense in the unit circle.  The normalized derivative
vectors are $(u(s),i)$ and $(R_0u(t),i)$.

Reality of the row--column coefficient, by comparison of its constant,
linear, and bilinear circle components, gives
\[
 K^{ij}=\begin{pmatrix}A&i b\\i d^{\mathsf T}&e\end{pmatrix},
 \qquad A,b,d,e\text{ real}.
\]
Its sign condition is
\begin{equation}\label{dynspin:cross-lr}
 u^{\mathsf T}AR_0v-b\cdot u-d\cdot R_0v-e\leq0
 \qquad (u,v\in S^1).
\end{equation}
For completeness, the comparison follows by averaging the imaginary
part first alone, then against $u$, $R_0v$, and $u(R_0v)^{\mathsf T}$;
the circle moments are $\mathbb E u=0$ and
$\mathbb E uu^{\mathsf T}=I_2/2$.
Reality of the row--row coefficient similarly gives $h_{\perp\perp}=0$
and $h_{zz}=0$.  With $k=h_{\perp z}$ and
$\ell^{\mathsf T}=h_{z\perp}$, its remaining sign condition is
\begin{equation}\label{dynspin:cross-ll}
 e-u^{\mathsf T}Av+k\cdot u+\ell\cdot v\leq0.
\end{equation}
Independent circle averages of \eqref{dynspin:cross-lr} and
\eqref{dynspin:cross-ll} are $-e$ and $e$, so $e=0$.
Each continuous nonpositive function then has zero average and vanishes
everywhere.  The same circle moments give $A=b=d=k=\ell=0$.
Thus all cross Gram blocks and pair Hamiltonians vanish.  The diagonal
Gram blocks are PSD and define local GKLS generators, proving locality.
\end{proof}

\begin{proof}[Proof of Theorem~\ref{dyn:arbitrary-spin}]
We prove (2)$\Rightarrow$(4) first.  The preceding lemmas give locality.
Each local semigroup preserves its pure Lee--Yang inputs near zero:
pad such an input with arbitrary pure Lee--Yang states at the other
sites.  Its evolved global polynomial is a product of the local output
polynomials, each nonzero since its density has trace one.  Stability
of that product implies stability of each factor.

Fix a site of capacity $\kappa\geq2$, retaining the normalization
$T=2S$.  Let $q(s,t)$ denote its row--column principal coefficient.
Use the pure vector polynomial $p(z)=(z-s)(z+t)$, with $s+t\ne0$.
At row coordinate $s$ and column coordinate $t$ both factors in
\eqref{dynspin:reflection} have simple roots.  A positive-slope line
has a double root, and only $q(s,t)$ contributes at the contact.
The common reflection sign cancels on dividing by the initial leading
coefficient.  Thus $q(s,t)$ is real and nonpositive; continuity adds
the omitted line $s+t=0$.

For a general Hermitian Gram matrix write
\[
 K=\begin{pmatrix}
 a&u+i\eta&r+i\mu\\
 u-i\eta&b&v+i\nu\\
 r-i\mu&v-i\nu&c
 \end{pmatrix}.
\]
Multiplication of \eqref{dynspin:vectors} gives
\[
 \Im q(s,t)=2(s+t)[-\eta(1-st)+r(1+st)+v(t-s)].
\]
Hence $\eta=r=v=0$, and
\begin{equation}\label{dynspin:local-shape}
 K=\begin{pmatrix}B&i d\\-i d^{\mathsf T}&c\end{pmatrix},
 \qquad B=B^{\mathsf T}\in\R^{2\times2},\quad d\in\R^2.
\end{equation}
With
\[
 r(s)=\frac{(2s,1-s^2)}{1+s^2},\qquad
 v(t)=\frac{(2t,t^2-1)}{1+t^2},
\]
the mixed coefficient satisfies
\begin{equation}\label{dynspin:circle}
 \frac{q(s,t)}{(1+s^2)(1+t^2)}
 =r(s)^{\mathsf T}Bv(t)+d\cdot(v(t)-r(s))-c\leq0.
\end{equation}
Both circle variables range densely and independently.  Taking them
equal gives $B\preceq cI_2$.

There is a second, independent contact available at this capacity.
Take $p(z)=(z-s)^2$ and a column coordinate away from its reflected
root.  Its pure-row principal coefficient must be real and nonpositive.
This follows either by real specialization of the column, or by a
positive-slope line through the contact; the nonzero complementary
factor cancels from the root-cluster coefficient.
Write the quadratic Hamiltonian as $H_2=T^{\mathsf T}QT$, with $Q$
real symmetric.  Put $u=r(s)$.  The pure-row principal coefficient,
divided by $(1+s^2)^2$, is
\begin{equation}\label{dynspin:row-quadratic}
 \tfrac12(c-u^{\mathsf T}Bu)+2Q_{\perp z}\cdot u
       -i(u^{\mathsf T}Q_{\perp\perp}u-Q_{zz}).
\end{equation}
Indeed its dissipative part is
$-\mathbf v(s)^{\mathsf T}(\Re K)\mathbf v(s)/2$, and its Hamiltonian part is
$-i\mathbf v(s)^{\mathsf T}Q\mathbf v(s)$, using the three-component vector in
\eqref{dynspin:vectors}.  Reality gives
$Q_{\perp\perp}=Q_{zz}I_2$.  Remove the corresponding Casimir scalar
and set $k=Q_{\perp z}$.  Nonpositivity at $u$ and $-u$ now gives
$c-u^{\mathsf T}Bu\leq0$.  Together with $B\preceq cI_2$, this yields
$B=cI_2$; the remaining inequality $2k\cdot u\leq0$ gives $k=0$.
The quadratic Hamiltonian is consequently scalar.

Equation~\eqref{dynspin:circle} reduces to
\[
 c(r\cdot v-1)+d\cdot(v-r)\leq0\qquad(r,v\in S^1).
\]
If $d\ne0$, take $r\perp d$ and
$v=(\cos\theta)r+(\sin\theta)d/|d|$ for small positive $\theta$.
The left side is $c(\cos\theta-1)+|d|\sin\theta>0$, a contradiction.
Thus $d=0$, $K=cI_3$, and positivity gives $c\geq0$.

The transformed isotropic dissipator has the explicit form
\begin{align}\label{dynspin:cayley-diffusion}
 \widetilde{\mathcal D}={}&-2c\kappa(\kappa zw+1)
 +2c\kappa z(zw-1)\partial_z
 +2c\kappa w(zw-1)\partial_w\notag\\
 &-2c(zw-1)^2\partial_z\partial_w.
\end{align}
To verify it directly, the left actions of $(T_x,T_y,T_z)$ are
\[
 2z\partial_z-\kappa,\quad
 (1-z^2)\partial_z+\kappa z,\quad
 i[(1+z^2)\partial_z-\kappa z],
\]
and the right actions are the same formulas in $w$ except that the
middle action has its sign reversed.  Substitute them into
$c\sum_\alpha(L_\alpha R_\alpha
 -(L_\alpha^2+R_\alpha^2)/2)$.
In particular this is a real operator and its pure-row and pure-column
second-order coefficients vanish.
For the remaining Hamiltonian $H=\widehat h_xT_x+\widehat h_yT_y+
\widehat h_zT_z$, use $p(z)=z-s$ and keep the reflected column off its
root.  The imaginary simple-root velocity comes only from the row
Hamiltonian and is
\[
 \widehat h_y s^2-2\widehat h_xs-\widehat h_y\geq0\qquad(s\in\R),
\]
by Lemma~\ref{gen:cluster-moments}.  Hence
$\widehat h_x=\widehat h_y=0$.  Returning to $S=T/2$ gives
\eqref{dynspin:isotropic}, with $h=2\widehat h_z$ and $\gamma=4c$.

It remains to recover the full qubit condition in a mixed-spin system.
Fix a qubit and a high-spin partner whose local form has just been
proved.  Their mixed principal coefficients will be denoted by
$q_{\rm q}(s,t)$ and $q_{\rm h}(u,v)=-2c(uv-1)^2$.
For arbitrary real $s,t$ with $s+t\ne0$, choose
\begin{equation}\label{dynspin:mask-parameters}
 \begin{gathered}
 x=(s+t)/2,\qquad a=(s-t)/2,\qquad c_0>|x|,\\
 b=\sqrt{(c_0/x)^2-1},\qquad
 u_*=b+c_0/x,\qquad v_*=-b+c_0/x.
 \end{gathered}
\end{equation}
Then $u_*v_*=1$.  The polynomial
\begin{equation}\label{dynspin:mixed-test}
 p(z,u)=(z-a)(u-b)-c_0
\end{equation}
is upper-half-plane stable, since for $z\in\HH$ its root in $u$ is
$b+c_0/(z-a)$, which lies in the lower half-plane.
It is a valid physical vector in capacities $(1,\kappa)$, including
when its degree is below $\kappa$.  Apart from the common parity sign,
its reflected factor is $(w+a)(v+b)-c_0$.

At $(z,u,w,v)=(s,u_*,t,v_*)$ both factors vanish.  Along the line with
all four slopes one, the initial quadratic coefficient, apart from
the same common sign, is $(c_0/x+x)^2>0$.  Locality makes its first
variation at the contact equal, with that sign, to
\begin{equation}\label{dynspin:mask-isolation}
 \frac{c_0^2}{x^2}q_{\rm q}(s,t)
       +x^2q_{\rm h}(u_*,v_*)
 =\frac{c_0^2}{x^2}q_{\rm q}(s,t).
\end{equation}
All lower-order terms vanish because one input factor remains zero;
the high-spin pure-row and pure-column coefficients are zero by
\eqref{dynspin:cayley-diffusion}.  Other sites can be padded by pure
stable factors kept nonzero at the contact.  The root-cluster condition
thus makes $q_{\rm q}(s,t)$ real and nonpositive.  Continuity again adds
$s+t=0$.  Consequently \eqref{dynspin:local-shape} and
\eqref{dynspin:circle} hold at every qubit as well.  For this Gram shape
the dissipator is real after Cayley transformation, as in the proof of
Proposition~\ref{dyn:local-complete}.  The same simple-root test excludes
its transverse Hamiltonian.  Lemma~\ref{dyn:lmi-reduction} gives precisely
\eqref{dynspin:qubit-cone}, including $\Tr B=0$.  This proves
(2)$\Rightarrow$(4), without treating the evolving partner as inert.

Conversely assume (4).  Formula~\eqref{dynspin:cayley-diffusion} is a
real endomorphism of the doubled capacity box whose only second-order
coefficient is $-2c(zw-1)^2\leq0$.  The axial Hamiltonian contributes
only a real first-order endomorphism.  At qubit sites, the proof of
Proposition~\ref{dyn:local-complete} gives a real box endomorphism with
its only mixed coefficient nonpositive by
Lemma~\ref{dyn:lmi-reduction}.  The sum therefore satisfies
Theorem~\ref{gen:main-real} on the full doubled capacity box.  Its
exponential preserves all complex stable polynomials after arbitrary
inert-variable extensions.  Undoing the coefficient Cayley transforms
gives the asserted physical complete preservation, in particular for
arbitrary inert finite spins.  No convexity of the stable class is used.
This proves (4)$\Rightarrow$(3).  Finally (3)$\Rightarrow$(1) and
(1)$\Rightarrow$(2) are immediate.
\end{proof}

The all-qubit ordinary classification is Theorem~\ref{dyn:main}; it is
not asserted there or here to follow from pure inputs alone.  On a single
qubit the ordinary cone is instead Proposition~\ref{dyn:one-qubit}, and
ordinary preservation need not be complete.  Theorem~\ref{dyn:arbitrary-spin}
also makes no classification claim for isolated noncovariant high-spin
channels: its semigroup hypothesis is essential to the argument.

\subsection{An application to rotation convolution semigroups}
\label{dynspin:brownian-subsection}

Gaussian processes on compact groups can already be recognized through
a faithful finite-dimensional representation and its tensor products;
see \cite[Theorem~2.8]{Voit03}, which restates the earlier compact-group
characterization.  For the real defining representation of $SO(3)$,
the spin-one conjugation action is its tensor square and contains the
defining representation as its antisymmetric part.  The following
application combines that recognition mechanism with the physical
classification above.  We give the short character proof in this case
using Hunt's formula \cite{Hunt56} and its twirling realization
\cite[Theorem~5.1]{Aniello10}.

\begin{corollary}[Lee--Yang rotation dynamics]\label{dyn:brownian-recognition}
Let $(\mu_t)_{t\geq0}$ be a weakly continuous probability convolution
semigroup on $SO(3)$ with $\mu_0=\delta_e$.  Let $U$ be its spin-one
representation and set
\[
 \Phi_t(\rho)=\int_{SO(3)}U(g)\rho U(g)^*\,\mu_t(dg).
\]
Then $\Phi_t$ preserves every pure coherently Lee--Yang input in a
neighborhood of zero if and only if the group generator is
\begin{equation}\label{dynspin:group-generator}
 \mathcal A=hX_z+\frac\gamma2(X_x^2+X_y^2+X_z^2),
 \qquad h\in\R,\quad\gamma\geq0.
\end{equation}
Here $X_\alpha$ are the left-invariant fields normalized by
$U(\exp(tE_\alpha))=e^{-itS_\alpha}$ and
$X_\alpha f(g)=\left.\frac{d}{dt}f(g\exp(tE_\alpha))\right|_{t=0}$.
In particular the L\'evy measure is zero, without a prior assumption of
continuous sample paths or rotation invariance of $\mu_t$.
\end{corollary}

\begin{proof}
The twirling maps form a continuous CPTP semigroup.  Under the stated
preservation hypothesis Theorem~\ref{dyn:arbitrary-spin} gives
\eqref{dynspin:isotropic}.  The spin-one operator space decomposes as
the angular-momentum representations $0\oplus1\oplus2$.  On its
$\ell$ component the double-commutator Casimir is $\ell(\ell+1)I$;
the axial commutator has trace zero.  Thus the generator restrictions
satisfy
\[
 \Tr\mathcal L_1=-3\gamma,\qquad \Tr\mathcal L_2=-15\gamma.
\]
Let $\chi_1,\chi_2$ be the characters of the three- and five-dimensional
representations.  For a rotation of angle $\theta$, put
\begin{equation}\label{dynspin:character-witness}
 F(g)=2-\chi_1(g)+\tfrac15\chi_2(g)
     =\tfrac45(1-\cos\theta)^2.
\end{equation}
This smooth function is strictly positive off the identity and has
zero jet through order two at the identity.  Since
$\int\chi_\ell\,d\mu_t=\Tr(e^{t\mathcal L_\ell})$, the displayed
traces give $\mathcal AF(e)=3\gamma-15\gamma/5=0$.

Write Hunt's generator with drift, diffusion matrix $a\succeq0$, and
L\'evy measure $\nu$, taking its second-order part to be
$\tfrac12\sum_{\alpha,\beta}a_{\alpha\beta}X_\alpha X_\beta$.
At the identity all local terms and the linear jump compensation
vanish on $F$.  Therefore
\[
 0=\mathcal AF(e)=\int_{SO(3)\setminus\{e\}}F(g)\,\nu(dg).
\]
The integral is finite: near the identity $F=O(\dist(g,e)^4)$, while
a L\'evy measure has the requisite local second-moment integrability;
away from the identity compactness applies.  Positivity of $F$ off
$e$ implies $\nu=0$, including for an initially infinite-activity
L\'evy measure.

On the dipole space the symmetric part of the remaining generator is
$\tfrac12(a-(\Tr a)I_3)$.  Comparing with $-\gamma I_3$, first taking
traces and then substituting back, gives $a=\gamma I_3$.  Its
antisymmetric part identifies the drift as $hX_z$.  This proves
\eqref{dynspin:group-generator}.  Conversely, that group generator
induces \eqref{dynspin:isotropic} under twirling, so
Theorem~\ref{dyn:arbitrary-spin} gives complete preservation.
\end{proof}

The group and the tested spin matter.  Conjugation cannot detect jumps
in the center $\{\pm I\}$ of $SU(2)$, so the zero-L\'evy-measure
conclusion is stated on $SO(3)$.  Moreover, Poisson jumps of rate
$\lambda$ by angle $\pi$ about the $z$ axis induce on a qubit
\[
 \Phi_t(\rho)=\tfrac12(1+e^{-2\lambda t})\rho
              +\tfrac12(1-e^{-2\lambda t})Z\rho Z.
\]
This is an allowed dephasing semigroup, but its rotation process has
nonzero jumps.  Thus spin one cannot be replaced uniformly by spin
one-half in Corollary~\ref{dyn:brownian-recognition}.

% END SECTION high_spin_dynamics

% BEGIN SECTION positive_capacities
\section{Positive coefficients and particle capacities}
\label{pos:section}

The restriction to nonnegative coefficients changes the generator problem substantially.
In particular, collective particle loss is compatible with stability, whereas its extension
by an arbitrary unchanged coordinate need not be.  We first classify the unrestricted
positive problem and then give a capacity-sensitive criterion.  Throughout this section,
coefficients are ordinary monomial coefficients, and matrix entries use the convention
\(Az^x=\sum_y A_{y,x}z^y\).

Write \(\mathcal S_+\) for the nonzero real stable polynomials with nonnegative
coefficients.  Write \(\mathcal{LP}_+\) for their nonzero locally uniform entire limits.
When taking closed cones we adjoin zero explicitly.  On a countable state space we use
the coefficient space \(\ell^1\).  Thus every member of \(\mathcal{LP}_+\) belongs to
\(\ell^1\), although a general member of \(\ell^1\) need not be entire.
\begin{lemma}[Positive Laguerre--P\'olya coefficient closure, \cite{LVR10}]\label{pos:lp-closure}
Coefficientwise limits of members of \(\mathcal{LP}_+\), when the limiting
coefficient sum is finite, belong to \(\mathcal{LP}_+\cup\{0\}\).
Every member of \(\mathcal{LP}_+\) has stable Jensen-polynomial approximants
converging in \(\ell^1\), whose coefficients are bounded by those of the
limiting series.
\end{lemma}
This is the positive Laguerre--P\'olya closure and Jensen characterization.
We use this precise standard theorem when passing from coefficient norm convergence
to an entire stable limit.

\subsection{Root contacts and Poisson tangency}

\begin{lemma}[A repeated-root first variation]\label{pos:root-contact}
Suppose \(p_t\) is holomorphic near zero, has real Taylor coefficients, and
\[
p_t(u)=c u^m+t h(u)+o(t),\qquad c>0,
\]
locally uniformly as \(t\downarrow0\).  Suppose that all zeros sufficiently near zero
are real.  Then \(h(0)=0\) if \(m\ge3\), and \(h(0)\le0\) if \(m=2\).
\end{lemma}
\begin{proof}
Fix a small circle isolating the zero of \(p_0\).  The argument principle and
\(p_t'/p_t\) express the power sums of the enclosed \(m\) roots as contour integrals.
Consequently their monic root polynomial has a right first variation, and its first
two nonleading coefficients are \(O(t)\).  If its roots are \(r_1(t),\ldots,r_m(t)\),
Newton's identity gives \(\sum_jr_j(t)^2=O(t)\).  Since these roots are real, each is
\(O(\sqrt t)\).  Their product is \(o(t)\) when \(m\ge3\).  The complementary analytic
factor tends to \(c\), so this proves the first assertion.  For \(m=2\), the
discriminant of the root polynomial is nonnegative.  Its linear coefficient is
\(O(t)\), hence its constant coefficient has upper first derivative at most zero.
Multiplication by the positive limiting factor \(c\) proves the second assertion.
\end{proof}

For a capacity vector \(\kappa\in(\N\cup\{\infty\})^d\), let
\(\Omega_\kappa=\{x\in\N^d:x_i\le\kappa_i\}\).  Coordinates of capacity zero are
removed.  Every linear map on the allowed polynomials has a unique locally finite
normal-ordered expression
\begin{equation}\label{pos:normal-recursion}
 A=\sum_{\alpha\le\kappa}P_\alpha(z)\partial^\alpha,
 \qquad
 P_\alpha=\frac1{\alpha!}\left(Az^\alpha-
 \sum_{\substack{\beta\le\alpha\\\beta\ne\alpha}}
 (\alpha)_\beta z^{\alpha-\beta}P_\beta\right).
\end{equation}
Here the \(P_\alpha\) may have degrees outside the original coordinate box.
The expression is locally finite because only finitely many derivatives act on a
fixed polynomial.

\begin{lemma}[Forced differential order]\label{pos:order}
Suppose the polynomial action of a positive stability-preserving semigroup has a
locally uniform generator first variation near a real open negative box.
Then \(P_\alpha=0\) whenever \(|\alpha|\ge3\).
For unrestricted coordinates it is enough to assume preservation of every
\((c+a\cdot z)^m\), where \(c\ge0\), \(a>0\), and \(m\ge0\).
\end{lemma}
\begin{proof}
For an allowed \(\alpha\), use
\(f(z)=\prod_i(z_i+s_i)^{\alpha_i}\), with \(s_i>0\) small.
At \(z=-s\), every normal-ordered term vanishes except the term with derivative
\(\alpha\); thus \(Af(-s)=\alpha!P_\alpha(-s)\).
Along a strictly positive direction through \(-s\), the initial polynomial has a
root of multiplicity \(|\alpha|\), with positive leading coefficient.
Lemma~\ref{pos:root-contact} annihilates \(P_\alpha\) on an open box, hence
identically by analytic continuation.  For the last assertion choose \(y<0\),
\(c=-a\cdot y\), and evaluate \(A(c+a\cdot z)^m\) at \(y\).  This gives
\(m!\sum_{|\alpha|=m}a^\alpha P_\alpha(y)=0\) for every \(a>0\).
Coefficient comparison in \(a\) gives the same conclusion.
\end{proof}

The local first variation in this lemma is automatic in finite dimension.  For
an existing positive \(C_0\)-semigroup, the required domain and local analytic
statements are supplied by Theorem~\ref{ana:domain}; the subsequent extension
to arbitrary real contacts is kept separate from this initial local argument.

\begin{lemma}[Poisson tangency]\label{pos:poisson}
Let \(G_t\in\mathcal{LP}_+\), \(G_0(z)=c\exp(a\cdot z)\), where \(c>0\), \(a\ge0\).
Suppose \(G_t=G_0+tH+o(t)\) locally uniformly near the origin.  Then
\(H/G_0\) is a polynomial of total degree at most two.  Its homogeneous quadratic
part is nonpositive on every strictly positive vector.
\end{lemma}
\begin{proof}
The one-variable positive Laguerre--P\'olya factorization, including its exponential
factor, is
\[
 g(s)=g(0)e^{\sigma s}\prod_j(1+\rho_js),
 \qquad \sigma,\rho_j\ge0,\quad \sum_j\rho_j<\infty.
\]
It follows from the standard Laguerre--P\'olya product representation
\cite{LVR10}; positive coefficients put every finite zero on the negative real axis.
If \(\log(g(s)/g(0))=\sum_{k\ge1}\ell_ks^k\), then for \(k\ge2\)
\[
 \ell_k=\frac{(-1)^{k-1}}k\sum_j\rho_j^k,
 \qquad
 |\ell_k|\le\frac{(-2\ell_2)^{k/2}}k.
\]
The second inequality follows from
\(\sum\rho_j^k\le(\sum\rho_j^2)^{k/2}\).
Apply this to \(G_t(sv)\) for fixed \(v>0\).  Its logarithm has a first variation,
\(\ell_2(t)=O(t)\), and \(\ell_2(t)\le0\).  Therefore
\(\ell_k(t)=o(t)\) for every \(k\ge3\).
The first variation of its logarithm is \(H(sv)/G_0(sv)\).
Each homogeneous Taylor coefficient of degree at least three vanishes on the open
positive orthant, and is therefore the zero polynomial.  The degree-two coefficient
is nonpositive there, as claimed.
\end{proof}

\subsection{Collective loss and an exact Euler step}

Put \(E=\sum_i z_i\partial_i\), \(D=\sum_i\partial_i\), and
\begin{equation}\label{pos:coal-operator}
 C=ED-E(E-1)=\sum_{i,j}z_j(1-z_i)\partial_i\partial_j.
\end{equation}
For \(|x|=n\),
\(Cz^x=(n-1)\sum_i x_i z^{x-e_i}-n(n-1)z^x\).
Thus \(C\) is a Markov generator of single-particle losses, with the rate of losing
a type \(i\) particle equal to \(x_i(|x|-1)\).

\begin{proposition}[Sharp positive Euler step]\label{pos:coal}
On polynomials of total degree at most \(N\ge2\), the operator \(I+hC\) preserves
\(\mathcal S_+\) and evaluation at \(\boldsymbol1\) whenever
\(0\le h\le1/[N(N-1)]\).  This is the exact interval on which it is a Markov Euler
step on the whole cutoff.  Consequently \(e^{tC}\) preserves \(\mathcal S_+\).
\end{proposition}
\begin{proof}
Write \(f=\sum_{n=0}^Nf_n\) in homogeneous parts and homogenize positively:
\(F(w,z)=\sum_nw^{N-n}f_n(z)\).  Positive homogenization preserves stability
\cite[Theorem 4.5]{BBL09}.  Hence
\[
 G(\xi,\eta,z)=F(1+\xi+\eta,z+\eta\boldsymbol1)
\]
is stable.  Set \(a=1-hN(N-1)\), \(b=h(N-1)\), and
\(\mathcal L=a+b(\partial_\xi+\partial_\eta)-h\partial_\xi\partial_\eta\).
For \(h>0\) its negative differential symbol is
\(a-b(u+v)-huv\).  Solving for \(v\) gives
\[
 \Im\frac{a-bu}{b+hu}
 =-\frac{b^2+ah}{|b+hu|^2}\Im u<0,
 \qquad b^2+ah=h[1-h(N-1)]>0.
\]
Thus this symbol is stable.  The constant-coefficient differential-symbol theorem
\cite{BB09,BB10} shows that \(\mathcal LG\) is stable or zero, including in the
unchanged \(z\)-variables.
For \(k=N-n\), direct differentiation of
\((1+\xi+\eta)^kf_n(z+\eta\boldsymbol1)\) gives, at \(\xi=\eta=0\),
\[
 \mathcal LG_n=
 [a+2bk-hk(k-1)]f_n+(b-hk)Df_n
 =[1-hn(n-1)]f_n+h(n-1)Df_n.
\]
Therefore \((\mathcal LG)|_{\xi=\eta=0}=(I+hC)f\).  Real specialization preserves
stability or zero.  The displayed monomial formula for \(C\) proves nonnegativity
of the coefficients and preservation of their sum, so the output is nonzero.
For \(h>1/[N(N-1)]\), the coefficient of \(z_1^N\) in \((I+hC)z_1^N\) is negative.
Finally take Euler limits on each fixed finite total-degree cutoff; coefficient
closure and preservation of mass give the assertion for \(e^{tC}\).
\end{proof}

\subsection{The unrestricted logarithmic symbol}

For \(a\in\R^d\), define the logarithmic symbol, when the indicated exponential
belongs to the generator domain, by
\(\Psi_A(z,a)=e^{-a\cdot z}A(e^{a\cdot z})\).

\begin{theorem}[Unrestricted positive classification]\label{pos:symbol}
Let \((T_t)_{t\ge0}\) be a coefficientwise-positive \(C_0\)-semigroup of bounded
operators on \(\ell^1(\N^d)\), with generator \(A\).  The following are equivalent:
\begin{enumerate}
\item \(T_t\mathcal{LP}_+\subseteq\mathcal{LP}_+\) for every \(t\ge0\).
\item \(T_t\mathcal S_+\subseteq\mathcal{LP}_+\) for every \(t\ge0\).
\item All polynomials belong to \(D(A)\), and on polynomials
\begin{equation}\label{pos:canonical}
 \begin{aligned}
 A={}&cI+\lambda\cdot z+
 \sum_i\left(b_i+\sum_jm_{ij}z_j\right)\partial_i
 \\
 &+ED_\nu-\beta E(E-1)-\sum_i\eta_i z_i^2\partial_i^2,
 \qquad D_\nu=\sum_i\nu_i\partial_i .
 \end{aligned}
\end{equation}
where \(\lambda_i,b_i,\nu_i,\beta,\eta_i\ge0\), \(m_{ij}\ge0\) for \(i\ne j\),
and \(c,m_{ii}\in\R\).
\end{enumerate}
For \(d\ge2\) the parameters are unique.  For \(d=1\), combine \(\beta+\eta_1\)
into one nonnegative parameter.  Equivalently,
\begin{equation}\label{pos:log-symbol}
 \Psi_A(z,a)=c+\lambda\cdot z+b\cdot a+a^{\mathsf T}Mz
 +(\nu\cdot a)(z\cdot a)-\beta(z\cdot a)^2
 -\sum_i\eta_i(z_i a_i)^2.
\end{equation}
The theorem concerns an existing semigroup.  The separate existence criterion for
a formal expression is Theorem~\ref{ana:weighted-generation}.
\end{theorem}
\begin{proof}
The first implication is immediate.  Assume the second statement.  Jensen
approximation and Lemma~\ref{pos:lp-closure} give preservation of all
\(\mathcal{LP}_+\), initially allowing zero.  Zero is impossible: for each basis
vector, strong continuity gives \((T_{t/n})_{x,x}>0\) for sufficiently large
\(n\), and positivity and the semigroup law give
\((T_t)_{x,x}\ge((T_{t/n})_{x,x})^n>0\).
Every nonzero positive input dominates a positive multiple of a basis vector.
The domain conclusion follows from Corollary~\ref{ana:polynomial-domain}.  Lemma~\ref{pos:order} gives
\(A=\sum_{|\alpha|\le2}P_\alpha\partial^\alpha\), with \(P_\alpha\in\ell^1\).
Consequently
\begin{equation}\label{pos:graph-bound}
 \|Az^x\|_1\le\sum_{|\alpha|\le2}(x)_\alpha\|P_\alpha\|_1
 \le C(1+|x|)^2.
\end{equation}
The polynomial partial sums of \(e^{a\cdot z}\) converge in graph norm by this
bound, since Poisson factorial moments are finite.  Closedness of \(A\) gives
\[
 e^{a\cdot z}\in D(A),\qquad
 A(e^{a\cdot z})=e^{a\cdot z}\sum_{|\alpha|\le2}a^\alpha P_\alpha(z).
\]
The stable approximants \((1+a\cdot z/m)^m\) converge in coefficient norm.
Their images and positive coefficient closure imply
\(T_te^{a\cdot z}\in\mathcal{LP}_+\).
Lemma~\ref{pos:poisson}, followed by coefficient comparison in \(a>0\), now gives
\(\deg P_\alpha\le2\) and
\begin{equation}\label{pos:poisson-inequality}
 P_0^{[2]}(v)+\sum_i a_iP_i^{[2]}(v)
 +\sum_{|\alpha|=2}a^\alpha P_\alpha^{[2]}(v)\le0
 \quad(a,v>0).
\end{equation}

Positivity of the semigroup makes the off-diagonal entries of \(A\) nonnegative.
The nonconstant coefficients of \(P_0=A1\) are therefore nonnegative.  Let
\(a\to0\) in \eqref{pos:poisson-inequality}; its quadratic part must vanish, giving
\(P_0=c+\lambda\cdot z\), \(\lambda\ge0\).
If \(|\rho|=2\), the coefficient of \(z^{(n-1)e_i+\rho}\) in \(Az_i^n\) is
\(n[z^\rho]P_i+O(1)\).  No second-order term has sufficiently high output degree
to contribute.  Positivity for arbitrarily large \(n\) gives
\([z^\rho]P_i\ge0\).  Replacing \(a\) by \(\varepsilon a\) in
\eqref{pos:poisson-inequality} and taking \(\varepsilon\downarrow0\) eliminates
these quadratic coefficients.  Thus
\(P_i=b_i+\sum_jm_{ij}z_j\), with the asserted signs obtained from \(Az_i\).

At this point \(A-M_\lambda\), where \(M_\lambda f=(\lambda\cdot z)f\), preserves
every finite total-degree cutoff, and \(M_\lambda\) is bounded on \(\ell^1\).
Theorem~\ref{ana:weighted-generation} gives compatible weighted semigroups:
if the cutoff-preserving semigroup has bound \(Me^{\omega t}\), then
\(\|T_tf\|_R\le M R^N e^{(\omega+M\Lambda R)t}\|f\|_1\) for
\(\deg f\le N\), \(R\ge1\), where \(\Lambda=\sum_i\lambda_i\).
Since \(A^2f\) has degree at most \(N+2\), the Taylor integral gives
\[
 \|T_tf-f-tAf\|_R\le
 \frac{t^2}{2}M R^{N+2}e^{(\omega+M\Lambda R)\tau}\|A^2f\|_1
 \quad(0\le t\le\tau).
\]
Thus generator first variations are legitimate at every real point, not only
in the unit polydisc.  We now determine the second-order part algebraically.

For \(|\alpha|=2\),
\begin{equation}\label{pos:reaction-positive}
 [z^\rho]P_\alpha\ge0\qquad(\rho\ne\alpha).
\end{equation}
For \(\alpha=2e_i\), divide the relevant off-diagonal coefficient in \(Az_i^n\)
by \(n(n-1)\) and let \(n\to\infty\).  For \(\alpha=e_i+e_j\), at least one of
\(\rho_i,\rho_j\) is zero if \(|\rho|\le2\) and \(\rho\ne\alpha\).
If \(\rho_i=0\), use source \(e_i+ne_j\) and target \((n-1)e_j+\rho\).
Terms differentiating only \(z_j\) retain \(z_i\), and \(\partial_i^2\) vanishes.
The coefficient is therefore \(n[z^\rho]P_\alpha+O(1)\).  This proves
\eqref{pos:reaction-positive}.

The square test \((c+a\cdot z)^2\) gives
\[
 H_a(z):=\sum_{|\alpha|=2}a^\alpha P_\alpha(z)\le0
 \quad\text{whenever }a>0,\ a\cdot z<0.
\]
Scaling \(z\) to zero first removes the nonnegative constant terms.
The linear part \(\ell(a)\cdot z\) is nonpositive on the half-space
\(a\cdot z<0\), so \(\ell(a)=k(a)a\) with \(k(a)\ge0\).
The polynomial identities \(a_j\ell_i=a_i\ell_j\) show that every \(\ell_i\)
is divisible by \(a_i\), and the linear quotients agree.  Thus
\(k(a)=\nu\cdot a\), with \(\nu_i\ge0\).
Scaling \(z\) to infinity shows that the quadratic part of \(H_a\) is nonpositive
on the same half-space.  Being even in \(z\), it is nonpositive everywhere.

Here is the full elimination of off-diagonal quadratic outputs.  Set \(a=e_i\).
On \(z_i=0\), the nonnegative coefficients in \eqref{pos:reaction-positive}
force every quadratic term not containing \(z_i\) to vanish.  A surviving
\(z_i z_j\) term is also impossible, because \(z_j\) can have either sign and
arbitrarily large magnitude.  For mixed sources, the coefficient of \(z_k^2\)
in the nonpositive quadratic form is
\(-h_{kk}a_k^2+\sum_{i<j}[z_k^2]P_{e_i+e_j}a_i a_j\).
Put \(a_k=\varepsilon\), all other entries equal to one, and let
\(\varepsilon\downarrow0\).  This first removes all coefficients whose
source avoids \(k\).  Divide the remaining inequality by \(\varepsilon\)
and take the limit again to remove sources containing \(k\).  Thus every
displayed nonnegative mixed-source coefficient vanishes.  Finally, a mixed output different from its mixed source
has an index outside the source.  Support \(a\) on the two source indices.
The corresponding quadratic form has zero square coefficient at the outside
index; a nonpositive quadratic form with such a zero has an identically zero
row there.  This removes that output as well.  We have proved
\[
 P_{2e_i}=\nu_i z_i-h_{ii}z_i^2,\qquad
 P_{e_i+e_j}=\nu_jz_i+\nu_i z_j-2h_{ij}z_i z_j,
\]
where \(H=(h_{ij})\) is positive semidefinite.

Write the principal matrix as
\(Q(z)=(\nu z^{\mathsf T}+z\nu^{\mathsf T})/2-
\diag(z)H\diag(z)\).  Different affine factors supply the additional restriction.
At a real common zero of
\((c_1+a\cdot z)(c_2+b\cdot z)\), with nonnegative factors,
Lemma~\ref{pos:root-contact} gives \(a^{\mathsf T}Q(z)b\le0\).
For distinct \(i,j,k\), apply it to \((z_i+z_k)(1+z_j)\) at
\((z_i,z_k,z_j)=(s,-s,-1)\).  The result is
\[
 -\frac{\nu_i+\nu_k}{2}+s(h_{ij}-h_{kj})\le0\quad(s\in\R).
\]
Hence all off-diagonal entries of \(H\) are a common value \(\beta\) when
\(d\ge3\); for \(d=2\), put \(\beta=h_{12}\).
The contact \((s+z_i)(s+z_j)\), at \(z_i=z_j=-s\), gives
\(-s(\nu_i+\nu_j)/2-\beta s^2\le0\), whence \(\beta\ge0\).
The contact \((z_i+z_j)(s+z_i)\), at \((z_i,z_j)=(-s,s)\), gives
\(-s(\nu_i+\nu_j)/2-s^2(h_{ii}-\beta)\le0\), whence \(h_{ii}\ge\beta\).
Thus \(H=\beta\boldsymbol1\boldsymbol1^{\mathsf T}+\diag(\eta)\),
\(\beta,\eta\ge0\).  This proves necessity and uniqueness.

For sufficiency, remove the bounded multiplication \(M_\lambda\) and work first
on a fixed total-degree cutoff.  The affine first-order part has flow
\[
 f(z)\longmapsto e^{ct}
 f\left(e^{tM}z+\int_0^te^{sM}b\,ds\right).
\]
The matrix \(M\) is Metzler, so \(e^{tM}\) is nonnegative with positive diagonal;
this substitution preserves \(\mathcal S_+\).
For \(\beta>0\) and \(\nu_i>0\), conjugate \(\beta C\) by positive coordinate
scaling \(S_sf(z)=f(s_1z_1,\ldots,s_dz_d)\), with \(s_i=\nu_i/\beta\):
\[
 S_s^{-1}(\beta C)S_s=ED_\nu-\beta E(E-1).
\]
Proposition~\ref{pos:coal} proves preservation.  Zero parameters follow by
parameter limits on the same finite-dimensional cutoff.
For individual damping \(K_i=-z_i^2\partial_i^2\), the active algebraic symbol of
\(I+hK_i\) at coordinate degree \(N\) factors as
\[
 (z+w)^{N-2}
 [w+(1-\sqrt\theta)z][w+(1+\sqrt\theta)z],
 \qquad\theta=hN(N-1)\in[0,1].
\]
It is stable, and the coefficients of the operator on monomials are nonnegative.
The finite-degree symbol theorem and Euler limits give preservation by
\(e^{t\eta_iK_i}\).  Finite-dimensional Lie products combine these elementary
flows.  Stability follows by closure, and invertibility of each finite
exponential excludes zero outputs.
Finally the bounded-perturbation product formula, justified in
Theorem~\ref{ana:weighted-generation}, adds \(M_\lambda\); its own flow is
multiplication by \(e^{t\lambda\cdot z}\).  Jensen approximation and coefficient
closure extend the result to \(\mathcal{LP}_+\).  Nonzero outputs follow from
\(T_tz^x\ge e^{A_{x,x}t}z^x\), the variation-of-constants inequality for a positive
semigroup and a basis vector in its domain.
\end{proof}

\begin{corollary}[Arbitrary-jump Markov rigidity]\label{pos:markov}
Let \(T_t\) be the transition semigroup of a conservative nonexplosive continuous-time
Markov chain on \(\N^d\), with finite exit rate at every state.  No restriction on
jump sizes or rate dependence is imposed.  The following are equivalent:
\begin{enumerate}
\item Every input in \(\mathcal S_+\) is mapped into \(\mathcal{LP}_+\) at every time.
\item The same assertion holds for all \((c+a\cdot z)^m\), where
\(c\ge0\), \(a>0\), and \(m\ge0\).
\item The generator is
\begin{equation}\label{pos:markov-generator}
 A=\sum_i\lambda_i(z_i-1)+
 \sum_i\left[\delta_i(1-z_i)+\sum_{j\ne i}r_{ij}(z_j-z_i)\right]\partial_i
 +\gamma C,
\end{equation}
with \(\lambda_i,\delta_i,r_{ij},\gamma\ge0\).
\end{enumerate}
In particular, the only transitions are constant immigration, independent
migration, and the single deaths
\(x\to x-e_i\) at rate \(x_i[\delta_i+\gamma(|x|-1)]\).
If all outputs from polynomial inputs must remain polynomials, then
\(\lambda=0\), and this condition is sufficient.
\end{corollary}
\begin{proof}
Finite exit rates put every basis vector in the \(\ell^1\) generator domain:
conditioning on the first jump gives
\(\|(T_tz^x-z^x)/t-Az^x\|_1\to0\) and
\(\|T_tz^x-z^x\|_1\le2tq(x)\).
Under the second assertion, the multinomial version of Lemma~\ref{pos:order}
applies.  The diagonal coefficient gives
\[
 q(x)=-\sum_{|\alpha|\le2}(x)_\alpha[z^\alpha]P_\alpha
 \le C(1+|x|)^2.
\]
Normalized multinomial approximants to a product-Poisson law converge in total
variation.  Markov contraction and positive coefficient closure therefore imply
preservation of those Poisson inputs.  The graph-norm and Poisson-tangency argument
in Theorem~\ref{pos:symbol} applies unchanged.  It yields constant immigration,
affine first-order coefficients, and a second-order part
\[
 P_{2e_i}=\nu_i z_i-h_{ii}z_i^2,\qquad
 P_{e_i+e_j}=\nu_jz_i+\nu_i z_j-2h_{ij}z_i z_j,
 \qquad H\succeq0.
\]
Only square multinomial contacts were used to reach this point.
Mass preservation gives \(P_\alpha(\boldsymbol1)=0\), hence
\(h_{ii}=\nu_i\) and \(2h_{ij}=\nu_i+\nu_j\).
The quadratic part of the square-contact expression is consequently
\(-(a\cdot z)(\sum_i\nu_i a_i z_i)\), nonpositive for every real \(z\).
Two nonzero real linear forms whose product is everywhere nonnegative are
nonnegative multiples of one another: the second must vanish on the kernel
of the first.  Taking \(a=\boldsymbol1\) shows that every \(\nu_i\) equals a
common \(\gamma\ge0\).  The constant and linear mass identities give exactly
\eqref{pos:markov-generator}.  This proves the difficult implication without
assuming preservation of products of different affine factors.
Sufficiency follows from Theorem~\ref{pos:symbol}, or from the same elementary
flows and bounded immigration construction.  The chain in
\eqref{pos:markov-generator} is nonexplosive: its population is bounded by the
initial population plus a rate \(\sum_i\lambda_i\) Poisson immigration process,
and every bounded-population set is finite.
If \(\lambda_i>0\), from the empty state a path with any prescribed number of
type \(i\) immigrations and no other jumps has positive probability before any
fixed positive time.  The output therefore has infinite support.  Conversely
\(\lambda=0\) leaves every total-degree cutoff invariant.
\end{proof}

Modulo the lineality generated by \(I,z_1\partial_1,\ldots,z_d\partial_d\),
Theorem~\ref{pos:symbol} describes, for \(d\ge2\), the formal cone generated by
\[
 z_i,\quad\partial_i,\quad z_j\partial_i\ (i\ne j),\quad
 E\partial_i,\quad-E(E-1),\quad-z_i^2\partial_i^2.
\]
This is a statement about the formal cone; it does not assert that each ray
individually generates bounded operators on every fixed coefficient space.

\begin{corollary}[Diagonal selection]\label{pos:diagonal}
Suppose \(v:\N^d\to\R\) is bounded above.  The positive semigroup
\(T_tz^x=e^{tv(x)}z^x\) preserves \(\mathcal{LP}_+\) if and only if
\[
 v(x)=c+\mu\cdot x-\beta|x|(|x|-1)-\sum_i\eta_i x_i(x_i-1),
 \qquad \beta,\eta_i\ge0,
\]
with the same one-dimensional merging convention as above.
\end{corollary}
\begin{proof}
The upper bound gives a bounded-operator positive \(C_0\)-semigroup on \(\ell^1\).
In \eqref{pos:canonical}, every nondiagonal parameter must vanish because its
monomial matrix entry is nonnegative and cannot cancel another off-diagonal
entry.  The remaining diagonal terms are exactly the displayed expression.
The converse is Theorem~\ref{pos:symbol}.
\end{proof}

\begin{example}[The square-symbol test is insufficient]\label{pos:psd-counterexample}
The diagonal contraction generated by
\(A=-z_1^2\partial_1^2-4z_1z_2\partial_1\partial_2-4z_2^2\partial_2^2\)
has logarithmic symbol \(-(a_1z_1+2a_2z_2)^2\).
Nevertheless, at \(t=(\log2)/2\), it sends the stable polynomial
\((z_1+z_2)(1+z_1)\) to
\(z_1+z_2+z_1^2/2+z_1z_2/4\), which vanishes at
\((z_1,z_2)=(-3+i,2+6i)\).
Thus positive semidefiniteness of the preliminary damping matrix is strictly
weaker than the common-damping condition in Theorem~\ref{pos:symbol}.
By contrast, \(E(\partial_1+2\partial_2)-2E(E-1)\) is a positive contraction
preserving stability: its column sum at \(z^x\) is
\(-x_1(|x|-1)\le0\), and it has the form \eqref{pos:canonical}.
It represents unequal collective loss together with killing.
\end{example}

\subsection{Finite capacities and quadratic contacts}

Now let every \(\kappa_i\) be finite and positive, put
\(V_\kappa=\operatorname{span}\{z^x:x\le\kappa\}\), and let
\(J=\{i:\kappa_i=1\}\).  The numerical capacities enter the boundary factors of
the rates.  The intrinsic contact domain remembers exactly the binary coordinates.
With \(e=\boldsymbol1\), define \(\mathfrak Q_J\) by
\begin{equation}\label{pos:contact-domain}
 \begin{gathered}
 y\in\R^d,\quad H=H^{\mathsf T},\quad H_{ij}\ge0,
 \quad H_{ii}=0\ (i\in J),\quad e^{\mathsf T}He=1,\\
 (He)(He)^{\mathsf T}-H\succeq0,\qquad
 Hy\le0,\qquad y^{\mathsf T}Hy\ge0.
 \end{gathered}
\end{equation}
Here \(H_{ij}\ge0\) is entrywise nonnegativity, and is not a positive-semidefinite
condition on \(H\).  Under the normalization, the displayed matrix inequality
is equivalent to \(H\) having exactly one positive eigenvalue.  Indeed, decompose
\(u=(e^{\mathsf T}Hu)e+v\), with \(e^{\mathsf T}Hv=0\).  A quadratic form with
one positive direction is nonpositive on this orthogonal complement, which gives
\(u^{\mathsf T}Hu\le(e^{\mathsf T}Hu)^2\); the converse follows by restriction.

\begin{theorem}[Finite-capacity classification]\label{pos:finite}
A real operator \(A:V_\kappa\to V_\kappa\) generates a positive semigroup preserving
\(\mathcal S_+\cap V_\kappa\) if and only if its only off-diagonal entries are
\begin{equation}\label{pos:finite-rates}
\begin{aligned}
 A_{x-e_i,x}&=x_i\left[\delta_i+\alpha_i(x_i-1)
                         +\sum_{j\ne i}a_{ij}x_j\right],\\
 A_{x+e_i,x}&=(\kappa_i-x_i)\left[b_i+\chi_i(\kappa_i-x_i-1)
                         +\sum_{j\ne i}g_{ij}(\kappa_j-x_j)\right],\\
 A_{x-e_i+e_j,x}&=c_{ij}x_i(\kappa_j-x_j)\qquad(i\ne j),
\end{aligned}
\end{equation}
all the displayed parameters are nonnegative, and the diagonal is
\begin{equation}\label{pos:diagonal-rate}
 A_{x,x}=\omega+\sum_i\mu_ix_i+\sum_i\tau_ix_i(x_i-1)
                      +\sum_{i<j}\theta_{ij}x_ix_j,
\end{equation}
with real coefficients.  Set \(\alpha_i=\chi_i=\tau_i=0\) for \(i\in J\).
Writing the second-order part as \(\sum_{i,j}Q_{ij}(z)\partial_i\partial_j\),
where \(Q\) is symmetric, its entries are
\begin{equation}\label{pos:principal-matrix}
\begin{aligned}
 Q_{ii}(z)&=\alpha_i z_i+\tau_i z_i^2+\chi_i z_i^3,\\
 2Q_{ij}(z)&=a_{ij}z_j+a_{ji}z_i+\theta_{ij}z_iz_j
 -c_{ij}z_j^2-c_{ji}z_i^2
 +g_{ij}z_i^2z_j+g_{ji}z_iz_j^2.
\end{aligned}
\end{equation}
The additional necessary and sufficient condition is
\begin{equation}\label{pos:contact-test}
 \operatorname{tr}(Q(y)H)\le0\qquad((y,H)\in\mathfrak Q_J).
\end{equation}
Within the structural family \eqref{pos:finite-rates}--\eqref{pos:diagonal-rate},
every failure of preservation has a quadratic stable input witnessing failure
at first order in time.
\end{theorem}

We give the structural and geometric parts of the proof separately.

\begin{lemma}[Boundary factors and finite rate reduction]\label{pos:finite-structure}
Every positive stability-preserving generator on \(V_\kappa\) has the rate form
\eqref{pos:finite-rates}--\eqref{pos:diagonal-rate}, with the stated signs.
\end{lemma}
\begin{proof}
Lemma~\ref{pos:order} gives order at most two.  For a net shift \(\sigma\), write
\(p_\sigma(x)=[z^{x+\sigma}]Az^x\).  Normal ordering gives a polynomial of total
degree at most two in \(x\), represented in reduced degree on every finite
coordinate grid.  Put \(r_i=(-\sigma_i)_+\), \(s_i=(\sigma_i)_+\).  The rate
vanishes when a source has fewer than \(r_i\) particles or fewer than \(s_i\)
holes.  Reduced interpolation on each boundary slice makes these vanishings
polynomial identities.  Successive division gives
\begin{equation}\label{pos:boundary-factors}
 p_\sigma(x)=\prod_i(x_i)_{r_i}(\kappa_i-x_i)_{s_i}\,h_\sigma(x),
 \qquad\deg h_\sigma\le2-\sum_i(r_i+s_i).
\end{equation}
A shift with no admissible source has zero reduced polynomial.  Thus a nonzero
shift has \(\|\sigma\|_1\le2\).

To exclude double births, set \(g_t(s)=(e^{tA}1)(s,\ldots,s)\).
It has nonnegative coefficients, negative real zeros, and \(g_0=1\).
The coefficient of \(s^2\) in \(\log g_t\) is nonpositive.  Its derivative at
zero is \([s^2](A1)(s,\ldots,s)\), the sum of the nonnegative double-birth
rates at the empty state.  Each is zero.  Formula~\eqref{pos:boundary-factors}
shows that a double-birth rate is a constant times its two hole factors, so
it vanishes everywhere.  The involution
\(\mathcal Jf(z)=z^\kappa f(1/z)\) preserves nonnegative coefficients and real
stability: reciprocal upper-half-plane variables lie in the lower half-plane,
where a real stable polynomial is also nonvanishing.  Conjugating by \(\mathcal J\)
therefore excludes double deaths.  Only single births, single deaths, transfers,
and a quadratic diagonal remain.

The single-death affine factor has a unique expansion in
\(1,x_i-1,x_j\ (j\ne i)\), and the birth factor has a unique expansion in the
corresponding hole coordinates; ineffective binary terms are set to zero.
The transfer boundary factors exhaust the degree budget, giving \(c_{ij}\ge0\).
The pair test \((z_i+s)(z_j+t)\) at \((-s,-t)\) gives
\[
 c_{ji}s^2+c_{ij}t^2+a_{ji}s+a_{ij}t-\theta_{ij}st
 +g_{ij}s^2t+g_{ji}st^2\ge0\qquad(s,t>0).
\]
Let \(s\downarrow0\) and then \(t\downarrow0\) to obtain \(a_{ij}\ge0\);
exchange indices for \(a_{ji}\).  Let \(s\to\infty\) with \(t\) fixed, then
\(t\to\infty\), to obtain \(g_{ij}\ge0\); exchange indices again.
For \(\kappa_i\ge2\), the square test gives
\(\alpha_i-\tau_i s+\chi_i s^2\ge0\) for \(s>0\), so
\(\alpha_i,\chi_i\ge0\).
A state with one particle at site \(i\) and no others gives \(\delta_i\ge0\).
A state with exactly one hole at \(i\) and no other holes gives \(b_i\ge0\).
These are all asserted signs.  Direct normal ordering of the rates yields
\eqref{pos:principal-matrix}.
\end{proof}

\begin{lemma}[Quadratic witnesses]\label{pos:quadratic-witness}
For \((y,H)\in\mathfrak Q_J\), the polynomial
\(f_{y,H}(z)=(z-y)^{\mathsf T}H(z-y)\) belongs to
\(\mathcal S_+\cap V_\kappa\), and
\(Af_{y,H}(y)=2\operatorname{tr}(Q(y)H)\).
\end{lemma}
\begin{proof}
Nonnegative quadratic, linear, and constant coefficients are respectively the
conditions \(H\ge0\) entrywise, \(-2Hy\ge0\), and \(y^{\mathsf T}Hy\ge0\).
The binary restrictions are exactly \(H_{ii}=0\).
For \(v>0\), \(v^{\mathsf T}Hv>0\).  Since \(H\) has one positive eigenvalue,
its form is nonpositive on the \(H\)-orthogonal complement of \(v\).
If \((u+iv)^{\mathsf T}H(u+iv)=0\), the imaginary part makes \(u\) belong
to this complement, whereas the real part requires
\(u^{\mathsf T}Hu=v^{\mathsf T}Hv>0\), a contradiction.
Thus \(z^{\mathsf T}Hz\), and its real translate, are stable.
At \(y\) both the value and the gradient vanish, leaving precisely the asserted
second-order trace.  Lemma~\ref{pos:root-contact} proves necessity of
\eqref{pos:contact-test}; a strict violation proves the witness assertion.
\end{proof}

For sufficiency we need a precise finite-dimensional boundary argument.  We first
record explicitly the homogeneous phase property used in its localizations.

\begin{lemma}[Homogeneous common phase]\label{pos:phase}
A nonzero homogeneous complex stable polynomial is a nonzero complex scalar times
a polynomial with nonnegative real coefficients.
\end{lemma}
\begin{proof}
Let its degree be \(m\), and fix \(e>0\).  Homogeneity gives
\(P(ie)=i^mP(e)\ne0\).  For every real \(x\), the degree-\(m\) polynomial
\(P(x+te)\) has no upper-half-plane zero by stability and no lower-half-plane zero,
since \(P(x+te)=(-1)^mP(-x-te)\).  Dividing by its leading coefficient \(P(e)\)
gives a monic real-rooted polynomial.  In particular \(P(x)/P(e)\) is real for
every real \(x\), so \(P_0=P/P(e)\) has real coefficients.  It does not vanish
on the positive orthant, and equals one at \(e\), hence is positive there.
Fix positive real values of all coordinates except one.  The resulting
univariate polynomial is real-rooted, positive on the positive half-line,
and has positive leading coefficient.  Its roots are nonpositive, so all its
coefficients are nonnegative.  Each of these coefficient polynomials in the
remaining variables is homogeneous and stable or zero, by differentiation and
real specialization, and is nonnegative on the remaining positive orthant.
Induction on the number of variables therefore makes every coefficient of
\(P_0\) nonnegative.
\end{proof}

\begin{lemma}[Projective boundary and interior density]\label{pos:boundary}
Let \(K=(\mathcal S_+\cap V_\kappa)\cup\{0\}\), and let \(O\) be the open set
of polynomials with every allowed coefficient strictly positive.
A stable section of multidegree \(\kappa\) on \((\mathbb{CP}^1)^d\) having no
real projective singularity is interior to the real-stable set.
Moreover \(K=\overline{\operatorname{int}K\cap O}\).
\end{lemma}
\begin{proof}
Use real M\"obius charts of positive determinant at infinity; they preserve the
upper half-plane.  Suppose a stable section has a zero with a nonempty set \(I\)
of real coordinates and a nonempty complementary set \(J\) strictly in the upper
half-plane.  Real specialization and Hurwitz's theorem imply
\(f(x_I,z_J)\equiv0\).
Let \(h(u_I,z_J)\) be the first nonzero homogeneous term of
\(f(x_I+u_I,z_J)\) in \(u_I\).  Positive rescaling and coefficient limits show
that \(h\) is jointly stable.  For each \(z_J\in\HH^J\), its homogeneous
coefficients in \(u_I\) have one common phase.  Every nonzero such coefficient
is itself stable, by derivative and specialization closure, and is nonvanishing
on \(\HH^J\).  Ratios of any two of them are therefore holomorphic and real-valued
on this open connected set, hence constant.  Polynomial continuation gives
\(h(u_I,z_J)=h_0(u_I)g(z_J)\), where \(h_0\) is real homogeneous up to a phase
and \(g\) is a real stable section up to the opposite phase.
If the order in \(u_I\) is at least two, every point of this real slice is a
singularity.  If it is one, the first derivatives along the slice are multiples
of \(g\).  This section has a real projective zero: if it depends genuinely on
some coordinate, specialize the others generically to real values and use a
nonconstant real-rooted polynomial; if a prescribed positive degree is missing,
use infinity in that coordinate.  At such a zero all first derivatives of \(f\)
vanish.  Thus a partially nonreal boundary zero forces a real singularity.

At an all-real regular zero, linear localization shows that the nonzero gradient
components have the same sign.  If the components in a nonempty set \(J\) vanish,
the slice in those coordinates must vanish identically.  To justify this last
assertion, otherwise choose a positive direction \(v_J\) on which the slice is
not identically zero and a coordinate \(i\) with nonzero partial derivative.
The real specialization
\(q(s,t)=f(p+s e_i+t v_J)\) is stable,
\(q_s(0,0)\ne0\), and \(q(0,t)=ct^m+O(t^{m+1})\), with \(m\ge2\), \(c\ne0\).
The analytic implicit function theorem gives a zero branch
\(s=\phi(t)=-(c/q_s(0,0))t^m+O(t^{m+1})\).
For an appropriate ray \(t=\rho e^{i\vartheta}\), \(0<\vartheta<\pi\),
its leading term has positive imaginary part.  For small \(\rho>0\), both
\(s\) and \(t\) lie in the upper half-plane, contradicting stability.
The slice therefore vanishes identically and the preceding slice argument
again supplies a real projective singularity.

In the absence of singularities, every real zero thus has all gradient components
nonzero and of the same sign.  Choose that sign \(\sigma\) in a small complex
coordinate neighborhood, so that \(\Re(\sigma\partial_i f)>0\) there.
The inequality persists for sufficiently small real coefficient perturbations \(g\).
For real \(u\) and positive \(v\) within this neighborhood,
\[
 \Im(\sigma g(u+iv))=
 \sum_i v_i\int_0^1\Re(\sigma\partial_i g(u+itv))\,dt>0.
\]
Outside these neighborhoods the section is nonzero on the product of closed upper
hemispheres, by the exclusion of partially nonreal zeros.  Compactness gives a
finite cover and a positive lower bound on the complementary compact set.
Every sufficiently small real perturbation is therefore stable.  This proves
the first assertion.

To prove density, apply small independent birth--death flows, with strictly positive
birth and death rates, to each finite site.  Their two-state M\"obius matrices have
positive entries and positive determinant, so the induced degree-\(\kappa_i\)
action preserves stability.  Their transition matrices are strictly positive at
positive times, since every state communicates through nearest-neighbor moves.
Thus every nonzero member of \(K\) is approximated by members of \(K\cap O\).
For such an \(f\), apply \(D_s=e^{-sE(E-1)}\), which preserves positive stability
by Theorem~\ref{pos:symbol}'s finite-cutoff proof and leaves the box invariant.
For any real \(x\), the radial polynomial \(r\mapsto f(rx)\) is real-rooted.
Indeed, positive homogenization \(F(w,z)=w^Nf(z/w)\) is stable, and
\(t\mapsto F((0,x)+t(1,\varepsilon,\ldots,\varepsilon))\) is real-rooted for
\(\varepsilon>0\).  Its limit \(F(t,x)\) has nonzero leading coefficient
\(F(1,0)=f(0)>0\); reciprocity gives the assertion for \(f(rx)\).

Put \(p_s(r)=(D_sf)(rx)\).  It remains real-rooted and satisfies
\(\partial_s p_s=-r^2p_s''\).
Every nonzero root is simple when \(s>0\).  Otherwise let \(r_0\ne0\) have
multiplicity \(m\ge2\) at an interior time \(s_0>0\), and factor the nearby root
cluster as
\(u^m+a_1(s)u^{m-1}+a_2(s)u^{m-2}+\cdots\), \(u=r-r_0\).
Contour power sums make its coefficients differentiable.  Dividing the evolution
identity by the nonvanishing complementary factor shows that the derivative
of this monic factor has no terms below \(u^{m-2}\), and that
\(a_2'(s_0)=-r_0^2m(m-1)\).
The sum of squares of the real root displacements is \(a_1^2-2a_2\), nonnegative
and zero at \(s_0\), but with derivative \(2r_0^2m(m-1)>0\), a contradiction.
A finite critical zero of \(D_sf\) would give a repeated radial root at \(r=1\),
so none exists.  At a projective face with coordinates \(I\) at infinity, let
\(m=\sum_{i\in I}\kappa_i\) and let \(g\) be the leading-coefficient face of \(f\).
The corresponding face of \(D_sf\) is
\[
 e^{-sm(m-1)}(D_sg)(e^{-2sm}z_{I^c}).
\]
The face \(g\) is stable, has all coefficients positive and has positive constant
coefficient; the same radial argument excludes a critical zero there.  A zero
with all coordinates at infinity is excluded by the positive top coefficient.
The first assertion makes \(D_sf\) an interior point.  Letting the two smoothing
times tend to zero, and finally scaling a positive interior polynomial to zero,
proves the density assertion.
\end{proof}

\begin{lemma}[Ordinary boundary contacts]\label{pos:ordinary}
Let \(M=\dim V_\kappa\).  Except for a semialgebraic subset of dimension at most
\(M-2\), every point of \(\partial K\cap O\) has a unique real projective
singularity.  It is finite, has nonsingular Hessian, and admits a local defining
function
\[
 h(g)=\sigma g(y(g)),\qquad K\text{ locally }=\{h\le0\},
\]
where \(y(g)\) is its continued critical point and \(\sigma\in\{1,-1\}\).
At the original point \(f\), the matrix
\(H=\sigma\operatorname{Hess}f(y)/2\), after positive normalization, belongs
to \(\mathfrak Q_J\).
\end{lemma}
\begin{proof}
The univariate capacity-one case is immediate: the strictly positive affine
polynomials are interior.  In all other cases, for fixed projective \(p\), the
conditions \(f(p)=0\), \(df(p)=0\) impose \(d+1\) independent linear conditions.
First jets are surjective since every capacity is positive.  The incidence space
therefore has dimension \(M-1\).
Hessian degeneracy imposes a further nontrivial polynomial equation in each fiber:
the allowed local quadratic \(\sum_{i<j}u_iu_j\) has invertible Hessian for
\(d\ge2\); in one variable use \(u^2\).
A projective coordinate at infinity likewise reduces the dimension of the point
parameter by at least one.  At a finite nondegenerate singularity, the implicit
function theorem continues the critical point as \(y(g)\); the differential of
its critical value is evaluation at \(y\).  Evaluations at two distinct finite
points are linearly independent, as is seen on \(1,z_1,\ldots,z_d\).
Two distinct ordinary singularities thus lie on transverse critical-value
hypersurfaces and impose codimension at least two.
All these sets are projections of real algebraic incidence sets in finitely many
projective charts, hence semialgebraic.  Their dimension bounds give finite
smooth stratifications, each covered by countably many charts of dimension at
most \(M-2\).  This proves the exceptional-set assertion.

At a remaining boundary point, Lemma~\ref{pos:boundary} supplies its singularity.
The quadratic localization of the stable positive homogenization at \((1,y)\)
is, with its common sign chosen,
\((v-sy)^{\mathsf T}H(v-sy)\).
The homogeneous common-phase property and quadratic stability yield
\(H\ge0\) entrywise, \(Hy\le0\), \(y^{\mathsf T}Hy\ge0\), and exactly one
positive eigenvalue.  The box imposes \(H_{ii}=0\) for binary coordinates.
Its positive normalization is therefore in \(\mathfrak Q_J\).

The Hessian remains invertible nearby, and compactness excludes other projective
singularities outside the chosen neighborhood.  Lemma~\ref{pos:boundary} implies
that the local boundary is contained in \(h(g)=\sigma g(y(g))=0\).
Here \(dh\ne0\), since its differential is signed evaluation at \(y\).
If \(h(g)>0\), a positive-direction restriction through \(y(g)\) has a positive
critical value and a positive quadratic leading coefficient near the former
double root; its two nearby zeros are nonreal.  That side is unstable.
Interior density from Lemma~\ref{pos:boundary} gives stable interior points on
the side \(h<0\).  Shrink to an implicit-function chart in which this side is
connected.  Membership in \(K\) cannot change within the side without another
boundary point, and every boundary point has \(h=0\).
Thus the whole side \(h<0\) is stable.  Closedness then includes \(h=0\), giving
exactly the stated local half-space description.
\end{proof}

\begin{lemma}[Integration through ordinary boundary points]\label{pos:integration}
Let \(K\subset\R^M\) be closed, \(O\) open, and
\(K=\overline{\operatorname{int}K\cap O}\).
Suppose a complete ambient flow generated by a \(C^1\) vector field leaves
\(O\) forward invariant.
Suppose the boundary in \(O\), except for countably many \(C^1\) images of
dimension at most \(M-2\), has \(C^2\) local defining functions \(h\) with
\(dh\ne0\) and \(K=\{h\le0\}\).
Then the flow preserves \(K\) if and only if \(dh(V)\le0\) at every such
ordinary boundary point, where \(V\) is its vector field.
\end{lemma}
\begin{proof}
Necessity is differentiation from a boundary point.
For sufficiency fix a time horizon \(T\).  If \(\Sigma\) is the exceptional set,
the initial points whose trajectories meet \(\Sigma\) before time \(T\) lie in
the union of the images
\((t,u)\mapsto\Phi_{-t}(\psi(u))\), \(0\le t\le T\), where \(\psi\) ranges
over its charts.  Each parameter space has dimension at most \(M-1\), and
countable compact covers make these images Lipschitz images of lower-dimensional
parameter sets relative to \(M\)-dimensional volume.  Their union has measure zero.
Only the ambient inverse flow is used; \(O\) need not be backward invariant.

In a regular chart, the ordinary boundary points are dense in its smooth
hypersurface, since \(\Sigma\) has zero \((M-1)\)-dimensional measure.
Continuity therefore extends \(dh(V)\le0\) to that entire hypersurface.
A local projection \(\pi\) to it satisfies
\(\|x-\pi(x)\|\le C|h(x)|\).
The function \(a(x)=dh(x)V(x)\) is locally Lipschitz, so for \(h(x)\ge0\),
\(a(x)\le LC h(x)\).
Along a trajectory the positive part of \(h\) consequently satisfies the upper
right derivative inequality \(D^+h_+\le LC h_+\).
Gronwall's inequality forbids a crossing from \(h\le0\) to \(h>0\).
Every initial point of \(\operatorname{int}K\cap O\) outside the null hitting set
therefore remains in \(K\) until \(T\), because a first exit would occur in a
regular chart.  Such initial points are dense in \(K\).
Continuity of the flow and closedness of \(K\) finish the proof, since \(T\)
was arbitrary.
\end{proof}

\begin{proof}[Completion of the proof of Theorem~\ref{pos:finite}]
Necessity was proved in Lemmas~\ref{pos:finite-structure} and
\ref{pos:quadratic-witness}.  Conversely the rate conditions make \(A\) Metzler.
Its ambient flow \(e^{tA}\) is complete, and the strictly positive coefficient
orthant is forward invariant, since
\((e^{tA}f)_x\ge e^{tA_{x,x}}f_x>0\) for \(f\in O\).
Lemmas~\ref{pos:boundary} and \ref{pos:ordinary} verify the geometric hypotheses
of Lemma~\ref{pos:integration}.
At an ordinary boundary point, \(f(y)=0\) and \(\nabla f(y)=0\), so
\[
 dh_f(Af)=\sigma(Af)(y)=2\operatorname{tr}(Q(y)H)\le0.
\]
The normalization of \(H\) changes only a positive factor.
Lemma~\ref{pos:integration} now proves preservation of \(K\).
Invertibility of the finite matrix exponential excludes zero outputs from
nonzero inputs.  This proves sufficiency.
\end{proof}

\begin{remark}[A compact certificate]\label{pos:compact}
Write \(Q=Q^{[1]}+Q^{[2]}+Q^{[3]}\) by homogeneous degree and put
\(\widehat Q(r,y)=r^2Q^{[1]}(y)+rQ^{[2]}(y)+Q^{[3]}(y)\).
Condition~\eqref{pos:contact-test} is equivalent to
\(\operatorname{tr}(\widehat Q(r,y)H)\le0\) under
\eqref{pos:contact-domain}, \(r\ge0\), and \(r^2+\|y\|^2=1\).
For \(r>0\), this is the identity \(\widehat Q(r,y)=r^3Q(y/r)\);
the face \(r=0\) follows by continuity.  The set is compact, since
\(H\ge0\) entrywise and \(e^{\mathsf T}He=1\) bound every matrix entry.
The number of auxiliary variables is \(O(d^2)\), independently of the numerical
capacities.  This does not remove the capacities from the actual rates.
\end{remark}

\subsection{Mixed finite and infinite capacities}

Let \(F=\{i:\kappa_i<\infty\}\), \(U=\{i:\kappa_i=\infty\}\), and
\(J=\{i:\kappa_i=1\}\).  Let \(\mathcal{LP}_{+,\kappa}\) denote the positive
Laguerre--P\'olya class supported on \(\Omega_\kappa\).

\begin{theorem}[Mixed-capacity classification]\label{pos:mixed}
Let \(T_t\) be a coefficientwise-positive \(C_0\)-semigroup of bounded operators
on \(\ell^1(\Omega_\kappa)\) with generator \(A\).
Preservation of \(\mathcal{LP}_{+,\kappa}\) at every time is equivalent to
preservation for every allowed stable polynomial input, and to the following
structural assertion together with \eqref{pos:contact-test}:
all allowed polynomials belong to \(D(A)\), the only off-diagonal entries are
\begin{equation}\label{pos:mixed-rates}
\begin{aligned}
 A_{x-e_i,x}&=x_i\left[\delta_i+\alpha_i(x_i-1)+\sum_{j\ne i}a_{ij}x_j\right],\\
 A_{x+e_i,x}&=(\kappa_i-x_i)\Bigl[b_i+\chi_i(\kappa_i-x_i-1)\\
 &\hspace{3em}+\sum_{j\in F\setminus\{i\}}g_{ij}(\kappa_j-x_j)\Bigr] &&(i\in F),\\
 A_{x+e_j,x}&=\lambda_j &&(j\in U),\\
 A_{x-e_i+e_j,x}&=c_{ij}x_i(\kappa_j-x_j) &&(j\in F,\ i\ne j),\\
 A_{x-e_i+e_j,x}&=r_{ij}x_i &&(j\in U,\ i\ne j).
\end{aligned}
\end{equation}
Every displayed off-diagonal parameter is nonnegative.  The diagonal is
\eqref{pos:diagonal-rate}.  In \eqref{pos:principal-matrix} and
\eqref{pos:contact-domain} use the conventions
\[
 \begin{gathered}
 \alpha_i=\tau_i=0\ (i\in J),\qquad
 \chi_i=0\ (i\in U\cup J),\\
 g_{ij}=0\ \text{unless }i,j\in F,\qquad
 c_{ij}=0\ \text{unless }j\in F.
 \end{gathered}
\]
Entries with targets outside \(\Omega_\kappa\) are omitted.
The existence criterion for the formal operator is separate and is given in
Theorem~\ref{ana:weighted-generation}.
\end{theorem}
\begin{proof}
Assume polynomial-input preservation.  Jensen approximation, positive
coefficient closure, and the strictly positive diagonal entries of \(T_t\),
proved in Theorem~\ref{pos:symbol}, extend it to all
\(\mathcal{LP}_{+,\kappa}\).
Corollary~\ref{ana:polynomial-domain} supplies the polynomial domain and the
local first variation needed for
Lemma~\ref{pos:order}.  The order is at most two.
The growth-reduction theorem, Theorem~\ref{ana:growth-reduction}, applies with
Lemma~\ref{pos:poisson}: it identifies the bounded part that strictly increases
the total occupation of the unrestricted sites as
\[
 R=\sum_{j\in U}\lambda_jz_j+
   \sum_{i\in F,\,j\in U}r_{ij}z_j\partial_i,
 \qquad \lambda_j,r_{ij}\ge0,
\]
and supplies the locally uniform first variation at arbitrary real points.
In particular, \(L=A-R\) cannot increase that unrestricted occupation.
These analytic statements precede the signed contacts used below.

For completeness, we give the remaining algebraic reduction.
For a net shift \(\sigma\), the rate is a reduced polynomial of total degree
at most two.  Boundary interpolation, exactly as in
\eqref{pos:boundary-factors}, gives
\begin{equation}\label{pos:mixed-boundaries}
 p_\sigma(x)=\prod_i(x_i)_{(-\sigma_i)_+}
 \prod_{i\in F}(\kappa_i-x_i)_{(\sigma_i)_+}\,h_\sigma(x),
 \quad
 \deg h_\sigma\le2-\sum_i(-\sigma_i)_+-\sum_{i\in F}(\sigma_i)_+.
\end{equation}
For unrestricted coordinates, vanishing on every integer point is a polynomial
identity; for finite ones, use reduced interpolation.  Thus at most two source
or finite-target boundary factors can occur.
Double finite births have constant times two hole factors and are excluded by
\(A1\): its simultaneous two-birth coefficient is nonnegative but the first
variation of the second logarithmic coefficient along a positive diagonal is
nonpositive.  Hence that coefficient is zero.

For every allowed multi-index \(\alpha\) of total degree two, apply the product-root
test to \(\prod_i(z_i+s_i)^{\alpha_i}\).  It gives
\[
 P_\alpha(-s_{\supp\alpha},z_{\mathrm{spectators}})\le0
\]
for every real spectator vector.  We classify all contributions to this complete
coefficient polynomial, including catalytic one-particle shifts.
Quadratic monomials in unrestricted destinations outside the source
have nonnegative coefficients, by \eqref{pos:mixed-boundaries} and positivity.
Set these spectators equal and let them tend to positive infinity.  Their
quadratic coefficients must vanish.  A remaining free destination enters linearly.
For two distinct source sites its coefficient has the form
\[
 d+t_i z_i+t_jz_j-b_i z_i^2-b_jz_j^2;
\]
for a repeated source it has the form \(d+t_i z_i\).
These are all possibilities because \eqref{pos:mixed-boundaries} leaves at most
two total factors: a two-to-one replacement gives the constant \(d\),
a source-dependent transfer gives a term \(t_i z_i\), and a simultaneous finite
birth gives its negative hole term \(-b_i z_i^2\).
The destination can have either real sign, so its coefficient vanishes identically
on the negative source orthant.  Coefficient comparison removes every such
replacement, every source-self or third-site catalytic transfer, and every
simultaneous finite birth and unrestricted transfer.  A pure double death is the
remaining nonnegative constant in \(P_\alpha\); sending the negative source roots
to zero removes it as well.

A transfer to a finite site already has its two boundary factors and is
\(c_{ij}x_i(\kappa_j-x_j)\), \(c_{ij}\ge0\).
After the preceding eliminations a transfer to an unrestricted site can only be
\(x_i(r_{ij}+h_{ij}x_j)\), with \(r_{ij},h_{ij}\ge0\).
The pair contact at \((z_i,z_j)=(-s,-t)\) contains the term
\(h_{ij}t^2\) in its second-order coefficient, while the other remaining terms
are at most linear in \(t\) for fixed \(s\), apart from a possible reverse
finite-target contribution \(-c_{ji}s^2\).
Letting \(t\to\infty\) forces \(h_{ij}=0\).
Finally a finite birth with unrestricted catalyst has rate
\(h(\kappa_i-x_i)x_\ell\), \(h\ge0\), and contributes
\(-h z_i^2z_\ell\) to the mixed principal coefficient.
At \((-s,-t)\), divide the contact inequality by \(t\) and let \(t\to\infty\).
Its leading dependence on \(s\) is \(hs^2\); the remaining terms have degree
at most one.  Letting \(s\to\infty\) forces \(h=0\).
The surviving shifts are precisely \eqref{pos:mixed-rates}.
The small-root, large-root, and boundary-state arguments in
Lemma~\ref{pos:finite-structure} give all their asserted signs, also when the
other occupancies are unrestricted.  The diagonal is the reduced polynomial
of degree at most two given by normal ordering.
Lemma~\ref{pos:quadratic-witness} proves necessity of the contact inequality.

Conversely, impose artificial capacities \(N_j\ge2\) at the unrestricted sites.
Keep the original diagonal, and replace immigration and transfer into those
sites by
\(\lambda_j(1-x_j/N_j)\) and
\(r_{ij}x_i(1-x_j/N_j)\), respectively.
The principal trace changes by
\[
 \operatorname{tr}(Q_N(y)H)-\operatorname{tr}(Q(y)H)
 =-\sum_{j\in U,\,i\ne j}\frac{r_{ij}}{N_j}y_j^2H_{ij}\le0.
\]
The immigration correction is first order and the binary set is unchanged.
Every approximant therefore satisfies Theorem~\ref{pos:finite}.
Theorem~\ref{ana:killed-approximation} proves convergence of these finite flows,
started from their coordinate projections, to the given semigroup; its proof
uses the bounded increasing part and weighted estimates, not an unbounded
similarity on ordinary \(\ell^1\).
Theorem~\ref{ana:mixed-sufficiency} and Lemma~\ref{pos:lp-closure} pass preservation
to all \(\mathcal{LP}_{+,\kappa}\), and exclude zero outputs.
This proves all implications.
\end{proof}

\subsection{One-site cones, conservation, and exact obstructions}

\begin{corollary}[One finite site]\label{pos:one-site}
For capacity \(K\ge2\), the exact positive stability-preserving cone is
\[
 Az^n=d_nz^{n-1}+u_nz^{n+1}+v_nz^n,
\]
where
\[
 d_n=n[\delta+\alpha(n-1)],\quad
 u_n=(K-n)[b+\chi(K-n-1)],\quad
 v_n=c+\mu n+\tau n(n-1),
\]
\(\delta,b,\alpha,\chi\ge0\), \(c,\mu\in\R\), and
\(\tau\le2\sqrt{\alpha\chi}\).
For \(K=1\), every Metzler \(2\times2\) matrix qualifies.
For every \(K\ge2\), the cone modulo \(I,z\partial_z\) is nonpolyhedral.
\end{corollary}
\begin{proof}
In one variable the normalized contact has \(H=1\), \(y\le0\).
The condition \(\alpha y+\tau y^2+\chi y^3\le0\) is therefore
\(\alpha-\tau s+\chi s^2\ge0\) for all \(s>0\).
With \(\alpha,\chi\ge0\), this is equivalent to
\(\tau\le2\sqrt{\alpha\chi}\), including either zero parameter by limits.
For \(K=1\), every nonzero affine polynomial with nonnegative coefficients is
stable, and every positive invertible matrix exponential preserves this class.
For nonpolyhedrality, for each \(u>0\) the functional
\[
 \alpha/u+u\chi-\tau+\delta+b
\]
is nonnegative on the quotient cone.  Equality forces \(\delta=b=0\),
\(\alpha=u^2\chi\), and \(\tau=2u\chi\), by the arithmetic--geometric mean
inequality and the cone inequality.  Its zero set is exactly the ray generated
by \((\alpha,\chi,\tau)=(u^2,1,2u)\).
These distinct exposed rays, for every \(u>0\), exclude polyhedrality.
\end{proof}

\begin{corollary}[Conserving mixed systems]\label{pos:conserving}
A finite-exit, particle-number-conserving Markov generator on \(\Omega_\kappa\)
preserves all allowed positive stable polynomials if and only if its only rates are
\[
 q(x,x-e_i+e_j)=
 \begin{cases}
 c_{ij}x_i(\kappa_j-x_j),&i,j\in F,\quad c_{ij}=c_{ji}\ge0,\\
 r_{ij}x_i,&j\in U,\quad r_{ij}\ge0,
 \end{cases}
\]
and there is no transfer from an unrestricted to a finite site.
The preservation quantifier ranges over the entire allowed polynomial space,
including mixtures of different total particle numbers.
\end{corollary}
\begin{proof}
Every fixed-total state space is finite, so the chain is nonexplosive.
Conservation and nonnegative off-diagonal entries remove all births and deaths
from Theorem~\ref{pos:mixed}.
For two finite sites the pair part of the principal coefficient is
\[
 (c_{ij}+c_{ji})z_i z_j-c_{ij}z_j^2-c_{ji}z_i^2.
\]
The contact at \((-s,-t)\) requires
\((t-s)(c_{ij}t-c_{ji}s)\ge0\) for every \(s,t>0\).
Taking \(t\) on either side of \(s\) forces \(c_{ij}=c_{ji}\).
For \(i\in U,j\in F\), a rate \(c_{ij}x_i(\kappa_j-x_j)\) gives the pair
coefficient \(c_{ij}(z_i z_j-z_j^2)\).  Taking \(z_j=-s,z_i=-t\) and
\(t\to\infty\) forces \(c_{ij}=0\).
No second-order term restricts independent transfer into an unrestricted site.

For sufficiency, a finite pair contributes
\begin{equation}\label{pos:exclusion-generator}
 G_{ij}^\kappa=(z_j-z_i)(\kappa_j\partial_i-\kappa_i\partial_j)
                  -(z_j-z_i)^2\partial_i\partial_j.
\end{equation}
Normalized polarization replaces each site by \(\kappa_i\) binary slots.
Summing the swap generators over every cross-block pair of slots intertwines
with \(G_{ij}^\kappa\): on a state with occupancies \(x_i,x_j\), exactly
\(x_i(\kappa_j-x_j)\) swaps move a particle from \(i\) to \(j\).
A binary swap \(\tau\) has flow
\[
 e^{t(\tau-I)}=\frac{1+e^{-2t}}2I+\frac{1-e^{-2t}}2\tau,
\]
which preserves stability by partial symmetrization \cite{BBL09}.
Polarization, specialization, and finite-dimensional Lie products prove
preservation for finite exclusion.  Independent transfers into unrestricted
sites have the affine-substitution flow from Theorem~\ref{pos:symbol}.
On each fixed total-degree cutoff all these operators act in finite dimension,
so a final Lie product combines them.
\end{proof}

\begin{corollary}[Reversibility and conserving time-one embeddings]\label{pos:embedding}
On a fully finite box put \(w(x)=\prod_i\binom{\kappa_i}{x_i}\) and
\(D=\diag(\sqrt{w(x)})\).
Every generator in Corollary~\ref{pos:conserving} is reversible with respect to
\(w\).  If its site graph is connected, its unique stationary law at total
population \(n\) is
\(w(x)/\binom{\sum_i\kappa_i}{n}\).
A stochastic, number-conserving operator \(T\) is the time-one map of a positive
stability-preserving Markov semigroup if and only if
\(S=D^{-1}TD\) is symmetric positive definite and
\[
 D(\log S)D^{-1}=\sum_{\{i,j\}}c_{ij}G_{ij}^\kappa,
 \qquad c_{ij}\ge0.
\]
The embedding generator in this class is unique.
\end{corollary}
\begin{proof}
For \(y=x-e_i+e_j\), the binomial identities give
\[
 \frac{w(y)}{w(x)}=
 \frac{x_i}{\kappa_i-x_i+1}\,
 \frac{\kappa_j-x_j}{x_j+1}.
\]
This is exactly detailed balance for the symmetric rates.  Connectivity permits
moving particles along site paths until any prescribed occupancy of the same
total is reached, so the fixed-total chain is irreducible.  The normalization
is the coefficient of \(s^n\) in \(\prod_i(1+s)^{\kappa_i}\).

If \(T=e^A\) with \(A\) Markov and positive, choose \(a\) so that \(B=A+aI\)
is entrywise nonnegative.  If an entry of \(e^A\) between different total layers
is zero, every term in \(e^{-a}\sum_{m\ge0}B^m/m!\) at that entry is zero,
in particular the corresponding off-diagonal entry of \(A\).
Thus time-one conservation forces generator conservation.  The previous
corollary applies, and \(D^{-1}AD\) is symmetric.  Its exponential is symmetric
positive definite, and its unique symmetric logarithm is \(\log S\).
This proves necessity and uniqueness.  Conversely the displayed logarithm is
a nonnegative combination of the conserving generators
\eqref{pos:exclusion-generator}; its exponential is \(T\) and the previous
corollary proves preservation.
\end{proof}

For two binary sites, \(T_p=(1-p)I+p\tau\) preserves stability for all
\(0\le p\le1\), but its conserving Markov embedding exists exactly for
\(0\le p<1/2\).  Indeed the antisymmetric eigenvalue is \(1-2p\), and the
embedding generator is
\(-\tfrac12\log(1-2p)(\tau-I)\).
Furthermore the one-particle rates of any finite conserving system obey every
cycle identity: the product of forward rates along a labeled cycle equals
the product along its reverse, since \(a_{ij}=c_{ij}\kappa_j\).
Pointwise limits preserve these identities.  The independent directed
three-cycle on unrestricted sites violates them.  Thus unrestricted conserving
systems need not be dilute limits of finite conserving systems on the same
labeled sites.  This does not apply to the killed approximants used in
Theorem~\ref{pos:mixed}.

\begin{corollary}[Support restriction for universal conserving mixing]\label{pos:support-mixing}
On a full binary box, suppose a universally positive stability-preserving
conserving Markov flow leaves a fixed-size family \(\mathcal B\) invariant
and is irreducible on it.  Then, for a partition into site-graph components
\(C_1,\ldots,C_r\) and integers \(k_1,\ldots,k_r\),
\[
 \mathcal B=\{S:|S\cap C_\ell|=k_\ell\ \text{for every }\ell\}.
\]
\end{corollary}
\begin{proof}
Corollary~\ref{pos:conserving} makes the flow symmetric exclusion on a site graph.
Every edge transposition must preserve \(\mathcal B\), since a positive rate
cannot leave an invariant family.  Edge transpositions generate the full
symmetric group of each connected component.  Their orbits are exactly the
sets with prescribed component occupancies.  Irreducibility makes
\(\mathcal B\) one such orbit.
\end{proof}

\begin{proposition}[Exact compensation for a capped transfer]\label{pos:compensation}
Let \(z\) have finite capacity \(K\ge1\), and let \(w\) be unrestricted.
Add to common coalescence of rate \(\gamma\ge0\) the transition
\((m,n)\to(m+1,n-1)\) at rate \(cn(K-m)\), where \(c\ge0\).
The resulting Markov flow preserves positive stability if and only if
\(c\le2\gamma\).
\end{proposition}
\begin{proof}
The additional differential operator is
\(cK(z-w)\partial_w+cz(w-z)\partial_z\partial_w\).
For a contact matrix \(H=\left(\begin{smallmatrix}h&a\\a&k\end{smallmatrix}\right)\)
and \(y=(z,w)\), common coalescence has principal trace
\(\gamma(e^{\mathsf T}Hy-y^{\mathsf T}Hy)\le0\).
At \(c=2\gamma\), the total trace simplifies exactly to
\[
 \gamma e^{\mathsf T}Hy-\gamma(h+2a)z^2-\gamma k w^2\le0.
\]
The endpoints \(c=0\) and \(c=2\gamma\) therefore pass every contact, and
linearity gives the entire interval.  Theorem~\ref{pos:mixed} proves sufficiency.
For necessity take \(h=k=0,a=1/2\), \(z=-s,w=-t\), \(s,t>0\).
Twice the trace is
\(-\gamma(s+t)-(2\gamma-c)st-cs^2\).
If \(c>2\gamma\), choose \(s>\gamma/(c-2\gamma)\) and then sufficiently large
\(t\); the trace is positive.  This is an admissible contact, including when
\(K=1\), and proves necessity.
For example \(K=1,\gamma=1,c=3\), \(f=(z+2)(w+15)\), and \(t=\log2\) give
\[
 T_tf=30+\frac{139}{8}z+\frac38w+\frac14zw,
 \qquad
 \frac{139}{8}\frac38-30\frac14=-\frac{63}{64}<0.
\]
The last inequality is the bivariate multiaffine stability discriminant.
\end{proof}

\begin{proposition}[Capacity and unchanged coordinates]\label{pos:spectators}
The operator \(A=(z-z^2)\partial_z^2\) preserves positive stability on
\(V_{(2,1)}\), with the binary coordinate \(w\) unchanged.
It fails on \(V_{(2,1,1)}\) with two unchanged binary coordinates.
Moreover there is no positive stability-preserving generator \(\widetilde A\)
on three binary variables \((u,v,w)\) satisfying the exact compression
\(D\widetilde A P=A\), where
\[
 Pz=(u+v)/2,\quad Pz^2=uv,\quad Pw=w,
 \qquad Dg(z,w)=g(z,z,w).
\]
\end{proposition}
\begin{proof}
For one unchanged coordinate, the contact matrix is
\(H=\left(\begin{smallmatrix}h&a\\a&0\end{smallmatrix}\right)\), \(h,a\ge0\).
If \(h>0\), the condition \(Hy\le0\) forces \(y_z\le0\): use the second
row when \(a>0\), and the first row when \(a=0\).
Thus \(\operatorname{tr}(Q(y)H)=h(y_z-y_z^2)\le0\); if \(h=0\) it is zero.
Theorem~\ref{pos:finite} proves preservation.
With two unchanged binary variables, the stable input
\((z+w_1)(z+w_2)\) evolves into
\[
 \theta z^2+z(w_1+w_2)+w_1w_2+(1-\theta)z,
 \qquad\theta=e^{-2t}.
\]
Identify \(w_1=w_2=w\) and specialize \(z=r\in(0,1)\).
Its discriminant in \(w\) is
\(-4(1-\theta)r(1-r)<0\) for \(t>0\), proving failure.

Suppose the asserted lift exists.  Average it with the swap of \(u,v\).
The positive preserving generator cone is convex, by finite-dimensional Lie
products, and this averaging preserves the compression identity.
The averaged lift commutes with the swap.  Since \(PD\) is the projection onto
its symmetric subspace and \(D\) kills the antisymmetric subspace, we have
\(D\widetilde A=D\widetilde A PD=AD\).
Evaluation at all ones makes \(\widetilde A\) Markov, because \(A\) is Markov.
Positivity and this lumping identity forbid any change of \(w\), any jump from
zero \(u,v\)-occupancy, or any occupancy change from a singleton.
From the double occupancy \(uv\), symmetry forces deaths to \(u\) and \(v\)
at rate one each.  The only remaining transitions are swaps of the singleton
states, at rates \(r_0,r_1\ge0\) according to the value of \(w\).

For \(F_s=(u+w)(v+s)\), direct application gives at
\((u,v,w)=(2s,-s,-2s)\)
\[
 (\widetilde A F_s)(2s,-s,-2s)
 =s+(4-3r_0-6r_1)s^2.
\]
Put \(c=3r_0+6r_1\) and \(s=1/[2(1+c)]\).  This value is
\((c+6)/[4(c+1)^2]>0\), although both nonnegative affine factors vanish at
the chosen real point.  Lemma~\ref{pos:root-contact} gives a contradiction.
\end{proof}

For unrestricted \(z\)-coordinates, the Markov flow
\eqref{pos:markov-generator} with \(\lambda=0\) preserves all joint positive
stable inputs after adjoining one unrestricted unchanged coordinate if and only
if \(\gamma=0\).  Sufficiency is affine substitution.  For necessity, apply
Corollary~\ref{pos:markov} in the enlarged system: the common quadratic rate
must vanish because the unchanged coordinate has zero second-degree action.
It is enough to test \((w+a\cdot z)^m\), since positive scaling and translation
of \(w\) supply all enlarged multinomial tests.
This quantifier is different from the one-binary-spectator assertion above.


\begin{example}[The semigroup assumption in multinomial testing]\label{pos:isolated-tests}
Multinomial tests do not characterize arbitrary positive operators.
Define a stochastic linear map from two variables to one by
\[
 K(z_1z_2)=(1+w^2)/2,\qquad K(z^\alpha)=w\quad(\alpha\ne(1,1)).
\]
For \(F=(c+az_1+bz_2)^m\), put \(M=(c+a+b)^m\),
\(U=[z_1z_2]F\).  Then
\(KF=U(1+w^2)/2+(M-U)w\).
For \(m\ge2\),
\[
 M\ge\binom m2c^{m-2}(a+b)^2
 \ge 2m(m-1)c^{m-2}ab=2U.
\]
For \(m<2\), \(U=0\).  Thus the quadratic discriminant is
\(M(M-2U)\ge0\), and every such output is a nonzero positive stable polynomial.
Nevertheless the stable monomial \(z_1z_2\) is sent to a polynomial vanishing at
\(w=i\).
Even adjoining all monomial tests does not suffice if zero outputs are allowed.
The projection retaining only
\(z_1z_2,z_1z_4,z_2z_3,z_3z_4\) sends every multinomial, for \(m\ge2\), to
\[
 m(m-1)c^{m-2}(a_1z_1+a_3z_3)(a_2z_2+a_4z_4),
\]
and sends every monomial to a stable monomial or zero.
It sends \((z_1+z_2)(z_3+z_4)\) to \(z_1z_4+z_2z_3\), which vanishes at
\(z_1=z_4=1+i\), \(z_2=z_3=-1+i\).
These examples concern isolated maps and do not contradict
Corollary~\ref{pos:markov}.
\end{example}


\begin{corollary}[A finite-degree witness under an order bound]\label{pos:finite-witness}
Suppose a particle-number-conserving Markov generator on unrestricted sites is
known to have differential order at most \(R<\infty\).  If it is not an
independent-particle migration generator, then some multinomial of degree at
most \(\max(2,R)\) fails stability preservation at arbitrarily small positive
times.
\end{corollary}
\begin{proof}
If every first variation at a multinomial root contact of degree
\(3,\ldots,R\) passed Lemma~\ref{pos:root-contact}, coefficient comparison as
in Lemma~\ref{pos:order} would remove all derivatives of order greater than two.
Degree conservation already supplies the output-degree bound.
If all degree-two square contacts also passed, the positivity and mass arguments
in Corollary~\ref{pos:markov} would give common coalescence and independent
migration.  Conservation removes coalescence.  This contradicts the hypothesis.
Thus one of the stated root-contact conditions fails; its local root argument
produces nonreal zeros for arbitrarily small positive times.
\end{proof}

For example, the positive conserving binary-consensus generator
\(A=(z_1-z_2)^2\partial_1\partial_2\) sends \((z_1+z_2)^2\) at time \(t\) to
\((2-\theta)(z_1^2+z_2^2)+2\theta z_1z_2\), where \(\theta=e^{-2t}\).
Specializing \(z_2=1\) gives discriminant \(16(\theta-1)<0\) at every
positive time.

% END SECTION positive_capacities

\appendix
% BEGIN SECTION hermite
\section{Hermite compression and boundary localization}
\label{herm:section}

The Hermite--Sylvester criterion represents real-rootedness by a positive
semidefinite matrix of Newton sums; see, for example, \cite{Nathanson21}.
At a polynomial with repeated roots, the differential of this matrix must
be tested only on its kernel.  We show that this necessary condition is
also sufficient for coefficient tangency, at every root multiplicity.
The proof identifies the compressed matrix in independent root-jet
coordinates and constructs a real-rooted curve realizing each admissible
direction.  Root-cluster tangent descriptions belong to the variational
theory developed in \cite{BurkeOverton01,BurkeEaton12}; the matrix
formulation also records exactly when semidefinite normal multipliers
have a closed image.

\subsection{Newton sums and the compressed differential}

Fix an integer $n\geq1$.  Let
\[
 \mathcal R_n=
 \left\{z^n+a_1z^{n-1}+\cdots+a_n:
                  \text{all roots are real}\right\}.
\]
We regard $\mathcal R_n$ as a closed subset of the affine space of real
monic degree-$n$ polynomials.  Its ambient tangent space is
$V_{n-1}=\R[z]_{\leq n-1}$, with the Euclidean coefficient inner product
in the basis $1,z,\ldots,z^{n-1}$.  The tangent cone throughout this section
is the contingent cone
\begin{equation}
 T_{\mathcal R_n}(p)=
 \left\{h\in V_{n-1}:\ 
   \frac{p_j-p}{t_j}\longrightarrow h
   \text{ for some }p_j\in\mathcal R_n,
   \ t_j>0,\ t_j\downarrow0\right\}.
 \label{herm:tangent-definition}
\end{equation}
For any monic real polynomial $p$ of degree $n$, let $s_k(p)$ be the sum
of the $k$th powers of its roots, counted with multiplicity, and set
$s_0(p)=n$.  Newton's identities make every $s_k$ a polynomial function
of the coefficients of $p$.  Define
\begin{equation}
 H(p)=\bigl(s_{i+j}(p)\bigr)_{0\leq i,j\leq n-1}.
 \label{herm:matrix}
\end{equation}
Thus $H(p)$ has size $n$, whereas the vector space of polynomials of
degree at most $n$ has dimension $n+1$.  The derivative $DH_p[h]$ is
taken in the monic affine space, so that $\deg h<n$ and $Ds_0[h]=0$.

\begin{lemma}[Hermite representation]
\label{herm:representation}
A monic real degree-$n$ polynomial $p$ belongs to $\mathcal R_n$ if and
only if $H(p)\succeq0$.  If
\begin{equation}
 p(z)=\prod_{a\in\mathcal A}(z-a)^{q_a},
 \qquad \mathcal A\subset\R,\qquad
 q_a\geq1,\qquad \sum_{a\in\mathcal A}q_a=n,
 \label{herm:factorization}
\end{equation}
then, upon identifying a coefficient vector with its polynomial $v$,
\begin{equation}
 v^{\mathsf T}H(p)v=\sum_{a\in\mathcal A}q_a v(a)^2,
 \qquad
 \ker H(p)=\{v\in V_{n-1}:v(a)=0\text{ for every }a\in\mathcal A\}.
 \label{herm:kernel}
\end{equation}
\end{lemma}

\begin{proof}
Expanding the quadratic form in \eqref{herm:matrix} gives the sum of
$v(\lambda)^2$ over all roots $\lambda$ of $p$, counted with their
multiplicities.  This proves \eqref{herm:kernel} and the forward
implication.  Conversely, suppose that $p$ has a nonreal root $\lambda$.
Its conjugate has the same multiplicity.  Interpolate the values $i$ and
$-i$ at $\lambda$ and $\overline\lambda$, respectively, and the value
zero at every other distinct root.  The interpolating polynomial has
degree less than the number of distinct roots, hence less than $n$.
Conjugating its coefficients preserves the prescribed data, so uniqueness
of interpolation shows that this polynomial has real coefficients.
Its Hermite quadratic form equals minus twice the multiplicity of
$\lambda$, contradicting positive semidefiniteness.
\end{proof}

\begin{theorem}[Exact Hermite tangent criterion]
\label{herm:tangent}
For every $p\in\mathcal R_n$, put $K_p=\ker H(p)$.  Then
\begin{equation}
 \boxed{\displaystyle
 T_{\mathcal R_n}(p)=
 \left\{h\in V_{n-1}:
       \left.DH_p[h]\right|_{K_p}\succeq0\right\}.}
 \label{herm:tangent-formula}
\end{equation}
The restriction denotes compression of the quadratic form to $K_p$;
positive semidefiniteness of the entire matrix $DH_p[h]$ is not required.
Every direction on the right of \eqref{herm:tangent-formula} is realized
by a curve $p_t\in\mathcal R_n$, $t\geq0$, whose coefficients are
polynomials in $t$ and which satisfies $p_0=p$ and $p'_0=h$.
\end{theorem}

We first compute the compressed differential.  For a fixed polynomial
$F$, write $L_F(p)$ for the sum of $F$ over the roots of $p$, with
multiplicities.  This is again a polynomial coefficient functional.

\begin{lemma}[Residue formula and independent jet blocks]
\label{herm:jet-blocks}
Let $p$ have the factorization \eqref{herm:factorization}, and let
$h\in V_{n-1}$.  Write its unique partial-fraction expansion as
\begin{equation}
 \frac{h(z)}{p(z)}=
       \sum_{a\in\mathcal A}\sum_{\ell=1}^{q_a}
          \frac{b_{a,\ell}}{(z-a)^\ell}.
 \label{herm:principal-parts}
\end{equation}
The map $h\mapsto(b_{a,\ell})$ is a linear isomorphism onto $\R^n$.
For every fixed polynomial $F$,
\begin{equation}
 DL_F(p)[h]=
       -\sum_{a\in\mathcal A}
         \operatorname{Res}_{z=a}\frac{F'(z)h(z)}{p(z)}.
 \label{herm:residue}
\end{equation}
In the independent jet coordinates
\begin{equation}
 v(a+w)=\sum_{r\geq1}u_{a,r}w^r,
 \qquad
 (u_{a,1},\ldots,u_{a,q_a-1})_{a\in\mathcal A},
 \label{herm:jet-coordinates}
\end{equation}
on $K_p$, the form $\left.DH_p[h]\right|_{K_p}$ is the direct sum of
the blocks
\begin{equation}
 M^{(a)}_{rs}(b)=
 \begin{cases}
       -(r+s)b_{a,r+s},&r+s\leq q_a,\\
       0,&r+s>q_a,
 \end{cases}
 \qquad 1\leq r,s\leq q_a-1.
 \label{herm:blocks}
\end{equation}
A simple root contributes a block of size zero.
\end{lemma}

\begin{proof}
The rational function $h/p$ is proper.  Existence and uniqueness of
partial fractions give \eqref{herm:principal-parts}.  Conversely, any
choice of its $n$ real coefficients yields a polynomial of degree less
than $n$ after multiplication by $p$.  This proves the first assertion.

Choose a positively oriented circle enclosing all roots of $p$ and all
roots of its sufficiently small coefficient perturbations.  The argument
principle gives
\[
 L_F(p)=\frac{1}{2\pi i}\int F(z)\frac{p'(z)}{p(z)}\,dz.
\]
Differentiation in the coefficient direction $h$, followed by integration
by parts on the closed contour, gives
\[
 DL_F(p)[h]
 =\frac{1}{2\pi i}\int F(z)
                   \left(\frac{h(z)}{p(z)}\right)'dz
 =-\frac{1}{2\pi i}\int F'(z)\frac{h(z)}{p(z)}\,dz.
\]
Taking residues proves \eqref{herm:residue}.  This calculation does not
require differentiable labels for repeated roots.

The map assigning to $v\in V_{n-1}$ all Taylor coefficients of orders
$0,\ldots,q_a-1$ at each $a$ is an isomorphism.  Indeed, a polynomial in
its kernel is divisible by $p$ and has degree less than $n$, so it is
zero; the source and target both have dimension $n$.  By
\eqref{herm:kernel}, restriction to $K_p$ sets precisely the zeroth
coefficient at each root to zero.  This proves the independence in
\eqref{herm:jet-coordinates}.

Apply \eqref{herm:residue} to $F=v^2$.  Since $v(a)=0$,
\[
 (v^2)'(a+w)=
       \sum_{r,s\geq1}(r+s)u_{a,r}u_{a,s}w^{r+s-1}.
\]
The residue against \eqref{herm:principal-parts} is consequently
\[
 v^{\mathsf T}DH_p[h]v
 =-\sum_{a\in\mathcal A}
       \sum_{\substack{r,s\geq1\\r+s\leq q_a}}
         (r+s)b_{a,r+s}u_{a,r}u_{a,s}.
\]
These are exactly the blocks \eqref{herm:blocks}.
\end{proof}

\begin{lemma}[Positive semidefinite block elimination]
\label{herm:block-elimination}
For a root of multiplicity $q\geq2$, the block in
\eqref{herm:blocks} is positive semidefinite if and only if
\begin{equation}
 b_\ell=0\quad(3\leq\ell\leq q),
 \qquad b_2\leq0.
 \label{herm:block-condition}
\end{equation}
There is no condition on $b_1$.
\end{lemma}

\begin{proof}
Suppose the block is positive semidefinite and some $b_\ell$ with
$\ell\geq3$ is nonzero.  Choose the largest such index $\ell$.
The diagonal entry in row $r=\ell-1$ is zero: its coefficient has index
$2\ell-2>\ell$ and is either outside the range of the block or zero by
maximality.  For a positive semidefinite matrix, a zero diagonal entry
forces its row and column to vanish.  For example, positivity of each
$2$-by-$2$ principal minor gives
$|M_{rs}|^2\leq M_{rr}M_{ss}=0$.  But the entry in column $1$ of this
row is $-\ell b_\ell\ne0$, a contradiction.  Once all $b_\ell$ with
$\ell\geq3$ vanish, the only potentially nonzero entry is
$M_{11}=-2b_2$.  This proves both implications, including $q=2$.
\end{proof}

\begin{proof}[Proof of Theorem~\ref{herm:tangent}]
Suppose $p_j=p+t_jh+o(t_j)\in\mathcal R_n$, with $t_j\downarrow0$.
For any $v\in K_p$, Lemma~\ref{herm:representation} and polynomial
differentiability of $H$ give
\[
 0\leq v^{\mathsf T}H(p_j)v
       =t_jv^{\mathsf T}DH_p[h]v+o(t_j).
\]
Division by $t_j$ proves necessity.

For sufficiency, assume that the compression is positive semidefinite.
Lemmas~\ref{herm:jet-blocks} and \ref{herm:block-elimination} give
\eqref{herm:block-condition} independently at every repeated root.
For a root $a$ of multiplicity $q=q_a\geq2$, put $w=z-a$ and set
\begin{equation}
 P_{a,t}(z)=
 \left[\left(w+\frac{b_{a,1}}{q}t\right)^2+b_{a,2}t\right]
 \left(w+\frac{b_{a,1}}{q}t\right)^{q-2}.
 \label{herm:realizing-factor}
\end{equation}
For a simple root set $P_{a,t}(z)=z-a+b_{a,1}t$.  Each factor is monic
and has only real roots for every $t\geq0$.  In
\eqref{herm:realizing-factor}, these roots are
\[
 a-\frac{b_{a,1}}{q}t\pm\sqrt{-b_{a,2}t},
 \qquad
 a-\frac{b_{a,1}}{q}t
       \quad\text{with multiplicity }q-2.
\]
When $q=2$, the final factor in \eqref{herm:realizing-factor} is $1$.
Direct differentiation of its polynomial coefficients gives
\[
 \left.\frac{d}{dt}P_{a,t}(z)\right|_{t=0}
       =b_{a,1}(z-a)^{q-1}+b_{a,2}(z-a)^{q-2}
       \qquad(q\geq2).
\]
Now define $p_t=\prod_{a\in\mathcal A}P_{a,t}$.  This is a monic,
real-rooted degree-$n$ polynomial for every $t\geq0$, and its
coefficients are polynomials in $t$.  By the product rule,
\[
 p'_0=p\left(
       \sum_{a\in\mathcal A}\frac{b_{a,1}}{z-a}
       +\sum_{\substack{a\in\mathcal A\\q_a\geq2}}
                       \frac{b_{a,2}}{(z-a)^2}\right)=h,
\]
where the last equality is \eqref{herm:principal-parts} together with
\eqref{herm:block-condition}.  This supplies the claimed curve and
proves sufficiency.  The case $n=1$ has no nonempty jet block; the
simple-root construction proves the assertion in that case as well.
\end{proof}

\begin{corollary}[Root-contact form]
\label{herm:root-contact}
For a root $a$ of $p$ with multiplicity $q_a\geq2$, write
$p=(z-a)^{q_a}R_a$ with $R_a(a)\ne0$.  A direction $h\in V_{n-1}$
belongs to $T_{\mathcal R_n}(p)$ if and only if, at every such root,
\begin{equation}
 h^{(j)}(a)=0\quad(0\leq j\leq q_a-3),
 \qquad
 \frac{h^{(q_a-2)}(a)}{(q_a-2)!R_a(a)}\leq0.
 \label{herm:root-contact-formula}
\end{equation}
Simple roots impose no condition, and the list of equalities is empty
when $q_a=2$.
\end{corollary}

\begin{proof}
The coefficient $b_{a,\ell}$ is the Taylor coefficient of order
$q_a-\ell$ in $h/R_a$ at $a$.  Vanishing of all the coefficients with
$\ell\geq3$ is equivalent to divisibility of $h$ by $(z-a)^{q_a-2}$.
Under that divisibility, the next coefficient is
$b_{a,2}=h^{(q_a-2)}(a)/((q_a-2)!R_a(a))$.
Apply \eqref{herm:block-condition}.
\end{proof}

\subsection{Normal multipliers and their closedness}

Use the coefficient inner product on $V_{n-1}$ and the trace inner
product on real symmetric matrices to define the adjoint $DH_p^*$.
The regular normal cone is
\[
 \widehat N_{\mathcal R_n}(p)=
 \left\{\eta\in V_{n-1}:
 \limsup_{\substack{q\to p,\ q\in\mathcal R_n\\q\ne p}}
       \frac{\langle\eta,q-p\rangle}{\|q-p\|}\leq0\right\}.
\]
For a cone $C$, write
$C^\circ=\{h:\langle h,c\rangle\leq0\text{ for all }c\in C\}$.
In finite dimensions the definition of the contingent cone implies
$\widehat N_{\mathcal R_n}(p)=T_{\mathcal R_n}(p)^\circ$.
Indeed, normality gives the inequality on every contingent limit;
conversely, a sequence violating normality has a convergent subsequence
of normalized displacement vectors, producing a contingent direction
with strictly positive pairing.

\begin{theorem}[Exact normal formula and closed multiplier images]
\label{herm:closed}
For $p\in\mathcal R_n$, define
\begin{equation}
 C_p=\left\{-DH_p^*(Y):
       Y\succeq0,\ \operatorname{range}Y\subseteq K_p\right\}.
 \label{herm:multiplier-image}
\end{equation}
Then
\begin{equation}
 \widehat N_{\mathcal R_n}(p)=\overline{C_p}.
 \label{herm:normal-closure}
\end{equation}
The following conditions are equivalent:
\begin{enumerate}
 \item $C_p$ is closed;
 \item every root of $p$ has multiplicity at most two;
 \item there exists $h\in V_{n-1}$ such that
       $\left.DH_p[h]\right|_{K_p}$ is positive definite on $K_p$.
\end{enumerate}
Positive definiteness on the zero-dimensional space is interpreted
vacuously.  In particular, the closure in
\eqref{herm:normal-closure} cannot be omitted at a polynomial having a
root of multiplicity at least three.
\end{theorem}

\begin{proof}
For any $h\in V_{n-1}$,
\begin{align*}
 h\in C_p^\circ
 &\quad\Longleftrightarrow\quad
   -\Tr\bigl(YDH_p[h]\bigr)\leq0
      \text{ for all }Y\succeq0
      \text{ with }\operatorname{range}Y\subseteq K_p\\
 &\quad\Longleftrightarrow\quad
   \left.DH_p[h]\right|_{K_p}\succeq0.
\end{align*}
The last equivalence follows by testing rank-one matrices $vv^{\mathsf T}$,
$v\in K_p$, and then using the spectral decomposition of $Y$.
Theorem~\ref{herm:tangent} gives $C_p^\circ=T_{\mathcal R_n}(p)$.
The set $C_p$ is a convex cone, so the finite-dimensional bipolar
theorem gives
\[
 \widehat N_{\mathcal R_n}(p)
       =T_{\mathcal R_n}(p)^\circ
       =(C_p^\circ)^\circ=\overline{C_p}.
\]
This proves \eqref{herm:normal-closure} without a strict-feasibility
assumption.

We next describe the multiplier coordinates precisely.  Let
$E=\bigoplus_{a\in\mathcal A}\R^{q_a-1}$ be the jet space in
\eqref{herm:jet-coordinates}, and let $L:E\to V_{n-1}$ be the inverse
jet interpolation map with image $K_p$.  Lemma~\ref{herm:jet-blocks}
says
\begin{equation}
 L^{\mathsf T}DH_p[h]L
       =M(b):=\bigoplus_{a\in\mathcal A}M^{(a)}(b).
 \label{herm:congruence}
\end{equation}
Every positive semidefinite $Y$ with range in $K_p$ can be written
uniquely as $Y=LZL^{\mathsf T}$ with $Z\succeq0$ on $E$.
To see this, use a left inverse $L^+$ of $L$ and set
$Z=L^+Y(L^+)^{\mathsf T}$; conversely congruence preserves positive
semidefiniteness.  Consequently the normal functional associated with
$Y$ is, in the partial-fraction coordinates,
\begin{equation}
 h\longmapsto-\Tr(ZM(b))
       =-\sum_{a\in\mathcal A}\Tr\bigl(Z^{(a)}M^{(a)}(b)\bigr),
 \label{herm:dual-block-pairing}
\end{equation}
where $Z^{(a)}$ is the corresponding diagonal principal block of $Z$.
Each $Z^{(a)}$ is positive semidefinite.  Conversely, any tuple of such
positive semidefinite blocks is realized by their block-diagonal direct
sum.  Thus no orthogonality of the jet coordinates has been assumed,
and off-diagonal blocks between different roots do not affect the
multiplier image.

Suppose first that every $q_a\leq2$.  Every nonempty block $M^{(a)}$ is
the scalar $-2b_{a,2}$.  Formula \eqref{herm:dual-block-pairing} then
identifies $C_p$ with the cone of functionals
\[
 h\longmapsto
       \sum_{\substack{a\in\mathcal A\\q_a=2}}
                         2u_a b_{a,2},\qquad u_a\geq0.
\]
The $b$'s are independent linear coordinates by
Lemma~\ref{herm:jet-blocks}.  This cone is therefore a product of
nonnegative rays in a coordinate subspace, carried back by an
invertible linear map.  It is closed.  The same independence allows
one to choose $b_{a,2}=-1$ at every double root, producing a positive
definite compression.  If there are no double roots, $K_p=\{0\}$ and
$C_p=\{0\}$.

Now suppose that some root $a$ has $q_a\geq3$.  The functional
$\ell(h)=b_{a,3}$ vanishes on the tangent cone, by
\eqref{herm:block-condition}; hence $\ell\in\overline{C_p}$ by
\eqref{herm:normal-closure}.  If it belonged to $C_p$, represent it by
\eqref{herm:dual-block-pairing}.  Its coefficient at $b_{a,2}$ would be
$2Z^{(a)}_{11}=0$.  Positive semidefiniteness would then force the first
row of $Z^{(a)}$ to vanish.  But its coefficient at $b_{a,3}$ would be
$6Z^{(a)}_{12}=0$, instead of $1$.  Independence of the partial-fraction
coordinates prevents contributions from other roots from changing
either coefficient.  Thus $\ell\notin C_p$, and $C_p$ is not closed.
Finally, in this root block the diagonal entry with index $q_a-1$ is
identically zero because $2q_a-2>q_a$.  The compression can therefore
never be positive definite.  All three conditions are equivalent.
\end{proof}

\begin{example}[Diverging multipliers at a triple root]
\label{herm:cubic-multiplier}
Take $p=z^3$ and write $h=c_2z^2+c_1z+c_0$.  In the kernel basis
$z,z^2$, the compressed differential is
\[
 \begin{pmatrix}-2c_1&-3c_0\\-3c_0&0\end{pmatrix}.
\]
Thus the tangent conditions are $c_1\leq0$ and $c_0=0$, with $c_2$
free.  A multiplier
$Z=\begin{psmallmatrix}u&v\\v&w\end{psmallmatrix}\succeq0$
produces the functional $h\mapsto2u c_1+6v c_0$.  Its image consists of
all $A c_1+B c_0$ with $A>0$, together with the zero functional, whereas
its closure allows $A=0$ and arbitrary $B$.  For fixed $B\ne0$,
\[
 Z_\varepsilon=
 \begin{pmatrix}
   \varepsilon&B/6\\
   B/6&B^2/(36\varepsilon)
 \end{pmatrix}\succeq0
\]
has image $2\varepsilon c_1+B c_0$, which tends to $B c_0$ as
$\varepsilon\downarrow0$.  The multiplier norms diverge.  In fact,
every multiplier sequence representing this limiting functional must
be unbounded: a bounded sequence would have a convergent positive
semidefinite subsequence representing $B c_0$, contrary to the exact
image just computed.  Up to coordinate scaling, this is the elementary
nonclosed positive semidefinite projection in
\cite[Example~1]{JiangSturmfels21}.
\end{example}

\begin{remark}[Scope of the coefficient chart]
\label{herm:scope}
All tangent and normal formulas above use fixed degree and monic
normalization.  If $f_t$ has actual degree $n$ near $t=0$, has positive
leading coefficient $\lambda(t)$, and $f_0=\lambda(0)p$, its normalized
direction is
\[
 \left.\frac{d}{dt}\frac{f_t}{\lambda(t)}\right|_{t=0}
       =\frac{f'_0}{\lambda(0)}
        -\frac{\lambda'(0)}{\lambda(0)}p.
\]
This direction has degree less than $n$ and is the argument to use in
$DH_p$.  Degree drops and additional coefficient inequalities require
separate arguments.  In particular, the exact compression formula is
not a general rule asserting that tangents commute with arbitrary
intersections of polynomial or matrix constraints.
\end{remark}

\subsection{A finite-dimensional positive-flow criterion}

The monic formula also gives a criterion on a closed cone containing
polynomials of every actual degree.  The passage to lower degree uses
density rather than a change in the size of the Hermite matrix.
Write $V_m=\R[z]_{\leq m}$ and define
\[
 \mathcal K_m=\{0\}\cup
 \{f\in V_m\setminus\{0\}: f\text{ has nonnegative coefficients
                                      and only real roots}\}.
\]
A real linear map on $V_m$ is called Metzler when its matrix in the
ordinary coefficient basis has nonnegative off-diagonal entries.

\begin{theorem}[Hermite criterion for positive polynomial flows]
\label{herm:positive-flow}
Let $m\geq1$, and let $A:V_m\to V_m$ be Metzler.  Then
\[
 e^{tA}\mathcal K_m\subseteq\mathcal K_m\qquad(t\geq0)
\]
if and only if every monic degree-$m$ polynomial $p\in\mathcal R_m$
whose coefficients are all strictly positive satisfies
\begin{equation}
 \left.DH_p\bigl[Ap-([z^m]Ap)p\bigr]\right|_{\ker H(p)}\succeq0.
 \label{herm:positive-flow-test}
\end{equation}
The matrix in this test has size $m$; the bracketed direction has degree
less than $m$.
\end{theorem}

\begin{proof}
Let $\mathcal O\subset V_m$ be the open cone of strictly positive
coefficient vectors.  Metzler positivity implies, for $f\in\mathcal O$,
\begin{equation}
 (e^{tA}f)_i\geq e^{tA_{ii}}f_i>0\qquad(t\geq0).
 \label{herm:metzler-lower-bound}
\end{equation}
For example, choose $c$ so that $A+c\Id$ is entrywise nonnegative and
expand its exponential.  Keeping the paths that remain at index $i$
gives the stated diagonal lower bound.  Thus $\mathcal O$ is forward
invariant, and the leading coefficient of $e^{tA}p$ remains positive
for every polynomial tested in \eqref{herm:positive-flow-test}.

If the semigroup preserves $\mathcal K_m$, normalize this orbit by its
leading coefficient.  Its derivative at zero is
$Ap-([z^m]Ap)p$, and Theorem~\ref{herm:tangent} proves necessity.

Conversely, assume \eqref{herm:positive-flow-test}.  For
$f\in\mathcal K_m\cap\mathcal O$, write $f=sp$ with
$s=[z^m]f>0$ and $p$ monic, and put $a=[z^m]Ap$.
Theorem~\ref{herm:tangent} supplies a monic real-rooted curve $P_t$
with $P_0=p$ and $P'_0=Ap-ap$.  The curve
\[
 f_t=s(1+ta)P_t
\]
has derivative $Af$ at zero and lies in $\mathcal K_m\cap\mathcal O$
for all sufficiently small $t\geq0$.  Indeed its scalar factor and
all its coefficients are positive by continuity.  Hence
\begin{equation}
 Af\in T_{\mathcal K_m}(f)
       \qquad(f\in\mathcal K_m\cap\mathcal O).
 \label{herm:positive-local-tangent}
\end{equation}

We use the finite-dimensional contingent-cone form of Nagumo's theorem:
if a locally Lipschitz vector field takes its values in the contingent
tangent cone of a closed set at every point of that set, its unique
forward trajectories remain in the set throughout their existence;
see \cite{Nagumo42,GhorbalSogokonPlatzer17}.  To apply it using only
\eqref{herm:positive-local-tangent}, fix
$f\in\mathcal K_m\cap\mathcal O$ and a finite time $T>0$.
The known ambient linear orbit
\[
 \mathcal C=\{e^{tA}f:0\leq t\leq T\}
\]
is compact and lies in $\mathcal O$ by
\eqref{herm:metzler-lower-bound}.  Choose a nonnegative smooth function
$\chi$, compactly supported in $\mathcal O$, equal to one on a
neighborhood of $\mathcal C$.  The vector field
$B(x)=\chi(x)Ax$ is smooth with compact support, hence globally
Lipschitz.  At points of $\mathcal K_m\cap\mathcal O$ it is tangent
by \eqref{herm:positive-local-tangent}; at all other points of
$\mathcal K_m$ it vanishes.  The cone $\mathcal K_m$ is closed, as
verified below, so Nagumo's theorem makes it invariant under $B$.
The curve $e^{tA}f$ solves the same equation on $[0,T]$, since
$\chi=1$ along that curve.  Uniqueness therefore proves
$e^{tA}f\in\mathcal K_m$ throughout $[0,T]$.  This argument uses only
forward invariance of $\mathcal O$ and does not assume real-rootedness
of the ambient orbit when choosing $\chi$.

For completeness, nonnegative coefficients persist under coefficient
limits.  If a nonzero coefficient limit of real-rooted polynomials had
a nonreal zero, a small disk around that zero disjoint from the real
axis would, by local zero continuity, contain a zero of every
sufficiently close polynomial.  This is impossible.  Including the
zero polynomial proves that $\mathcal K_m$ is closed, including across
degree drops.

Finally, $\mathcal K_m\cap\mathcal O$ is dense in $\mathcal K_m$.
For a nonzero $f\in\mathcal K_m$ of actual degree $d$, all roots are
nonpositive, so
\[
 f(z)=c\prod_{j=1}^d(z+r_j),\qquad c>0,\quad r_j\geq0.
\]
For $\varepsilon>0$, the polynomials
\[
 f_\varepsilon(z)=
 c\prod_{j=1}^d(z+r_j+\varepsilon)
                 (1+\varepsilon z)^{m-d}
\]
have degree $m$, strictly positive coefficients, and only negative
roots, and tend coefficientwise to $f$.  Approximate zero by
$\varepsilon(1+z)^m$.  Continuity of $e^{tA}$ and closedness of
$\mathcal K_m$ now extend preservation to the entire cone.
\end{proof}

For $m=0$, the cone consists of nonnegative constants and every real
scalar generator preserves it.  No Hermite test is needed in that case.

\subsection{Exact localization at products of binary linear forms}

Fix integers $m\geq2$ and $r\geq1$, and write the coordinates on
$\R^{r+2}$ as $(T,X,Y)$, where $Y=(Y_1,\ldots,Y_r)$.
Let $\mathcal H_{m,r}$ be the set of homogeneous real polynomials of
degree $m$ that are monic in $T$ and hyperbolic with respect to
$e=(1,0,0)$.  Equivalently, $f\in\mathcal H_{m,r}$ if
$T\mapsto f(T,x,y)$ has only real roots for every $(x,y)\in\R^{r+1}$.
All tangent cones in this subsection are coefficient-space Bouligand
tangent cones.  Thus $h\in T_{\mathcal H_{m,r}}(p)$ means that there are
$\varepsilon_j\downarrow0$ and $p_j\in\mathcal H_{m,r}$ with
$p_j=p+\varepsilon_jh+o(\varepsilon_j)$ coefficientwise.  In particular,
$h$ is homogeneous of degree $m$ and has zero $T^m$ coefficient.

\begin{theorem}[Binary-factor localization]\label{herm:binary-localization}
Let $\lambda_1<\cdots<\lambda_m$ be real numbers, and set
\[
 \ell_i=T-\lambda_iX,\qquad
 p=\prod_{i=1}^m\ell_i,\qquad
 q_i=\frac{p}{\ell_i\ell_{i+1}}\quad(1\leq i<m).
\]
Then
\begin{equation}\label{herm:binary-tangent-formula}
\begin{aligned}
 T_{\mathcal H_{m,r}}(p)
 =\bigg\{&-\sum_{i=1}^m d_i(X,Y)\frac{p}{\ell_i}
           -\sum_{i=1}^{m-1}Q_i(Y)q_i(T,X):\\
         & d_i\in\R[X,Y]_1,\quad Q_i\in\R[Y]_2,
             \quad Q_i\geq0\bigg\}.
\end{aligned}
\end{equation}
Here the subscripts specify homogeneous degree, and $Q_i\geq0$
means positive semidefiniteness of the quadratic form on $\R^r$.
The forms $d_i$ and $Q_i$ in this representation are unique.  Consequently
the tangent cone, in its linear span, is linearly isomorphic to
\[
 \R^{m(r+1)}\times(\mathbb S_+^r)^{m-1},
\]
where $\mathbb S_+^r$ denotes the cone of real symmetric positive
semidefinite $r\times r$ matrices.  Every tangent direction has a realizing
curve $\varepsilon\mapsto p_\varepsilon\in\mathcal H_{m,r}$ whose
coefficients are polynomials in $\varepsilon\geq0$ and such that
$p_\varepsilon=p+\varepsilon h+O(\varepsilon^2)$.
\end{theorem}

The external representation theorem used in the necessity argument is
the following precise ternary form of the Helton--Vinnikov theorem:
if $f(T,X,Y)$ is homogeneous of degree $m$, $f(1,0,0)=1$, and $f$ is
hyperbolic with respect to $(1,0,0)$, then there are real symmetric
$m\times m$ matrices $A,B$ with
\begin{equation}\label{herm:ternary-determinant}
 f(T,X,Y)=\det(TI-XA-YB).
\end{equation}
This statement does not require irreducibility or smoothness.  The
normalization and the matrix size, including the case in which
$f(1,X,Y)$ has degree less than $m$, are established in
\cite[Eq.~(5) and pp.~2497--2498]{LewisParriloRamana05}.

\begin{proof}[Proof of Theorem~\ref{herm:binary-localization}]
We first prove necessity when $r=1$.  Choose a tangent sequence
$p_j=p+\varepsilon_jh+o(\varepsilon_j)$.
By \eqref{herm:ternary-determinant}, write
\[
 p_j(T,X,Y)=\det(TI-XA_j-YB_j),
\]
with $A_j,B_j$ real symmetric.  Orthogonally diagonalize $A_j$, applying
the same conjugation to $B_j$, and arrange its eigenvalues increasingly:
\[
 A_j=\diag(a_{1,j},\ldots,a_{m,j}).
\]
These eigenvalues are the roots of $p_j(T,1,0)$.  Since the roots of
$p(T,1,0)$ are distinct, the implicit function theorem for a simple root
and the first-order coefficient expansion give real numbers $\alpha_i$
such that
\begin{equation}\label{herm:binary-eigenvalue-expansion}
 a_{i,j}=\lambda_i+\varepsilon_j\alpha_i+o(\varepsilon_j).
\end{equation}

On setting $X=0$, the coefficients of $T^{m-1}Y$ and $T^{m-2}Y^2$
in $\det(TI-YB_j)$ show that
\[
 \Tr B_j=O(\varepsilon_j),\qquad e_2(B_j)=O(\varepsilon_j),
\]
where $e_2(B_j)$ is the second elementary symmetric function of the
eigenvalues.  Hence
\[
 \|B_j\|_{\mathrm F}^2
 =\Tr(B_j^2)=(\Tr B_j)^2-2e_2(B_j)=O(\varepsilon_j).
\]
Here symmetry is essential: it identifies $\Tr(B_j^2)$ with a sum of
squares.  Thus every entry of $B_j$ is $O(\sqrt{\varepsilon_j})$.

The coefficient of $Y$ in the determinant is exactly
\begin{equation}\label{herm:binary-linear-spectator}
 -\sum_{i=1}^m(B_j)_{ii}
       \prod_{k\ne i}(T-a_{k,j}X).
\end{equation}
The limiting forms $p/\ell_i$ are a basis of the binary forms of
degree $m-1$: evaluation at $(T,X)=(\lambda_i,1)$ annihilates all
but the $i$th form, whose value is nonzero.  Accordingly, the basis
matrices in \eqref{herm:binary-linear-spectator} converge to an
invertible matrix.  Since the coefficient of $Y$ in $p_j$ divided by
$\varepsilon_j$ converges coefficientwise, there are uniquely determined
real numbers $\beta_i$ such that
\begin{equation}\label{herm:binary-diagonal-expansion}
 (B_j)_{ii}=\varepsilon_j\beta_i+o(\varepsilon_j).
\end{equation}
After passing to a subsequence, the bounds on $B_j$ also give limits
\[
 \frac{(B_j)_{ik}}{\sqrt{\varepsilon_j}}\longrightarrow c_{ik}
 \qquad(1\leq i<k\leq m).
\]
No continuous choice of representations or diagonalizing matrices is
needed for these subsequence limits.

Expand the determinant by permutations.  The identity permutation,
together with \eqref{herm:binary-eigenvalue-expansion} and
\eqref{herm:binary-diagonal-expansion}, contributes
\[
 p-\varepsilon_j\sum_i(\alpha_iX+\beta_iY)\frac p{\ell_i}
 +o(\varepsilon_j).
\]
A permutation moving exactly two indices $i,k$ is a transposition;
its contribution to first order is
$-\varepsilon_jY^2c_{ik}^2p/(\ell_i\ell_k)$.
Every other nonidentity permutation moves at least three indices.
Its off-diagonal factors are each $O(\sqrt{\varepsilon_j})$, so its
contribution is $O(\varepsilon_j^{3/2})$.  Replacing a remaining
diagonal factor in a transposition term by its perturbation changes
that term by $O(\varepsilon_j^2)+o(\varepsilon_j^2)$.
These estimates are coefficientwise, since there are finitely many
permutations and all remaining linear factors have bounded coefficients.
We obtain
\begin{equation}\label{herm:binary-all-pairs}
 h=-\sum_{i=1}^m(\alpha_iX+\beta_iY)\frac p{\ell_i}
   -Y^2\sum_{i<k}c_{ik}^2\frac p{\ell_i\ell_k}.
\end{equation}
In particular, no term of $Y$-degree at least three occurs in $h$.

To reduce the pair terms to adjacent pairs, use the polynomial identity
\begin{equation}\label{herm:adjacent-telescoping}
 \frac p{\ell_i\ell_k}
 =\sum_{j=i}^{k-1}
    \frac{\lambda_{j+1}-\lambda_j}{\lambda_k-\lambda_i}\,q_j
 \qquad(i<k).
\end{equation}
For $X\ne0$ away from the finitely many vanishing factors, it follows
by telescoping
\[
 \frac1{\ell_{j+1}}-\frac1{\ell_j}
 =\frac{(\lambda_{j+1}-\lambda_j)X}{\ell_j\ell_{j+1}};
\]
multiplying by $p$ then proves the identity everywhere.
All coefficients in \eqref{herm:adjacent-telescoping} are positive.
Moreover, $q_1,\ldots,q_{m-1}$ are linearly independent: in a relation
$\sum_jb_jq_j=0$, evaluation at $(\lambda_1,1)$ forces $b_1=0$;
evaluation successively at $(\lambda_2,1),\ldots,(\lambda_{m-1},1)$
forces each remaining coefficient to vanish.  They therefore form a
basis of the binary forms of degree $m-2$.  Equation
\eqref{herm:binary-all-pairs} now proves the required necessity for
one spectator variable, with nonnegative coefficients of $Y^2q_j$.

For arbitrary $r$, restrict the tangent sequence to $Y=sy$, where
$y\in\R^r$ is fixed and $s$ is a single real variable.  Restriction
preserves monicity in $T$ and hyperbolicity in $e$, so the preceding
ternary conclusion applies to $h(T,X,sy)$ for every $y$.
Decompose $h$ into its components homogeneous in $Y$.
Every component of $Y$-degree at least three vanishes after all such
restrictions, and is therefore identically zero.
The $Y$-degree zero component is a binary form of degree $m$ with
zero $T^m$ coefficient; it is divisible by $X$ and consequently has a
unique expression $-\sum_i\alpha_iX p/\ell_i$.
Each coefficient of $Y_a$ has a unique expansion in the basis
$p/\ell_1,\ldots,p/\ell_m$.  Together these expressions give the
unique linear forms $d_i(X,Y)$.
Finally, since the $q_i$ are a basis, the quadratic component has a
unique expansion
\[
 -\sum_{i=1}^{m-1}Q_i(Y)q_i(T,X)
\]
with real homogeneous quadratic forms $Q_i$.
For every $y$, the ternary conclusion and uniqueness in the adjacent
basis give $Q_i(y)\geq0$ for every $i$.  Thus each $Q_i$ is positive
semidefinite.  This proves necessity and uniqueness in all dimensions
without a determinantal representation in more than three variables.

For sufficiency, start with any linear forms $d_i$ and positive
semidefinite forms $Q_i$ appearing in
\eqref{herm:binary-tangent-formula}.  For $\varepsilon\geq0$, put
\[
 L_i=T-\lambda_iX-\varepsilon d_i(X,Y)
\]
and define polynomials recursively by
\begin{equation}\label{herm:weighted-path-recurrence}
 P_0=1,\qquad P_1=L_1,\qquad
 P_k=L_kP_{k-1}-\varepsilon Q_{k-1}(Y)P_{k-2}\quad(2\leq k\leq m).
\end{equation}
The recurrence shows that $P_k$ is homogeneous of degree $k$, monic
in $T$, and has coefficients polynomial in $\varepsilon$.
For each fixed real $(X,Y)$, it is the characteristic-polynomial
recurrence for the real symmetric tridiagonal matrix with diagonal
entries $\lambda_iX+\varepsilon d_i(X,Y)$ and adjacent entries
$\sqrt{\varepsilon Q_i(Y)}$.
All these square roots are real, so all roots of $P_m$ in $T$ are
real.  Thus $P_m\in\mathcal H_{m,r}$ for every $\varepsilon\geq0$.
The matrix entries need not be linear functions of $Y$; the polynomial
itself is defined by \eqref{herm:weighted-path-recurrence} and contains
no square roots.

For completeness, expanding the recurrence over matchings $M$ of
the path with vertices $1,\ldots,m$ gives
\[
 P_m=\sum_M(-\varepsilon)^{|M|}
       \prod_{\{i,i+1\}\in M}Q_i(Y)
       \prod_{j\notin V(M)}L_j.
\]
This identity follows inductively by separating matchings according
to whether they contain the final edge.
The empty matching contributes the first-order changes of the $L_j$;
the single-edge matchings contribute the quadratic terms; every
other matching has a factor $\varepsilon^2$.
Consequently
\[
 P_m=p-\varepsilon\sum_i d_i\frac p{\ell_i}
       -\varepsilon\sum_iQ_i\frac p{\ell_i\ell_{i+1}}
       +O(\varepsilon^2)=p+\varepsilon h+O(\varepsilon^2).
\]
This proves sufficiency and polynomial realization.  The linear
isomorphism asserted in the theorem follows from the uniqueness of
the displayed coordinates.
\end{proof}

\begin{corollary}[Positive stability in the binary-factor stratum]
\label{herm:binary-positive-stability}
Let $N\geq3$, and let
\[
 F(z)=\prod_{i=1}^m L_i(z),\qquad z\in\R^N,
\]
where each $L_i$ has strictly positive coefficients, the $L_i$ are
pairwise nonproportional, and their linear span has dimension two.
Fix $v\in\R_{>0}^N$.
There is a real linear coordinate isomorphism
$\Phi:z\mapsto(T,X,Y)$, with $Y\in\R^{N-2}$ and $\Phi(v)=e$, such
that
\[
 \frac{F}{F(v)}=p\circ\Phi,\qquad
 p(T,X,Y)=\prod_i(T-\lambda_iX),\qquad
 \lambda_1<\cdots<\lambda_m.
\]
Let $\mathcal S^+_{m,N}$ be the homogeneous degree-$m$ real stable
polynomials with nonnegative coefficients.  Then
\begin{equation}\label{herm:binary-positive-tangent}
 T_{\mathcal S^+_{m,N}}(F)
 =\{aF+F(v)(h\circ\Phi):
        a\in\R,\ h\in T_{\mathcal H_{m,N-2}}(p)\}.
\end{equation}
Every direction on the right has a realizing curve with coefficients
polynomial in time, lying in $\mathcal S^+_{m,N}$ for all sufficiently
small nonnegative times.
\end{corollary}

\begin{proof}
Normalize the factors by $\widetilde L_i=L_i/L_i(v)$.
Their span is two-dimensional, and each has value one at $v$.
Choose a linear form $T$ in this span with $T(v)=1$, and a nonzero
linear form $X$ in this span with $X(v)=0$.
Then $\widetilde L_i=T-\lambda_iX$, with distinct $\lambda_i$;
reorder the factors increasingly.  Complete $X$ to a basis
$X,Y_1,\ldots,Y_{N-2}$ of the linear forms vanishing at $v$.
Together with $T$ these forms define the asserted isomorphism $\Phi$.

We use the following standard consequence of G\aa rding's component
theorem
\cite[Corollary~2.7 and Theorem~2.9]{HarveyLawson09}: a homogeneous real polynomial hyperbolic in
$v$ is hyperbolic in every direction in the connected component of
$\{f\ne0\}$ containing $v$.
In particular, if $v>0$, $f$ is hyperbolic in $v$, and $f$ has
nonnegative coefficients and is nonzero, then it is hyperbolic in
every positive direction.  Indeed, $f$ is strictly positive throughout
the positive orthant, which is connected and contains $v$.
Such an $f$ is real stable: if $z=x+iy$ with $y>0$, then
$t\mapsto f(x+ty)$ is a nonzero real-rooted polynomial, and so does
not vanish at $t=i$.  Conversely, a homogeneous real stable polynomial
with $f(v)>0$ is hyperbolic in $v$, since the real polynomial
$t\mapsto f(x+tv)$ has no root in the upper half-plane, hence by
conjugation none in the lower half-plane.

Suppose first that
$F_j=F+\varepsilon_jg+o(\varepsilon_j)\in\mathcal S^+_{m,N}$.
For large $j$, $F_j(v)>0$.  Set $a=g(v)/F(v)$ and normalize:
\[
 \frac{F_j}{F_j(v)}
 =\frac{F}{F(v)}
  +\varepsilon_j\frac{g-aF}{F(v)}+o(\varepsilon_j).
\]
Every normalized polynomial is hyperbolic in $v$ by the preceding
observation.  Pulling back by $\Phi^{-1}$ gives a tangent sequence
in $\mathcal H_{m,N-2}$ at $p$, proving necessity in
\eqref{herm:binary-positive-tangent}.

Conversely, let $a\in\R$ and let $h$ be as in the right side of
\eqref{herm:binary-positive-tangent}.  Use the polynomial curve
$P_\varepsilon$ from Theorem~\ref{herm:binary-localization} and put
\[
 F_\varepsilon(z)=F(v)(1+a\varepsilon)
                         P_\varepsilon(\Phi(z)).
\]
This is homogeneous, its coefficients are polynomials in
$\varepsilon$, and
\[
 F_\varepsilon=F+\varepsilon\bigl(aF+F(v)(h\circ\Phi)\bigr)
                      +O(\varepsilon^2).
\]
Every degree-$m$ monomial has a strictly positive coefficient in
$F$, because every coefficient of every $L_i$ is strictly positive.
There are finitely many such coefficients, so all remain strictly
positive in $F_\varepsilon$ for sufficiently small
$\varepsilon\geq0$.  Also $1+a\varepsilon>0$ for such times, and
$F_\varepsilon$ is hyperbolic in $v$ by construction.
The component argument above therefore makes $F_\varepsilon$ real
stable, proving sufficiency.
\end{proof}

\begin{remark}\label{herm:binary-localization-scope}
The restriction to a two-dimensional span of pairwise nonproportional
factors is part of the theorem.  Neither repeated factors nor products
whose factors span three or more dimensions are covered by
\eqref{herm:binary-tangent-formula}.
When $r=0$, the simple-root binary form lies in the open set of
monic real-rooted binary forms, and its tangent consists of all
homogeneous degree-$m$ perturbations with zero $T^m$ coefficient.
When $m=1$, the monic homogeneous forms are all hyperbolic and their
tangent is the full tangent space of the monic affine slice.
\end{remark}

\subsection{Why fixed univariate restrictions do not determine the tangent}

The exactness of Theorem~\ref{herm:tangent} does not allow one to
interchange taking tangents with the family of affine-line restrictions
that characterizes multivariate stability.  We give two explicit
examples.  Both have strictly positive homogeneous coefficients after
an invertible real change of coordinates.

We shall use the following standard form of the hyperbolicity component
theorem.  If a real homogeneous polynomial is hyperbolic in a direction
$e$ with $p(e)>0$, it is hyperbolic in every direction in the connected
component of $\{p>0\}$ containing $e$; see
\cite[Corollary~2.7 and Theorem~2.9]{HarveyLawson09}.
Here hyperbolicity means that $p(e)\ne0$ and every polynomial
$s\mapsto p(x+se)$, $x$ real, has only real roots.
For a homogeneous real polynomial, hyperbolicity in every strictly
positive direction is equivalent to real stability: a zero
$p(x+iy)=0$ with $y>0$ violates the corresponding line test, and a
nonreal root of a real line restriction has a conjugate in the upper
half-plane.  We will repeatedly use the following elementary consequence
of the component theorem.  If $p_j\to p$ coefficientwise, each $p_j$ is
hyperbolic in $d$, and a fixed path from $d$ to $e$ has compact image
contained in $\{p>0\}$, then $p_j$ is hyperbolic in $e$ for all
sufficiently large $j$.  Indeed, the positive minimum of $p$ on that
path persists under coefficient convergence.

\subsubsection{An exact cubic comparison}

In three variables set
\begin{equation}
 T=z_1+z_2+z_3,\qquad
 X=\frac{z_1+2z_2+3z_3}{4},\qquad Y=z_2,
 \label{herm:cubic-coordinates}
\end{equation}
and define
\begin{equation}
 f=T(T-X)(T+X),\qquad
 g_{a,b}=-(aT+bX)Y^2,\qquad a,b\in\R.
 \label{herm:cubic-family}
\end{equation}
The coordinate change is invertible.  Each of $T,T-X,T+X$ has strictly
positive coefficients in $z$, so $f$ is real stable and all its
degree-three coefficients are strictly positive.  Let $\mathcal S_3$
be the set of homogeneous real stable ternary cubics with nonnegative
coefficients, together with zero.

For an affine line $z=u+sv$ with $u\in\R^3$ and $v>0$, put
$p(s)=f(u+sv)$ and $h(s)=g_{a,b}(u+sv)$.  Its leading coefficient is
$c=f(v)>0$.  The normalized base and direction are
\begin{equation}
 \overline p=\frac pc,\qquad
 \overline h=\frac hc-\frac{[s^3]h}{c}\overline p.
 \label{herm:normalized-slice}
\end{equation}
Saying that this line passes the Hermite tangent test means
$DH_{\overline p}[\overline h]|_{\ker H(\overline p)}\succeq0$.

\begin{theorem}[Positive-direction slice tests are insufficient]
\label{herm:cubic-slices}
For the family \eqref{herm:cubic-family}, every affine restriction in
every strictly positive direction passes the Hermite tangent test if
and only if
\begin{equation}
 a+\frac b4\geq0,
 \qquad a+\frac{3b}{4}\geq0.
 \label{herm:cubic-slice-cone}
\end{equation}
In contrast,
\begin{equation}
 g_{a,b}\in T_{\mathcal S_3}(f)
       \quad\Longleftrightarrow\quad a\geq|b|.
 \label{herm:cubic-true-cone}
\end{equation}
When $a\geq|b|$, the straight perturbation $f+\varepsilon g_{a,b}$
already belongs to $\mathcal S_3$ for every sufficiently small
$\varepsilon\geq0$.
\end{theorem}

\begin{proof}
All three affine factors have strictly positive slopes along $v>0$.
If two restricted roots coincide, then $T=X=0$ there and the third
root coincides as well.  Thus a restriction has either three simple
roots or one triple root.  At a simple-root base the monic real-rooted
set is open and its tangent is unrestricted.  At a triple root $s=r$,
Corollary~\ref{herm:root-contact} gives precisely
$\overline h(r)=0$ and $\overline h'(r)\leq0$.
The leading-normalization term in \eqref{herm:normalized-slice}
vanishes to order three, whereas direct differentiation gives
\[
 h(r)=0,\qquad
 h'(r)=-(aT(v)+bX(v))Y(u+rv)^2.
\]
For every fixed $v>0$, a line in direction $v$ can pass through a point
with $T=X=0$ and any specified value of $Y$.  Hence all line tests are
equivalent to $aT(v)+bX(v)\geq0$ for every $v>0$.
The ratio $X(v)/T(v)$ ranges over $(1/4,3/4)$, so this is equivalent
to \eqref{herm:cubic-slice-cone}.

We prove the necessity in \eqref{herm:cubic-true-cone} directly for
arbitrary contingent sequences.  Suppose
$f_j=f+\varepsilon_jg_{a,b}+o(\varepsilon_j)\in\mathcal S_3$, where
$\varepsilon_j\downarrow0$.  In the coordinates $(T,X,Y)$ the base is
\[
 p=T^3-X^2T,
\]
which is hyperbolic in $e=(1,0,0)$, with positive component
$\{T>|X|\}$.  The original positive direction $(1,1,1)$ becomes
$d=(3,3/2,1)$.  The transformed $f_j$ are hyperbolic in $d$, and the
segment from $d$ to $e$ stays in $\{p>0\}$.  The component argument
above therefore makes them hyperbolic in $e$ for all sufficiently
large $j$.

Normalize the coefficient of $T^3$ to one.  This coefficient was
$1+o(\varepsilon_j)$, so normalization does not change the prescribed
first variation.  Write the normalized polynomials as
\[
 p_j=T^3+A_j(X,Y)T^2+B_j(X,Y)T+C_j(X,Y),
\]
where coefficientwise
\[
 A_j=o(\varepsilon_j),\qquad
 B_j=-X^2-\varepsilon_j aY^2+o(\varepsilon_j),\qquad
 C_j=-\varepsilon_j bXY^2+o(\varepsilon_j).
\]
The real translation eliminating the quadratic term produces
$T^3-Q_jT+R_j$, with
\begin{align}
 Q_j&=A_j^2/3-B_j
       =X^2+\varepsilon_j aY^2+o(\varepsilon_j),\notag\\
 R_j&=C_j-A_jB_j/3+2A_j^3/27
       =-\varepsilon_j bXY^2+o(\varepsilon_j).
 \label{herm:cubic-depressed}
\end{align}
All these estimates take place in fixed finite-dimensional spaces of
homogeneous coefficient forms.  The real-root criterion for a depressed
cubic is
\begin{equation}
 Q_j(X,Y)\geq0,\qquad
 |R_j(X,Y)|\leq\frac{2}{3\sqrt3}Q_j(X,Y)^{3/2}
       \quad(X,Y\in\R).
 \label{herm:cubic-discriminant-bound}
\end{equation}
For completeness, its discriminant is $4Q_j^3-27R_j^2$; equivalently,
evaluation at the two critical points
$\pm\sqrt{Q_j/3}$ gives the same necessary and sufficient inequalities.
At $(X,Y)=(0,1)$, the first inequality implies $a\geq0$.

Fix $s>0$ and evaluate at $(X,Y)=(\pm\sqrt{\varepsilon_j}s,1)$.
Then
\[
 Q_j(\pm\sqrt{\varepsilon_j}s,1)
       =\varepsilon_j(s^2+a)+o(\varepsilon_j).
\]
The difference of the corresponding $R_j$ values has the sharper
expansion
\begin{equation}
 R_j(\sqrt{\varepsilon_j}s,1)
       -R_j(-\sqrt{\varepsilon_j}s,1)
       =-2b\varepsilon_j^{3/2}s+o(\varepsilon_j^{3/2}).
 \label{herm:cubic-odd-cancellation}
\end{equation}
Indeed, $R_j$ is a homogeneous cubic in $X,Y$.  Its even powers of $X$
cancel in the difference; its $X^3$ coefficient is $o(\varepsilon_j)$,
and its $XY^2$ coefficient is
$-\varepsilon_jb+o(\varepsilon_j)$.  This justifies
\eqref{herm:cubic-odd-cancellation} even when the coefficients of the
even terms are only $o(\varepsilon_j)$.
Apply \eqref{herm:cubic-discriminant-bound} at both points and use the
triangle inequality.  After division by $2\varepsilon_j^{3/2}$ and
passage to the limit, we obtain
\begin{equation}
 |b|s\leq\frac{2}{3\sqrt3}(s^2+a)^{3/2}\qquad(s>0).
 \label{herm:cubic-optimized-bound}
\end{equation}
If $a=0$, let $s\downarrow0$ to obtain $b=0$.  If $a>0$, the minimum
of $(s^2+a)^3/s^2$ is $27a^2/4$, attained when $s^2=a/2$.
Squaring \eqref{herm:cubic-optimized-bound} at that value gives
$b^2\leq a^2$.  This proves necessity.

For sufficiency, assume $a\geq|b|$.  In transformed coordinates the
straight perturbation is
\[
 p_\varepsilon
       =T^3-(X^2+\varepsilon aY^2)T-\varepsilon bXY^2.
\]
Its depressed coefficient $Q=X^2+\varepsilon aY^2$ is nonnegative.
Its discriminant is also nonnegative because
\[
 4(X^2+\varepsilon aY^2)^3
       -27\varepsilon^2b^2X^2Y^4\geq0.
\]
When $\varepsilon Y^2>0$, divide this expression by
$\varepsilon^3Y^6$ and put $u=X^2/(\varepsilon Y^2)$; the inequality
follows from
\[
 4(u+a)^3-27b^2u
 \geq4(u+a)^3-27a^2u
       =4(u+4a)(u-a/2)^2\geq0.
\]
The case $\varepsilon Y^2=0$ is immediate.  Thus $p_\varepsilon$ is
hyperbolic in $e$ for every $\varepsilon\geq0$.
For sufficiently small $\varepsilon$, all its coefficients in the
original $z$ coordinates remain strictly positive, and the fixed segment
from $e$ to $d$ remains in its positive set.  The component theorem
then gives hyperbolicity in $d$.  Positivity of the coefficients makes
the entire original positive orthant a connected subset of the same
positive component.  Hyperbolicity in every positive direction follows,
and hence the original polynomial is real stable.  This proves
sufficiency and the straight-path assertion.
\end{proof}

For example, $a=0,b=1$ satisfies every fixed positive-direction Hermite
test but is not a coefficient tangent.  In this cubic family, enlarging
the directions to the full base hyperbolicity cone $T>|X|$ would recover
the condition $a\geq|b|$, since $X(v)/T(v)$ then ranges over $(-1,1)$.
The next example shows that this enlargement does not give a general
tangent criterion either.

\subsubsection{Failure even for all directions in the base cone}

Let $\mathcal H_5$ denote the monic-in-$T$ real homogeneous quintics
hyperbolic in $e=(1,0,0)$.  Set
\begin{equation}
 p=T(T^2-X^2)(T^2-4X^2)
       =T^5-5X^2T^3+4X^4T,
 \qquad h=-TX^2Y^2.
 \label{herm:quintic-example}
\end{equation}
The positive hyperbolicity component of $p$ containing $e$ is
$\Gamma=\{T>2|X|\}$.

\begin{theorem}[Full-base-cone slice obstruction]
\label{herm:quintic-slices}
Every affine restriction of \eqref{herm:quintic-example} in every
direction $v\in\Gamma$ satisfies the exact univariate tangent criterion
after monic normalization.  Nevertheless
\[
 h\notin T_{\mathcal H_5}(p).
\]
The same failure occurs for the tangent to homogeneous real stability
at a base with strictly positive coefficients, after the change
\begin{equation}
 T=z_1+z_2+z_3,\qquad
 X=\frac{z_1+2z_2+3z_3}{8},\qquad Y=z_2.
 \label{herm:quintic-coordinates}
\end{equation}
\end{theorem}

\begin{proof}
The five factors $T-\lambda X$,
$\lambda\in\{-2,-1,0,1,2\}$, all have positive slopes in direction
$v\in\Gamma$.  If any two restricted roots coincide, then $T=X=0$
at that point, so all five coincide.  Simple restrictions impose no
tangent condition.  At a quintuple root $s=r$, the restricted variation
vanishes to order at least three, and its coefficient of $(s-r)^3$ is
\[
       -T(v)X(v)^2Y(u+rv)^2\leq0.
\]
The base leading coefficient is positive, and the monic-normalization
term vanishes to order five.  Corollary~\ref{herm:root-contact} therefore
proves every asserted slice test.

Suppose, to the contrary, that
$p_j=p+\varepsilon_jh+o(\varepsilon_j)\in\mathcal H_5$, where
$\varepsilon_j\downarrow0$.  Write
\[
 p_j=T^5+a_{1,j}(X,Y)T^4+a_{2,j}(X,Y)T^3
             +\cdots+a_{5,j}(X,Y).
\]
For fixed real $(X,Y)$, let $M_{2,j}$ be the sum of squares of the five
real roots in $T$, and let $E_{4,j}$ be their fourth elementary
symmetric function.  Newton's identity gives, coefficientwise,
\begin{equation}
 M_{2,j}=a_{1,j}^2-2a_{2,j}=10X^2+o(\varepsilon_j),
 \qquad
 E_{4,j}=a_{4,j}=4X^4-\varepsilon_jX^2Y^2+o(\varepsilon_j).
 \label{herm:quintic-moments}
\end{equation}
For any five real numbers,
\begin{equation}
       |E_4|\leq5M_2^2.
 \label{herm:fourth-moment-bound}
\end{equation}
There are five products in $E_4$, and each absolute product is bounded
by $M_2^2$, since each individual absolute root is at most $\sqrt{M_2}$.

At $(X,Y)=(0,1)$, the first equation in
\eqref{herm:quintic-moments} gives $M_{2,j}=o(\varepsilon_j)$.
Inequality \eqref{herm:fourth-moment-bound} therefore improves the
coefficient of $Y^4$ in $E_{4,j}$ to $o(\varepsilon_j^2)$.
For fixed $s>0$, average the values at
$(X,Y)=(\pm\sqrt{\varepsilon_j}s,1)$.  The odd powers of $X$ cancel;
the improved $Y^4$ coefficient and the remaining even coefficients in
\eqref{herm:quintic-moments} give
\[
 \frac{E_{4,j}(\sqrt{\varepsilon_j}s,1)
            +E_{4,j}(-\sqrt{\varepsilon_j}s,1)}{2}
       =\varepsilon_j^2(4s^4-s^2)+o(\varepsilon_j^2).
\]
At both evaluation points,
$M_{2,j}=10\varepsilon_j s^2+o(\varepsilon_j)$.
Apply \eqref{herm:fourth-moment-bound} at each point, take the average
of the absolute bounds, divide by $\varepsilon_j^2$, and pass to the
limit.  It follows that
\[
       |4s^4-s^2|\leq500s^4\qquad(s>0).
\]
This is false whenever $0<s^2<1/504$.  Thus no contingent sequence
exists.

Finally, \eqref{herm:quintic-coordinates} is invertible, and all five
linear factors $T-\lambda X$ have strictly positive coefficients in
$z$.  Their product consequently has every degree-five coefficient
strictly positive.  The original direction $(1,1,1)$ becomes
$d=(3,3/4,1)$, and the segment from $d$ to $e$ lies in $\Gamma$.
A hypothetical stable tangent sequence in the original variables
would be hyperbolic in $d$.  The component argument would make it
hyperbolic in $e$ for large $j$; normalization of its $T^5$ coefficient,
which is $1+o(\varepsilon_j)$, would produce the impossible tangent
sequence just excluded.  This proves the positive-coefficient version.
\end{proof}

Both failures concern coefficient tangency at a fixed base polynomial.
Their proofs use lines whose locations vary on the scale
$X=\sqrt{\varepsilon}\,s$.  They do not give a counterexample to a
global generator criterion tested at every stable input; such a
criterion has additional compatibility across base polynomials.

\begin{remark}[Recovery from the adjacent quotient cone]
\label{herm:adjacent-recovery}
The localization formula gives a second explanation of these examples.
For the cubic roots $-1,0,1$, the adjacent quotients are $T-X,T+X$,
and
\[
 aT+bX=\frac{a-b}{2}(T-X)+\frac{a+b}{2}(T+X).
\]
Their nonnegative cone gives exactly $a\geq|b|$.  For the quintic
roots $-2,-1,0,1,2$, let $q_i=p/(\ell_i\ell_{i+1})$ in that order.
Direct expansion gives
\[
 TX^2=\frac1{12}(q_1-q_2-q_3+q_4).
\]
Uniqueness of the adjacent quotient coordinates excludes this form from
their nonnegative cone, although $TX^2\geq0$ on $\Gamma$.
The cubic discriminant proof and quintic moment proof above establish
the same conclusions independently of the determinantal representation
used in Theorem~\ref{herm:binary-localization}.
\end{remark}

% END SECTION hermite

% BEGIN SECTION analytic_generation
\section{Domains, generation, and polynomial-valued dynamics}
\label{ana:section}

The algebraic restrictions on a stability-preserving generator do not by
themselves specify a semigroup on an infinite-dimensional space.  We first
resolve the domain and approximation questions for summable positive
coefficients.  We then treat polynomial-valued semigroups, for which the
topology forces a different and particularly rigid conclusion.

\subsection{An automatic domain theorem}

Let \(E\) be countable and let \(r:E\to\N\) have finite fibres
\(E_k=r^{-1}(k)\).  Write \(e_x\) for the standard basis of \(\ell^1(E)\).
A nonnegative sequence has interval support if its positive entries form an
interval of integers, possibly a singleton or an unbounded interval.

\begin{theorem}[Automatic finite-support domain]\label{ana:domain}
Let \(T(t)\) be a positive strongly continuous semigroup on \(\ell^1(E)\),
with generator \(A\).  Suppose that, for each \(x\), and all sufficiently small
\(t>0\), the sequence
\[
 b_k(t)=\sum_{y\in E_k}(T(t)e_x)_y
\]
is log-concave and has interval support.  Then \(e_x\in D(A)\), and
\[
 \supp(Ae_x)\subseteq E_{r(x)-1}\cup E_{r(x)}\cup E_{r(x)+1},
 \qquad E_{-1}=\varnothing .
\]
In particular, every finitely supported vector belongs to \(D(A)\), and
\(A\) maps finitely supported vectors to finitely supported vectors.
\end{theorem}

We give the transition-function argument needed in the proof, to distinguish
existence of coordinate derivatives from an a priori finite-rate assumption.
The classical result is due to Kolmogorov; see \cite{Chung67,Kendall55}.

\begin{lemma}\label{ana:transition-derivatives}
Let \(S(t)\) be a standard countable substochastic transition kernel.  Then
every off-diagonal limit \(S_{yx}(t)/t\), as \(t\downarrow0\), exists and is
finite.  The limit \((1-S_{xx}(t))/t\) exists in \([0,\infty]\).
\end{lemma}

\begin{proof}
Transpose to the row convention and add an absorbing cemetery state.
We obtain a stochastic kernel \(P(t)\), with
\(P_{ij}(t)\to\delta_{ij}\) at zero.  Its entries are continuous: the
semigroup identity gives
\[
 |P_{ij}(t+h)-P_{ij}(t)|\le 1-P_{ii}(h),
\]
and the corresponding left estimate follows by writing \(t=(t-h)+h\).
Fix \(i\ne j\) and \(0<\varepsilon<1/3\).  Choose \(\delta_0>0\) so that
\(P_{ii}(s),P_{jj}(s)\ge1-\varepsilon\) for \(0\le s\le\delta_0\).
Consider the discrete chain with transition matrix \(P(h)\), started at
\(i\).  Let \(\tau\) be its first visit to \(j\), let
\(f_l=\Pr(\tau=l)\), and put
\(a_k=\Pr(X_k=i,\tau>k)\).
For \(kh\le\delta_0\), first-visit decomposition and
\(P_{ji}(s)\le\varepsilon\) give
\[
 P_{ii}(kh)
 =a_k+\sum_{l=1}^k f_lP_{ji}((k-l)h)\le a_k+\varepsilon .
\]
Consequently \(a_k\ge1-2\varepsilon\), and
\(f_{k+1}\ge a_kP_{ij}(h)\).  If \(nh\le\delta_0\), then
\[
 P_{ij}(nh)
 =\sum_{l=1}^n f_lP_{jj}((n-l)h)
 \ge n(1-\varepsilon)(1-2\varepsilon)P_{ij}(h)
 \ge n(1-3\varepsilon)P_{ij}(h).
\]
Fix \(0<\delta\le\delta_0\) and take \(n=\lfloor\delta/h\rfloor\).
Continuity implies
\[
 \limsup_{h\downarrow0}\frac{P_{ij}(h)}h
 \le \frac{P_{ij}(\delta)}{\delta(1-3\varepsilon)}<\infty .
\]
First let \(\delta\downarrow0\) along a sequence attaining the lower limit,
and then let \(\varepsilon\downarrow0\).  The off-diagonal assertion follows.

Subdivision of time, retaining only the path which stays at \(i\), gives
\(P_{ii}(t)>0\) for every \(t\).  Thus
\(\phi(t)=-\log P_{ii}(t)\) is continuous, nonnegative, and subadditive.
Writing \(\delta=nt+s\), \(0\le s<t\), yields
\(\phi(\delta)\le n\phi(t)+\phi(s)\).  Hence
\(\liminf_{t\downarrow0}\phi(t)/t\ge\phi(\delta)/\delta\).
Taking the supremum over \(\delta>0\) proves existence of the limit, possibly
infinite.  Since \((1-e^{-u})/u\to1\) at zero, this is also the asserted
diagonal limit.
\end{proof}

\begin{proof}[Proof of Theorem~\ref{ana:domain}]
Write \(p_{yx}(t)=(T(t)e_x)_y\).  Choose \(M,\omega\) with
\(\|T(t)\|\le Me^{\omega t}\), and choose
\(\lambda>\max(\omega,0)\).  Define the strictly positive bounded weights
\[
 w_x=\int_0^\infty e^{-\lambda s}\|T(s)e_x\|_1\,ds .
\]
No uniform positive lower bound on \(w_x\) is required.  Positivity,
Tonelli's theorem, and the semigroup law give
\[
 \sum_yw_yp_{yx}(t)
 =e^{\lambda t}
 \left(w_x-\int_0^t e^{-\lambda s}\|T(s)e_x\|_1\,ds\right).
\]
Therefore
\[
 S_{yx}(t)=e^{-\lambda t}\frac{w_y}{w_x}p_{yx}(t)
\]
is a standard substochastic kernel, with deficit
\[
 1-\sum_yS_{yx}(t)
 =w_x^{-1}\int_0^t e^{-\lambda s}\|T(s)e_x\|_1\,ds
 =t/w_x+o(t).
\]
Lemma~\ref{ana:transition-derivatives} supplies finite off-diagonal
derivatives \(q_{yx}=\lim_{t\downarrow0}p_{yx}(t)/t\).

Fix \(x\), put \(N=r(x)\), and use the \(b_k\) of the theorem.
Strong continuity gives \(b_N(t)\to1\).  Finiteness of adjacent fibres and
the coordinate derivatives imply \(b_{N-1}(t),b_{N+1}(t)=O(t)\).
Log-concavity and interval support give the geometric estimates
\[
 b_{N+j}(t)\le b_N(t)
 \left(\frac{b_{N+1}(t)}{b_N(t)}\right)^j,\qquad
 b_{N-j}(t)\le b_N(t)
 \left(\frac{b_{N-1}(t)}{b_N(t)}\right)^j .
\]
These remain valid when an adjacent coefficient vanishes, since the
corresponding tail then vanishes.  It follows that
\[
 \sum_{|k-N|\ge2}b_k(t)=O(t^2),\qquad
 \sum_{y\ne x}p_{yx}(t)=O(t).
\]
The second estimate also uses finiteness of the three central fibres.
Since \(w\) is bounded and \(w_x>0\), the off-diagonal column sum of \(S(t)\)
is \(O(t)\).  Its displayed deficit is \(O(t)\); hence
\(1-S_{xx}(t)=O(t)\).  The diagonal limit in the preceding lemma is thus
finite, and so is
\[
 q_{xx}=\lim_{t\downarrow0}\frac{p_{xx}(t)-1}{t}.
\]
On the finite union of the three central fibres, the difference quotient
converges coordinatewise and therefore in norm.  Outside that union its
\(\ell^1\) norm is \(O(t^2)/t=O(t)\).  This is precisely strong
differentiability of \(T(t)e_x\) at zero, with the stated support.
\end{proof}

\begin{corollary}\label{ana:polynomial-domain}
On any mixed-capacity space with finitely many coordinates, a positive
\(C_0\)-semigroup preserving \(\mathcal{LP}_{+,\kappa}\) automatically has all
polynomials in its generator domain.  Its action on each monomial has finite
support and changes total degree by at most one.
\end{corollary}

\begin{proof}
Use the proper grading \(r(x)=|x|\).  The diagonal specialization of
\(T(t)z^x\) belongs to the univariate positive Laguerre--P\'olya class.
Its coefficient sequence is log-concave with interval support: this follows
from negative-real-root factorization for its polynomial approximants and
passage to coefficient limits.  Theorem~\ref{ana:domain} applies.
\end{proof}

\subsection{Consistent cutoffs and weighted spaces}

Let \(\mathcal F\) denote the finitely supported vectors.  Suppose \(L\) is
a Metzler operator on \(\mathcal F\), meaning that its off-diagonal matrix
entries are nonnegative, and preserves
\[
 W_N=\operatorname{span}\{e_x:r(x)\le N\}.
\]
Suppose \(R\) is a bounded nonnegative operator whose entries raise the
grading by exactly one.  Write \(\rho=\|R\|_1\).

\begin{theorem}[Cutoff generation criterion]\label{ana:weighted-generation}
The operator \(L+R\) on \(\mathcal F\) has a positive \(C_0\)-semigroup
realization on \(\ell^1(E)\), with this action on \(\mathcal F\), if and only
if there are \(\tau>0\) and \(B<\infty\) such that
\[
 \sup_N\sup_{0\le t\le\tau}\|\exp(tL|_{W_N})\|_1\le B.
\]
When it exists, the realization is unique, its generator is the closure of
\(L+R\), and \(\mathcal F\) is a graph core.

For \(s\ge1\), put
\[
 X_s=\left\{f:\|f\|_s=\sum_x|f_x|s^{r(x)}<\infty\right\}.
\]
If \(\|U(t)\|_1\le Me^{\omega t}\) for the semigroup generated by the closure
of \(L\), with \(M\ge1\) and \(\omega\ge0\), then compatible semigroups satisfy
\[
 \|U(t)\|_{X_s}\le Me^{\omega t},\qquad
 \|R\|_{X_s}=s\rho,\qquad
 \|T^{(s)}(t)\|_{X_s}\le
 M\exp\bigl((\omega+Ms\rho)t\bigr).
\]
\end{theorem}

\begin{proof}
The finite flows are consistent by invariance of the \(W_N\).  The asserted
bound, followed by subdivision of time, bounds these flows on every compact
time interval.  They extend from \(\mathcal F\) to bounded operators \(U(t)\)
on \(\ell^1(E)\).  Density proves both the semigroup law and strong
continuity.  Finite-dimensional positivity proves positivity of the
extension.  The bounded positive perturbation \(R\) can be added by the
Dyson series, giving a positive semigroup \(T(t)\).

Conversely, subtract \(R\) from any proposed generator by bounded
perturbation.  Each \(W_N\) is contained in its domain and is invariant under
its action.  Uniqueness of classical solutions identifies the resulting
semigroup on \(W_N\) with \(\exp(tL|_{W_N})\).  In particular, it is positive
on \(\mathcal F\), hence positive everywhere, and its local boundedness
gives the necessary cutoff estimate.  It also proves uniqueness.

Let \(G\) denote the generator of the constructed semigroup \(U(t)\).
For a sufficiently large resolvent parameter \(\lambda\), its resolvent
preserves each \(W_N\), by the Laplace-transform formula.
Given \(f\in D(G)\), approximate
\((\lambda-G)f\) in \(\ell^1\) by vectors in \(\mathcal F\) and apply that
resolvent.  The resulting vectors lie in \(\mathcal F\) and converge to
\(f\) in the graph norm of \(G\).  Since \(G=L\) on \(\mathcal F\), this
proves \(G=\overline L\).  A bounded perturbation has the same domain and an
equivalent graph norm, proving the core assertion.

All entries of \(U(t)\) vanish unless the output grading is no larger than
the input grading.  Its weighted column sums are therefore bounded by the
unweighted column sums after dividing by the input weight.  This proves
the bound for \(U(t)\) on \(X_s\).  Density proves strong continuity there.
The identity \(\|R\|_{X_s}=s\rho\) follows directly from its one-step grading.
The \(k\)-th Dyson term has norm at most
\[
 M^{k+1}e^{\omega t}(s\rho)^k t^k/k!,
\]
and summation gives the stated bound.  These constructions are compatible
under \(X_s\hookrightarrow X_1\), term by term.
\end{proof}

\begin{corollary}\label{ana:column-sum}
If the column sums of \(L+R\) are bounded above, the realization in
Theorem~\ref{ana:weighted-generation} exists.
\end{corollary}
\begin{proof}
The column sums of each finite restriction of \(L\) are no larger than those
of \(L+R\).  The positive matrix exponential has norm at most \(e^{ct}\)
when its column sums are at most \(c\).  Apply the cutoff criterion.
\end{proof}

\subsection{Growth reduction and the legitimate scaling limit}

For mixed capacities, write \(F\) and \(U\) for the finite and unrestricted
coordinates and use \(r(x)=|x_U|\).  All spaces \(W_N\) below are finite
dimensional.  The next proposition supplies the analytic bridge between
local normal-order tests and tests with arbitrary real spectators.

\begin{proposition}\label{ana:growth-reduction}
Let \(T(t)\) be a positive \(C_0\)-semigroup on \(X_\kappa=\ell^1(\Omega_\kappa)\)
preserving \(\mathcal{LP}_{+,\kappa}\).  Suppose its polynomial action has
already been reduced to canonical differential order at most two.  Then
every canonical coefficient has unrestricted-coordinate degree at most two.
The part that increases \(|x_U|\) is the bounded operator
\[
 R=\sum_{j\in U}\lambda_jz_j+
   \sum_{\substack{i\in F\\j\in U}}r_{ij}z_j\partial_i,
 \qquad \lambda_j,r_{ij}\ge0.
\]
For \(L=A-R\), the conclusions of
Theorem~\ref{ana:weighted-generation} hold.  In particular, for every
polynomial \(f\), the expansion
\(T(t)f=f+tAf+O(t^2)\) holds locally uniformly on every bounded polydisc.
\end{proposition}

\begin{proof}
Corollary~\ref{ana:polynomial-domain} shows that all canonical coefficients
are polynomials: the triangular normal-order recursion involves only
images of finitely many monomials.  Since there are finitely many
coefficients after order reduction, and finite-coordinate exponents are
bounded, there is a constant \(C\) such that
\[
 \|Az^{(m,x_U)}\|_1\le C(1+|x_U|)^2.
\]
Thus the Taylor truncations of \(z_F^m e^{a\cdot z_U}\), \(a\ge0\),
converge in graph norm.  Closedness of \(A\) puts this exponential in its
domain and identifies its derivative by the differential expression.

Specialize finite variables to \(b>0\); this specialization is bounded
because their degrees remain bounded.  Dividing the first variation by
the initial exponential gives
\[
 \sum_{|\alpha|\le2}(m)_{\alpha_F}b^{-\alpha_F}
       a^{\alpha_U}P_\alpha(b,z_U).
\]
Lemma~\ref{pos:poisson} makes this polynomial quadratic in
\(z_U\).  Comparing powers of \(a\), and then the linearly independent
falling factorials on the finite \(m\)-grid, shows
\(\deg_{z_U}P_\alpha\le2\).

Represent the operator as a matrix on the finite-coordinate state space:
\[
 A_{nm}=\sum_{\substack{|p|\le2\\|q|\le2}}
              C_{nm;pq}z_U^q\partial_U^p .
\]
Terms increasing \(|x_U|\) have \((|p|,|q|)=(0,2),(1,2)\), or \((0,1)\).
The coefficients with \((0,2)\) are nonnegative, as coefficients of the
image of the input with no unrestricted particles.  Poisson tangency at
\(a=0\) makes their quadratic part nonpositive on the positive orthant;
therefore they vanish.

For \((1,2)\), fix \(p=e_l\).  If \(q_l>0\), its coefficient is a slope in
the nonnegative affine rate for adding one unrestricted particle.
Letting \(x_l\to\infty\) proves that the slope is nonnegative.  If \(q_l=0\),
the shift removes a particle at \(l\) and adds two elsewhere, and its rate
is precisely that coefficient times \(x_l\); the coefficient is again
nonnegative.  After the \((0,2)\) elimination, apply Poisson tangency at
\(a=\varepsilon v\), \(v>0\), divide its quadratic inequality by
\(\varepsilon\), and let \(\varepsilon\downarrow0\).  Every remaining
linear-in-\(a\), quadratic-in-\(z_U\) coefficient is nonnegative and their
sum is nonpositive.  They all vanish.  The increasing part is consequently
\[
 R=\sum_{j\in U}z_jK_j,
\]
with \(K_j\) a nonnegative matrix on the finite state space.  This proves
boundedness before any scaling is performed.

Subtracting this bounded operator gives \(L=A-R\).  Its finite restrictions
are Metzler and preserve \(W_N\); the necessity part of
Theorem~\ref{ana:weighted-generation} supplies the positive semigroup
\(U(t)\) and its compatible weighted realizations.  Put
\[
 S_s:X_s\longrightarrow X_1,\qquad
 (S_sf)(z_F,z_U)=f(z_F,sz_U).
\]
This is an isometric isomorphism between the indicated different spaces.
For
\[
 \widetilde U_s(t)=S_sU^{(s)}(t/s)S_s^{-1}
\]
we have \(\|\widetilde U_s(t)\|\le Me^{\omega t/s}\).
On a fixed \(W_N\), every generator entry is
\[
 s^{-1}s^{|y_U|-|x_U|}L_{yx}=O(s^{-1}).
\]
Thus \(\widetilde U_s(t)\to I\) strongly, uniformly on compact time
intervals, first on \(\mathcal F\), then by density on \(X_1\).
The perturbation of this rescaled semigroup is exactly \(R\), since
\(S_sRS_s^{-1}=sR\).  Its \(k\)-th Dyson term is bounded by
\[
 M^{k+1}e^{\omega t/s}\rho^k t^k/k!,
\]
and converges strongly to \(t^kR^k/k!\).  Dominated summation proves
\[
 S_sT^{(s)}(t/s)S_s^{-1}\longrightarrow e^{tR}
\]
strongly and locally uniformly in time.  Positive scaling and coefficient
closure of \(\mathcal{LP}_+\), Lemma~\ref{pos:lp-closure}, show that
\(e^{tR}\) preserves that class.

Set all unrestricted variables except \(z_j\) to zero.  We find that
\(e^{tz_jK_j}\) preserves positive stability.  Since this one-step-raising
operator is bounded on every \(X_s\), its exponential has a first variation
on every bounded polydisc.  Finite-site double contacts
with the real variable \(z_j\) as a spectator force every second-order
normal coefficient of \(K_j\) to vanish: its product with \(z_j\) has to be
nonpositive for both signs of \(z_j\).  The first-order specialization of
Lemma~\ref{gen:box-decomposition} then expresses \(K_j\) by independent deaths, capacity
births, and an affine diagonal.  Corollary~\ref{ana:polynomial-domain}
applied to \(1\) excludes the births: a birth in \(K_j1\) would create a
degree-two term in \(R1\).  Positivity gives
\[
 K_j=\lambda_j I+\sum_{i\in F}(r_{ij}\partial_i+m_{ij}z_i\partial_i),
 \qquad r_{ij}\ge0.
\]
Its remaining diagonal slopes need not yet have a sign.

If \(m=m_{ij}\ne0\), use the input \(c+u\) in the corresponding finite
variable \(u\), with \(c>\max(0,r_{ij}/m)\).  Up to a nonvanishing
exponential factor, the output is
\[
 c+e^{mtz}u+\frac{r_{ij}}m(e^{mtz}-1).
\]
For \(z=iv\), choose \(v>0\) with \(\sin(mtv)>0\).  Its root
\[
 u=-r_{ij}/m-(c-r_{ij}/m)e^{-mtz}
\]
then has positive imaginary part, a contradiction.  Hence every
\(m_{ij}=0\), proving the displayed form of \(R\), including
\(\lambda_j\ge0\).

Finally, the generator on every \(X_s\) agrees with \(A\) on polynomials.
If \(f\in W_N\), then \(A^2f\in W_{N+2}\).  Twice integrating the
semigroup equation and using the weighted estimate gives, for
\(0\le t\le\tau\),
\[
 \|T(t)f-f-tAf\|_s
 \le \frac{t^2}{2}M s^{N+2}
        e^{(\omega+Ms\rho)\tau}\|A^2f\|_1 .
\]
Finite-coordinate radii add only a fixed factor.  This proves the final
assertion on arbitrary bounded polydiscs.
\end{proof}

\begin{remark}
The order reduction used in the preceding proposition needs only the
first variation on a small polydisc, which follows directly from the
\(\ell^1\) generator derivative.  Negative collision points can be chosen
inside that polydisc; polynomial identity then gives the global
normal-order conclusion.  Arbitrary-radius first variations are proved
only after the bounded growth decomposition.  Thus no weighted
regularity is presupposed in deriving that decomposition.
\end{remark}

\subsection{Finite approximation with unchanged diagonals}

\begin{theorem}\label{ana:killed-approximation}
Let a mixed-capacity operator have the rate form and quadratic-contact
inequality of the positive classification, and suppose it has the
realization in Theorem~\ref{ana:weighted-generation}.  Impose artificial
capacities \(N_j\ge2\) at unrestricted sites.  Keep its diagonal entries
unchanged and replace the rates for immigration and transfer into such a
site by
\[
 \lambda_j(1-x_j/N_j),\qquad
 r_{ij}x_i(1-x_j/N_j),
\]
respectively.  Denote the finite operator by \(A_N\), and the coordinate
projection onto its box by \(P_N\).  Then every \(A_N\) satisfies the finite
positive stability criterion.  As all \(N_j\to\infty\), the following
convergence is strong on \(X_\kappa\), uniformly on compact time intervals:
\[
 e^{tA_N}P_N\longrightarrow T(t).
\]
\end{theorem}

\begin{proof}
In the convention that the second-order expression is
\(\sum_{i,j}Q_{ij}\partial_i\partial_j\) with \(Q\) symmetric, the corrected
principal matrix is
\[
 Q_N(y)=Q(y)-
 \frac12\sum_{\substack{j\in U\\i\ne j}}
 \frac{r_{ij}}{N_j}y_j^2(E_{ij}+E_{ji}).
\]
For each admissible contact matrix \(H\), whose entries are nonnegative,
\[
 \Tr(Q_N(y)H)=\Tr(Q(y)H)
       -\sum_{\substack{j\in U\\i\ne j}}
                 \frac{r_{ij}}{N_j}y_j^2H_{ij}\le0.
\]
The immigration correction is first order, and the binary sites are
unchanged because \(N_j\ge2\).  The finite contact theorem applies.

Split \(A_N=L_N+R_N\) by unrestricted-degree increase.  The increasing
parts satisfy \(\|R_N\|\le\rho\) and converge strongly to \(R\).
For a nonnegative input in the finite box, variation of constants gives
\[
 U(t)f-e^{tL_N}f
 =\int_0^tU(t-s)(L-L_N)e^{sL_N}f\,ds\ge0.
\]
The finite trajectory is polynomial-valued and belongs to \(D(L)\).
The integrand is nonnegative because only off-diagonal entries were
reduced; their diagonals agree exactly.  Therefore
\(\|e^{tL_N}P_N\|\le Me^{\omega t}\).

On every fixed \(W_m\), the matrices \(L_N\) converge to \(L\) once
\(N_j\ge m\).  Finite-dimensional exponential convergence and density,
with the preceding uniform bound, imply
\(e^{tL_N}P_N\to U(t)\) strongly and uniformly on compact intervals.
The Dyson expansions adding \(R_N\) have uniformly summable majorants
\[
 M^{k+1}e^{\omega t}\rho^k t^k/k!.
\]
Each fixed term converges to the corresponding term for \(U,R\);
summation proves the claim.  The projection \(P_N\), including at time
zero, is essential in this full-space statement.
\end{proof}

\begin{corollary}\label{ana:mixed-sufficiency}
The rate and contact conditions in the positive mixed-capacity
classification, together with the cutoff generation criterion, imply
preservation of \(\mathcal{LP}_{+,\kappa}\).
\end{corollary}
\begin{proof}
For stable polynomial input, sufficiently large \(P_N\) act as the
identity.  The finite outputs are stable and nonnegative.  Their
\(\ell^1\) limit belongs to \(\mathcal{LP}_{+,\kappa}\) by
Lemma~\ref{pos:lp-closure}, normalizing at \(\mathbf1\) if the probability-measure
form of that theorem is used.  A nonzero positive input has nonzero
output: the diagonal entry \(T(t)_{xx}\) is positive by subdivision of
time and strong continuity.  Positive Jensen approximants converge in
\(\ell^1\) to every member of \(\mathcal{LP}_{+,\kappa}\).
Boundedness of \(T(t)\) and the same closure theorem complete the proof.
\end{proof}

\subsection{Further generation criteria and topology}

For comparison with other analytic topologies, let
\[
 Y_R=\left\{f=\sum_xf_xz^x:\sum_x|f_x|R^{|x|}<\infty\right\},
 \qquad \mathcal E_{\rm ent}=\bigcap_{R\ge1}Y_R .
\]
These coefficient seminorms give the usual compact-open topology on
entire functions in finitely many variables.  The map \(f(z)\mapsto f(Rz)\)
is an isometry \(Y_R\to Y_1\), rather than a bounded similarity of \(Y_1\)
when \(R>1\).  In the unrestricted positive canonical expression below,
this conjugation replaces
\[
 \lambda\mapsto R\lambda,\qquad b\mapsto b/R,\qquad
 \nu\mapsto\nu/R,
\]
and leaves \(c,M,\beta,\eta\) unchanged.  With coordinate weights \(R_i\),
the off-diagonal coefficient \(m_{ij}\) becomes \(m_{ij}R_j/R_i\).
These assertions follow by conjugating multiplication and differentiation
on each monomial.

\begin{proposition}\label{ana:entire-cutoff}
Let \(L\) be a degree-nonincreasing polynomial operator, and let its
consistent finite-degree flows be \(U_N(t)\).  They extend to a locally
equicontinuous \(C_0\)-semigroup on \(\mathcal E_{\rm ent}\) if and only if
for every \(R\ge1\) and \(\tau>0\) there are \(S\ge1\) and \(C<\infty\)
such that
\[
 \|U_N(t)p\|_{Y_R}\le C\|p\|_{Y_S}
 \quad(N\ge0,\ 0\le t\le\tau,\ p\in\mathcal P_N).
\]
Such an extension is unique.
\end{proposition}
\begin{proof}
Necessity is the definition of local equicontinuity, since finitely many
coefficient seminorms can be dominated by their largest radius.
For sufficiency, Taylor polynomials are dense in every seminorm.
The estimate makes their evolved sequences Cauchy, uniformly in
\(0\le t\le\tau\), in every target seminorm.  Their limits define
continuous operators satisfying the same estimate.  Density transfers
the semigroup law, and uniform approximation by a polynomial orbit
proves strong continuity.  The same density proves uniqueness.
\end{proof}

\begin{example}\label{ana:topology-warning}
The Euler flow \(e^{tE}f(z)=f(e^tz)\) acts continuously on all entire
functions, but for every fixed \(R\) and \(t>0\) its norm on \(Y_R\)
is infinite, as the degree-\(n\) monomial has norm ratio \(e^{tn}\).
Conversely, the positive stability-preserving polynomial flow
generated by \(z\partial_z^2\) cannot extend globally to all positive
entire inputs.  Indeed,
\[
 [z]e^{tz\partial_z^2}z^n=n!t^{n-1}\qquad(n\ge1).
\]
This follows by applying the lowering operator \(n-1\) times and dividing
by \((n-1)!\) in its exponential.  The coefficient of \(z\) in the evolved
Taylor truncations of \(e^{az}\), with \(a>0\), is therefore
\(\sum_{n=1}^N a^nt^{n-1}\), which diverges when \(at\ge1\).
Polynomial stability of this flow can be seen on any fixed degree bound
\(N\): it is the dilute limit of the one-site generators
\((z-z^2/K)\partial_z^2\), \(K\ge N\), which satisfy the finite positive
criterion.  Their restrictions to degree at most \(N\) converge as
\(K\to\infty\).  This example does not satisfy coefficient-\(\ell^1\)
generation and does not contradict Theorem~\ref{ana:unary-l1}.
\end{example}

For the unrestricted positive canonical operator, write
\[
 A=cI+\lambda\cdot z+
 \sum_i\left(b_i+\sum_jm_{ij}z_j\right)\partial_i
 +E D_\nu-\beta E(E-1)-\sum_i\eta_i z_i^2\partial_i^2,
\]
where \(E=\sum_i z_i\partial_i\), \(D_\nu=\sum_i\nu_i\partial_i\),
and the signs are those in the positive classification.  On the probability
simplex define
\[
 F(p)=\nu\cdot p-\beta-\sum_i\eta_i p_i^2.
\]

\begin{proposition}\label{ana:quadratic-growth}
Existence of a positive coefficient-\(\ell^1\) realization forces
\(F(p)\le0\) throughout the simplex.  The strict inequality
\(\max F<0\) is sufficient for such a realization.
\end{proposition}
\begin{proof}
For a finite set of basis states \(S\), let \(B\) be the principal
matrix of the generator on \(S\).  Positivity gives
\[
 P_ST(t)|_S\ge e^{tB}
\]
entrywise.  Indeed, the left side, denoted \(G(t)\), satisfies
\(G'(t)=G(t)B+R(t)\), where \(R(t)\ge0\), by differentiating on the
finite set of generator-domain columns.  Variation of constants proves
the comparison.

Choose \(x^{(N)}/N\to p\), with zero coordinates when \(p_i=0\), and retain
the finite states \(x^{(N)}-\alpha\), where \(0\le\alpha\le x^{(N)}\)
and \(|\alpha|\le K\).  For sufficiently large \(N\), this is the
fixed index set of such \(\alpha\) supported on \(\{i:p_i>0\}\).
Divide their principal matrix by \(N^2\).  It converges to the matrix with diagonal
\[
 -D(p),\qquad D(p)=\beta+\sum_i\eta_i p_i^2,
\]
and rates \(\nu_i p_i\) from \(\alpha\) to \(\alpha+e_i\) when
\(|\alpha|<K\).  All linear terms vanish in this scaling.  Its exponential
has column mass at \(\alpha=0\) equal to
\[
 e^{-sD(p)}\sum_{j=0}^K\frac{(s\nu\cdot p)^j}{j!}.
\]
If \(T(t)\) is locally bounded by \(M\), apply the comparison at
\(t=s/N^2\) and let \(N\to\infty\).  Then let \(K\to\infty\).
We obtain \(M\ge e^{sF(p)}\) for every \(s>0\), forcing \(F(p)\le0\).

Conversely, the column sum at \(x\) is
\[
 (Az^x)(\mathbf1)
 =|x|^2 F(x/|x|)+O(|x|)+O(1),
\]
with a uniform remainder.  If \(F\le-\varepsilon\), these column sums
are bounded above.  Corollary~\ref{ana:column-sum} applies.
\end{proof}

\begin{proposition}[Exact unary generation criterion]\label{ana:unary-l1}
For
\[
 A=c+\lambda z+(b+mz)\partial_z+
             \nu z\partial_z^2-\chi z^2\partial_z^2,
 \qquad \lambda,b,\nu,\chi\ge0,\quad c,m\in\R,
\]
the positive coefficient-\(\ell^1\) realization exists if and only if
\[
 \chi>\nu
 \quad\hbox{or}\quad
 \chi=\nu\ \hbox{and}\ m+b\le0 .
\]
\end{proposition}
\begin{proof}
The terms \(cI+\lambda z\) are bounded.  Removing them from an existing
generator preserves generation; positivity of the resulting semigroup
follows from its Metzler restrictions to finite degree cutoffs and
density, rather than from an assertion that subtracting a positive
operator always preserves positivity.  Conversely, add \(\lambda z\)
by positive bounded perturbation and \(cI\) by multiplying the semigroup
by \(e^{ct}\).  We may set \(c=\lambda=0\).
The column sum is
\[
 (\nu-\chi)n(n-1)+(m+b)n.
\]
This proves sufficiency, while Proposition~\ref{ana:quadratic-growth}
excludes \(\chi<\nu\).

Suppose \(\chi=\nu\) and \(a=m+b>0\).  Then \(A=Q+aE\), where \(Q\)
is pure death with rates \(d_k=bk+\nu k(k-1)\).
On the finite cutoff, the elementary finite-state Feynman--Kac formula
(obtained by conditioning on the first holding time) gives
\[
 \|e^{tA}z^n\|_1
 =\mathbb E_n\exp\left(a\int_0^t N_s\,ds\right).
\]
If \(\nu>0\), fix \(t>0\) and choose \(K\) with
\(\sum_{k>K}d_k^{-1}<t/4\).
The independent exponential holding times \(\tau_k\) give
\[
 H_n=\sum_{k=K+1}^n\tau_k,\qquad
 Z_n=\sum_{k=K+1}^n k\tau_k.
\]
Markov's inequality yields \(\Pr(H_n>t)\le1/4\), while
\[
 \mathbb EZ_n=\nu^{-1}\log n+O(1),\qquad
 \operatorname{Var}Z_n=\sum_{k=K+1}^n k^2/d_k^2=O(1).
\]
For large \(n\), Chebyshev's inequality gives
\(\Pr(Z_n\ge(\log n)/(2\nu))\ge3/4\).
On an event of probability at least \(1/2\), both estimates hold, and
the integrated population up to time \(t\) is at least \(Z_n\).
Thus \(\|e^{tA}z^n\|_1\ge\tfrac12 n^{a/(2\nu)}\), contradicting boundedness.
If \(\nu=0\), direct solution instead gives
\[
 \|e^{tA}z^n\|_1
 =\left(e^{mt}+b\int_0^t e^{ms}\,ds\right)^n .
\]
The parenthesis is \(1+at+O(t^2)>1\) for small positive \(t\), giving
the same contradiction.
\end{proof}

\subsection{Automatic filtration of polynomial-valued semigroups}

Let \(\mathcal P=\C[z_1,\ldots,z_d]\), let
\(\mathcal P_N=\{f:\deg f\le N\}\), and give \(\mathcal P\) the strict
locally convex direct-limit topology of these finite-dimensional spaces.
This is the finest locally convex topology on \(\mathcal P\).

\begin{lemma}[A degree bound from the constant polynomial]
\label{ana:degree-defect}
An injective complex-linear stability preserver
\(T:\mathcal P\to\mathcal P\) satisfies
\[
 \deg Tp\le\deg p+\deg T1\qquad(p\ne0).
\]
\end{lemma}
\begin{proof}
Fix a monomial \(z^\alpha\), a coordinate \(i\), and put
\(G=Tz^\alpha\), \(F=T(z_i z^\alpha)\).  For every \(w\in\HH\),
the input \(z^\alpha(z_i+w)\) is stable, and its nonzero image is
\(F+wG\).  Thus \(F/G\) is holomorphic on \(\HH^d\) and has nonnegative
imaginary part there: a value with negative imaginary part would permit
the choice \(w=-F/G\in\HH\).

Choose \(v>0\) where the leading homogeneous parts of both \(F,G\)
are nonzero.  The rational function \(F(sv)/G(sv)\) maps \(\HH\) into
its closure.  Its numerator degree minus denominator degree is at most
one.  Indeed, if that difference were \(m\ge2\), its asymptotic form
on \(s=Re^{i\theta}\) would be \(aR^m e^{im\theta}(1+o(1))\),
with \(a\ne0\).  As \(0<\theta<\pi\), some fixed angle makes the leading
imaginary part negative, contradicting the mapping property for large
\(R\).  Cancellation does not change the degree difference.
Consequently \(\deg F\le\deg G+1\).  Iteration along a monomial
multiplication chain, followed by linearity, proves the assertion.
\end{proof}

\begin{theorem}[Automatic polynomial filtration]\label{ana:filtration}
Every strongly continuous complex-linear semigroup on \(\mathcal P\)
preserving upper-half-plane stability has an everywhere-defined locally
finite generator \(A\).  Its operators are linear automorphisms of the
vector space \(\mathcal P\), and
\[
 A\mathcal P_N\subseteq\mathcal P_N,\qquad
 T(t)\mathcal P_N=\mathcal P_N
 \quad(N\ge0,\ t\ge0).
\]
The same conclusion holds algebraically if \(A\) is assumed locally finite
and its finite-orbit exponentials preserve stability.
\end{theorem}
\begin{proof}
A bounded subset of this direct limit lies in some \(\mathcal P_N\).
For the needed implication, otherwise choose polynomials with selected
nonzero monomials of strictly increasing degrees.  Assign successively
positive monomial weights making their weighted coefficient seminorms
at least \(1,2,\ldots\).  Every such weighted seminorm is continuous in
the finest locally convex topology, contradicting boundedness.

For fixed \(f\), its orbit over \(0\le t\le1\) is compact, hence contained
in a finite-dimensional space.  Put
\[
 W_r=\operatorname{span}\{T(t)f:0\le t\le r\}.
\]
The finite integer dimensions have a minimum for \(r>0\).  Choose
\(\varepsilon>0\) sufficiently small that
\(W_\varepsilon=W_{2\varepsilon}\).
Then \(T(s)W_\varepsilon\subseteq W_\varepsilon\) for
\(0\le s\le\varepsilon\), and iteration proves invariance for all \(s\ge0\).
Finite-dimensional continuous semigroup theory gives an exponential
on this space.  Its generator agrees on overlapping invariant spaces
because it is the derivative of the same orbit.  These derivatives define
a locally finite linear \(A\) on all of \(\mathcal P\), and the finite-space
inverses \(e^{-tA}\) agree as well.  Thus every \(T(t)\) is injective.

Recursive evaluation on monomials gives the canonical locally finite
normal expansion \(A=\sum_\alpha P_\alpha(z)\partial^\alpha\), with
polynomial coefficients but no assumed order bound.  Apply the flow to
\(\prod_i(z_i-y_i)^{\alpha_i}\) and restrict to \(z=y+sv\), \(v>0\).
This one orbit lies in a finite-dimensional polynomial space, so the
restricted degrees are uniformly bounded.  At the contact \(y\),
the first variation is exactly \(\alpha!P_\alpha(y)\).
Lemma~\ref{gen:cluster-moments} makes it zero when \(|\alpha|\ge3\).
Polynomial identity proves \(P_\alpha=0\) for those \(\alpha\).
Thus \(A\) is a finite-order polynomial differential operator.

The entire orbit of \(1\) lies in one finite invariant space; choose
\(D\) with \(\deg T(t)1\le D\) for all \(t\ge0\).
Lemma~\ref{ana:degree-defect} gives
\[
 T(t)\mathcal P_N\subseteq\mathcal P_{N+D}.
\]
Analyticity on finite orbit spaces allows differentiation to every order:
\[
 A^k\mathcal P_N\subseteq\mathcal P_{N+D}
 \quad\hbox{for every }k,N.
\]
Suppose the maximal degree shift of a normal-ordered term in \(A\)
is \(s>0\), and let \(B\ne0\) be the homogeneous component of shift \(s\).
The component of shift \(ks\) in \(A^k\) is \(B^k\).
For \(ks>D\), the preceding inclusions force \(B^k\) to vanish on
every polynomial.  This is impossible: if \(B\) has differential order
\(r\), the principal symbol of \(B^k\) is the \(k\)-th power of its
nonzero symbol in the integral domain \(\C[z,\xi]\).
Consequently all shifts are nonpositive.  This proves invariance of
each \(\mathcal P_N\); its exponential is invertible on that space.
The algebraic locally finite case uses the same argument after the first
paragraphs, since finite invariant orbit spaces are already available.
\end{proof}

\begin{remark}
The topological necessity remains true for any Hausdorff vector-space
topology on \(\mathcal P\) for which the given maps are continuous and the
semigroup is strongly continuous.  Shkarin's orbit theorem
\cite[preprint Lemma 5.4]{Shkarin11} states that an orbit span has either finite
dimension or continuum dimension.  Countable Hamel dimension forces the
first alternative, after which the algebraic proof applies.  The converse
classification below is asserted for the specified strict direct-limit
topology; continuity in another topology is a separate requirement.
\end{remark}

\begin{theorem}[All-polynomial classification]\label{ana:polynomial-classification}
The semigroups in Theorem~\ref{ana:filtration} are precisely the
polynomial exponentials of
\[
 A=cI+\sum_i\left(b_i+i\sigma_i+\sum_jm_{ij}z_j\right)\partial_i
   -\sum_i q_i(z_i)\partial_i^2
   -\sum_{i<j}c_{ij}\partial_i\partial_j,
\]
where \(c\in\C\), \(b_i\in\R\), \(\sigma_i\ge0\), \(M=(m_{ij})\)
is real Metzler, each \(q_i\) is a real nonnegative polynomial of degree
at most two, and \(c_{ij}\ge0\).  Their flows preserve stability after
adjoining any finite collection of inert polynomial variables.
\end{theorem}
\begin{proof}
Filtration and canonical recursion give \(\deg P_\alpha\le|\alpha|\).
The preceding contact argument gives order at most two.
Double contacts make every second-order coefficient real and nonpositive
on \(\R^d\).  A simple-root contact makes
\(\Im P_{e_i}(y)\ge0\) everywhere.  Since \(P_{e_i}\) is affine, its
imaginary part is a constant \(\sigma_i\ge0\); \(P_0=c\) is constant.

We require the following homogeneous sign fact.  A homogeneous complex
stable polynomial has all coefficients in one common phase, and after
removing that phase they are nonnegative.  Indeed, homogeneity also gives
nonvanishing on the lower half-plane.  For \(u,v>0\),
\(s\mapsto f(u+sv)\) has only real roots and no nonnegative root;
its leading coefficient is \(f(v)\ne0\).
Thus \(f(u)/f(v)>0\).  Fixing \(v\) gives a common phase on the positive
orthant.  The negative real roots give nonnegative directional
derivatives there.  Letting directions tend to coordinate vectors, and
repeating this argument for stable homogeneous derivatives
\cite{BB09}, proves
nonnegativity of all coefficients.

After removal of the scalar phase \(i\Im c\), the highest homogeneous
block of the generator is real.  Its exponential applied to a monomial
has a highest homogeneous part which is stable or zero by scaling and
Hurwitz.  Near zero time its original monomial coefficient is positive.
The homogeneous sign fact makes all other coefficients of that block
nonnegative.  Its generator matrix is therefore Metzler in every degree.

Put \(U_i=P_{2e_i}\).  Its homogeneous quadratic part is negative
semidefinite.  For \(j\ne i\), the shift \(2e_j-2e_i\) can only be produced
by the \(z_j^2\) coefficient of \(U_i\).  Homogeneous Metzler positivity
forces that coefficient to be nonnegative; nonpositivity of \(U_i\)
forces it to vanish.  A semidefinite quadratic form with zero diagonal
entry has its corresponding row zero.  Hence the quadratic part of
\(U_i\) depends only on \(z_i\), and global nonpositivity also eliminates
linear coefficients in other variables.  Thus \(U_i=-q_i(z_i)\).

For \(U_{ij}=P_{e_i+e_j}\), the same unique-shift argument at every
coordinate outside \(\{i,j\}\) eliminates its outside square coefficients
and their rows.  Its quadratic part depends only on \(z_i,z_j\).
Take an input with \(x_i=N,x_j=1\), all other exponents zero, and inspect
the shift \(e_i-e_j\).  Its rate is
\[
 N[z_i^2]U_{ij}+m_{ji}.
\]
Here all other possible second-order contributions have already been
eliminated by the preceding variable-dependence conclusions.
Positivity for every \(N\) gives \([z_i^2]U_{ij}\ge0\), hence equality.
Interchange \(i,j\).  The negative semidefinite quadratic part now has
all diagonal entries zero, so vanishes.  An affine polynomial nonpositive
on all of \(\R^d\) is constant; write it as \(-c_{ij}\), \(c_{ij}\ge0\).
The degree-one block finally yields \(m_{ij}\ge0\) when \(i\ne j\).

Conversely, the affine first-order flow is substitution by
\[
 z\longmapsto e^{tM}z+\int_0^t e^{sM}(b+i\sigma)\,ds.
\]
The matrix \(e^{tM}\) is nonnegative with strictly positive diagonal, so
this maps \(\HH^d\) into itself.  Each unary term
\(-q_i(z_i)\partial_i^2\) and each binary term
\(-c_{ij}\partial_i\partial_j\) preserves every finite coordinate box and
satisfies Theorem~\ref{gen:main-real}.  Their flows therefore preserve
complex stability, with arbitrary inert variables.  On every
\(\mathcal P_N\), Lie products of these flows and the scalar exponential
converge to \(e^{tA}\).  Hurwitz and invertibility on \(\mathcal P_N\)
give nonzero stable outputs.  The consistent finite-dimensional flows
are strongly continuous in the strict direct-limit topology, which
proves sufficiency.
\end{proof}

\subsection{A common space of entire functions}

Polynomial-valued dynamics must be distinguished both from
coefficient-\(\ell^1\) generation and from continuous dynamics on all entire
functions.  A common space for the preceding full complex cone is
\[
 \mathcal E_{2,0}
 =\left\{f(z)=\sum_\alpha f_\alpha z^\alpha:
   \|f\|_R=\sum_\alpha |f_\alpha|\sqrt{\alpha!}\,R^{|\alpha|}
                    <\infty\ \hbox{for every }R>0\right\}.
\]
Positive integer radii suffice to define its Fr\'echet topology.
In particular
\[
 |f(z)|\le\|f\|_R
       \exp\left(\frac{\sum_i|z_i|^2}{2R^2}\right),
\]
and Taylor polynomials converge in every seminorm.

\begin{theorem}\label{ana:minimal-quadratic}
Every generator in Theorem~\ref{ana:polynomial-classification} has a unique
locally equicontinuous \(C_0\)-semigroup extension to \(\mathcal E_{2,0}\).
Its generator is the same continuous differential operator, defined on
all of \(\mathcal E_{2,0}\).  For every \(R>0\) there are \(C_R,K_R\ge0\)
such that
\[
 \|T(t)f\|_R\le e^{C_Rt}\|f\|_{R e^{K_Rt}}
 \quad(t\ge0).
\]
It preserves the complex Laguerre--P\'olya class intersected with
\(\mathcal E_{2,0}\), with zero output allowed.
\end{theorem}
\begin{proof}
Write \(q_i=a_i z_i^2+d_i z_i+e_i\).  Work first on any finite degree
cutoff, with weight \(w_R(\alpha)=\sqrt{\alpha!}R^{|\alpha|}\).
For a complex matrix flow, the upper right derivative of the weighted
absolute coefficient sum is bounded by the weighted column sums formed
from the real diagonal and absolute off-diagonal entries.  Their
contributions are bounded by
\[
\begin{array}{ll}
 (b_i+i\sigma_i)\partial_i:
       &|b_i+i\sigma_i|\sqrt{\alpha_i}/R,\\
 m_{ij}z_j\partial_i\quad(i\ne j):
       &m_{ij}\sqrt{\alpha_i(\alpha_j+1)},\\
 -d_i z_i\partial_i^2:
       &(|d_i|/R)\sqrt{\alpha_i}(\alpha_i-1),\\
 -e_i\partial_i^2:
       &(e_i/R^2)\sqrt{\alpha_i(\alpha_i-1)},\\
 -c_{ij}\partial_i\partial_j:
       &(c_{ij}/R^2)\sqrt{\alpha_i\alpha_j}.
\end{array}
\]
The diagonal contains \(-a_i\alpha_i(\alpha_i-1)\).  When \(a_i>0\),
\[
 \sup_{n\ge0}\{-a_i n(n-1)+(|d_i|/R)\sqrt n(n-1)\}<\infty.
\]
When \(a_i=0\), nonnegativity of \(q_i\) forces \(d_i=0\).
All remaining contributions, including the scalar and linear diagonal,
are bounded by \(C_R+K_R|\alpha|\), using elementary arithmetic-geometric
mean inequalities.  The constants may be chosen to work at every radius
\(r\ge R\), since the radius-dependent terms decrease.

For the finite polynomial flow put
\(H(t,r)=\sum_\alpha |f_\alpha(t)|\sqrt{\alpha!}r^{|\alpha|}\).
The upper derivative inequality is
\[
 D_t^+H(t,r)\le C_RH(t,r)+K_Rr\partial_r H(t,r),\qquad r\ge R.
\]
Along \(r(s)=R e^{K_R(t-s)}\), the radius derivative cancels the last term.
Gronwall proves the asserted estimate, uniformly in every cutoff.

Taylor approximation now constructs \(T(t)f\) in each seminorm,
uniformly on compact time intervals.  The same estimate transfers the
semigroup law and strong continuity from polynomials.  Every
differential monomial in \(A\) is continuous between a larger-radius
seminorm and a smaller one, since an exponential ratio absorbs all
polynomial factors in \(\alpha\).  Polynomial approximation in
\[
 T(t)p-p=\int_0^tT(s)Ap\,ds
\]
therefore proves the identity for all \(f\), identifies the full generator
domain, and proves uniqueness of the extension.

For preservation use the multivariate Jensen polynomials
\[
 J_nf(z)=\sum_{\alpha_i\le n}f_\alpha
               \prod_i\frac{(n)_{\alpha_i}}{n^{\alpha_i}}z^\alpha .
\]
They are stable or zero for complex Laguerre--P\'olya inputs by
\cite[Theorem 5.1 and Lemma 5.2]{BB09}.  Their multipliers lie in \([0,1]\)
and converge pointwise to one, so dominated convergence gives
\(J_nf\to f\) in every displayed seminorm.  The flow estimate and Hurwitz
now prove preservation.
\end{proof}

\begin{proposition}[A sharp full-class obstruction]\label{ana:gaussian-obstruction}
For \(a,t>0\) with \(4at\ge1\), the polynomial operator
\(e^{-t\partial_z^2}\) has no compact-open continuous extension even to
the full real Laguerre--P\'olya class that agrees with it on real-rooted
polynomials.
\end{proposition}
\begin{proof}
The real-rooted polynomials \(g_N(z)=(1-az^2/N)^N\) converge locally
uniformly to \(e^{-az^2}\).  Direct differentiation gives
\[
 (e^{-t\partial_z^2}g_N)(0)
 =\sum_{k=0}^N\frac{(N)_k}{N^k}\binom{2k}{k}(at)^k .
\]
All summands are nonnegative and each fixed multiplier tends to one.
The infinite series without those multipliers diverges for \(4at\ge1\),
including equality since \(\binom{2k}{k}\sim4^k/\sqrt{\pi k}\).
Keeping any finite initial segment and then increasing its length proves
divergence of the displayed values.  This contradicts continuity of
evaluation at zero for the proposed extension.
\end{proof}

\subsection{Exactly when a unary flow acts on all entire functions}

\begin{theorem}\label{ana:entire-unary}
Let
\[
 A=\gamma I+(u+mz)\partial_z-(az^2+bz+c)\partial_z^2,
\]
where \(\gamma\in\C\), \(\Im u\ge0\), \(m\in\R\), and
\(q(z)=az^2+bz+c\) is real and nonnegative on \(\R\).
Its polynomial semigroup extends to a continuous \(C_0\)-semigroup
on the compact-open space \(\mathcal H(\C)\) of all entire functions if
and only if \(q=0\) or \(a>0\).  In these cases the extension is unique,
its generator is \(A\) on all entire functions, and it preserves the
compact-open closure of stable polynomials.  If \(q\) is a nonzero
constant, no individual positive-time operator has a continuous
extension to \(\mathcal H(\C)\).
\end{theorem}
\begin{proof}
The seminorms \(\|f\|_R=\sum_n|f_n|R^n\) define the compact-open topology:
they dominate the supremum on the radius-\(R\) disk and are at most twice
the supremum on the radius-\(2R\) disk.
If \(a>0\), choose \(R\) so large that
\(\delta_R=a-|b|/R-c/R^2>0\).
On every finite degree cutoff, absolute weighted column sums with the
real diagonal give logarithmic norm at most
\[
 M_R=\sup_{n\ge0}\{
 \Re\gamma+(m+|u|/R)n-\delta_R n(n-1)\}<\infty.
\]
Thus \(\|e^{tA}p\|_R\le e^{M_Rt}\|p\|_R\), uniformly in the degree.
Every target radius is dominated by a sufficiently large such \(R\).
Taylor approximation constructs the unique locally equicontinuous
extension and transfers its semigroup law and strong continuity.
The differential expression is continuous on \(\mathcal H(\C)\);
passing to the limit in the integrated polynomial equation identifies
its generator on all of that space.  Hurwitz gives preservation of the
polynomial-closure class.
If \(q=0\), all assertions follow directly from
\[
 T(t)f(z)=e^{\gamma t}
 f\left(e^{mt}z+u\int_0^t e^{ms}\,ds\right).
\]

The only remaining case is \(a=b=0,c>0\).
If \(T(t)\) extended continuously for a fixed \(t>0\), evaluation at zero
after \(T(t)\) would be bounded by \(C\) times the supremum on some disk
of radius \(R\).  Polynomial moments give
\[
 (T(t)e^{wz})(0)
 =\exp\left(\gamma t+uB_tw-cC_tw^2\right),
\]
where
\[
 B_t=\int_0^t e^{ms}\,ds,\quad C_t=\int_0^t e^{2ms}\,ds>0.
\]
This identity does not presuppose existence on exponentials: expanding
the right side in \(w\) gives the unique polynomial solutions for all
monomial data, as is verified with
the ansatz \(e^{v(t)z+d(t)}\), \(v'=mv\),
\(d'=\gamma+uv-cv^2\).
Continuity of the proposed extension then permits passage from the
Taylor polynomials of \(e^{wz}\).
At \(w=iY\) the displayed modulus grows like
\(\exp(cC_tY^2+O(|Y|))\), whereas continuity bounds it by
\(C e^{R|Y|}\).  This contradiction proves necessity.
\end{proof}

\begin{proposition}\label{ana:all-entire-separation}
For every \(t>0\), the continuous entire-function operator
\(e^{-tz^2\partial_z^2}\) fails to preserve the class of all entire
functions nonvanishing on \(\HH\), although it preserves the
compact-open closure of stable polynomials.
\end{proposition}
\begin{proof}
Put
\[
 G_t(z)=\sum_{n\ge0}e^{-tn(n-1)}\frac{z^n}{n!}.
\]
This is \(T(t)e^z\).  Its order is zero, since
\[
 \limsup_{n\to\infty}
 \frac{n\log n}{tn(n-1)+\log n!}=0.
\]
It is nonconstant, as \(G_t'(0)=1\). A zero-free entire function of
order zero is constant, by the zero-order case of Hadamard
factorization \cite{Boas54}. Here is a direct proof of the fact used.
If such a function is \(e^h\) with \(h\) entire, its growth gives
\(\Re h(z)\le R^{1/2}\) on \(|z|\le R\) for sufficiently large \(R\).
For a fixed center \(a\), the nonnegative function
\((R+|a|)^{1/2}-\Re h(a+Re^{i\theta})\) has average
\((R+|a|)^{1/2}-\Re h(a)\). Bounding its first Fourier coefficient
by its average gives
\[
 |h'(a)|\le
 \frac{2((R+|a|)^{1/2}-\Re h(a))}{R}\longrightarrow0.
\]
Thus \(h\) is constant. Consequently \(G_t\) has a zero.
The polynomials \(T(t)(1+z/N)^N\) are real-rooted with nonnegative
coefficients, and Theorem~\ref{ana:entire-unary} makes them converge
locally uniformly to \(G_t\).
Hurwitz off the nonpositive real axis, and \(G_t(0)=1\), show that all
zeros of \(G_t\) are strictly negative real.  Choose one, \(-r\).
The entire function \(e^{iz}\) is globally zero-free, but its image
\(G_t(iz)\) vanishes at \(z=ir\in\HH\).
\end{proof}

% END SECTION analytic_generation

% BEGIN SECTION static_channels
\section{Static channels and stable encodings}
\label{stat:section}

This section classifies invertible Hermiticity-preserving, trace-preserving
maps with Lee--Yang full Choi tensors, and gives structural results for
singular channels. The
variables associated with different physical qubits are always distinct.
In particular, permutation of physical wires is allowed, whereas an arbitrary
change of the computational basis is not.

\subsection{Coefficient tensors and complete preservation}

For a matrix \(M:(\C^2)^{\otimes m}\to(\C^2)^{\otimes n}\), put
\begin{equation}
 G_M(z,w)=\sum_{a\in\{0,1\}^n}\sum_{b\in\{0,1\}^m}
 M_{ab}z^aw^b.
 \label{stat:matrix-polynomial}
\end{equation}
We call \(M\) \emph{tensor-LY} if this polynomial is nonzero whenever all
its variables belong to the open unit disk \(\D\).  Zero is excluded from
this definition.  For a complex-linear map
\(\Phi:M_{2^m}(\C)\to M_{2^n}(\C)\), our Choi convention is
\begin{equation}
 J(\Phi)=\sum_{c,d}\Phi(|c\rangle\langle d|)
                 \otimes |c\rangle\langle d|.
 \label{stat:choi}
\end{equation}
Thus full Choi--LY means stability in all \(2(m+n)\) variables of
\(G_{J(\Phi)}\).  Hermiticity-preserving and trace-preserving maps are
abbreviated HP and TP, respectively.

We use the disk tensor-contraction theorem in the following form: contracting
two binary indices of a tensor-LY polynomial with unit radii produces a
tensor-LY polynomial or zero.  Tensor products, permutation of indices, and
coefficientwise conjugation preserve tensor-LY.  These facts follow from
the classical Asano contraction theorem; the tensor and channel formulation
is described in \cite[Section~2]{WBG26}.  The possible zero result at unit
radius will be excluded explicitly where necessary.

\begin{lemma}[The Choi bridge]
\label{stat:bridge}
Let \(\Phi:M_{2^m}(\C)\to M_{2^n}(\C)\) be complex linear.
If \(J(\Phi)\) is tensor-LY, then \(\Phi\), and each of its extensions
by identity maps, sends tensor-LY matrices to tensor-LY matrices or zero.
Its Hilbert--Schmidt adjoint \(\Phi^\dagger\) is also full Choi--LY.
For a CPTP map, full Choi--LY is equivalent to preservation of tensor-LY
density matrices with every finite qubit reference system left inert.
\end{lemma}

\begin{proof}
The superoperator and Choi arrays satisfy
\[
 S(\Phi)_{(a,b),(c,d)}=\Phi(|c\rangle\langle d|)_{ab},
 \qquad
 J(\Phi)_{(a,c),(b,d)}=S(\Phi)_{(a,b),(c,d)}.
\]
They therefore differ only by a permutation of tensor indices.  Applying
\(\Phi\) is contraction of this array with the input coefficient array.
Input transposes occurring in the matrix convention also only permute
indices.  The contraction theorem proves the first assertion, including
inert indices.  The superoperator of \(\Phi^\dagger\) is obtained by
conjugating entries and interchanging the input and output index pairs,
which proves the adjoint assertion.

For a CPTP map the output of a density matrix has trace one, so cannot
be the zero alternative.  Conversely, take the normalized maximally
entangled vector between the input and an equally large reference:
\[
 |\Omega_m\rangle=2^{-m/2}\sum_c|c\rangle|c\rangle.
\]
Its amplitude polynomial is
\(2^{-m/2}\prod_i(1+x_i u_i)\), and its density-matrix polynomial is
the product of this polynomial and its coefficientwise conjugate in
independent variables.  Both are stable.  Applying \(\Phi\) to one
side produces \(2^{-m}J(\Phi)\), proving the converse.
\end{proof}

Ordinary multiplication of tensor-LY matrices is likewise a tensor
contraction, with the same zero alternative.  In particular, multiplication
on either side by an invertible tensor-LY matrix whose inverse is
tensor-LY preserves the class of tensor-LY matrices.

\subsection{An elementary matrix edge and hyperbolic pencils}

Write
\[
 X=\begin{pmatrix}0&1\\1&0\end{pmatrix},\qquad
 Y=\begin{pmatrix}0&-i\\i&0\end{pmatrix},\qquad
 Z=\begin{pmatrix}1&0\\0&-1\end{pmatrix}.
\]
Subscripts denote physical sites, with identity factors suppressed.  The
real vector spaces needed below are
\begin{equation}
 \mathfrak g=\operatorname{span}_{\R}\{X,Y,iZ\},\qquad
 \mathcal E=\R I+\mathfrak g,
 \qquad
 \mathfrak e_n=\C I+\bigoplus_{j=1}^n\mathfrak g_j.
 \label{stat:edge-spaces}
\end{equation}

\begin{lemma}[Two-sided tangency]
\label{stat:edge}
Suppose \(M(t)\) is a matrix-valued curve defined for real \(t\) near
zero, differentiable at zero, with \(M(0)=I\) and \(M'(0)=B\).
If \(M(t)\) is tensor-LY for both signs of all sufficiently small
\(t\), then \(B\in\mathfrak e_n\).
Conversely, \(B\in\mathfrak e_n\) implies that \(e^{tB}\) is
tensor-LY for every real \(t\).
\end{lemma}

\begin{proof}
Fix a site \(j\) and \(|\zeta|=1\).  Specializing \(z_j=\zeta\)
is justified by first using \(r\zeta\), \(r<1\), and then applying
Hurwitz's theorem.  The specialization is stable or zero.  At \(t=0\)
its value with all remaining variables zero is one, so it is not the
zero polynomial for sufficiently small \(|t|\).
Fix interior values \(u=(z_k,w_k)_{k\ne j}\), and put
\(D(u)=\prod_{k\ne j}(1+z_kw_k)\).  The resulting affine polynomial
in \(w_j\) is
\[
 D(u)(1+\zeta w_j)+tG_B(\zeta,w_j,u)+o(t).
\]
Its root \(r(t)\) is differentiable at zero, with
\(r(0)=-\zeta^{-1}\), and stability implies \(|r(t)|\ge1\) for
both signs of small \(t\).  Differentiating the root equation gives
\[
 \frac{r'(0)}{r(0)}
 =\frac{G_B(\zeta,-\zeta^{-1},u)}{D(u)}.
\]
The derivative of \(|r(t)|^2\) at its two-sided minimum is zero.
Consequently the holomorphic function on the remaining polydisk
\begin{equation}
 u\longmapsto\frac{G_B(\zeta,-\zeta^{-1},u)}{D(u)}
 \label{stat:contact-quotient}
\end{equation}
has purely imaginary values, and is therefore a purely imaginary
constant.

Expand \(B\) in the Pauli tensor basis, separating site \(j\):
\(B=\sum_{\sigma,\tau}b_{\sigma,\tau}\sigma_j\otimes\tau\).
At \((z_j,w_j)=(\zeta,-\zeta^{-1})\), the one-site polynomials are
\[
 G_I=0,\qquad G_X=\zeta-\zeta^{-1},\qquad
 G_Y=i(\zeta+\zeta^{-1}),\qquad G_Z=2.
\]
Since the Pauli polynomials on the other sites are linearly independent,
constancy of \eqref{stat:contact-quotient} gives, for every nonidentity
word \(\tau\),
\[
 b_{X,\tau}(\zeta-\zeta^{-1})
 +ib_{Y,\tau}(\zeta+\zeta^{-1})+2b_{Z,\tau}=0.
\]
Taking \(\zeta=1,-1,i\) forces all three coefficients to vanish.
As \(j\) was arbitrary, every Pauli coefficient of weight at least
two is zero.  Thus
\(B=cI+\sum_j(x_jX_j+y_jY_j+d_jZ_j)\).
The remaining scalar contact expression is
\[
 x_j(\zeta-\zeta^{-1})+iy_j(\zeta+\zeta^{-1})+2d_j.
\]
Its real part is zero.  The same three values of \(\zeta\) imply
\(x_j,y_j\in\R\) and \(d_j\in i\R\), as asserted.  The scalar
\(c\) is unrestricted.

For the converse, if
\(N=aI+xX+yY+izZ\in\mathcal E\), direct multiplication gives
\begin{equation}
 NZN^*=\delta Z,
 \qquad \delta=a^2+z^2-x^2-y^2=\det N.
 \label{stat:unit-identity}
\end{equation}
When \(\delta>0\), write
\(G_N(s,w)=A(s)+C(s)w\).  Equation~\eqref{stat:unit-identity}
implies
\[
 |A(s)|^2-|C(s)|^2=\delta(1-|s|^2)>0\qquad(|s|<1),
\]
so \(N\) is tensor-LY.  Its inverse has the same property.
For \(b=xX+yY+izZ\), one has \(b^2=(x^2+y^2-z^2)I\).
The exponential \(e^{tb}\) belongs to \(\mathcal E\) for real
\(t\) and has determinant one.  Therefore it is tensor-LY.  Finally,
for \(B=cI+\sum_jb_j\),
\(e^{tB}=e^{tc}\bigotimes_j e^{tb_j}\), which proves the assertion.
\end{proof}

The space \(\mathcal E\) is closed under multiplication.  Its members
sufficiently close to \(I\), and their inverses, are tensor-LY by
\eqref{stat:unit-identity}.  We call these matrices local stability
units.  A traceless element \(b=xX+yY+izZ\in\mathfrak g\) is
\emph{hyperbolic} if \(x^2+y^2-z^2>0\).  A diagonal phase
conjugation followed by conjugation with \(e^{rY}\), where
\(\tanh(2r)=z/\sqrt{x^2+y^2}\), carries it to
\(\sqrt{x^2+y^2-z^2}\,X\).  The conjugating matrices and their
inverses are local stability units.

\begin{lemma}[Hyperbolic affine pencil]
\label{stat:hyperbolic}
Let \(B=cI+\sum_jb_j\in\mathfrak e_n\), and suppose one
\(b_k\) is hyperbolic.  If \(I+tB\) is tensor-LY for some nonzero
real \(t\), then \(b_j=0\) for all \(j\ne k\) and \(c\in\R\).
\end{lemma}

\begin{proof}
Conjugating by a local stability unit, assume \(b_k=aX_k\),
\(a>0\).  For a one-site matrix \(b\) define
\(q_b(z,w)=G_b(z,w)/(1+zw)\).  The function
\[
 q_X(z,z)=\frac{2z}{1+z^2}
\]
attains every nonreal complex number for \(|z|<1\).  Indeed, the
inverse equation has two reciprocal roots; neither is on the unit
circle for a nonreal target, because the displayed quotient is real
there wherever finite.  Exactly one inverse root is therefore in
the disk.

If \(b_j\ne0\) for another site, \(q_{b_j}\) is a nonconstant
holomorphic function, since \(b_j\) is nonzero and traceless.
Freeze all other variable pairs at zero, and divide the polynomial
of \(I+tB\) by \(\prod_i(1+z_iw_i)\).  The result has the form
\(C+taq_X(z_k,w_k)+tq_{b_j}(z_j,w_j)\).  The open image of
\(q_{b_j}\) allows its variables to be chosen so that the target
needed from \(q_X\) is nonreal.  The preceding inverse construction
then gives an interior zero, a contradiction.  Once the other
components vanish, an imaginary part of \(c\) would make
\(-(1+tc)/(ta)\) nonreal and give the same contradiction.
\end{proof}

\subsection{Routing and factorization of visible outputs}

Let \(\Phi:M_{2^m}(\C)\to M_{2^n}(\C)\) be HP, TP and full
Choi--LY, and put \(\phi=\Phi^\dagger\).  An output site \(i\)
is called \emph{transversely visible} if at least one of
\(\phi(X_i),\phi(Y_i)\) is nonscalar.  Denote the set of such
sites by \(A\).

\begin{lemma}[Local routing]
\label{stat:routing}
For each \(i\in A\) there is a unique input site \(\pi(i)\) such
that, writing \(k=\pi(i)\),
\begin{equation}
 \begin{split}
 \phi(X_i)&=b_{ix}I+D_{i,11}X_k+D_{i,21}Y_k,\\
 \phi(Y_i)&=b_{iy}I+D_{i,12}X_k+D_{i,22}Y_k,\\
 \phi(Z_i)&=\lambda_i Z_k,
 \end{split}
 \label{stat:local-block}
\end{equation}
where all coefficients are real and \(D_i\ne0\).  No injectivity
of \(\Phi\) or \(D_i\) is required.
\end{lemma}

\begin{proof}
For a local \(U\in\mathfrak g_i\), the identity
\(U^2=\delta I\), \(\delta\in\R\), gives
\(e^{tU}=c_\delta(t)I+s_\delta(t)U\), with real analytic scalar
coefficients satisfying \(c_\delta(0)=1\), \(s_\delta(0)=0\),
and \(s_\delta'(0)=1\).  These exponentials are stability units.
Since \(\phi\) is unital, their images, divided by
\(c_\delta(t)\), give a two-sided stable pencil
\(I+r\phi(U)\) near zero.  A zero image is excluded near \(I\).
Lemma~\ref{stat:edge} puts \(\phi(U)\) in \(\mathfrak e_m\).
HP then implies
\[
 \phi(X_i),\phi(Y_i)\in\R I+
       \sum_k\operatorname{span}_{\R}\{X_k,Y_k\},
 \qquad
 \phi(Z_i)\in\R I+\sum_k\R Z_k.
\]
Every nonscalar transverse image has a hyperbolic component, so
Lemma~\ref{stat:hyperbolic} confines it to one input site.  If two
transverse directions used different sites, their sum would have
two nonzero hyperbolic components.  Thus the whole transverse plane
uses the same site whenever its nonscalar image is nonzero.

Choose a transverse \(U_i\) with nonzero traceless image on this
site, and apply the argument to \(U_i+i\varepsilon Z_i\) for a
small nonzero real \(\varepsilon\).  The routed component remains
hyperbolic.  Lemma~\ref{stat:hyperbolic} eliminates all other
longitudinal components, and its real-scalar conclusion eliminates
the scalar coefficient of \(\phi(Z_i)\).  This proves
\eqref{stat:local-block}.
\end{proof}

\begin{theorem}[Visible-output decomposition]
\label{stat:visible}
Let \(\Phi:M_{2^m}(\C)\to M_{2^n}(\C)\) be HP, TP and full
Choi--LY.  With the preceding notation, put \(K=\pi(A)\) and
\(G_k=\pi^{-1}(k)\).  There are HP, TP, full Choi--LY maps
\(\Xi_k:M_2(\C)\to M_{2^{|G_k|}}(\C)\) such that, after
ordering the physical output groups,
\begin{equation}
 \Tr_{A^c}\circ\Phi
 =\left(\bigotimes_{k\in K}\Xi_k\right)\circ\Tr_{K^c}.
 \label{stat:visible-form}
\end{equation}
The routing function is unique.  If \(\Phi\) is CP, every
\(\Xi_k\) is CP.  The trace on the right is the ordinary,
unnormalized partial trace.
\end{theorem}

\begin{proof}
For each \(i\in A\), choose a small real open set
\(\mathcal O_i\subset\mathfrak g_i\) about a transverse direction
whose image is nonscalar.  Shrink it so that every \(U_i\in
\mathcal O_i\), and the traceless part of \(\phi(U_i)\), is
hyperbolic.  By \eqref{stat:local-block},
\(\phi(U_i)\in\mathcal E_{\pi(i)}\).  These open sets have
complex span equal to the entire local traceless matrix space,
even when \(D_i\) has rank one.

Fix choices \(U_i\in\mathcal O_i\).  For \(S\subseteq A\) set
\[
 F_S(t_S)=\phi\left(\prod_{i\in S}(I+t_iU_i)\right).
\]
We prove by induction that, for all sufficiently small real parameters,
\begin{equation}
 F_S(t_S)=\prod_{k\in\pi(S)} F_{k,S}(t_{S\cap G_k}),
 \qquad F_{k,S}\in\mathcal E_k,\qquad F_{k,S}(0)=I,
 \label{stat:group-induction}
\end{equation}
where each factor is multiaffine in its own parameters.  Distinct
factors act on distinct sites.  They and their inverses are stability
units near zero.  The empty and singleton cases follow from unitality
and Lemma~\ref{stat:routing}.

Suppose \eqref{stat:group-induction} holds, and add \(i\notin S\),
with route \(k\).  Put
\[
 \Delta(t_S)=\phi\left(U_i\prod_{j\in S}(I+t_jU_j)\right),
 \qquad B(t_S)=F_S(t_S)^{-1}\Delta(t_S).
\]
The new input product is tensor-LY.  Its image is
\(F_S+t_i\Delta\), which is nonzero near the identity.  Left
multiplication by the stable inverse gives the two-sided stable
pencil \(I+t_iB(t_S)\).  Lemma~\ref{stat:edge} applies to this
pencil.  At \(t_S=0\) its coefficient is \(\phi(U_i)\), whose
component on site \(k\) is hyperbolic.  Continuity and
Lemma~\ref{stat:hyperbolic} therefore imply
\begin{equation}
 B(t_S)\in\mathcal E_k
 \label{stat:localized-increment}
\end{equation}
on a sufficiently small parameter neighborhood.

Fix a parameter \(t_j\) belonging to a different group \(G_\ell\),
and fix all the other parameters.  Write
\(F_{\ell,S}=C_0+t_jC_1\).  The matrices \(C_0,C_1\) are
linearly independent for sufficiently small fixed values of the
other parameters: at zero they are \(I,\phi(U_j)\), with a
nonzero hyperbolic traceless part in the second matrix.
Remove the fixed other-site factors from \(\Delta=F_SB\).
An affine function of \(t_j\) is thereby written as
\[
 (C_0+t_jC_1)\otimes R(t_j),\qquad R(t_j)=F_{k,S}B(t_S).
\]
Choose linear functionals \(\alpha,\beta\) with
\[
 \alpha(C_0)=1,\quad\alpha(C_1)=0,\qquad
 \beta(C_0)=0,\quad\beta(C_1)=1.
\]
Applying them shows that both \(R(t_j)\) and \(t_jR(t_j)\) are
affine.  Hence \(R\), and therefore \(B\), is constant in
\(t_j\).  Repeating this for every parameter outside \(G_k\)
gives the update
\[
 F_{S\cup\{i\}}=
 \prod_{\ell\ne k}F_{\ell,S}
 \left(F_{k,S}+t_iF_{k,S}B(t_{S\cap G_k})\right).
\]
The last factor lies in the real algebra \(\mathcal E_k\), which
is closed under multiplication.  It is multiaffine, as is seen by
setting every other group's parameters to zero in the original
definition of \(F_{S\cup\{i\}}\).  This completes the induction;
only finitely many neighborhood shrinkings are needed.

Setting the other groups' parameters to zero identifies each
\(F_{k,S}\) as the image of its own group product.  Extracting
coefficients in \eqref{stat:group-induction} consequently gives
\[
 \phi\left(\prod_{i\in S}U_i\right)=
 \prod_{k\in\pi(S)}
 \phi\left(\prod_{i\in S\cap G_k}U_i\right).
\]
This is a complex multilinear identity.  Its validity on a product
of real open sets extends first to their real spans and then to their
complex spans.  Adjoining identity factors spans the entire matrix
algebra on \(A\).  Thus the restriction of \(\phi\) to that
algebra is the tensor product of maps on the groups, each with image
on its routed input site.

For an observable \(O\) on \(G_k\), extract the one-site image of
\(\phi(O\otimes I)\) by normalized partial trace over the other
input sites.  This gives a unital HP map \(\psi_k\), CP if
\(\phi\) is CP.  Its adjoint \(\Xi_k\) is TP and has the same
positivity properties.  Equivalently, \(\Xi_k\) is obtained by
preparing every other input maximally mixed, applying \(\Phi\),
and tracing the other outputs.  These operations have tensor-LY
Choi arrays.  Their contraction is nonzero because \(\Xi_k\) is
TP, so \(\Xi_k\) is full Choi--LY.  Taking adjoints proves
\eqref{stat:visible-form}.  The adjoint of the partial trace in
that formula tensors with identity, which also verifies its
normalization.
\end{proof}

\begin{theorem}[Invertible static classification]
\label{stat:invertible}
An invertible complex-linear HP, TP map
\(\Phi:M_{2^n}(\C)\to M_{2^n}(\C)\) is full Choi--LY if and
only if
\begin{equation}
 \Phi=(\Phi_1\otimes\cdots\otimes\Phi_n)\circ\mathcal P,
 \label{stat:invertible-form}
\end{equation}
where \(\mathcal P\) permutes the physical input wires and the
\(\Phi_i\) are invertible one-qubit HP, TP, full Choi--LY maps.
The map \(\Phi\) is CP if and only if all its local factors are CP.
\end{theorem}

\begin{proof}
Apply the preceding results to \(\phi=\Phi^\dagger\).  The induced
map on the real space of transverse fields modulo scalars is injective.
Indeed, if \(\phi(U)=bI\), unitality gives
\(\phi(U-bI)=0\); injectivity forces \(U=bI\), and a traceless
transverse field is then zero.  In particular every output is visible,
each of its two-dimensional transverse planes maps injectively to its
routed two-dimensional input plane, and two such planes cannot route
to the same input plane.  With \(n\) planes on each side, the routes
form a permutation and all groups have one site.

Theorem~\ref{stat:visible} now gives a tensor product of local factors.
Each factor is invertible: a nonzero kernel vector of one factor,
tensored with suitable nonzero vectors for the other factors, would
give a kernel vector of the full map.  The same theorem proves
inheritance of CP.  Conversely, tensor products and physical wire
permutations preserve full Choi--LY and all the stated properties.
\end{proof}

\begin{remark}
The proof of Theorem~\ref{stat:invertible} does not require a
semigroup representation, complete positivity, or injectivity of an
unproved local restriction.  For completeness, its adjoint local
longitudinal coefficient in \eqref{stat:local-block} is positive:
stability of \(I+s\lambda_iZ\), obtained from \(I+sZ\) for small
\(s>0\), gives \(|1+s\lambda_i|\ge|1-s\lambda_i|\), hence
\(\lambda_i\ge0\); injectivity excludes zero.
\end{remark}

\begin{example}[Ordinary static preservation is weaker]
\label{stat:ordinary-counterexample}
Let \(n\ge2\), \(d=2^n\),
\(\tau=\diag(3/4,1/4)\), and \(\sigma=\tau^{\otimes n}\).
For
\[
 0<a<\frac1{4^n+1},\qquad
 \Phi_a(M)=aM+(1-a)\Tr(M)\sigma,
\]
the map \(\Phi_a\) is invertible and CPTP, and sends every density
matrix to a matrix with a polynomial nonzero on the closed unit
polydisk.  Nevertheless it is neither a product of local channels
up to a physical permutation nor full Choi--LY.

Indeed, it fixes \(\sigma\) and multiplies the trace-zero subspace
by \(a\), proving invertibility.  On the closed polydisk,
\[
 |G_\sigma|\ge2^{-n}=d^{-1},\qquad |G_\rho|\le d
\]
for every density matrix \(\rho\).  The second inequality follows
by bounding the two monomial-vector norms by \(\sqrt d\) and
\(\|\rho\|_{\mathrm{op}}\) by one.  Hence
\[
 |G_{\Phi_a(\rho)}|\ge(1-a)d^{-1}-ad>0.
\]
On input \(|0^n\rangle\langle0^n|\), the occupation indicators at
any two distinct output sites have covariance \(a(1-a)/16>0\).
The output is not product, whereas a product of local channels,
with or without a wire permutation, sends product inputs to products.

One may also exhibit the Choi zero directly.  Its polynomial is
\[
 a\prod_i(1+z_ix_i)(1+w_iy_i)
 +(1-a)G_\sigma(z,w)\prod_i(1+x_iy_i).
\]
Set all output variables except \(z_1,w_1\) to zero, and set
\(z_1=w_1=t\), \(x_1=y_1=v\).  Define
\[
 p(t)=(3/4)^{n-1}(3/4+t^2/4),\quad
 C(t)=\frac{1-a}{a}p(t),\quad C_0=C(1).
\]
Choose \(v_0=e^{i\theta}\) with
\(\cos\theta=-1/(1+C_0)\), \(\pi/2<\theta<\pi\).  Then
\((1+v_0)^2/(1+v_0^2)=-C_0\).
For small \(\varepsilon>0\), put
\[
 t=1-\varepsilon,\qquad v=(1-\varepsilon)v_0,
 \qquad q=-\frac{(1+tv)^2}{C(t)(1+v^2)}-1.
\]
As \(\varepsilon\downarrow0\), \(q\to0\), so all of
\(|t|,|v|,|q|\) are smaller than one for sufficiently small
positive \(\varepsilon\).  Set \(x_2=y_2=\sqrt q\) and all
remaining reference variables to zero.  The Choi polynomial becomes
\(a(1+tv)^2+(1-a)p(t)(1+v^2)(1+q)=0\), at an interior point.
Lemma~\ref{stat:bridge} turns this into an ancillary physical witness.
The same example is the late-time channel
\(e^{t(\mathcal R_\sigma-\Id)}=e^{-t}\Id+(1-e^{-t})\mathcal R_\sigma\),
where \(\mathcal R_\sigma(M)=\Tr(M)\sigma\).  Its ordinary
strict-preservation property for \(t>\log(4^n+1)\) therefore does
not imply preservation at times approaching zero.
\end{example}

\subsection{Saturated columns and reflected polynomial factors}

For a polynomial \(p\), let \(\|p\|_2\) be the Euclidean norm of
its coefficient vector.  If \(p\) is multiaffine in a specified
finite set \(S\), define its full reflection by
\begin{equation}
 p^*(z)=z^S\overline{p(1/\bar z)}.
 \label{stat:reflection}
\end{equation}
The set \(S\) is part of the definition, and may include variables
absent from \(p\).  In contrast, \(\overline p(w)=
\overline{p(\bar w)}\) below denotes coefficientwise conjugation
without reflection.

\begin{lemma}[Saturation and disjoint factors]
\label{stat:saturation}
Let \(F(z,x)=\sum_{b\in\{0,1\}^m}A_b(z)x^b\) be a stable
multiaffine polynomial.  Suppose
\(\|A_{e_i}\|_2=\|A_0\|_2\) for every input site \(i\).
Then, over the rational function field in \(z\),
\begin{equation}
 F(z,x)=A_0(z)\prod_{i=1}^m
       \left(1+x_i\frac{A_{e_i}(z)}{A_0(z)}\right).
 \label{stat:saturation-product}
\end{equation}
Writing the quotients in coprime form \(q_i/p_i\), there is a
polynomial factorization
\begin{equation}
 F(z,x)=C(z)\prod_{i=1}^m(p_i(z)+x_iq_i(z))
 \label{stat:primitive-product}
\end{equation}
in which the output-variable supports of the displayed factors are
pairwise disjoint.  All these factors are stable, and
\(\|p_i\|_2=\|q_i\|_2\).
\end{lemma}

\begin{proof}
Setting all input variables except \(x_i\) to zero gives a stable
affine polynomial \(A_0+x_iA_{e_i}\), whence
\(|A_{e_i}(z)|\le|A_0(z)|\) in the output polydisk.
The inequality extends continuously to the output torus.  Parseval
and norm equality imply equality of these moduli everywhere on that
torus.  At a torus point where \(A_0\ne0\), boundary Hurwitz makes
the input polynomial stable; its nonzero constant coefficient
excludes the zero alternative.  Write this input polynomial as
\(P_i(x_{\ne i})+x_iQ_i(x_{\ne i})\).  The holomorphic ratio
\(Q_i/P_i\) is contractive on the remaining input polydisk, and
has modulus one at its origin.  It is therefore constant by the
maximum modulus principle.  Applying this to every input gives
\eqref{stat:saturation-product} at these torus points.
The zero set of the nonzero polynomial \(A_0\) has torus measure
zero.  Clearing denominators and using uniqueness of Fourier
coefficients proves the rational identity everywhere.

Each \(p_i+x_iq_i\) is primitive over \(\C[z]\) and irreducible
over the rational function field, where it is linear in its own
input variable.  Gauss's lemma applied to the primitive part of
\(F\), considered as a polynomial in \(x\), gives
\eqref{stat:primitive-product} with \(C\) polynomial.  Degrees in
any one output variable add under multiplication in this integral
domain.  Since every such degree of \(F\) is at most one, each
output variable belongs to at most one displayed factor.  Stability
of \(F\) implies stability of every factor and, by setting
\(x_i=0\), of \(p_i\).  Finally coefficient norms multiply on
disjoint supports.  Comparing \(A_0\) with \(A_{e_i}\) gives the
claimed local norm equality.
\end{proof}

\begin{lemma}[Reduced reflected pairs]
\label{stat:reflected-pair}
Let \(p+xq\) be stable and multiaffine, where \(p,q\) are
coprime, have equal coefficient norms, and their union of output
supports is \(S\).  Then
\begin{equation}
 q=e^{i\theta}p^*
 \label{stat:pair-reflection}
\end{equation}
for some real \(\theta\).  Conversely, for any stable multiaffine
\(p\) on a specified set \(S\), the polynomial
\(p+e^{i\theta}xp^*\) is stable.
\end{lemma}

\begin{proof}
Stable affine slices and equality of norms imply \(|p|=|q|\) on
the output torus.  Multiplication by \(z^S\) gives
\(pp^*=qq^*\).  Coprimality implies \(p\mid q^*\); write
\(q^*=ph\), so \(p^*=qh\).
The constant coefficient of the stable polynomial \(p\) is nonzero.
Thus \(p^*\) has degree one in every variable of \(S\).
If \(h\) depends on \(z_j\), degree additivity makes \(q\)
independent of \(z_j\), and the support-union property makes
\(p\) depend on \(z_j\).  Independence of \(q\) implies
\(z_j\mid q^*=ph\).  Since \(p\) has nonzero constant term,
\(z_j\mid h\), and then \(z_j\mid p^*=qh\).  The latter
says \(p\) is independent of \(z_j\), a contradiction.  Hence
\(h\) is constant.  Reflecting twice shows \(|h|=1\), proving
\eqref{stat:pair-reflection}.

For the converse, \(p^*/p\) is holomorphic on the polydisk and
has modulus at most one.  To justify this assertion even in the
presence of boundary zeros, use \(p_r(z)=p(rz)\), \(r<1\).
The polynomial \(p_r\) has no closed-polydisk zero, and
\(p_r^*/p_r\) has unit modulus on the torus.  Successive maximum
modulus principles make it contractive inside.  Letting \(r\uparrow1\)
gives the asserted bound.  Hence
\(1+e^{i\theta}x p^*/p\ne0\) for \(x\in\D\).
\end{proof}

Reflection descriptions of rational inner functions are classical;
see \cite{Knese24}.  The elementary proof above specifies
the full support set and the coprime multiaffine reduction used here.

\subsection{All stable isometries and their coding obstruction}

\begin{theorem}[Stable isometries]
\label{stat:isometry}
Let \(V:(\C^2)^{\otimes m}\to(\C^2)^{\otimes n}\) be an
isometry.  Its coefficient polynomial is stable if and only if,
after a physical output permutation,
\begin{equation}
 V=|\psi\rangle\otimes V_1\otimes\cdots\otimes V_m,
 \label{stat:isometry-form}
\end{equation}
where \(\psi\) is a normalized stable pure state on an optional
prepared block \(S_0\).  The remaining nonempty, disjoint blocks
\(S_i\) carry one-input isometries whose polynomials are
\begin{equation}
 G_{V_i}(z,x_i)=
 \frac{p_i(z)+e^{i\theta_i}x_i p_i^*(z)}{\|p_i\|_2},
 \qquad \gcd(p_i,p_i^*)=1,
 \qquad [z^{S_i}]p_i(z)^2=0.
 \label{stat:isometry-block}
\end{equation}
Each \(p_i\) is stable and multiaffine, and reflection uses the
whole block \(S_i\).  The sets \(S_0,S_1,\ldots,S_m\) partition
the physical outputs.  In particular, an isometric channel
\(M\mapsto VMV^*\) is full Choi--LY exactly for these isometries.
\end{theorem}

\begin{proof}
Write \(G_V=\sum_bA_b(z)x^b\).  Isometry gives
\(\langle A_b,A_c\rangle=\delta_{bc}\), so
Lemma~\ref{stat:saturation} applies.  In its factorization,
disjointness makes both norms and inner products multiplicative.
Comparing columns \(0\) and \(e_i\) yields
\(\|p_i\|_2=\|q_i\|_2>0\) and
\(\langle p_i,q_i\rangle=0\).  A factor with no output variable
could not satisfy the latter equality, so every logical block is
nonempty.  Normalize the common prepared factor and each pair.
There is no residual normalization, because the all-zero column
has norm one.  Variables absent from the polynomial are fixed-zero
prepared outputs.

Lemma~\ref{stat:reflected-pair} gives the reflected form.
Furthermore,
\[
 \langle p_i,p_i^*\rangle
 =\overline{\sum_{T\subseteq S_i}(p_i)_T(p_i)_{S_i\setminus T}}
 =\overline{[z^{S_i}]p_i(z)^2},
\]
so orthogonality is exactly the final condition in
\eqref{stat:isometry-block}.  Conversely that condition gives
orthogonal columns, reflection preserves their norms, and
Lemma~\ref{stat:reflected-pair} gives stability.  Tensor products
prove the converse for \(V\).  The isometric channel's Choi
polynomial is \(G_V(z,x)\overline{G_V}(w,y)\) in independent
variables, proving the final equivalence.
\end{proof}

\begin{lemma}[Strict longitudinal bias]
\label{stat:strict-bias}
If \(p\) is stable and multiaffine on a nonempty specified set
\(S\), and \(\gcd(p,p^*)=1\), then for every \(j\in S\),
\begin{equation}
 \|p|_{z_j=0}\|_2^2>\|\partial_jp\|_2^2.
 \label{stat:strict-slice}
\end{equation}
\end{lemma}

\begin{proof}
Write \(p=a+z_jb\).  Stable affine slices imply \(|b|\le|a|\)
on the remaining polydisk, hence \(\|b\|_2\le\|a\|_2\).
If \(b=0\), strictness follows from \(a\ne0\).
Otherwise suppose equality holds.  Torus saturation gives
\(aa^*=bb^*\), reflecting both polynomials over
\(R=S\setminus\{j\}\), while the full reflection is
\(p^*=b^*+z_ja^*\).
Let \(g=\gcd(a,b)\) and
\(h=a/g+z_jb/g\).  This polynomial is primitive over \(\C[z_R]\)
and has positive degree in \(z_j\).  Over \(\C(z_R)[z_j]\),
\[
 p^*=\frac{a^*}{b}(a+z_jb).
\]
Thus \(h\) divides both \(p\) and \(p^*\) over the fraction
field.  Gauss's lemma gives the same divisibility in the polynomial
ring, contradicting coprimality.  This proof also covers empty
\(R\) and variables absent from either slice.
\end{proof}

\begin{corollary}[Single-output erasure]
\label{stat:no-code}
For a reduced logical block in Theorem~\ref{stat:isometry}, every
physical output \(j\in S_i\) satisfies
\begin{equation}
 V_i^*Z_jV_i=\delta_j Z_i,
 \qquad
 \delta_j=
 \frac{\|p_i|_{z_j=0}\|_2^2-\|\partial_jp_i\|_2^2}
      {\|p_i\|_2^2}>0.
 \label{stat:bias-compression}
\end{equation}
Erasure of any such output is not exactly correctable on the
logical space.  Every nontrivial stable isometric quantum code
therefore has distance one.
\end{corollary}

\begin{proof}
Reflection complements occupied subsets, so the two diagonal
expectations of \(Z_j\) have opposite signs.  The cross term is
zero because complementary subsets cancel:
\[
 \langle p_i,Z_jp_i^*\rangle
 =\overline{\sum_{T\subseteq S_i}
       (-1)^{\mathbf{1}_{j\in T}}(p_i)_T(p_i)_{S_i\setminus T}}=0.
\]
Lemma~\ref{stat:strict-bias} now proves
\eqref{stat:bias-compression}.  For erasure of output \(j\),
the two Kraus operators may be taken as
\(E_a=\langle a|_j\otimes I\), \(a=0,1\).
The Knill--Laflamme criterion \cite{KnillLaflamme97} requires
\(V^*E_a^*E_bV\) to be scalar for all \(a,b\).
Subtracting the diagonal cases requires \(V^*Z_jV\) to be scalar,
contrary to \eqref{stat:bias-compression}.  Equivalently, this
weight-one phase observable is not detectable by the code.
\end{proof}

\begin{remark}
Outputs in the prepared block are exempt from the erasure conclusion.
Nor is there a uniform positive lower bound on the individual bias:
for \(p(z_1,z_2)=1+r z_1\), \(0\le r<1\), reflected over
\(\{1,2\}\), one has \(p^*=z_2(r+z_1)\), coprime orthogonal
columns, and \(\delta_1=(1-r^2)/(1+r^2)\to0\).  The conclusion
concerns exact correction only.
\end{remark}

\begin{example}[Coherent information invisible in one-site marginals]
\label{stat:repetition}
The repetition isometry \(V|0\rangle=|00\rangle\),
\(V|1\rangle=|11\rangle\) has polynomial \(1+xz_1z_2\).
The square channel \(\Phi(M)=V\Tr_2(M)V^*\) is full Choi--LY,
with polynomial
\[
 (1+x_1z_1z_2)(1+y_1w_1w_2)(1+x_2y_2).
\]
All its single-output transverse adjoints vanish, although
\(\Phi^\dagger(X_1X_2)=X_1\otimes I_2\).
Thus Theorem~\ref{stat:visible} cannot identify its omitted
outputs as a classical or input-independent sector.
\end{example}

\begin{proposition}[No stable Stinespring representation in general]
\label{stat:no-dilation}
There is a two-qubit CPTP full Choi--LY map with no Stinespring
isometry whose coefficient tensor is LY, for any finite qubit
environment.
\end{proposition}

\begin{proof}
Let \(\Delta\) be complete computational-basis dephasing on two
qubits, let \(P\) swap the qubits, and set
\(\Phi=\tfrac12\Delta+\tfrac12\operatorname{Ad}_P\circ\Delta\).
This is CPTP.  With \(r_j=z_jw_j\) and \(s_j=x_jy_j\), its
Choi polynomial is
\begin{equation}
 \begin{split}
 K(r,s)&=\tfrac12(1+r_1s_1)(1+r_2s_2)
       +\tfrac12(1+r_1s_2)(1+r_2s_1)\\
 &=1+\tfrac12(r_1+r_2)(s_1+s_2)+r_1r_2s_1s_2.
 \end{split}
 \label{stat:symmetrized-kernel}
\end{equation}
For fixed \(s_1,s_2\in\D\), this is the symmetric multiaffine
polarization in \(r_1,r_2\) of
\((1+rs_1)(1+rs_2)\).  The Grace--Walsh--Szeg\H{o} coincidence
theorem on the convex circular domain \(\D\) implies its
nonvanishing for \(r_1,r_2\in\D\); see \cite{BB10II}.
Thus the full Choi tensor is stable.

However,
\(\Phi^\dagger(Z_1)=\Phi^\dagger(Z_2)=(Z_1+Z_2)/2\).
If a stable Stinespring isometry existed, apply
Theorem~\ref{stat:isometry} to its physical and environmental
output qubits together.  Every one-site observable on a physical
output would pull back either to a scalar or to an operator on
one logical input, because its output belongs to one disjoint
encoding block or to the prepared block.  Tracing the environment
merely inserts identities at its sites.  The displayed two-input
longitudinal observable is incompatible with this conclusion.
\end{proof}

\subsection{All Schur channels}

\begin{lemma}[A scalar equal-diagonal criterion]
\label{stat:two-by-two}
For a nonzero PSD matrix
\(H=\begin{psmallmatrix}a&b\\\bar b&d\end{psmallmatrix}\),
its tensor polynomial is stable if and only if \(a>0\) and
\(a\ge d\).  For arbitrary complex \(\beta,\gamma\), the
polynomial \(1+\beta x+\gamma y+xy\) is stable if and only if
\(\gamma=\bar\beta\) and \(|\beta|\le1\).
\end{lemma}

\begin{proof}
A stable polynomial must have nonzero constant coefficient, so
\(a>0\).  The necessary boundary root-modulus inequality is
\[
 |a+\bar b\zeta|^2-|b+d\zeta|^2
 =(a-d)(a+d+2\Re(\bar b\zeta))\ge0,
 \qquad |\zeta|=1.
\]
Its average gives \(a\ge d\).  Conversely, PSD gives
\(|b|^2\le ad\), and \(a\ge d\) gives \(|b|\le a\).
If \(|b|<a\), the affine constant factor \(a+\bar b z\)
is nonzero on the closed disk, and the displayed inequality
and the maximum modulus principle bound the other coefficient
divided by this one by one.  This proves stability.  If
\(|b|=a\), PSD and \(a\ge d\) force \(d=a\), and the
polynomial factors as
\(a(1+(\bar b/a)z)(1+(b/a)w)\), also stable.

For the second assertion, boundary root-modulus inequalities give
\(|1+\beta\zeta|\ge|\gamma+\zeta|\), and the analogous
inequality with \(\beta,\gamma\) exchanged.  Squaring and
averaging gives \(|\beta|=|\gamma|\).  The remaining nonnegative
linear trigonometric polynomial has mean zero, so vanishes
identically, yielding \(\beta=\bar\gamma\).  Setting one
variable to zero gives \(|\beta|\le1\).  The first assertion
applied to the resulting equal-diagonal PSD matrix proves the
converse.
\end{proof}

\begin{theorem}[Stable correlation matrices and Schur channels]
\label{stat:schur}
Let \(C\succeq0\) be a \(2^m\)-by-\(2^m\) matrix with
\(C_{aa}=1\).  Its tensor polynomial is stable if and only if
\begin{equation}
 C=\bigotimes_{i=1}^m
 \begin{pmatrix}1&\eta_i\\\bar\eta_i&1\end{pmatrix},
 \qquad |\eta_i|\le1.
 \label{stat:correlation-product}
\end{equation}
Equivalently, a CPTP map fixing every computational-basis state is
full Choi--LY if and only if it is a tensor product of one-qubit
phase/dephasing channels.  Parameters \(\eta_i=0\) are allowed.
\end{theorem}

\begin{proof}
Separate one physical site and write
\(C=\begin{psmallmatrix}A&B\\B^*&D\end{psmallmatrix}\).
For arbitrary \(z\in\D^{m-1}\), compress the other sites against
the product vector \(v(z)=\sum_a z^a|a\rangle\).
The two-by-two compression is PSD.  Its tensor polynomial is an
interior specialization of \(G_C\), hence stable, and its
constant coefficient is strictly positive.  By
Lemma~\ref{stat:two-by-two},
\[
 \langle v(z),(A-D)v(z)\rangle\ge0.
\]
At any fixed positive interior coordinate moduli, the angular
average of this expression is
\(\sum_a(A_{aa}-D_{aa})|z|^{2a}=0\).
Continuity and nonnegativity force pointwise equality on every
such torus.  Varying the moduli proves the polynomial identity
\(\langle v(z),(A-D)v(z)\rangle=0\).  Independence of the
monomials \(\bar z^a z^b\) gives \(A=D\) entrywise.

Writing \(u\) for the remaining independent row and column
variables, the full polynomial has the form
\[
 a(u)(1+xy)+b(u)x+c(u)y.
\]
The specialization \(a(u)\) is nonzero throughout the polydisk.
Lemma~\ref{stat:two-by-two} gives
\(c(u)/a(u)=\overline{b(u)/a(u)}\) and
\(|b(u)/a(u)|\le1\) for every complex \(u\) there.
Both ratios are holomorphic on this connected domain, so they
are constant.  This factors off one matrix in
\eqref{stat:correlation-product}.  The remaining factor \(A\)
is PSD, has diagonal one, and is tensor-LY.  Induction proves
the assertion.  The converse follows from the one-site criterion
and multiplication on disjoint variables.

For a Schur multiplier \(\mathcal S_C(M)=C\circ M\), CP and
TP are equivalent to \(C\succeq0\) and diagonal one.  Its full
Choi polynomial is \(G_C(r,s)\), where \(r_i=z_iu_i\) and
\(s_i=w_iv_i\).  Each disk value is attainable by these
products, so full Choi stability is exactly tensor stability
of \(C\).  Finally, any CPTP map fixing all computational-basis
states has diagonal Kraus operators: each Kraus image of a basis
vector must lie in the one-dimensional range of the corresponding
pure output.  It is therefore a Schur map, proving the equivalent
formulation.
\end{proof}

\subsection{Channels with pure computational-basis outputs}

\begin{theorem}[Pure-basis-output classification]
\label{stat:pure-basis}
Let \(\Phi:M_{2^m}(\C)\to M_{2^n}(\C)\) be CPTP, and suppose
every computational-basis output \(\Phi(|a\rangle\langle a|)\)
is pure.  The outputs need not be orthogonal or distinct.  Then
\(\Phi\) is full Choi--LY if and only if it has the following
form.

Partition the inputs as \(E\sqcup K=[m]\), and the outputs into
an optional prepared block \(S_0\) and nonempty disjoint blocks
\(S_i\), \(i\in K\).  Prepare a normalized stable pure state
\(\psi\) on \(S_0\).  On each \(S_i\) take a reduced stable
two-column matrix \(V_i\) with normalized, linearly independent
columns \(v_{i0},v_{i1}\).  Reduced means that their amplitude
polynomials are coprime.  Equivalently,
\begin{equation}
 G_{V_i}(z,x_i)=
 \frac{p_i(z)+e^{i\theta_i}x_i p_i^*(z)}{\|p_i\|_2},
 \qquad \gcd(p_i,p_i^*)=1,
 \label{stat:pure-pair}
\end{equation}
where \(p_i\) is stable and multiaffine and reflection uses the
full block.  Put \(\mu_i=\langle v_{i0},v_{i1}\rangle\), and
choose \(|\eta_i|\le1\) subject to
\begin{equation}
 \mu_i\eta_i=0.
 \label{stat:overlap-noise}
\end{equation}
With
\[
 C_i=\begin{pmatrix}1&\eta_i\\\bar\eta_i&1\end{pmatrix},
 \qquad
 \mathcal B_i(M)=V_i(C_i\circ M)V_i^*,
\]
the channel, after output permutation, is exactly
\begin{equation}
 \Phi(\rho)=|\psi\rangle\langle\psi|\otimes
       \left(\bigotimes_{i\in K}\mathcal B_i\right)
                         (\Tr_E\rho).
 \label{stat:pure-basis-form}
\end{equation}
Empty \(K\) is allowed.  If \(\mu_i=0\), the block is a stable
isometric encoding preceded by arbitrary local phase/dephasing.
If \(\mu_i\ne0\), it measures the input in the computational
basis and prepares the corresponding nonorthogonal reflected
pure state.
\end{theorem}

\begin{proof}
Assume first that \(\Phi\) is full Choi--LY.  Choose normalized
output vectors \(\psi_a\), with amplitude polynomials \(p_a\).
Complete input dephasing is full Choi--LY; its Choi polynomial is
a product of factors \(1+x_iy_iu_iv_i\).  After composing with
this map and grouping each input row/column pair into its product,
we obtain the stable polynomial
\begin{equation}
 F(z,w,s)=\sum_a A_a(z,w)s^a,
 \qquad A_a(z,w)=p_a(z)\overline p_a(w).
 \label{stat:pure-density-kernel}
\end{equation}
Grouping is an equivalence of disk stability because disk-variable
products cover the disk.  Each \(A_a\) has coefficient norm one,
being the polynomial of a normalized rank-one projection.
Lemma~\ref{stat:saturation} gives, for every nonempty subset
\(T\subseteq[m]\),
\begin{equation}
 A_T A_0^{|T|-1}=\prod_{i\in T}A_{e_i}.
 \label{stat:pure-column-identity}
\end{equation}

Substitute \eqref{stat:pure-density-kernel} into this identity.
The rational function
\[
 R_T(z)=\frac{p_T(z)p_0(z)^{|T|-1}}
                   {\prod_{i\in T}p_{e_i}(z)}
\]
satisfies \(R_T(z)\overline{R_T(\bar w)}=1\) in independent
variables.  Fixing a generic \(w\) shows that it is a constant
\(\xi_T\) of modulus one.  We may therefore align the output
column phases so that the polynomial
\begin{equation}
 G(z,x)=p_0(z)\prod_i
              \left(1+x_i\frac{p_{e_i}(z)}{p_0(z)}\right)
 \label{stat:phase-aligned-amplitude}
\end{equation}
has exactly those phase-aligned columns.  More explicitly, its
coefficient at \(x^T\) is \(\xi_T^{-1}p_T\); hence it is a
multiaffine polynomial, not merely a rational expression.
The stable constant slice \(A_0=p_0\overline p_0\) makes
\(p_0\) zero-free in the output polydisk.  The singleton modulus
domination for \eqref{stat:pure-density-kernel}, specialized to
\(w=\bar z\), gives \(|p_{e_i}/p_0|\le1\).
Thus \eqref{stat:phase-aligned-amplitude} is stable.  It defines
a matrix \(V\) with normalized columns; no orthogonality is
being asserted.

Apply Lemma~\ref{stat:saturation} to \(G\), whose column norms
are all one, and then Lemma~\ref{stat:reflected-pair} to its
primitive factors.  We obtain a stable prepared factor and
disjoint reduced reflected pairs with equal local norms.  A
dependent primitive pair \(q_i=c_ip_i\) must have both
polynomials constant, so has empty output support; conversely
every nonempty primitive pair has independent columns.  Assign
the empty factors to \(E\) and the others to \(K\).  Their
equal norms imply \(|c_i|=1\).  Replacing each erased input
variable by \(c_i^{-1}x_i\) is a disk-preserving phase change
of columns, and makes the two erased-input columns identical.
Normalize the prepared factor and all retained pairs.  The total
normalization is one because the all-zero full column has norm
one.  All aligned output vectors now have the form
\begin{equation}
 v_a=\psi\otimes\bigotimes_{i\in K}v_{i,a_i},
 \label{stat:aligned-vectors}
\end{equation}
independent of \(a_E\).  Common prepared factors and variables
absent from \(G\) belong to \(S_0\).  This proves all the
polynomial and support assertions about the pairs in the theorem.

Choose Kraus operators \(K_\ell\) of \(\Phi\).  Since the
positive sum of their output projections on input \(a\) has
rank one, every \(K_\ell|a\rangle\) is proportional to \(v_a\).
Writing \(K_\ell|a\rangle=c_{\ell a}v_a\) gives
\begin{equation}
 \Phi(|a\rangle\langle b|)=C_{ab}|v_a\rangle\langle v_b|,
 \qquad C_{ab}=\sum_\ell c_{\ell a}\bar c_{\ell b},
 \label{stat:kraus-correlation}
\end{equation}
where \(C\succeq0\) and \(C_{aa}=1\).  Original trace
preservation is exactly
\begin{equation}
 C_{ab}\langle v_b,v_a\rangle=\delta_{ab}.
 \label{stat:original-trace}
\end{equation}
In particular, if \(a_K=b_K\), the aligned vectors in
\eqref{stat:aligned-vectors} are identical, so
\begin{equation}
 C_{ab}=\delta_{ab}\quad\text{whenever }a_K=b_K.
 \label{stat:erased-trace}
\end{equation}

For each retained block, its independent normalized columns have
positive definite Gram matrix
\[
 H_i=V_i^*V_i=\begin{pmatrix}1&\mu_i\\\bar\mu_i&1\end{pmatrix},
 \qquad |\mu_i|<1.
\]
Both \(V_i^*\) and
\[
 H_i^{-1}=\frac1{1-|\mu_i|^2}
       \begin{pmatrix}1&-\mu_i\\-\bar\mu_i&1\end{pmatrix}
\]
are tensor-LY, the latter by Lemma~\ref{stat:two-by-two}.
Consequently \(L_i=H_i^{-1}V_i^*\) is tensor-LY and satisfies
\(L_iV_i=I\).  Contraction cannot give zero because of this
identity.  The rectangular matrix
\(L=\langle\psi|\otimes\bigotimes_{i\in K}L_i\) is stable
and obeys \(Lv_a=|a_K\rangle\).

The CP map \(\Psi(M)=L\Phi(M)L^*\) is full Choi--LY or zero,
and satisfies
\[
 \Psi(|a\rangle\langle b|)=C_{ab}|a_K\rangle\langle b_K|.
\]
Its trace is \(C_{ab}\delta_{a_K,b_K}=\delta_{ab}\) by
\eqref{stat:erased-trace}.  Thus \(\Psi\) is TP and nonzero.
The compression itself need not be TP on arbitrary output
matrices; this argument proves TP only for its composition with
\(\Phi\), which is all that is required.

The Choi polynomial of \(\Psi\) is
\[
 \sum_{a,b}C_{ab}z^{a_K}w^{b_K}u^av^b.
\]
For \(i\in K\), replace \(z_iu_i,w_iv_i\) by \(r_i,s_i\);
for \(i\in E\), put \(r_i=u_i,s_i=v_i\).
The variables \(r_i,s_i\) range independently over the disk,
so stability of this Choi polynomial is precisely tensor
stability of the entire correlation matrix \(C\).
Theorem~\ref{stat:schur} gives
\(C=\bigotimes_i\begin{psmallmatrix}1&\eta_i\\\bar\eta_i&1\end{psmallmatrix}\),
with \(|\eta_i|\le1\).  Taking strings differing only at an
erased input in \eqref{stat:erased-trace} forces \(\eta_i=0\)
there.  Taking strings differing at a retained input in
\eqref{stat:original-trace} gives
\(\eta_i\bar\mu_i=0\), equivalently
\eqref{stat:overlap-noise}.  Combining these tensor-product
coefficients with \eqref{stat:aligned-vectors} proves
\eqref{stat:pure-basis-form}; the erased coefficient is exactly
\(\delta_{a_E,b_E}\), hence ordinary partial trace.

Conversely, start from the displayed data.  Each \(C_i\) is a
PSD correlation matrix, so its Schur map is CP.  Following it
by \(M\mapsto V_iMV_i^*\) preserves CP.  Normalization of
the columns gives the correct diagonal traces, while
\eqref{stat:overlap-noise} makes the off-diagonal traces zero.
Thus each \(\mathcal B_i\) is TP and has pure basis outputs.
Its two component Choi tensors are stable by
Lemma~\ref{stat:reflected-pair} and
Theorem~\ref{stat:schur}; their composition is nonzero because
it is TP.  Stable pure preparation, partial trace, tensor
products and physical permutations prove the converse.
The descriptions of the two cases follow directly from
\(\mu_i\eta_i=0\).
\end{proof}

Theorems~\ref{stat:visible}, \ref{stat:schur}, and
\ref{stat:pure-basis} concern different singular strata.
In particular, none asserts a classification of general channels
with mixed computational-basis outputs.  Proposition~\ref{stat:no-dilation}
prevents extending the isometry theorem to that class by an
assumption of a stable Stinespring representation.

\subsection{Pauli-diagonal maps without invertibility}

\begin{lemma}[Factorization on slit domains]
\label{stat:slit-factorization}
Let \(\Omega_1,\ldots,\Omega_n\) be nonempty connected open
subsets of \(\C\) with \(\C\setminus\Omega_j\subseteq\R\).
A multiaffine polynomial nonvanishing on \(\prod_j\Omega_j\)
is a nonzero scalar times a product of one-variable affine
polynomials, with constant factors permitted.
\end{lemma}

\begin{proof}
Write \(p(r_1,r')=a(r')+r_1b(r')\).  If \(b=0\), omit that
coordinate.  Otherwise, on a nonempty connected open neighborhood
where \(b\ne0\), the holomorphic function \(-a/b\) avoids
\(\Omega_1\) and is therefore real-valued.  It is constant,
say \(c\notin\Omega_1\).  The polynomial identity principle
gives \(a+cb=0\) everywhere, so \(p=(r_1-c)b\).
The remaining factor is nonvanishing on the remaining domain.
Induction proves the assertion.
\end{proof}

\begin{theorem}[All-rank Pauli-diagonal rigidity]
\label{stat:pauli}
Suppose a complex-linear HP, TP map on \(n\) qubits is diagonal
in the Pauli-string basis:
\(\Phi(\sigma_\alpha)=\lambda_\alpha\sigma_\alpha\), where
\(\sigma_\alpha=\bigotimes_j\sigma_{\alpha_j}\) and
\(\sigma_0=I\).  If \(\Phi\) is full Choi--LY, then
\begin{equation}
 \lambda_\alpha=\prod_j\lambda_{\alpha_j^{(j)}},
 \qquad \Phi=\bigotimes_j\Phi_j,
 \label{stat:pauli-product}
\end{equation}
where \(\alpha_j^{(j)}\) denotes the string with the stated
Pauli at site \(j\) and identities elsewhere.  Each local
factor is HP, TP and full Choi--LY.  If \(\Phi\) is CP, all
factors are CP, and the converse holds.  Zero eigenvalues are
allowed.
\end{theorem}

\begin{proof}
HP makes the \(\lambda_\alpha\) real, and TP gives
\(\lambda_0=1\).  Pauli orthogonality yields
\[
 J(\Phi)=2^{-n}\sum_\alpha\lambda_\alpha
                 \sigma_\alpha\otimes\sigma_\alpha^{\mathrm T}.
\]
The four one-site polynomials are
\(f_I=1+zw\), \(f_X=z+w\), \(f_Y=i(z-w)\),
and \(f_Z=1-zw\); the transpose contributes a minus sign
only for \(Y\).

At every site select a desired axis, independently of the others.
For axis \(X\), set the output variables to \(z=w=1\).
Only the identity and \(X\) terms remain, and their normalized
ratio is
\[
 r_X=\frac{x+y}{1+xy}.
\]
Its image on the bidisk is
\(\Omega_\perp=\C\setminus(({-\infty},-1]\cup[1,\infty))\).
For a real \(|r|\ge1\), the inverse equation would require
\(y=(r-x)/(1-rx)\), whereas
\[
 |r-x|^2-|1-rx|^2=(r^2-1)(1-|x|^2)\ge0.
\]
The endpoint equations also require a boundary variable.
Conversely, setting \(x=y=u\) gives an inverse quadratic with
reciprocal roots.  A unit-circle root would place \(r\) on
the omitted rays.  For every nonzero \(r\in\Omega_\perp\),
exactly one inverse root is therefore in the disk; zero is
attained at \(u=0\).

For axis \(Y\), set \(z=i,w=-i\).  The output factor is
\(f_Y=-2\), and the input transpose changes the sign again.
The surviving ratio is
\(r_Y=i(x-y)/(1+xy)\), with the same image
\(\Omega_\perp\), by replacing \((x,y)\) with \((ix,-iy)\).
For axis \(Z\), set \(z=s,w=-s,x=t,y=t\), all interior.
The output specialization kills \(X\) and the input
specialization kills \(Y\).  The ratio is
\[
 r_Z=\frac{1+s^2}{1-s^2}\frac{1-t^2}{1+t^2}.
\]
Each factor ranges over the right half-plane, so their product
ranges exactly over \(\Omega_Z=\C\setminus(-\infty,0]\).
Indeed products have arguments strictly between \(-\pi\) and
\(\pi\), and the principal square root realizes every point
of this slit plane as such a product.

For axes \(a_1,\ldots,a_n\), divide out the nonzero identity
factors after these substitutions.  The resulting polynomial is
\begin{equation}
 p_a(r)=\sum_{S\subseteq[n]}\lambda_{a_S}\prod_{j\in S}r_j,
 \qquad p_a(0)=1,
 \label{stat:axis-polynomial}
\end{equation}
where \(a_S\) has the chosen axes on \(S\).
Boundary specialization at the \(X,Y\) sites gives stable or
zero by Hurwitz.  The zero alternative is impossible: the
ratio maps have the nonempty open product image just computed,
and the nonzero polynomial \eqref{stat:axis-polynomial} cannot
vanish identically there.  The \(Z\) substitutions are interior.
Thus \(p_a\) is nonvanishing on the corresponding product of
slit domains.

Lemma~\ref{stat:slit-factorization} and \(p_a(0)=1\) give
\(p_a(r)=\prod_j(1+\lambda_{a_j^{(j)}}r_j)\).
The normalization at zero is algebraic and does not require
zero to belong to \(\Omega_Z\).  Comparing coefficients, and
varying the axes, proves \eqref{stat:pauli-product} for every
Pauli string.  Local factors inherit the Choi and positivity
properties by the marginal construction in the proof of
Theorem~\ref{stat:visible}.  Tensor products prove the converse.
\end{proof}

For CPTP local Pauli factors the exact one-site condition is
\[
 \lambda_z\ge\max\{|\lambda_x|,|\lambda_y|\},
\]
or equivalently
\(\min(p_0,p_z)\ge\max(p_x,p_y)\) in the Pauli-error
probabilities.  This is the Pauli specialization of the local
classification in Theorem~\ref{loc:classification}; its
sufficiency is also \cite[Theorem~3]{WBG26}.
Because the Pauli transform between error probabilities and
eigenvalues is an invertible tensor-product transform,
\eqref{stat:pauli-product} says that the entire error
distribution is a product distribution.

\subsection{Classical channels and the purity boundary}

Consider the channel that measures \(m\) input bits in the
computational basis and prepares a diagonal distribution
\(P(b\mid a)\) on \(n\) output bits.  Its full Choi
polynomial reduces to
\begin{equation}
 K(r,s)=\sum_{a,b}P(b\mid a)r^bs^a=\sum_aA_a(r)s^a,
 \qquad K(1,s)=\prod_i(1+s_i).
 \label{stat:classical-kernel}
\end{equation}
The coefficients are nonnegative.  As above, reducing each
row/column pair to its product is an equivalence of disk
stability.  Define the collision purity of column \(a\) by
\(\mathcal C(a)=\sum_bP(b\mid a)^2=\|A_a\|_2^2\).

\begin{theorem}[The constant-purity boundary]
\label{stat:classical-purity}
Every stable stochastic kernel \eqref{stat:classical-kernel}
satisfies \(\mathcal C(a)\le\mathcal C(0)\) for every input.
Equality at every singleton input \(e_i\) holds if and only
if, for a partition \(S_0,S_1,\ldots,S_m\) of the output
variables, the kernel is
\begin{equation}
 K(r,s)=C(r_{S_0})\prod_{i=1}^m
                 \bigl(p_i(r_{S_i})+s_i p_i^*(r_{S_i})\bigr).
 \label{stat:classical-purity-form}
\end{equation}
Here \(C\) and the \(p_i\) are stable, multiaffine, have
nonnegative coefficients, satisfy \(C(1)=p_i(1)=1\), and
\(\gcd(p_i,p_i^*)=1\); reflection is over the full indicated
block.  Empty blocks \(S_i\) are allowed and contribute
\(1+s_i\).  Every column of such a kernel has the same
collision purity.
\end{theorem}

\begin{proof}
Fix the output variables in the polydisk.  For any input subset
\(T\), set its variables equal to \(t\) and the other input
variables to zero.  The resulting univariate polynomial is
nonvanishing for \(|t|<1\), has constant coefficient \(A_0\),
and top coefficient \(A_T\) in degree \(|T|\).  If that top
coefficient is nonzero, the product of root moduli is at least
one, whence \(|A_T|\le|A_0|\).  If it is zero the inequality
is immediate.  Continuity and torus Parseval prove the purity
inequality.

Under singleton equality, Lemma~\ref{stat:saturation} gives
disjoint primitive factors \(C\prod_i(p_i+s_iq_i)\).
The stable constant slice makes the constant coefficients of
\(C,p_i\) nonzero.  Fix their phases to make these constants
positive.  Each coefficient of each factor can be isolated
as a coefficient of \(K\) by selecting the constant monomial
in all the other factors; disjoint supports therefore make
all coefficients of \(C,p_i,q_i\) nonnegative.
Lemma~\ref{stat:reflected-pair} gives
\(q_i=e^{i\theta_i}p_i^*\), and nonnegativity forces the
phase to be one.  Stochasticity at output variables one gives
\(p_i(1)=q_i(1)>0\).  Normalizing these common values and
the prepared factor gives \eqref{stat:classical-purity-form},
with no scalar remainder because every column has mass one.
Variables absent from \(K\) are fixed-zero prepared outputs.

Conversely, Lemma~\ref{stat:reflected-pair} makes each factor
stable.  The displayed normalization makes the nonnegative
coefficient kernel stochastic.  Reflection preserves
coefficient norms and disjoint supports make them multiply,
so all columns have equal collision purity.
\end{proof}

\begin{corollary}[Deterministic classical channels]
\label{stat:deterministic}
A deterministic stochastic kernel is full Choi--LY if and
only if every output bit is identically zero or copies one
input bit.  An input may be copied to any disjoint block of
outputs or discarded.
\end{corollary}

\begin{proof}
Write \(P(b\mid a)=\mathbf{1}_{b=f(a)}\).  All column
purities equal one.  The rational product identity in
Lemma~\ref{stat:saturation} therefore gives
\[
 r^{f(a)}=r^{f(0)+\sum_{i:a_i=1}(f(e_i)-f(0))}.
\]
Stable evaluation at the full origin requires \(f(0)=0\).
Each output coordinate is consequently a sum of its
singleton Boolean values.  Since it must still be a bit
when all inputs are one, at most one singleton value can
be one.  This proves necessity.  Conversely the kernel of
such a map is
\(\prod_i(1+s_i\prod_{j\in S_i}r_j)\), with disjoint
copied-output blocks, and is manifestly stable.
\end{proof}

\begin{theorem}[One classical output]
\label{stat:classical-one-output}
A stochastic kernel from \(m\) input bits to one output bit
is full Choi--LY if and only if
\begin{equation}
 P(1\mid a)=\alpha+\sum_i\lambda_i a_i,
 \qquad \alpha\ge0,\quad\lambda_i\ge0,\quad
                   2\alpha+\sum_i\lambda_i\le1.
 \label{stat:randomized-copy}
\end{equation}
Thus it chooses to copy input \(i\) with probability
\(\lambda_i\), to output zero with probability
\(1-2\alpha-\sum_i\lambda_i\), or to output a fresh fair
bit with probability \(2\alpha\).
\end{theorem}

\begin{proof}
Write \(K(r,s)=A(s)+rB(s)\).  Stochasticity gives
\(A+B=\prod_i(1+s_i)\).  Put \(H=(A-B)/(A+B)\), so
\[
 K=\frac{\prod_i(1+s_i)}2\bigl((1+r)+(1-r)H(s)\bigr).
\]
The quotient \((1+r)/(1-r)\) ranges over the open right
half-plane; thus stability is exactly \(\Re H(s)\ge0\)
on the input polydisk.  Set
\(u_i=(1-s_i)/(1+s_i)\).  Then \(H\) becomes a real
multiaffine polynomial in the right-half-plane variables
\(u_i\), because every
\(s^a/\prod_i(1+s_i)\) is a product of factors
\((1\pm u_i)/2\).

A real polynomial with nonnegative real part on a product
of right half-planes has the form
\(c+\sum_i\lambda_i u_i\), with \(c,\lambda_i\ge0\).
To see this, restrict to \(u_i=v_i t\), \(v_i>0\),
\(\Re t>0\).  A real one-variable polynomial of degree
at least two has a leading term with negative real part
along some interior right-half-plane ray at large radius.
Every such restriction therefore has degree at most one.
Varying \(v\) in its open positive orthant forces all
higher homogeneous parts to vanish.  Positive real
evaluations near zero, or with one variable tending to
infinity, show that the remaining coefficients are
nonnegative.

Evaluating the transformed polynomial identity at
\(u_i=(-1)^{a_i}\) gives
\(1-2P(1\mid a)=c+\sum_i\lambda_i(-1)^{a_i}\).
This is polynomial evaluation, including at vertices
corresponding to poles of the original change of variables.
Consequently \eqref{stat:randomized-copy} holds with
\(\alpha=(1-c-\sum_i\lambda_i)/2\).  The original
probability at \(a=0\) makes \(\alpha\ge0\), and
\(c\ge0\) gives the last inequality.  Conversely these
conditions define the stated stochastic mixture and a
positive-real polynomial \(H\), proving stability.
\end{proof}

Every single-output marginal of a multibit classical
full Choi--LY channel therefore has the form
\eqref{stat:randomized-copy}.  Marginalization sets the
other output variables to one; boundary Hurwitz applies,
and stochastic normalization excludes zero.  This does
not determine the joint output correlations.

\begin{proposition}[Nonconvexity among reset channels]
\label{stat:reset-nonconvex}
For \(n\ge3\), the distribution
\(\nu_\theta=(1-\theta)\delta_{0^n}
+\theta\operatorname{Unif}(\{0,1\}^n)\),
\(0\le\theta\le1\), has a stable polynomial precisely when
\begin{equation}
 0\le\theta\le\frac1{1+\cos^n(\pi/n)}
 \quad\text{or}\quad\theta=1.
 \label{stat:reset-threshold}
\end{equation}
For \(n=1,2\), all \(\theta\in[0,1]\) are allowed.
Consequently the set of full Choi--LY channels is not
convex, even among classical reset channels.
\end{proposition}

\begin{proof}
The distribution polynomial is
\(f_\theta(r)=1-\theta+2^{-n}\theta\prod_j(1+r_j)\).
The endpoints \(\theta=0,1\) are stable.  For
\(0<\theta<1\), an interior zero means that a product
of \(n\) points in the open disk centered at one with
radius one equals \(-2^n(1-\theta)/\theta\).
Write these nonzero points as
\(\rho_j e^{i\alpha_j}\), where
\(|\alpha_j|<\pi/2\) and
\(0<\rho_j<2\cos\alpha_j\).
A negative real product requires
\(\sum_j\alpha_j\) to be an odd multiple of \(\pi\),
which is impossible for \(n=1,2\).
For \(n\ge3\), concavity of \(\log\cos\) implies
\[
 \prod_j\rho_j
 <2^n\prod_j\cos\alpha_j
 \le[2\cos(\pi/n)]^n.
\]
An odd multiple of larger absolute value can only reduce
the bound.  Every smaller positive product magnitude is
attained by setting all arguments to \(\pi/n\) and
choosing equal moduli.  Therefore an interior zero occurs
exactly when
\(2^n(1-\theta)/\theta<[2\cos(\pi/n)]^n\), the
complement of \eqref{stat:reset-threshold} away from the
separately treated endpoint.  Equality requires boundary
points and is stable for the open-disk definition.
Resetting to this distribution multiplies its polynomial
by the stable input trace factor \(\prod_i(1+s_i)\),
so the same criterion applies to the channel.
\end{proof}

% END SECTION static_channels

% BEGIN SECTION local_channels
\section{The complete one-qubit static classification}\label{loc:section}

We give exact parameter criteria for a single qubit, including every
singular stratum.  The stability classification first requires only
Hermiticity preservation and trace preservation.  Complete positivity is
then imposed by explicit matrix or scalar inequalities.

Write a complex-linear Hermiticity-preserving, trace-preserving map in
Bloch coordinates as
\begin{equation}\label{loc:bloch}
 \Phi(I)=I+t\cdot\sigma,\qquad
 \Phi(\sigma_j)=\sum_iT_{ij}\sigma_i,
 \qquad t\in\R^3,\quad T\in\R^{3\times3},
\end{equation}
where $\sigma=(X,Y,Z)$.  On trace-one Hermitian matrices its action is
$r\mapsto t+Tr$.  Invertibility means complex-linear invertibility on
$M_2(\C)$, equivalently $\det T\ne0$.  Full Choi stability always
refers to all four variables of $F_{J(\Phi)}$.

\begin{theorem}[All one-qubit HP/TP maps]\label{loc:classification}
The map \eqref{loc:bloch} has a full Lee--Yang Choi kernel if and only if
it belongs to exactly one of the following two classes.
\begin{enumerate}
\item There is a nonzero real $2\times2$ matrix $D$ such that
\begin{equation}\label{loc:transverse-form}
 T=\begin{pmatrix}D&0\\0&\lambda\end{pmatrix},\qquad
 \lambda>0,\qquad t=(\tau,0),\qquad \tau\in\R^2,\quad\|\tau\|<1,
\end{equation}
and, for every $n\in\R^2$ with $\|n\|=1$,
\begin{equation}\label{loc:angular}
 \|D^{\mathsf T}n\|^2\leq m_\lambda(\tau\cdot n),\qquad
 m_\lambda(r)=\min_{-1\leq x\leq1}
                  \bigl((x+r)^2+\lambda^2(1-x^2)\bigr).
\end{equation}
\item The transverse linear action vanishes and
\begin{equation}\label{loc:erasure-form}
 T=\diag(0,0,\lambda),\qquad
 \lambda\geq0,\quad t_z\geq0,\quad\|t_\perp\|\leq1.
\end{equation}
\end{enumerate}
For $|r|<1$, the minimum in \eqref{loc:angular} is
\begin{equation}\label{loc:minimum}
 m_\lambda(r)=
 \begin{cases}
 \lambda^2\bigl(1-r^2/(1-\lambda^2)\bigr),
   &0<\lambda<1,\ |r|\leq1-\lambda^2,\\
 (1-|r|)^2,&0<\lambda<1,\ |r|>1-\lambda^2,\\
 (1-|r|)^2,&\lambda\geq1.
 \end{cases}
\end{equation}
The complex-linear rank is $2+\rank D$ in the first class and
$1+\boldsymbol{1}_{\{\lambda>0\}}$ in the second.  In particular the
first class is invertible exactly when $D$ is invertible.
\end{theorem}

\subsection{Boundary kernels and the proof of the classification}

\begin{lemma}[Trace-one scalar kernels]\label{loc:scalar-kernel}
A Hermitian trace-one matrix with transverse Bloch norm $R$ and
longitudinal coordinate $h$ has a stable coefficient polynomial if and
only if $R\leq1$ and $h\geq0$.  For $0\leq R\leq1$, the function
\begin{equation}\label{loc:positive-real}
 K_R(z,w)=\frac{1+zw+R(z+w)}{1-zw}
\end{equation}
has strictly positive real part on $\D^2$, and the infimum of that real
part is zero.
\end{lemma}

\begin{proof}
A rotation about $Z$ makes the transverse component $(R,0)$ and only
rotates the disk variables.  The coefficient polynomial becomes
\[
 f(z,w)=\frac{1+h+R(z+w)+(1-h)zw}{2}
       =\frac{1-zw}{2}(K_R(z,w)+h).
\]
For $a=|z|<1$, $b=|w|<1$ and $0\leq R\leq1$, multiplication by
$|1-zw|^2$ shows that the real part of the numerator of $K_R$ is at
least
\begin{align*}
 1-a^2b^2-R\bigl(a(1-b^2)+b(1-a^2)\bigr)
 &\geq(1-ab)(1-a)(1-b)>0.
\end{align*}
This proves sufficiency.  Setting $z=w=is$ and letting $s\uparrow1$
shows that the infimum of $\Re K_R$ is zero.

For necessity, stability gives the boundary root-modulus inequality
obtained by treating $f$ as affine in $w$.  At $z=e^{i\theta}$ its
squared difference between the constant and linear coefficients is
$h(1+R\cos\theta)$.  Its average is $h$, so $h\geq0$.  If $h>0$,
the same inequality for all $\theta$ gives $R\leq1$.  If $h=0$ and
$R>1$, setting $w=0$ gives the interior zero $z=-1/R$.
\end{proof}

\begin{lemma}[The local block form]\label{loc:block}
Every map in Theorem~\ref{loc:classification} whose full Choi kernel is
stable satisfies
\begin{equation}\label{loc:block-intermediate}
 T=D\oplus\lambda,\qquad\lambda\geq0.
\end{equation}
If $D\ne0$, then $t_z=0$ and $\lambda>0$.
\end{lemma}

\begin{proof}
Let $\phi=\Phi^\dagger$.  It is unital and Hermiticity preserving.
Lemma~\ref{stat:bridge} makes it a preserver of complex matrix-LY
polynomials, with a possible zero output.  Near the identity that zero
alternative is excluded by continuity.  Apply
Lemma~\ref{stat:edge} to the two-sided stable curves
$\phi(e^{sX})$, $\phi(e^{sY})$, and $\phi(e^{siZ})$.
Hermiticity preservation, and its anti-Hermitian counterpart, give
\begin{equation}\label{loc:adjoint-block}
 \begin{aligned}
 \phi(X)&=t_xI+D_{11}X+D_{12}Y,\\
 \phi(Y)&=t_yI+D_{21}X+D_{22}Y,\\
 \phi(Z)&=t_zI+\lambda Z.
 \end{aligned}
\end{equation}
The stable matrix $I+sZ$, $s>0$ small, maps to
$(1+st_z)I+s\lambda Z$.  Its root-modulus inequality is
\[
 |1+st_z+s\lambda|^2-|1+st_z-s\lambda|^2
       =4s\lambda(1+st_z)\geq0,
\]
so $\lambda\geq0$.

If $D\ne0$, choose a real transverse field $A$ such that the transverse
component of $\phi(A)$ is nonzero.  For small nonzero real
$\varepsilon$, both $A+i\varepsilon Z$ and the traceless part of its
image are hyperbolic.  The input pencil
$I+s(A+i\varepsilon Z)$ is stable for both signs of sufficiently
small $s$: its exponentials are stability units by
Lemma~\ref{stat:edge}, and the identity
$(A+i\varepsilon Z)^2\in\R I$ rewrites these exponentials as scalar
multiples of the pencil near the identity.  Its image is a stable pencil
whose scalar part is $t_A+i\varepsilon t_z$.
Lemma~\ref{stat:hyperbolic} forces this scalar to be real, hence $t_z=0$.

Finally specialize the two input Choi variables to
$u=e^{i\theta}$, $v=-e^{-i\theta}$.  The input matrix is
$Z+i n\cdot\sigma_\perp$, with $n$ ranging over the real unit circle.
The output matrix is
$\lambda Z+i(Dn)\cdot\sigma_\perp$.  Hurwitz's theorem makes its
polynomial stable or zero.  If $\lambda=0$, choose $n$ with $Dn\ne0$.
The polynomial is then nonzero but its constant coefficient is zero,
contradicting stability.  This proves $\lambda>0$.
\end{proof}

\begin{proof}[Proof of Theorem~\ref{loc:classification}]
By Lemma~\ref{loc:block}, it remains to treat the two possible transverse
blocks.  Suppose first that $D=0$, and rotate about $Z$ so that
$t_\perp=(R,0)$.  Direct expansion of the Choi kernel gives
\begin{equation}\label{loc:erasure-kernel}
 F_{J(\Phi)}=
 \frac{(1+uv)(1-zw)}2
 \left(K_R(z,w)+t_z+\lambda\frac{1-uv}{1+uv}\right).
\end{equation}
The last quotient ranges over the right half-plane.  Thus
$R\leq1$, $t_z\geq0$, and $\lambda\geq0$ imply strict nonvanishing
by Lemma~\ref{loc:scalar-kernel}.
Conversely, specialization $u=v=0$ gives a trace-one Hermitian kernel
with transverse norm $R$ and longitudinal coordinate $t_z+\lambda$,
so $R\leq1$.  If $\lambda=0$, that same lemma gives $t_z\geq0$.
If $\lambda>0$ and $t_z<0$, choose $z,w$ with
$\Re(K_R+t_z)<0$.  A value of $(1-uv)/(1+uv)$ in the right
half-plane then cancels the parenthesis in
\eqref{loc:erasure-kernel}, a contradiction.  This proves the second
class completely, without positivity.

Now suppose $D\ne0$.  Lemma~\ref{loc:block} gives
$\lambda>0$ and $t_z=0$.  Set
\begin{equation}\label{loc:coefficients}
 \begin{gathered}
 A_0=(1+\lambda)/2,\qquad B_0=(1-\lambda)/2,
 \qquad \zeta=(t_x-it_y)/2,\\
 a=(D_{11}+D_{22}+i(D_{12}-D_{21}))/2,\qquad
 b=(D_{11}-D_{22}+i(D_{12}+D_{21}))/2.
 \end{gathered}
\end{equation}
The full polynomial has the reciprocal decomposition
\begin{equation}\label{loc:reciprocal}
 F_{J(\Phi)}=p+vp^\sharp,\qquad
 p=g+uh,\quad
 g=A_0+\zeta w+\bar\zeta z+B_0zw,\quad
 h=\bar b w+\bar a z,
\end{equation}
where
$p^\sharp(z,w,u)=zwu\,\overline{p(1/\bar z,1/\bar w,1/\bar u)}$.
This is an algebraic identity with no positivity assumption.

If the full polynomial is stable, its specializations $p$ and $g$ are
stable, since their constant coefficient is $A_0>0$.
Lemma~\ref{loc:scalar-kernel} applied to $g$ gives
$R=\|t_\perp\|\leq1$.  The root test in $u$ and continuity give
$|h|\leq|g|$ on the output two-torus.  Write
$z=e^{i(\theta+\varphi)}$, $w=e^{i(\theta-\varphi)}$ and
$n=(\cos\varphi,-\sin\varphi)$.  Direct expansion yields
\begin{equation}\label{loc:torus-norms}
 |h|^2=\|D^{\mathsf T}n\|^2,\qquad
 |g|^2=(\cos\theta+t_\perp\cdot n)^2+
                         \lambda^2\sin^2\theta.
\end{equation}
Minimization in $x=\cos\theta$ proves \eqref{loc:angular}.
If $R=1$, rotate so that $t_\perp=e_1$.  The endpoint $x=-1$
gives $\|D^{\mathsf T}n\|\leq1-n_1$ for every unit $n$.
At $n=e_1$ the first row of $D$ vanishes.  Setting
$n=(\cos\alpha,\sin\alpha)$ then gives
$|\sin\alpha|\|D^{\mathsf T}e_2\|\leq1-\cos\alpha$.
Division by $|\sin\alpha|$ and passage to $\alpha=0$ kills the
second row.  Since $D\ne0$, necessarily $R<1$.

Conversely, assume \eqref{loc:transverse-form} and
\eqref{loc:angular}.  The polynomial $g$ has no zero on the closed
bidisk.  Indeed its constant coefficient in $w$, namely
$A_0+\bar\zeta z$, has no zero on the closed disk, and at $|z|=1$
the squared root-modulus difference is
\[
 |A_0+\bar\zeta z|^2-|\zeta+B_0z|^2
       =\lambda(1+2\Re(\bar\zeta z))>0.
\]
Here $2|\zeta|=R<1$.  The maximum modulus principle in $z$ proves
the asserted closed-bidisk nonvanishing.  Equations
\eqref{loc:angular} and \eqref{loc:torus-norms}, followed by the
maximum modulus principle in each variable, imply $|h/g|\leq1$ on
the bidisk.  Hence $p=g+uh$ is stable.

For any stable multiaffine $p$, one has $|p^\sharp/p|\leq1$ in the
open polydisk.  To include boundary zeros, apply the maximum modulus
principle to $p_r(z,w,u)=p(rz,rw,ru)$, $r<1$, which has no zeros on
the closed polydisk.  On its distinguished boundary
$|p_r^\sharp|=|p_r|$, so $|p_r^\sharp/p_r|\leq1$ inside.
Let $r\uparrow1$.  This proves the assertion, and hence
$p+vp^\sharp\ne0$ when $|v|<1$.  Sufficiency follows.

The formula \eqref{loc:minimum} follows by minimizing the quadratic
$(1-\lambda^2)x^2+2rx+r^2+\lambda^2$: for $\lambda<1$ its critical
point is $-r/(1-\lambda^2)$, and otherwise its minimum on the interval
is at an endpoint.  Finally the Pauli transfer matrix is
$\begin{psmallmatrix}1&0\\t&T\end{psmallmatrix}$, whose rank is
$1+\rank T$.  This gives all rank assertions.
\end{proof}

\subsection{Complete positivity and finite inequalities}

\begin{theorem}[The CPTP subfamilies]\label{loc:cptp}
In the first class of Theorem~\ref{loc:classification}, complete
positivity is equivalent to
\begin{equation}\label{loc:choi-psd}
 J=\begin{pmatrix}
 A_0&0&\zeta&a\\
 0&B_0&\bar b&\zeta\\
 \bar\zeta&b&B_0&0\\
 \bar a&\bar\zeta&0&A_0
 \end{pmatrix}\succeq0,
\end{equation}
with coefficients \eqref{loc:coefficients}.  For a positive
trace-preserving map of this form with $0<\lambda<1$, the stability
criterion \eqref{loc:angular} is equivalent to the finite matrix
inequality
\begin{equation}\label{loc:lmi}
 \frac{DD^{\mathsf T}}{\lambda^2}
       +\frac{\tau\tau^{\mathsf T}}{1-\lambda^2}\preceq I_2.
\end{equation}
At $\lambda=1$, positivity forces $\tau=0$, and the stability
criterion is $\|D\|_{\mathrm{op}}\leq1$.

For the singular rank-three subclass, input and output rotations about
$Z$ put $D=\diag(d,0)$ with $d>0$ and $\tau=(p,q)$.
The complete CPTP and full-Choi-LY conditions are the two scalar
inequalities
\begin{equation}\label{loc:rank-three}
 q^2+\bigl(d+\sqrt{p^2+\lambda^2}\bigr)^2\leq1,
 \qquad
 \frac{d^2}{\lambda^2}
       +\frac{p^2}{1-\lambda^2-q^2}\leq1,
 \qquad d,\lambda>0.
\end{equation}
The first inequality makes the denominator in the second strictly
positive.

In the transverse-erasure class, the complete conditions are
\begin{equation}\label{loc:erasure-cptp}
 t_z,\lambda\geq0,\qquad
 \|t_\perp\|^2+(t_z+\lambda)^2\leq1.
\end{equation}
\end{theorem}

\begin{proof}
The matrix in \eqref{loc:choi-psd} is the Choi matrix in the order
$(00,01,10,11)$, with each pair ordered as (output,input).
Its partial trace over the output equals the input identity.  Choi's
criterion gives the first assertion.

For the finite inequality, let $k=1-\lambda^2>0$ and fix a transverse
unit vector $n$.  Reversing $n$ if necessary, set
$r=\tau\cdot n\geq0$ and $s^2=\|D^{\mathsf T}n\|^2$.
The desired scalar form of \eqref{loc:lmi} is
\begin{equation}\label{loc:lmi-scalar}
 Q=\lambda^2-s^2-\lambda^2r^2/k\geq0.
\end{equation}
If $r\leq k$, this is exactly the first branch of
\eqref{loc:minimum}.  If $r>k$, put $a=k/r\in(0,1)$.
Positivity means that the affine image of the Bloch unit ball lies in
the unit ball.  Its support in the unit output direction
$(an,\sqrt{1-a^2})$ therefore satisfies
\[
 ar+\sqrt{a^2s^2+\lambda^2(1-a^2)}\leq1.
\]
Squaring the nonnegative sides gives
\[
 0\leq k-2ar+a^2(r^2-s^2+\lambda^2)=\frac{k^2}{r^2}Q.
\]
Thus positivity itself implies \eqref{loc:lmi-scalar} in the clipped
region.  Conversely \eqref{loc:lmi-scalar} bounds the unconstrained
minimum over all real $x$, hence implies \eqref{loc:angular} on
$[-1,1]$.  Varying $n$ proves \eqref{loc:lmi}.  Positivity at the
input Bloch points $\pm e_z$ gives
$\|\tau\|^2+\lambda^2\leq1$, proving both $\lambda\leq1$ and the
stated endpoint assertion.

Suppose now $D$ has rank one.  Real rotations in the transverse input
and output planes put it in the form $\diag(d,0)$, $d>0$; these are
implemented by rotations about $Z$ and preserve both properties under
consideration.  In these coordinates the Choi matrix satisfies
$2J=I+H$, where
\[
 H=pX\otimes I+qY\otimes I+dX\otimes X+\lambda Z\otimes Z.
\]
Pauli multiplication gives
\[
 H^2=(p^2+q^2+d^2+\lambda^2)I
          +2pd\,I\otimes X-2d\lambda\,Y\otimes Y.
\]
The last two Pauli matrices anticommute.  The eigenvalues of $H^2$ are
$q^2+(\sqrt{p^2+\lambda^2}\pm d)^2$, each twice.  Also
$\Tr H=\Tr H^3=0$, by Pauli orthogonality.  Newton's identities make
the characteristic polynomial of $H$ even, so its spectrum is
symmetric about zero.  Its least eigenvalue is consequently
$-\sqrt{q^2+(\sqrt{p^2+\lambda^2}+d)^2}$.
This proves the first inequality in \eqref{loc:rank-three} as the
exact CP condition.  Since $d,\lambda>0$, it also implies
$\lambda<1$ and $1-\lambda^2-q^2>0$.
Taking the Schur complement of the lower-right entry in
\eqref{loc:lmi} now gives precisely the second inequality in
\eqref{loc:rank-three}.

Finally, a transverse-erasure map measures in the $Z$ basis and prepares
\begin{equation}\label{loc:prepared-states}
 \rho_\pm=\tfrac12\bigl(I+t_xX+t_yY+(t_z\pm\lambda)Z\bigr).
\end{equation}
Its Choi matrix is block diagonal on the input and is positive exactly
when both prepared matrices are positive, equivalently
$\|t_\perp\|^2+(t_z\pm\lambda)^2\leq1$.
Since $t_z,\lambda\geq0$, the two inequalities reduce to
\eqref{loc:erasure-cptp}.
\end{proof}

\begin{remark}\label{loc:singular-distinction}
Full Choi stability does not require both states in
\eqref{loc:prepared-states} to be individually Lee--Yang: the quantity
$t_z-\lambda$ may be negative.  The choice
$t_\perp=0$, $t_z=\lambda=1/3$ is a rank-two CPTP example with
positive longitudinal drift.  A rank-three example retaining transverse
information is $t=0$, $D=\diag(1/5,0)$, $\lambda=2/5$.
At rank one, the maps are precisely resets
$M\mapsto\Tr(M)\rho$; their CPTP full-Choi-LY condition is
$t_z\geq0$ and $\|t\|\leq1$.
These singular examples explain why the zero longitudinal drift conclusion for an
invertible map cannot be extended without a rank hypothesis.
\end{remark}

% END SECTION local_channels

% BEGIN SECTION higher_capacity
\section{Coherent stability at higher capacities}\label{spin:section}

The dynamical classification at arbitrary finite spins is
Theorem~\ref{dyn:arbitrary-spin}. This appendix gives two
complementary results. The infinitesimal classification determines all Hermitian
Hamiltonians whose coherent Gibbs kernels are stable at every temperature.
A separate, static argument classifies all independently rotation-covariant
trace-preserving maps, including maps between different spin representations.
The covariance assumption in the latter result is explicit throughout.

\subsection{Coherent normalization and the Hamiltonian cone}

At site $i$, fix a positive integer $\kappa_i$ and put $j_i=\kappa_i/2$.
Let $\mathcal H_{\kappa_i}=\C^{\kappa_i+1}$ have orthonormal basis
$|a_i\rangle$, $0\leq a_i\leq\kappa_i$.  For multi-indices $a$ set
\begin{equation}\label{spin:weights}
 b_a=\prod_i\binom{\kappa_i}{a_i},\qquad
 |a\rangle\longleftrightarrow \sqrt{b_a}\,z^a.
\end{equation}
Thus the monomial inner product is
$\langle z^a,z^b\rangle=\delta_{ab}/b_a$.  All adjoints in this section
refer to the physical orthonormal basis, or equivalently to this weighted
polynomial inner product.  The coherent kernel of a matrix $M$ is
\begin{equation}\label{spin:kernel}
 F_M(z,w)=\sum_{a,b}\sqrt{b_a b_b}\,M_{ab}z^a w^b.
\end{equation}
In particular, $F_I(z,w)=\prod_i(1+z_iw_i)^{\kappa_i}$.
We call $M$ coherently Lee--Yang if $F_M$ has no zero when all its variables
belong to $\D$.  Both weights in \eqref{spin:kernel} are essential.

The spin matrices have the following differential realization:
\begin{equation}\label{spin:differential}
 S_i^x=\frac{(1-z_i^2)\partial_i+\kappa_i z_i}{2},\qquad
 S_i^y=\frac{(1+z_i^2)\partial_i-\kappa_i z_i}{2i},\qquad
 S_i^z=\frac{\kappa_i}{2}-z_i\partial_i.
\end{equation}
They satisfy
$(S_i^x)^2+(S_i^y)^2+(S_i^z)^2=j_i(j_i+1)I$.
Write $S_i^\perp=(S_i^x,S_i^y)^{\mathsf T}$.

\begin{theorem}[Higher-spin Hamiltonians]\label{spin:hamiltonian}
For a Hermitian operator $H$ on $\bigotimes_i\mathcal H_{\kappa_i}$, the
following conditions are equivalent:
\begin{enumerate}
\item $e^{-\beta H}$ is coherently Lee--Yang for every $\beta\geq0$;
\item this holds for every $0\leq\beta<\varepsilon$, for some
      $\varepsilon>0$;
\item $H$ admits the expression
\begin{align}\label{spin:hamiltonian-form}
 H={}&h_0 I+\sum_i(h_i^xS_i^x+h_i^yS_i^y+h_i^zS_i^z)
       +\sum_{i:\,\kappa_i\geq2}(S_i^\perp)^{\mathsf T}B_iS_i^\perp
       \notag\\
    &+\sum_{i<j}\bigl((S_i^\perp)^{\mathsf T}J_{ij}S_j^\perp
                         +e_{ij}S_i^zS_j^z\bigr),
\end{align}
where all coefficients are real, $B_i=B_i^{\mathsf T}$, and
\begin{equation}\label{spin:hamiltonian-cone}
 h_i^z\leq0,\qquad B_i\succeq0,\qquad
 e_{ij}\leq-\|J_{ij}\|_{\mathrm{op}}.
\end{equation}
\end{enumerate}
The coefficient $(B_i)_{12}$ multiplies
$S_i^xS_i^y+S_i^yS_i^x$.  At a capacity-one site all single-site
quadrupoles reduce to scalars or zero and are omitted from
\eqref{spin:hamiltonian-form}.
\end{theorem}

We give the normalization and coordinate computations in the proof, so that
the finite-capacity infinitesimal theorem can be applied without changing
the physical adjoint.

\begin{lemma}[Polarization and the coherent Cayley transform]
\label{spin:cayley}
Define the Dicke isometry
\begin{equation}\label{spin:dicke}
 W_i|a\rangle=\binom{\kappa_i}{a}^{-1/2}
                  \sum_{|x|=a}|x\rangle,\qquad W=\bigotimes_iW_i.
\end{equation}
The binary matrix kernel of $WMW^*$ is the normalized polarization of
$F_M$, separately in each row and column block.  Moreover, let
\begin{equation}\label{spin:cayley-matrix}
 C_1=\frac1{\sqrt2}\begin{pmatrix}i&-i\\1&1\end{pmatrix},\qquad
 C_i=W_i^*C_1^{\otimes\kappa_i}W_i,\qquad C=\bigotimes_iC_i.
\end{equation}
For Hermitian $H$, coherent stability of $e^{-tH}$ for all sufficiently
small $t\geq0$ is equivalent to the condition that
$\exp(-tCHC^*)$ is a complete upper-half-plane stability preserver for
every $t\geq0$.
\end{lemma}

\begin{proof}
At binary words of weights $a,b$, the matrix entry of $WMW^*$ is
$M_{ab}/\sqrt{b_ab_b}$.  Identifying the variables in their respective
blocks multiplies this coefficient by $b_ab_b$, giving exactly
\eqref{spin:kernel}.  The disk polarization theorem therefore identifies
stability of the two kernels; see \cite[Theorem~2.1]{BB10II}.
Polarization, binary tensor contraction, and diagonal specialization show
that a coherently stable matrix sends each stable coherent vector to a
stable vector or zero.  An invertible matrix cannot produce the latter
alternative.  This argument does not require $WMW^*$ to be invertible
outside the symmetric subspace.

The matrix $C_i$ is unitary.  On polynomials it acts by
\begin{equation}\label{spin:cayley-action}
 f(z_i)\longmapsto
 \left(\frac{s_i+i}{\sqrt2}\right)^{\kappa_i}
 f\left(\frac{s_i-i}{s_i+i}\right),
\end{equation}
which identifies disk stability with upper-half-plane stability.
Consequently coherent stability of $e^{-tH}$ implies ordinary stability
preservation by $e^{-tCHC^*}$.  Preservation for a sufficiently small
time interval implies preservation for all times by composition.
Theorem~\ref{gen:main-complex} then gives complete preservation.

For the converse, put
$D_i=\diag_{0\leq a\leq\kappa_i}((-1)^{\kappa_i-a})$ and
$D=\bigotimes_iD_i$.  Symmetric powers of $C_1C_1^{\mathsf T}$ give
$CC^{\mathsf T}=D$, and
\begin{equation}\label{spin:universal-halfplane}
 F_D(s,t)=\prod_i(s_it_i-1)^{\kappa_i}
\end{equation}
is upper-half-plane stable.  Indeed $st=1$ is impossible for two points
of the upper half-plane.  A complete preserver
$N=e^{-tCHC^*}=Ce^{-tH}C^*$ applied in the $s$ variables to this kernel
produces the coherent coefficient matrix
$ND=Ce^{-tH}C^{\mathsf T}$.  This is precisely the separate Cayley
transform of both groups of variables in $F_{e^{-tH}}$.  It is nonzero
because $N$ is invertible.  Applying the inverse transforms proves the
required disk stability.
\end{proof}

\begin{proof}[Proof of Theorem~\ref{spin:hamiltonian}]
Put $A=-CHC^*$.  By Lemma~\ref{spin:cayley},
Theorem~\ref{gen:main-complex} applies to $A$.  We recall the parts of its conclusion
used here.  The generator belongs to the degree-at-most-two part of the
enveloping algebra of the local $\mathfrak{sl}_2$ actions.  Its
second-order differential coefficients are real and nonpositive on the
appropriate real lines or planes.  If $A=R+iS$ in the real monomial
basis, then its imaginary part has the form
\begin{equation}\label{spin:imaginary-wedge}
 S=cI+\sum_iJ_{i,\kappa_i}(q_i),\qquad
 J_{i,\kappa_i}(q)=q\partial_i-\frac{\kappa_i}{2}q',\qquad
 q_i\in\R[z_i]_{\leq2},\quad q_i\geq0\text{ on }\R.
\end{equation}
Conversely these conditions, together with preservation of the finite
degree box, characterize the generators.

The degree-two operator space of Proposition~\ref{gen:box-decomposition}
decomposes as the direct sum of the scalar
space, the local spin spaces, the local traceless symmetric quadrupoles
when $\kappa_i\geq2$, and the spaces spanned by
$S_i^aS_j^b$ for $i<j$.  Its dimension is
\begin{equation}\label{spin:degree-two-dimension}
 1+3n+5\#\{i:\kappa_i\geq2\}+9\binom n2.
\end{equation}
For clarity, the local assertion follows from the usual decomposition of
the square of the three-dimensional spin representation into its scalar,
antisymmetric, and traceless symmetric parts: the antisymmetric part is
the spin commutator, the scalar part is the Casimir, and the rank-two
part is nonzero precisely for $\kappa_i\geq2$.  Distinct sites give
independent tensor factors.  This also proves the stated directness.
Since the Cayley action is a local spin rotation, it preserves this
operator space.  Thus an arbitrary Hermitian $H$ satisfying condition (2)
already has total spin degree at most two, with the stated capacity-one
exceptions.

Hermiticity of $A$ and the real diagonal monomial inner product give
$R^*=R$ and $S^*=-S$.  For one site write
\[
 J^-=\partial,\qquad J^0=z\partial-\kappa/2,\qquad
 J^+=z^2\partial-\kappa z.
\]
In the weighted inner product,
$(J^-)^*=-J^+$ and $(J^0)^*=J^0$.
Consequently skew-adjointness of
$aJ^-+bJ^0+cJ^+$, for real $a,b,c$, forces $b=0$ and $a=c$.
The scalar in \eqref{spin:imaginary-wedge} vanishes, and hence
\begin{equation}\label{spin:imaginary-fields}
 q_i(s)=\alpha_i(1+s^2),\quad\alpha_i\geq0,
 \qquad iJ_{i,\kappa_i}(q_i)=-2\alpha_iS_i^y.
\end{equation}
The Cayley rotations are
\begin{equation}\label{spin:rotation}
 CS_i^xC^*=-S_i^z,\qquad CS_i^yC^*=S_i^x,\qquad
 CS_i^zC^*=-S_i^y.
\end{equation}
Thus the original longitudinal field contributes $h_i^zS_i^y$ to $A$,
and \eqref{spin:imaginary-fields} gives $h_i^z=-2\alpha_i\leq0$.
Each mixed longitudinal--transverse quadratic term in $H$ gives an
imaginary second-order coefficient in $A$.  The direct decomposition
above, or the linearly independent principal polynomials at each site,
prevents cancellation between distinct such terms.  They all vanish.

It remains to compute the real second-order inequalities.  The principal
coefficients of $(S^x,S^y,S^z)$ in \eqref{spin:differential} are
$((1-s^2)/2,-i(1+s^2)/2,-s)$.  Put
\[
 u(s)=\left(\frac{2s}{1+s^2},\frac{1-s^2}{1+s^2}\right)^{\mathsf T};
\]
its range is dense in the real unit circle.  For the pair term
$(S_i^\perp)^{\mathsf T}J_{ij}S_j^\perp+e_{ij}S_i^zS_j^z$,
the coefficient $q_{ij}$ of $\partial_i\partial_j$ in $A$ satisfies
\begin{equation}\label{spin:pair-principal}
 \frac{4q_{ij}(s,t)}{(1+s^2)(1+t^2)}
       =e_{ij}-u(s)^{\mathsf T}J_{ij}u(t).
\end{equation}
Its nonpositivity is therefore equivalent to
$e_{ij}\leq-\|J_{ij}\|_{\mathrm{op}}$.
Before fixing the local Casimir gauge, write a surviving local quadratic
term as $(S_i^\perp)^{\mathsf T}K_iS_i^\perp+c_i(S_i^z)^2$, with
$K_i$ real symmetric.  For $\kappa_i\geq2$ its coefficient $q_i$ of
$\partial_i^2$ satisfies
\begin{equation}\label{spin:local-principal}
 \frac{4q_i(s)}{(1+s^2)^2}=c_i-u(s)^{\mathsf T}K_iu(s).
\end{equation}
Thus $c_i\leq\lambda_{\min}(K_i)$.  The Casimir identity rewrites this
term as $(S_i^\perp)^{\mathsf T}B_iS_i^\perp$ plus a scalar, where
$B_i=K_i-c_iI_2\succeq0$.  This proves necessity.

Conversely, \eqref{spin:hamiltonian-cone} makes
\eqref{spin:pair-principal} and \eqref{spin:local-principal} real and
nonpositive and makes the imaginary part satisfy
\eqref{spin:imaginary-wedge}.  The expression preserves the degree box
because it is built from the spin operators.
Theorem~\ref{gen:main-complex} and Lemma~\ref{spin:cayley} give coherent stability at
every time.  This proves (3)$\Rightarrow$(1); (1)$\Rightarrow$(2) is
immediate.
\end{proof}

\begin{corollary}[Exact symmetric-qubit realization]\label{spin:qubit-realization}
Every Hamiltonian in Theorem~\ref{spin:hamiltonian} has a Hermitian qubit
extension $\widehat H$, invariant under permutations within the capacity
blocks, whose pair coefficients satisfy the qubit Lee--Yang cone and such
that
\begin{equation}\label{spin:intertwining}
 \widehat HW=WH,\qquad e^{-\beta\widehat H}W=We^{-\beta H}
 \quad(\beta\in\R).
\end{equation}
\end{corollary}

\begin{proof}
For each $B_i\succeq0$, let $\lambda_i=\lambda_{\max}(B_i)$ and use
the exact identity
\begin{equation}\label{spin:casimir-extension}
 (S_i^\perp)^{\mathsf T}B_iS_i^\perp
 =(S_i^\perp)^{\mathsf T}(B_i-\lambda_iI_2)S_i^\perp
   -\lambda_i(S_i^z)^2+\lambda_i j_i(j_i+1)I.
\end{equation}
Replace each spin in this expression and in
\eqref{spin:hamiltonian-form} by
$\widehat S_i^a=\frac12\sum_{r=1}^{\kappa_i}\sigma_{ir}^a$.
Retain the last term of \eqref{spin:casimir-extension} as the indicated
scalar.  Interblock pairs have transverse matrix $J_{ij}/4$ and
longitudinal coefficient $e_{ij}/4$.  Intrablock pairs have transverse
matrix $(B_i-\lambda_iI_2)/2$ and longitudinal coefficient
$-\lambda_i/2$.  If $b_i=\lambda_{\min}(B_i)$, then
\[
 \|B_i-\lambda_iI_2\|_{\mathrm{op}}=\lambda_i-b_i\leq\lambda_i.
\]
All pairs therefore satisfy the qubit cone.  Identical-slot symmetric
off-diagonal products vanish by Pauli anticommutation; the remaining
identical-slot terms are scalars.  Each collective spin preserves the
symmetric sector and intertwines the physical spin through $W_i$.
This proves both identities in \eqref{spin:intertwining}.
\end{proof}

The reduction of higher spins to spin one-half and the corresponding
anisotropic sufficient Lee--Yang theorems are classical; see, in
particular, \cite[Theorem~2 and equations~(2.38)--(2.44)]{SuzukiFisher71}.
The corollary supplies explicit coefficients
for every Hamiltonian in the necessary-and-sufficient cone, including
single-ion quadrupoles.  It is an invariant-sector identity and does not
assert a spectral comparison with the other qubit spin sectors.

\subsection{Independently covariant maps between different capacities}

Let $U_k$ denote the spin-$k/2$ representation on
$\mathcal H_k=\C^{k+1}$, now allowing $k=0$.
For input capacities $\kappa=(\kappa_1,\ldots,\kappa_n)$ and output
capacities $\ell=(\ell_1,\ldots,\ell_n)$, a complex-linear map
$\Phi:\operatorname{End}(\bigotimes_i\mathcal H_{\kappa_i})
\to\operatorname{End}(\bigotimes_i\mathcal H_{\ell_i})$
is independently covariant if
\begin{equation}\label{spin:covariance}
 \Phi\bigl(U_\kappa(g)MU_\kappa(g)^*\bigr)
 =U_\ell(g)\Phi(M)U_\ell(g)^*,\qquad
 U_k(g)=\bigotimes_iU_{k_i}(g_i),
\end{equation}
for every $g=(g_1,\ldots,g_n)\in SU(2)^n$.
The input and output sites in this definition are matched in advance.

Use the Choi convention
$J(\Phi)=\sum_{a,b}\Phi(|a\rangle\langle b|)\otimes
|a\rangle\langle b|$.  Its coherent polynomial is
\begin{equation}\label{spin:choi-weights}
 F_{J(\Phi)}(z,w,u,v)=
 \sum_{c,e,a,b}\sqrt{b^{\ell}_cb^{\ell}_e
                          b^{\kappa}_ab^{\kappa}_b}\,
 [\Phi(|a\rangle\langle b|)]_{ce}\,z^cw^eu^av^b.
\end{equation}
Full coherent Choi stability means stability in all four groups of
variables.  For one matched pair of sites put
\begin{equation}\label{spin:invariants}
 K_o=1+zw,\qquad K_i=1+uv,\qquad
 A=(1+zu)(1+wv),\qquad B=K_oK_i.
\end{equation}

\begin{theorem}[Covariant root classification]\label{spin:covariant}
Suppose $\Phi$ is complex linear, trace preserving, and independently
covariant as in \eqref{spin:covariance}.  Its full coherent Choi kernel is
Lee--Yang if and only if $\Phi=\bigotimes_i\Phi_i$, where each local
factor is specified by an integer $0\leq d_i\leq\min(\kappa_i,\ell_i)$
and an unordered list $s_{i1},\ldots,s_{id_i}\geq0$ as follows.  With
$q_i(r)=\prod_{a=1}^{d_i}(r+s_{ia})$, its local kernel is
\begin{equation}\label{spin:covariant-kernel}
 F_{J(\Phi_i)}=
 \frac{K_{o,i}^{\ell_i-d_i}K_{i,i}^{\kappa_i-d_i}
            \prod_{a=1}^{d_i}(A_i+s_{ia}B_i)}
      {(\ell_i+1)\int_0^1q_i(x)\,dx}.
\end{equation}
Every such map is completely positive, although neither positivity nor
Hermiticity preservation was assumed.  Its complex-linear rank is
\begin{equation}\label{spin:covariant-rank}
 \rank\Phi=\prod_i(d_i+1)^2.
\end{equation}
It is injective exactly when $d_i=\kappa_i$ at every site, surjective
exactly when $d_i=\ell_i$ at every site, and invertible exactly when
$d_i=\kappa_i=\ell_i$ at every site.
\end{theorem}

We prove the theorem together with a factorization through an intermediate
spin representation.  The representation theory of covariant maps and
the general covariant completely positive simplex are established in
\cite{AlNuwairan14}; coherent descriptions and multipliers are also
studied in \cite{ARS24}.  We use the multiplicity-free commutant
description, and derive the additional root restriction directly from
stability.

\begin{lemma}[Range of the coherent ratio]\label{spin:ratio}
The image of $A/B$ on $\D^4$ is precisely
$\Omega=\C\setminus(-\infty,0]$.
\end{lemma}

\begin{proof}
For $a\in\D$, the disk automorphisms
\[
 f_a(x)=\frac{x-a}{1-\bar a x},\qquad
 h_a(y)=\frac{y+\bar a}{1+ay}
\]
satisfy
\[
 1+f_a(x)h_a(y)=
 \frac{(1-|a|^2)(1+xy)}{(1-\bar a x)(1+ay)}.
\]
Apply $f_a$ to $z,v$ and $h_a$ to $u,w$.  The ratio $A/B$ is unchanged.
Choosing $a=z$ makes the new $z$ zero, and reduces the ratio to
$(1+w'v')/(1+u'v')$.  It is nonzero.  If it equalled $-t$ for some
$t>0$, then $1+t=-v'(w'+tu')$, whereas
$|v'|(|w'|+t|u'|)<1+t$.  This proves inclusion in $\Omega$.
Conversely, setting $u=z$ and $v=w=-z$ gives
\[
 \frac AB=\left(\frac{1+z^2}{1-z^2}\right)^2.
\]
The successive maps $z\mapsto z^2$, $x\mapsto(1+x)/(1-x)$, and
$y\mapsto y^2$ map onto the disk, the right half-plane, and $\Omega$,
respectively.  This proves surjectivity.
\end{proof}

\begin{lemma}[Polynomial factorization on a product of slit planes]
\label{spin:slit-product}
If a nonzero polynomial $Q\in\C[r_1,\ldots,r_n]$ has no zero on
$\Omega^n$, then
\begin{equation}\label{spin:slit-factor}
 Q(r)=c\prod_iq_i(r_i),\qquad c\ne0,
\end{equation}
where each $q_i$ is monic and all its roots belong to $(-\infty,0]$.
\end{lemma}

\begin{proof}
Write $Q(x,y)=\sum_{k=0}^da_k(y)x^k$, with $a_d\not\equiv0$.
On any connected open patch of $\Omega^{n-1}$ where $a_d\ne0$, every
root of the monic polynomial $Q(x,y)/a_d(y)$ is real and nonpositive.
Its coefficients $a_k/a_d$ are therefore real-valued holomorphic
functions, and hence constant on the patch.  Polynomial identity gives
$Q(x,y)=a_d(y)p(x)$ everywhere, for a fixed monic polynomial $p$ with
all roots in $(-\infty,0]$.  Since $p$ is nonzero on $\Omega$, the
polynomial $a_d$ has no zero on $\Omega^{n-1}$.  Induction proves the
claim, skipping variables of degree zero.  This reasoning includes
repeated roots and slices on which the degree might initially drop.
\end{proof}

For $0\leq r\leq k$, let $R_{k\to r}$ be partial trace from a
Dicke-encoded $k$-qubit input to $r$ qubits, identified with
$\operatorname{End}(\mathcal H_r)$.  The reduced operator is supported
on the symmetric sector: a vector antisymmetric in a retained pair is
orthogonal to the fully symmetric input sector.  Define the adjoint map
\begin{equation}\label{spin:raw-cloner}
 C_{r\to k}(N)=
 W_k^*\bigl(W_rNW_r^*\otimes I_2^{\otimes(k-r)}\bigr)W_k.
\end{equation}
These maps are completely positive and covariant; $R_{k\to r}$ is
trace preserving.  At $r=0$ they are the ordinary trace and scalar
multiplication of $I_{k+1}$, respectively.

\begin{lemma}[Coherent compression kernels]\label{spin:compression}
For $0\leq r\leq\min(\kappa,\ell)$, the map
$T_r^{\kappa,\ell}=C_{r\to\ell}R_{\kappa\to r}$ has coherent Choi
kernel
\begin{equation}\label{spin:basis-kernel}
 A^rK_o^{\ell-r}K_i^{\kappa-r}.
\end{equation}
These $\min(\kappa,\ell)+1$ maps form a basis of all complex-linear
$SU(2)$-covariant maps between the two matrix spaces.  Their trace
effects are
\begin{equation}\label{spin:basis-effect}
 (T_r^{\kappa,\ell})^*(I_{\ell+1})
            =\frac{\ell+1}{r+1}I_{\kappa+1}.
\end{equation}
\end{lemma}

\begin{proof}
The Dicke coefficients give, for $c=a+k$ and $e=b+k$,
\begin{equation}\label{spin:compression-coefficients}
 [C_{r\to\ell}(|a\rangle\langle b|)]_{c,e}
 =\frac{\sqrt{\binom ra\binom rb}\binom{\ell-r}{k}}
             {\sqrt{\binom\ell c\binom\ell e}},
\end{equation}
and zero if the two shifts do not agree.  In particular,
\begin{equation}\label{spin:kernel-multiplication}
 F_{C_{r\to\ell}(N)}(z,w)=(1+zw)^{\ell-r}F_N(z,w).
\end{equation}
Adjunction gives the coefficients of $R_{\kappa\to r}$, with
$\ell$ replaced by $\kappa$.  If its input indices are $a+m,b+m$,
the two square-root binomial weights from the output and input of the
Choi kernel cancel exactly the two denominators in these formulas.
What remains is
\[
 \binom ra\binom rb\binom{\ell-r}{k}\binom{\kappa-r}{m},
\]
the coefficient in the expansion of \eqref{spin:basis-kernel}.

Covariance is equivalent to the Choi matrix commuting with
$U_\ell\otimes\overline{U_\kappa}$.  The Clebsch--Gordan decomposition
is multiplicity free, with spins
$|\ell-\kappa|/2,|\ell-\kappa|/2+1,\ldots,(\ell+\kappa)/2$.
Its commutant therefore has dimension $\min(\kappa,\ell)+1$.
The displayed covariant kernels are independent, since division by
$K_o^\ell K_i^\kappa$ gives the successive powers of $A/B$, whose
range is open by Lemma~\ref{spin:ratio}.  They consequently span the
entire complex commutant, without a Hermiticity assumption.

Finally $C_{r\to\ell}(I_{r+1})=I_{\ell+1}$: the projection onto the
symmetric sector of the first $r$ qubits acts identically on the fully
symmetric $\ell$-qubit sector.  Covariance and irreducibility imply that
$C_{r\to\ell}^*(I_{\ell+1})$ is a scalar multiple of $I_{r+1}$.
Taking traces identifies this scalar as $(\ell+1)/(r+1)$.
Composing with the trace-preserving $R_{\kappa\to r}$ proves
\eqref{spin:basis-effect}.
\end{proof}

\begin{proof}[Proof of Theorem~\ref{spin:covariant}]
For the independent product group, the commutant is the tensor product
of the local commutants in Lemma~\ref{spin:compression}.  There is
therefore a unique polynomial $Q$ of separate degrees at most
$\min(\kappa_i,\ell_i)$ such that
\begin{equation}\label{spin:global-invariant}
 F_{J(\Phi)}=
 \prod_iK_{o,i}^{\ell_i}K_{i,i}^{\kappa_i}\,
 Q\left(\frac{A_1}{B_1},\ldots,\frac{A_n}{B_n}\right).
\end{equation}
The separate degree bounds make this a polynomial.  The prefactor
never vanishes in the polydisk.  The variables of different sites are
independent, so Lemma~\ref{spin:ratio} shows that full coherent stability
is equivalent to nonvanishing of $Q$ on $\Omega^n$.
Lemma~\ref{spin:slit-product} gives
$Q=c\prod_iq_i$, where
$q_i(r)=\prod_{a=1}^{d_i}(r+s_{ia})$ and $s_{ia}\geq0$.
The kernel thus factors, up to one nonzero scalar, into the numerators
of \eqref{spin:covariant-kernel}.

For one factor write $q(r)=\sum_{a=0}^dq_ar^a$ and set
$\Psi_{\kappa,\ell,q}=\sum_{a=0}^dq_aT_a^{\kappa,\ell}$.
Every coefficient $q_a$ is nonnegative, so this map is completely
positive.  By \eqref{spin:basis-effect} its trace effect is
\begin{equation}\label{spin:normalization}
 \Psi_{\kappa,\ell,q}^*(I_{\ell+1})
 =c_\ell(q)I_{\kappa+1},\qquad
 c_\ell(q)=(\ell+1)\sum_{a=0}^d\frac{q_a}{a+1}
          =(\ell+1)\int_0^1q(x)\,dx>0.
\end{equation}
In the global factorization, trace preservation forces
$c\prod_ic_{\ell_i}(q_i)=1$.  Thus the initial scalar is positive
real and the map is the tensor product of the individually normalized
completely positive maps in \eqref{spin:covariant-kernel}.
Conversely these maps are covariant and trace preserving by construction,
and their kernels are stable because $q_i$ has no root in $\Omega$.
This proves the classification and automatic complete positivity.

To prove the rank assertion, first consider the square case $d=\kappa
=\ell$.  The map $\Psi_{d,d,q}$ is also the Dicke compression
\begin{equation}\label{spin:core-compression}
 \Psi_{d,d,q}(M)=W_d^*\left[
     \left(\bigotimes_{a=1}^d\mathcal T_{s_a}\right)(W_dMW_d^*)\right]W_d,
 \qquad \mathcal T_s(N)=N+s\Tr(N)I_2,
\end{equation}
where the tensor product in brackets acts on the $d$ qubit matrix
factors.  Its kernel is $\prod_a(A+s_aB)$ by diagonalization of the
binary variables, which verifies the identity with $\Psi_{d,d,q}$.
As an operator on Hilbert--Schmidt space, $\mathcal T_s$ is positive
definite: its eigenvalues are $1$ on traceless matrices and $1+2s$ on
the scalar space.  The embedding $M\mapsto W_dMW_d^*$ is a
Hilbert--Schmidt isometry.  Its compression of the positive definite
tensor product is therefore positive definite and invertible.  Define
\begin{equation}\label{spin:core}
 \Lambda_{d,q}=\frac{\Psi_{d,d,q}}{c_d(q)},\qquad
 \Lambda_{0,1}=\Id_{\C}.
\end{equation}

The multiplication identity \eqref{spin:kernel-multiplication} gives
$C_{d\to\ell}C_{a\to d}=C_{a\to\ell}$, and adjunction gives
$R_{d\to a}R_{\kappa\to d}=R_{\kappa\to a}$.  Hence
\begin{equation}\label{spin:raw-factorization}
 \Psi_{\kappa,\ell,q}
 =C_{d\to\ell}\Psi_{d,d,q}R_{\kappa\to d}.
\end{equation}
The map $C_{d\to\ell}$ is injective by
\eqref{spin:kernel-multiplication}; its adjoint $R_{\ell\to d}$ is
therefore surjective, and the same holds for $R_{\kappa\to d}$.
The middle map in \eqref{spin:raw-factorization} is invertible on a
space of dimension $(d+1)^2$.  This proves local rank $(d+1)^2$.
Ranks multiply under tensor products, which gives
\eqref{spin:covariant-rank} and all assertions about injectivity,
surjectivity, and invertibility.
\end{proof}

\begin{corollary}[Trace, invertible core, and symmetric cloning]
\label{spin:trace-core-cloner}
Each local map in Theorem~\ref{spin:covariant} has the exact factorization
\begin{equation}\label{spin:normalized-factorization}
 \Phi_i=\mathcal C_{d_i\to\ell_i}\circ\Lambda_{d_i,q_i}
                                      \circ R_{\kappa_i\to d_i},
 \qquad
 \mathcal C_{d\to\ell}=\frac{d+1}{\ell+1}C_{d\to\ell}.
\end{equation}
All three maps are trace preserving and completely positive, and the
middle map is an invertible coherent Lee--Yang map.
For equal input and output capacities the local factors are also unital.
\end{corollary}

\begin{proof}
The first two assertions follow from \eqref{spin:raw-factorization},
\eqref{spin:normalization}, and
$c_d(q)/c_\ell(q)=(d+1)/(\ell+1)$.  For equal capacities,
covariance makes $\Phi(I)$ scalar, and trace preservation makes this
scalar one.
\end{proof}

The normalized map $\mathcal C_{d\to\ell}$ is the universal symmetric
cloner of \cite[Eq.~(3.3)]{Werner98}; its dimension-ratio normalization
and concatenation law are classical.  Here stability forces the entire
factorization \eqref{spin:normalized-factorization}.  The integer $d$
specifies an intermediate spin dimension and the exact algebraic rank;
it is not an assertion of noiseless recovery or an operational quantum
communication capacity.  If a specified covariance relation permutes the
matched output sites, the theorem applies after undoing that permutation;
an invertible map then matches equal capacities.

\begin{example}[The spin-one spectral region]\label{spin:spin-one}
For $\kappa=\ell=2$, let $\lambda_1,\lambda_2$ be the multipliers of a
trace-preserving covariant map on the dipole and quadrupole subspaces.
The exact coherent Choi condition is
\begin{equation}\label{spin:spin-one-region}
 0\leq\lambda_2\leq\lambda_1,\qquad
 3\lambda_1\leq2+\lambda_2,\qquad
 3\lambda_1^2+\lambda_2^2\geq4\lambda_2.
\end{equation}
Indeed, in the degree-two case write $q(r)=r^2+ar+b$.
Its roots are real and nonpositive exactly when
$a,b\geq0$ and $a^2\geq4b$.  The maps with kernels $A^2$, $AB$,
and $B^2$ have scalar multipliers $1,3/2,3$, respectively, dipole
multipliers $1,1,0$, and quadrupole multipliers $1,0,0$.
For the middle map, partial trace annihilates the quadrupole
$\diag(1,-2,1)$, while the difference of the two extreme projectors is
sent to $S^z$; covariance determines its action on each irreducible
operator space.  Therefore
\[
 \lambda_1=\frac{1+a}{1+3a/2+3b},\qquad
 \lambda_2=\frac1{1+3a/2+3b}.
\]
Solving for $a,b$ gives \eqref{spin:spin-one-region} when
$\lambda_2>0$.  Degree one gives
$\lambda_2=0$ and $0<\lambda_1\leq2/3$, and degree zero gives the
origin, proving all endpoint cases.  In particular, the ordinary
qutrit depolarizing map with $\lambda_1=\lambda_2=t$ fails coherent
Choi stability for every $0<t<1$, although both endpoints are stable.
\end{example}

\begin{remark}\label{spin:scope}
Theorem~\ref{spin:covariant} assumes independent local covariance, not
only covariance under a common rotation.  Its proof is static and does
not use Theorem~\ref{spin:hamiltonian}.  It also does not apply the
invertible qubit classification to a polarized higher-spin map: a
Dicke lift is generally singular on the full qubit matrix space.
No unrestricted noncovariant higher-spin static classification is
asserted here.
\end{remark}

% END SECTION higher_capacity

% BEGIN SECTION global_semigroups
\section{Connected preservers and the obstruction to flow generation}
\label{glob:section}

The exact infinitesimal cone does not generate every connected static
preserver, even after taking limits of products.  We prove this on the
fixed multiaffine space
$V_n=\C[z_1,\ldots,z_n]_{\mathrm{ma}}\simeq(\C^2)^{\otimes n}$.
Let $S_n$ be the semigroup of invertible complex-linear operators that
preserve disk stability.  For its ordinary monomial coefficient matrix
$M$, write
\begin{equation}\label{glob:symbol}
 G_M(z,w)=\sum_{a,b\in\{0,1\}^n}M_{ab}z^aw^b.
\end{equation}
For \(M\in\operatorname{GL}(V_n)\), the symbol criterion is
\begin{equation}\label{glob:symbol-equivalence}
 M\in S_n\quad\Longleftrightarrow\quad
 G_M\text{ is stable on }\D^{2n}.
\end{equation}
For necessity, fix $w\in\D^n$ and apply $M$ to the stable polynomial
$\prod_i(1+w_iz_i)$.  Invertibility excludes the zero output.
Conversely, tensor contraction of $G_M$ with a stable input proves
preservation, again excluding zero by invertibility.  This is the
multiaffine disk-symbol convention of the preceding sections.
Hurwitz's theorem shows that $S_n$ is closed relative to
$\operatorname{GL}(V_n)$: an invertible limit cannot have identically
zero symbol.  Nonzero scalar multiplication preserves $S_n$.

Define
\begin{equation}\label{glob:semigroups}
 \begin{split}
 W_n&=\{A:e^{tA}\in S_n\text{ for every }t\geq0\},\\
 E_n&=\langle e^A:A\in W_n\rangle,\qquad
 T_n=\overline{E_n}^{\,\operatorname{GL}(V_n)},\qquad
 H_n=S_n\cap S_n^{-1}.
 \end{split}
\end{equation}
Here $E_n$ consists of finite products.  The closure $T_n$ is a
different semigroup.  The group $H_n$ is the group of reversible
preservers, and $H_n^0$ denotes its identity component.

\begin{theorem}[Global connectedness and a strict flow obstruction]
\label{glob:main}
For every $n\geq1$, $S_n$ is path connected.  Nevertheless,
\begin{equation}\label{glob:reversible-parts}
 E_n\cap H_n=H_n^0,
 \qquad T_n\cap T_n^{-1}=H_n^0.
\end{equation}
For $n\geq2$, every nontrivial coordinate permutation $\pi$ satisfies
\begin{equation}\label{glob:permutation-obstruction}
 \pi\in S_n\setminus T_n.
\end{equation}
Thus a positive-radius matrix neighborhood of $\pi$ contains no finite
product of admissible one-parameter flows, even with unbounded product
lengths and generator norms.
\end{theorem}

\begin{proof}[Path connectedness]
Let $D_rz^a=r^{|a|}z^a$, $0<r\leq1$.  It is an invertible preserver,
and
\begin{equation}\label{glob:sandwich}
 G_{D_rMD_r}(z,w)=G_M(rz,rw).
\end{equation}
For $M\in S_n$, its constant coefficient $M_{00}$ is nonzero.
A path of nonzero scalars normalizes it to one.  With this normalization,
$D_rMD_r$ converges as $r\downarrow0$ to the vacuum projector
$|0\rangle\langle0|$.

Consider the convex open set in the complex affine hyperplane $X_{00}=1$
given by
\begin{equation}\label{glob:coefficient-ball}
 \mathcal B=\left\{X:X_{00}=1,
        \ \sum_{(a,b)\ne(0,0)}|X_{ab}|<1\right\}.
\end{equation}
Every member has a kernel nonvanishing on the closed polydisk, by strict
constant-term domination.  For small positive $r,s$, both $D_rMD_r$
and $D_s$ belong to $\mathcal B$ and are invertible.  Indeed the
nonconstant coefficient sum for $D_s$ is $(1+s)^n-1$.

The invertible part of $\mathcal B$ is path connected.  To see this,
take invertible $X,Y\in\mathcal B$ and set
$L(\zeta)=(1-\zeta)X+\zeta Y$.  Its determinant is a nonzero
polynomial in $\zeta$, since it is nonzero at zero.  The real segment
$L([0,1])$ is compactly contained in $\mathcal B$, so a sufficiently
thin complex neighborhood of $[0,1]$ still maps into $\mathcal B$.
A path from $0$ to $1$ in this neighborhood avoids the finitely many
determinant zeros.  Its image connects $X$ to $Y$ through invertible
members of $\mathcal B$.
Now connect $M$ to $D_rMD_r$ using \eqref{glob:sandwich}, connect this
point to $D_s$ inside $\mathcal B$, and connect $D_s$ to $I$ by
increasing $s$ to one.  Every point is in $S_n$.  The singular vacuum
projector is used only as a limit to enter $\mathcal B$; it is never a
point of the path.
\end{proof}

For reference, there is also an explicit two-variable path.  On
$\operatorname{span}(z_1,z_2)$ let $P_\pm$ be the symmetric and
antisymmetric projections, and let $U_\theta$ act as
$P_++e^{i\theta}P_-$ there and as the identity on
$\operatorname{span}(1,z_1z_2)$.  For $r=1/4$,
\begin{align}\label{glob:swap-path}
 G_{D_rU_\theta}
 ={}&1+ra_\theta(z_1w_1+z_2w_2)
       +rb_\theta(z_1w_2+z_2w_1)
       +r^2z_1z_2w_1w_2,\notag\\
 a_\theta={}&(1+e^{i\theta})/2,\qquad
 b_\theta=(1-e^{i\theta})/2.
\end{align}
Its nonconstant coefficient sum is at most
$2\sqrt2r+r^2<1$, and its determinant is $r^4e^{i\theta}\ne0$.
The paths $I\to D_r$, $D_rU_\theta$ for $0\leq\theta\leq\pi$,
and $D_s\pi$ for $r\leq s\leq1$ connect the identity to the
coordinate exchange.  These paths are not asserted to have preserving
forward increments.

\begin{lemma}[Reversible factors]\label{glob:unit-factors}
If $A,B\in S_n$ and $AB\in H_n$, then $A,B\in H_n$.
Moreover every element of $H_n^0$ is a nonzero scalar times a tensor
product of one-site invertible matrices, whereas no nontrivial coordinate
permutation has this form.
\end{lemma}

\begin{proof}
The identities $A^{-1}=B(AB)^{-1}$ and $B^{-1}=(AB)^{-1}A$ prove
the first assertion.  The group $H_n$ is a closed subgroup of
$\operatorname{GL}(V_n)$ and hence a real Lie group.  Its Lie algebra
is $W_n\cap(-W_n)$.  By Corollary~\ref{gen:gram-cones}, in
half-plane coordinates this algebra is
\begin{equation}\label{glob:unit-algebra}
 \C I+\bigoplus_i\mathfrak{sl}_2(\R)_i.
\end{equation}
The fixed tensor Cayley transformation preserves products and
coordinate permutations.  A connected Lie group is generated by the
exponentials of its Lie algebra.  Exponentiating
\eqref{glob:unit-algebra} and taking finite products therefore gives
only scalar multiples of local tensor products.
A nontrivial coordinate permutation cannot have this form: varying one
input factor varies a different output factor, whereas an invertible
local tensor product varies only the corresponding output factor.
Equivalently, for a two-factor exchange, an identity
$Ax\otimes By=c\,y\otimes x$ with fixed nonzero $x$ and varying
independent $y$ would force the fixed vector $Ax$ to be proportional
to every $y$.
\end{proof}

\begin{proof}[The exact-product assertion in Theorem~\ref{glob:main}]
If $e^{A_m}\cdots e^{A_1}\in H_n$, with every $A_j\in W_n$,
Lemma~\ref{glob:unit-factors} makes every exponential factor a unit.
For $0\leq s\leq1$, factor
$e^{A_j}=e^{(1-s)A_j}e^{sA_j}$.  Both factors preserve stability,
so the same lemma gives $e^{sA_j}\in H_n$.  Thus each factor belongs
to $H_n^0$, and so does their product.  Conversely $H_n^0$ is
generated by exponentials of its Lie algebra, which is contained in
$W_n$.  This proves $E_n\cap H_n=H_n^0$.
\end{proof}

For the closure assertion we use Neeb's classical connected-unit theorem:
if $G$ is a connected Lie group and $T\subseteq G$ is a closed
submonoid satisfying
\[
 T=\overline{\langle\exp L(T)\rangle}^{\,G},\qquad
 L(T)=\{A:\exp(tA)\in T\text{ for every }t\geq0\},
\]
then $T\cap T^{-1}$ is connected. This is
\cite[Definition~II.1(3) and Proposition~III.2, pp.~171, 178]{Neeb92};
see also \cite[Section~6.1]{SanMartinSantana02}.

\begin{proof}[The closure assertion in Theorem~\ref{glob:main}]
The ambient real Lie group $G=\operatorname{GL}_{2^n}(\C)$ is connected.
By definition $T_n$ is a closed submonoid and $T_n\subseteq S_n$.
Every $A\in W_n$ generates a one-parameter semigroup in $E_n$, while
$L(T_n)\subseteq L(S_n)=W_n$. Hence
\[
 L(T_n)=W_n,\qquad
 T_n=\overline{\langle\exp L(T_n)\rangle}^{\,G}.
\]
Neeb's theorem therefore makes $T_n\cap T_n^{-1}$ a connected
subgroup of $H_n$, so it is contained in $H_n^0$.
The reverse inclusion follows because the exponentials of both signs of
\eqref{glob:unit-algebra} belong to $T_n$ and generate $H_n^0$.

If a nontrivial coordinate permutation $\pi$ belonged to $T_n$, its
finite order would imply $\pi^{-1}\in T_n$.  It would therefore be
a unit of $T_n$, and hence belong to $H_n^0$, contrary to
Lemma~\ref{glob:unit-factors}.  This proves
\eqref{glob:permutation-obstruction}.  Relative closure in the general
linear group and ordinary matrix-norm closure agree at the invertible
point $\pi$, proving the neighborhood statement.
\end{proof}

\begin{corollary}[Exact local divisibility]\label{glob:local-divisibility}
Call $M\in S_n$ exactly locally divisible if, for every identity
neighborhood $U$ in $\operatorname{GL}(V_n)$, it is a finite product
of members of $S_n\cap U$.  The exactly locally divisible units of
$S_n$ are precisely $H_n^0$.  In addition, if a continuous path
$M:[0,1]\to S_n$ starts at $I$ and satisfies
$M(t)M(s)^{-1}\in S_n$ for $s\leq t$, then a unit endpoint must
belong to $H_n^0$.
\end{corollary}

\begin{proof}
The identity component of a Lie group is open, so choose $U$ with
$H_n\cap U\subset H_n^0$.  Any exact factorization of a unit into
$S_n\cap U$ consists entirely of units by
Lemma~\ref{glob:unit-factors}, hence has endpoint in $H_n^0$.
Conversely, subdividing a finite exponential representation of a member
of $H_n^0$ puts all its factors in any prescribed identity neighborhood.
For the path assertion, write
$M(1)=[M(1)M(t)^{-1}]M(t)$.  A unit endpoint makes $M(t)$ a unit
for every $t$, so the continuous path stays in $H_n$ and ends in
$H_n^0$.
\end{proof}

\begin{remark}\label{glob:scope}
The identity for the closure in \eqref{glob:reversible-parts} is
$T_n\cap T_n^{-1}=H_n^0$.  It does not assert
$T_n\cap H_n=H_n^0$ for arbitrary units of the larger semigroup;
finite order supplies the inverse inclusion needed for permutations.
The result concerns flows of the infinitesimal cone.  Static apolar maps
may lie on other branches: for $F(s,t)=s-t$,
$B_F=\pi-I$, so $I+B_F=\pi$, whereas
\[
 e^{tB_F}=\frac{1+e^{-2t}}2I+\frac{1-e^{-2t}}2\pi.
\]
Allowing such a static permutation as an additional primitive changes
the factorization question.  No description of the remaining
noninvertible boundary or the nonunit part of $T_n$ is claimed.
\end{remark}

% END SECTION global_semigroups

% BEGIN SECTION robustness
\section{Quantitative rigidity and correlated tests}
\label{rob:section}

The finite-test rigidity of Theorem~\ref{dyn:trine} is an exact statement.
Its quantitative behaviour depends on the allowed inputs.  We prove that
products of two GHZ states detect precisely the failure of a diagonal
Hamiltonian to have additive energies.  The approximate-additivity theorem
of Kalton and Roberts then gives a dimension-independent comparison with
distance from one-site dynamics.  In contrast, even all product Lee--Yang
inputs admit no such comparison.

Throughout this section, a density matrix acts on
$\mathcal H_n=(\C^2)^{\otimes n}$, and
\begin{align*}
 F_X(z,w)&=\sum_{x,y\in\{0,1\}^n}X_{xy}z^xw^y,\\
 \mathcal Y_n&=\{\rho\geq0:\Tr\rho=1,\ 
                 F_\rho\ne0\text{ on }\D^{2n}\}.
\end{align*}
The condition is zero-freeness, not merely that $F_\rho$ is a nonzero
polynomial.  We use the unhalved trace norm $\|\cdot\|_1$ and diamond norm
$\|\cdot\|_\diamond$, and write
$\operatorname{dist}_1(X,\mathcal Y_n)=
\inf_{\rho\in\mathcal Y_n}\|X-\rho\|_1$.
Let $\operatorname{LocHP}_n$ be the real vector space of maps
\[
 \mathcal M=\sum_{i=1}^n
       \mathcal M_i\otimes\Id_{\{1,\ldots,n\}\setminus\{i\}},
\]
where every $\mathcal M_i$ is complex linear and Hermiticity preserving.
Let $\operatorname{LocGKLS}_n$ be the cone obtained by requiring each
$\mathcal M_i$ to be a one-qubit GKLS generator.  In particular, local
Hamiltonian generators belong to both classes.

\subsection{The diagonal quotient norm}

Identify a subset $S\subseteq[n]$ with the computational basis vector
whose occupied sites are exactly $S$.  For a real function $h$ on $2^{[n]}$,
put
\begin{gather*}
 H=\diag(h(S):S\subseteq[n]),\qquad
 \mathcal L_H(X)=-i[H,X],\\
 \operatorname{osc}h=\max_S h(S)-\min_S h(S).
\end{gather*}

\begin{proposition}
\label{rob:distance}
For every diagonal Hermitian $H$,
\begin{align}
 d(H)
 &:=
 \operatorname{dist}_\diamond(\mathcal L_H,\operatorname{LocHP}_n)
 \notag\\
 &=\operatorname{dist}_\diamond(\mathcal L_H,\operatorname{LocGKLS}_n)
 \notag\\
 &=\min_{b\in\R^n}\operatorname{osc}
       \left(h(S)-\sum_{i\in S}b_i\right).
 \label{rob:distance-formula}
\end{align}
If $h(S)=h([n]\setminus S)$ for all $S$, then
$d(H)=\operatorname{osc}h=\|\mathcal L_H\|_\diamond$.
\end{proposition}

\begin{proof}
Fix $\mathcal M\in\operatorname{LocHP}_n$ and set
$E_{xy}=|x\rangle\langle y|$.  The coefficient of $E_{xy}$ in
$\mathcal M(E_{xy})$ equals
\[
 \langle x|\mathcal M(E_{xy})|y\rangle
 =\sum_i\langle x_i|\mathcal M_i(E_{x_i y_i})|y_i\rangle.
\]
If $x_i=y_i$, the summand is real.  The summands for $(x_i,y_i)=(0,1)$
and $(1,0)$ are conjugates.  Thus, for a vector $b\in\R^n$ independent
of $x,y$, its imaginary part is $b\cdot(y-x)$.  Since
$\|E_{xy}\|_1=1$ and every matrix entry is bounded by the trace norm,
\[
 \|\mathcal L_H-\mathcal M\|_\diamond
 \geq \max_{x,y}
       |(h(x)-b\cdot x)-(h(y)-b\cdot y)|
 =\operatorname{osc}(h-b\cdot x).
\]
Here and below $h(x)$ is the subset notation written in occupation
coordinates.

For a Hermitian matrix $D$ one has
\begin{equation}
 \|-i[D,\cdot]\|_\diamond
   =\lambda_{\max}(D)-\lambda_{\min}(D).
 \label{rob:commutator-norm}
\end{equation}
Indeed, subtract the midpoint of the spectral interval from $D$.
The trace-norm ideal inequality, also after adjoining an identity
ancilla, gives the upper bound $2\|D\|_{\mathrm{op}}$ for the resulting
commutator.  The matrix unit between eigenvectors at the two spectral
endpoints gives the reverse bound.  Apply this identity to
$D=H-H_b$, where
$H_b=\sum_i b_i|1\rangle\langle1|_i$.
The local Hamiltonian generator $\mathcal L_{H_b}$ attains the
corresponding oscillation.  Because
$\operatorname{LocGKLS}_n\subseteq\operatorname{LocHP}_n$,
the two distance infima agree with the infimum of these oscillations.
It is a minimum: the objective is continuous and
\[
 \operatorname{osc}(h-b\cdot x)
 \geq \operatorname{osc}(b\cdot x)-\operatorname{osc}h
 =\|b\|_{\ell^1}-\operatorname{osc}h.
\]

Suppose now that $h$ is invariant under complementation.  Replacing
$x$ by $1-x$ shows that
$\operatorname{osc}(h-b\cdot x)=\operatorname{osc}(h+b\cdot x)$.
The oscillation is a convex function of its argument, so
\[
 \operatorname{osc}h
 \leq \tfrac12\operatorname{osc}(h-b\cdot x)
       +\tfrac12\operatorname{osc}(h+b\cdot x)
 =\operatorname{osc}(h-b\cdot x).
\]
Taking $b=0$ proves the last assertion.
\end{proof}

\subsection{Two correlated blocks}

For disjoint nonempty $A,B\subseteq[n]$, define
\begin{equation}
\begin{aligned}
 |\psi_{A,B}\rangle
 &=\frac{|0_A\rangle+|1_A\rangle}{\sqrt2}
  \otimes\frac{|0_B\rangle+|1_B\rangle}{\sqrt2}
  \otimes|0_{[n]\setminus(A\cup B)}\rangle,\\
 \rho_{A,B}&=|\psi_{A,B}\rangle\langle\psi_{A,B}|.
\end{aligned}
 \label{rob:ghz-probe}
\end{equation}
The pure-state coefficient polynomial is
$\frac12(1+\prod_{i\in A}z_i)(1+\prod_{i\in B}z_i)$.
It is zero-free on $\D^n$, and hence $\rho_{A,B}\in\mathcal Y_n$.
Set
\begin{align*}
 \Delta_{A,B}(H)&=h(A\cup B)-h(A)-h(B)+h(\varnothing),\\
 \Delta(H)&=\max_{A\cap B=\varnothing}|\Delta_{A,B}(H)|.
\end{align*}

\begin{theorem}[Correlated diagonal rigidity]
\label{rob:diagonal}
Let $n\geq2$ and let $H$ be a diagonal Hermitian Hamiltonian.  For every
pair of disjoint nonempty subsets $A,B$ the following limit exists:
\begin{equation}
 \lim_{t\downarrow0}
 \frac{\operatorname{dist}_1
       (e^{-itH}\rho_{A,B}e^{itH},\mathcal Y_n)}{t}
 =\frac{|\Delta_{A,B}(H)|}{2}.
 \label{rob:exact-defect}
\end{equation}
If $\delta(H)$ denotes the maximum of these limits, then
\begin{equation}
 \delta(H)=\frac{\Delta(H)}2,
 \qquad
 \delta(H)\leq d(H)\leq4K\,\delta(H),
 \label{rob:diagonal-bounds}
\end{equation}
where $K$ is any universal constant in the real Kalton--Roberts
approximate-additivity theorem.  In particular, one may take $K=45$.
\end{theorem}

We first establish the polynomial and boundary facts used in the proof.

\begin{lemma}[Block compression]
\label{rob:compression}
Let $f(u_1,\ldots,u_k,v)$ be multiaffine and zero-free on the open unit
polydisk.  If $a(v)$ and $d(v)$ are its constant and full-degree
coefficients in the $u$ variables, then $a(v)+d(v)s$ is zero-free on
the open unit polydisk in $(s,v)$.

Consequently, if
$V:\C^2\otimes\C^2\longrightarrow\mathcal H_n$ is the isometry
\[
 V|a,b\rangle=|a_A,b_B,0_{[n]\setminus(A\cup B)}\rangle,
\]
then $F_{V^*\sigma V}$ is zero-free on $\D^4$ for every
$\sigma\in\mathcal Y_n$.
\end{lemma}

\begin{proof}
Fix the remaining variables $v$ in their polydisk.  The polynomial
$g(u)=f(u,\ldots,u,v)$ has no zero in $\D$, and $a(v)=g(0)\ne0$.
If $\deg g=k$, all its roots have modulus at least one, and the
product-of-roots identity gives $|d(v)|\leq|a(v)|$.
If $\deg g<k$, then $d(v)=0$.  Thus $a(v)+d(v)s\ne0$ for $|s|<1$.
This proves the first assertion; it is the unit-disk form of the
Asano coefficient contraction \cite{AsanoJapanese70}.

For the second assertion, first set all row and column variables
outside $A\cup B$ to zero.  Apply the first assertion successively to
the row variables in $A$, the row variables in $B$, and the corresponding
two groups of column variables.  The resulting coefficient polynomial
is exactly $F_{V^*\sigma V}$.  No normalization or zero-polynomial
exception occurs in this argument.
\end{proof}

The compression in Lemma~\ref{rob:compression} is completely positive
and trace nonincreasing, but need not preserve trace.  We use its
unnormalized output.  For every Hermitian $X$,
$\|V^*XV\|_1\leq\|X\|_1$, as follows either from trace-norm duality
or from the positive and negative parts of $X$.

\begin{lemma}[A two-qubit boundary contact]
\label{rob:contact}
Write $\rho_0=|++\rangle\langle++|$, where
$|+\rangle=(|0\rangle+|1\rangle)/\sqrt2$, and put
\begin{align*}
 u&=|00\rangle+|10\rangle+|01\rangle+|11\rangle,\\
 v&=|00\rangle-|10\rangle-|01\rangle+|11\rangle,\qquad
 \ell(X)=\Im\langle v|X|u\rangle.
\end{align*}
If $t_j\downarrow0$, $X_j$ is Hermitian, $F_{X_j}$ is zero-free on
$\D^4$, and $\|X_j-\rho_0\|_1=O(t_j)$, then
$\ell(X_j)=o(t_j)$.  Moreover,
\begin{equation}
 |\ell(X)|\leq2\|X\|_1
 \quad\text{for every Hermitian }X.
 \label{rob:functional-norm}
\end{equation}
\end{lemma}

\begin{proof}
The Hermitian observable representing $\ell$ is
\[
 R=\frac{|u\rangle\langle v|-|v\rangle\langle u|}{2i}.
\]
The vectors $u,v$ are orthogonal and both have norm two.  On their span
$R$ has eigenvalues $2,-2$, and it vanishes on the orthogonal
complement.  Thus $\|R\|_{\mathrm{op}}=2$, proving
\eqref{rob:functional-norm}.

Consider
\[
 g_j(s)=F_{X_j}(s-1,s-1,1,1)
       =a_{2,j}s^2+a_{1,j}s+a_{0,j}.
\]
For $X_j=\rho_0$ this polynomial is $s^2$.  Continuity of its
coefficients gives $a_{2,j}=1+O(t_j)$ and
$a_{1,j},a_{0,j}=O(t_j)$.  The column variables have been specialized
to boundary points, so we justify their use.  For fixed $j$, the
polynomials $F_{X_j}(s-1,s-1,r,r)$, with $0<r<1$, are zero-free on
$|s-1|<1$.  Hurwitz's theorem shows that their limit as $r\uparrow1$
is zero-free there or identically zero.  For all sufficiently large
$j$, $a_{2,j}\ne0$, which excludes the latter alternative.

Let $r_{1,j},r_{2,j}$ be the two roots of $g_j$, counted with
multiplicity.  Vieta's formulas give their sum and product as $O(t_j)$;
the quadratic formula then gives $|r_{k,j}|=O(\sqrt{t_j})$.
Because $|r_{k,j}-1|\geq1$,
\[
 \Re r_{k,j}\leq\tfrac12|r_{k,j}|^2=O(t_j).
\]
Their real parts have sum $O(t_j)$, so the same real parts are bounded
below by $-O(t_j)$.  Consequently
$\Re r_{k,j}=O(t_j)$ and $\Im r_{k,j}=O(\sqrt{t_j})$.
It follows that
\[
 \Im(r_{1,j}r_{2,j})=O(t_j^{3/2}),\qquad
 \Im a_{0,j}=O(t_j^{3/2})+O(t_j^2)=o(t_j).
\]
Finally $a_{0,j}=F_{X_j}(-1,-1,1,1)=\langle v|X_j|u\rangle$.
\end{proof}

\begin{proof}[Proof of Theorem~\ref{rob:diagonal}]
Fix $A,B$, let $V$ be as in Lemma~\ref{rob:compression}, and abbreviate
$\Delta=\Delta_{A,B}(H)$ and
$\rho(t)=e^{-itH}\rho_{A,B}e^{itH}$.
The image of $V$ is invariant under $H$, and
$V^*\rho(t)V=e^{-itD}\rho_0e^{itD}$, where the logical diagonal
Hamiltonian $D$ has energies
$h(\varnothing),h(A),h(B),h(A\cup B)$.
Since $\rho_0=|u\rangle\langle u|/4$ and $\langle v|u\rangle=0$,
differentiation gives
\begin{equation}
 \ell(V^*\rho(t)V)=-\Delta t+O(t^2).
 \label{rob:contact-derivative}
\end{equation}

Let $t_j\downarrow0$ be arbitrary, and choose
$\sigma_j\in\mathcal Y_n$ with
\[
 \|\rho(t_j)-\sigma_j\|_1
 \leq\operatorname{dist}_1(\rho(t_j),\mathcal Y_n)+t_j^2.
\]
Because $\rho_{A,B}\in\mathcal Y_n$ and
$\rho(t_j)=\rho_{A,B}+O(t_j)$, both the distance on the right and
$\|\sigma_j-\rho_{A,B}\|_1$ are $O(t_j)$.
Lemma~\ref{rob:compression} shows that
$X_j=V^*\sigma_j V$ has a zero-free coefficient polynomial.
It is Hermitian and $X_j=\rho_0+O(t_j)$, so
Lemma~\ref{rob:contact} applies without rescaling its trace.
Equations~\eqref{rob:functional-norm} and
\eqref{rob:contact-derivative} imply
\[
 2\|\rho(t_j)-\sigma_j\|_1
 \geq |\ell(V^*\rho(t_j)V-X_j)|
 =|\Delta|t_j+o(t_j).
\]
As the sequence was arbitrary, this proves the lower bound
$|\Delta|/2$ for the liminf in \eqref{rob:exact-defect}.

For a matching upper bound, subtract $(\Delta/4)Z_1Z_2$ from $D$,
where $Z=\diag(1,-1)$.  The remaining logical energy is affine in the
two occupation bits: its mixed difference is zero.  Its evolution of
the logical vector $|++\rangle$, embedded by $V$, is therefore a
product of two GHZ states with unit-modulus relative phases and a
vacuum on the remaining sites.  This comparison state belongs to
$\mathcal Y_n$.  Since all four logical configurations have equal
probability and the two values of $Z_1Z_2$ occur equally often, the
inner product between the actual and comparison pure vectors is
$\cos(t\Delta/4)$.  Their trace distance is
\[
 2|\sin(t\Delta/4)|\leq t|\Delta|/2\qquad(t\geq0).
\]
This proves \eqref{rob:exact-defect}.  Empty subsets have mixed
difference zero, and the family of pairs is finite, so
$\delta(H)=\Delta(H)/2$.

We now invoke the precise inherited input.  The Kalton--Roberts
theorem \cite[Theorem~4.1]{KaltonRoberts83} supplies a universal
$K\leq45$ with the following property: if a real set function $f$ on
a set algebra satisfies $f(\varnothing)=0$ and
\[
 |f(A\cup B)-f(A)-f(B)|\leq\eta
 \quad(A\cap B=\varnothing),
\]
then there is an additive real set function $\mu$ such that
$\sup_S|f(S)-\mu(S)|\leq K\eta$.
Apply it on $2^{[n]}$ to $f(S)=h(S)-h(\varnothing)$ and
$\eta=\Delta(H)$.  On this finite algebra,
$\mu(S)=\sum_{i\in S}b_i$ for some real $b_i$, and hence
\[
 \operatorname{osc}(h-b\cdot x)\leq2K\Delta(H).
\]
If $\Delta(H)=0$, the same conclusion follows directly from finite
additivity.  Proposition~\ref{rob:distance} now yields
$d(H)\leq2K\Delta(H)=4K\delta(H)$.

Finally choose a minimizer $b$ in \eqref{rob:distance-formula}.
The flow of $H_b$ consists of independent diagonal one-qubit
unitaries.  It rotates each polynomial variable by a unit complex
number and thus preserves $\mathcal Y_n$.
The Duhamel formula and unitary trace-norm invariance imply, for every
density matrix $\rho$ and $t\geq0$,
\[
 \|e^{t\mathcal L_H}\rho-e^{t\mathcal L_{H_b}}\rho\|_1
 \leq t\|\mathcal L_H-\mathcal L_{H_b}\|_\diamond=t\,d(H).
\]
Apply this estimate to each $\rho_{A,B}$, divide by $t$, and use
\eqref{rob:exact-defect}.  Taking the maximum gives
$\delta(H)\leq d(H)$.
\end{proof}

\begin{remark}
\label{rob:scope}
The universal constant in \eqref{rob:diagonal-bounds} is inherited
from Kalton--Roberts.  The bridge to that theorem is the exact
root-contact identity \eqref{rob:exact-defect}.
The probes may contain arbitrarily large blocks, their family is
exponential in $n$, and the defect uses distance to the full set
$\mathcal Y_n$.  No efficient certification procedure or corresponding
statement for non-diagonal Hamiltonians is asserted.
\end{remark}

\subsection{Product inputs do not give uniform robustness}

Let $\mathcal P_n\subseteq\mathcal Y_n$ be the set of all product
density matrices in $\mathcal Y_n$, including mixed products, and put
\[
 \delta_{\mathrm{prod}}(\mathcal L)
 =\sup_{\rho\in\mathcal P_n}\limsup_{t\downarrow0}
   \frac{\operatorname{dist}_1(e^{t\mathcal L}\rho,\mathcal Y_n)}{t}.
\]
We give two different mechanisms by which this defect can be small.
For use with mixed states, recall the elementary estimate
\begin{equation}
 \|e^{-itH}\rho e^{itH}-\rho\|_1
 \leq2|t|\sqrt{\Tr(\rho H^2)-(\Tr(\rho H))^2}.
 \label{rob:variance-bound}
\end{equation}
To prove it, purify $\rho$ to a unit vector $\Psi$ and evolve the
purification by $H\otimes I$.  For a pure-state projection
$P=|\Psi\rangle\langle\Psi|$, the two possibly nonzero eigenvalues
of $-i[H\otimes I,P]$ are
$\pm\sqrt{\operatorname{Var}_\rho H}$.
Its trace norm is therefore twice this square root.
The variance is constant along the unitary orbit, so integration gives
the same bound for purified states.  Partial trace gives
\eqref{rob:variance-bound}.

\begin{proposition}[A nearest-neighbour obstruction]
\label{rob:path}
On a path with $m\geq1$ edges and $n=m+1$ vertices, let
\[
 H_m=\frac1{2m}\sum_{j=1}^m Z_jZ_{j+1},
 \qquad \mathcal L_m=-i[H_m,\cdot].
\]
Then
\begin{gather*}
 \|\mathcal L_m\|_\diamond=d(H_m)=1,\\
 \operatorname{dist}_1(e^{t\mathcal L_m}\rho,\mathcal Y_n)
 \leq\sqrt{\frac3m}|t|
 \quad(\rho\in\mathcal P_n,\ t\in\R).
\end{gather*}
For the equatorial pure product
$|\psi_\theta\rangle=\bigotimes_{j=1}^n
(|0\rangle+e^{i\theta_j}|1\rangle)/\sqrt2$ and its density
$\rho_\theta$, the sharper formula is
\begin{equation}
 \|e^{t\mathcal L_m}\rho_\theta-\rho_\theta\|_1
 =2\sqrt{1-\cos(t/(2m))^{2m}}
 \leq\frac{|t|}{\sqrt m}.
 \label{rob:path-exact}
\end{equation}
\end{proposition}

\begin{proof}
The diagonal energy is invariant under simultaneous complementation.
Equal computational spins have energy $1/2$, alternating spins have
energy $-1/2$, and these are the extrema.  Thus
Proposition~\ref{rob:distance} gives the two norm identities.

For a product density, the computational spins are independent signs.
Write $W_j=Z_jZ_{j+1}$ for their edge products, regarded as random
variables in this diagonal distribution.  Disjoint edges have zero
covariance, each variance is at most one, and each adjacent covariance
has absolute value at most one by Cauchy--Schwarz.  Hence
\[
 \operatorname{Var}\left(\sum_{j=1}^m W_j\right)
 \leq m+2(m-1)\leq3m.
\]
Equation~\eqref{rob:variance-bound} gives the stated mixed-state bound.
The initial state itself belongs to $\mathcal Y_n$ and is an
admissible comparison state in the distance.

For $\psi_\theta$, the computational distribution is uniform.
The root sign and the $m$ edge products determine all vertex signs
bijectively, so the edge products are independent uniform signs.
It follows that
\[
 \langle\psi_\theta|e^{-itH_m}|\psi_\theta\rangle
 =\cos(t/(2m))^m.
\]
The trace distance of unit pure states $\psi,\phi$ is
$2\sqrt{1-|\langle\psi,\phi\rangle|^2}$; this follows by restricting
the difference of their projections to their two-dimensional span.
It proves the equality in \eqref{rob:path-exact}.
Independence gives $\operatorname{Var}_{\rho_\theta}H_m=1/(4m)$,
so \eqref{rob:variance-bound} proves its last inequality.
Finally, $\rho_\theta\in\mathcal Y_n$, since its coefficient polynomial
is a product of linear factors with roots on the unit circle.
\end{proof}

\begin{lemma}[One-qubit physical states]
\label{rob:onequbit}
A one-qubit density matrix
$\tau=\left(\begin{smallmatrix}a&\overline b\\ b&d\end{smallmatrix}\right)$
belongs to $\mathcal Y_1$ if and only if $a\geq d$.
\end{lemma}

\begin{proof}
Its polynomial is
$f(z,w)=a+bz+\overline b w+dzw$.
Suppose $a\geq d$.  Positivity gives $|b|^2\leq ad\leq a^2$,
and $a>0$ because $a+d=1$.  The denominator of
\[
 q(z)=\frac{\overline b+dz}{a+bz}
\]
has no zero in $\D$.  For $|z|=1$,
\[
 |a+bz|^2-|\overline b+dz|^2
 =(a-d)(a+d+2\Re(bz))\geq0.
\]
If $|b|<a$, the maximum principle gives $|q(z)|\leq1$ in $\D$.
The remaining case $|b|=a$ forces $a=d=|b|$, and cancellation makes
$q$ a constant of modulus one.  Thus in either case
$f(z,w)=(a+bz)(1+wq(z))$ is nonzero for $|z|,|w|<1$.

Conversely, suppose $a<d$.  Then
$|b|\leq\sqrt{ad}<(a+d)/2$, so the displayed boundary difference is
strictly negative on the unit circle.  By continuity, some $z$ inside
the circle satisfies
$|a+bz|<|\overline b+dz|$.
The denominator in
$w=-(a+bz)/(\overline b+dz)$ is nonzero and $|w|<1$;
this gives an interior zero of $f$.
\end{proof}

\begin{proposition}[An exponential product-input obstruction]
\label{rob:rankone}
For $n\geq2$, put $P_n=|1^n\rangle\langle1^n|$ and
$\mathcal K_n=-i[P_n,\cdot]$.  Then
\begin{equation}
 \|\mathcal K_n\|_\diamond=1,\qquad
 d(P_n)=\frac{n-1}{n}.
 \label{rob:rankone-distance}
\end{equation}
For every $\rho\in\mathcal P_n$ and every real $t$,
\begin{equation}
 \operatorname{dist}_1(e^{t\mathcal K_n}\rho,\mathcal Y_n)
 \leq\|e^{t\mathcal K_n}\rho-\rho\|_1
 \leq2^{1-n/2}|t|.
 \label{rob:rankone-bound}
\end{equation}
For equatorial pure products, the middle expression is exactly
\[
 4\sqrt{2^{-n}(1-2^{-n})}\,|\sin(t/2)|.
\]
\end{proposition}

\begin{proof}
The commutator norm is one by \eqref{rob:commutator-norm}.
To evaluate \eqref{rob:distance-formula}, average $b$ over coordinate
permutations.  Convexity of oscillation and permutation symmetry of
$h(x)=\mathbf1_{\{x=1^n\}}$ show that this cannot increase the objective.
We may therefore take $b_i=c$ for all $i$.  The residual energies are
\[
 \{-ck:0\leq k\leq n-1\}\ \cup\ \{1-cn\}.
\]
If $c\leq1/n$, the last value minus the value at weight $n-1$ is
$1-c\geq(n-1)/n$.  If $c\geq1/n$, the difference between weights
zero and $n-1$ has absolute value $c(n-1)\geq(n-1)/n$.
For $c=1/n$, all residual energies lie in
$[-(n-1)/n,0]$ and both endpoints occur.  This proves
\eqref{rob:rankone-distance}.

The coefficient polynomial of a product density is the product of
the individual coefficient polynomials in disjoint variables.  Thus
it is zero-free precisely when all factors are zero-free.
Lemma~\ref{rob:onequbit} implies that every occupied population is at
most $1/2$.  Consequently $p=\Tr(P_n\rho)\leq2^{-n}$.
For a purification $\Psi$ of $\rho$,
\[
 \langle\Psi|(e^{-itP_n}\otimes I)|\Psi\rangle
 =1-p+pe^{-it}.
\]
The pure-state trace-distance formula and partial-trace contractivity
therefore give
\[
 \|e^{t\mathcal K_n}\rho-\rho\|_1
 \leq4\sqrt{p(1-p)}\,|\sin(t/2)|
 \leq2\sqrt p\,|t|.
\]
This proves \eqref{rob:rankone-bound}, using the initial LY state as
comparison.  For an equatorial pure product $p=2^{-n}$, and no partial
trace is needed, giving the asserted equality.
\end{proof}

\begin{corollary}
\label{rob:product-obstruction}
There is no function $\omega:[0,\infty)\to[0,\infty)$ with
$\omega(0)=0$ and $\lim_{s\downarrow0}\omega(s)=0$, independent of $n$,
such that
\[
 \operatorname{dist}_\diamond
       (\mathcal L,\operatorname{LocGKLS}_n)
 \leq\omega(\delta_{\mathrm{prod}}(\mathcal L))
\]
for every Hamiltonian generator $\mathcal L$ of diamond norm one.
This remains false if the Hamiltonian is required to be a sum of
nearest-neighbour two-site terms on a path.
\end{corollary}

\begin{proof}
Proposition~\ref{rob:path} gives a sequence with distance one and
$0\leq\delta_{\mathrm{prod}}(\mathcal L_m)\leq\sqrt{3/m}\to0$,
contradicting the required behaviour of $\omega$.
Proposition~\ref{rob:rankone} separately gives exponentially small
product defect and distance tending to one.
\end{proof}

The normalization of the path interactions is by their total strength;
no positive lower bound on each edge coefficient is assumed.  These
examples concern a global diamond norm and do not obstruct estimation
of individual interaction coefficients in other access models.
Nor do they contradict Theorem~\ref{rob:diagonal}: correlated inputs
can detect what all product inputs detect only weakly.  For instance,
the GHZ state on all $n$ sites belongs to $\mathcal Y_n$ and has
$\Tr(P_n\rho)=1/2$, so the product-population bound used above does not
extend to arbitrary physical Lee--Yang states.

% END SECTION robustness

% BEGIN SECTION counterexamples
\section{An exact counterexample to a stronger ground-state radius}\label{cx:section}

The qualitative strict radius in Theorem~\ref{eq:strict-field} cannot
in general be replaced by the proposed explicit radius for deformed
EPR ground states. We address the stronger clause of Conjecture~1 in
\cite[Section 5.4]{WBG26}, with its stated unweighted-graph hypothesis.

\begin{theorem}\label{cx:epr}
Let $G$ have vertices $\{0,1,2,3,4,5\}$ and edges
\[
 \{01,02,03,14,15\}.
\]
Put
\[
 H=-\sum_{ij\in E(G)}
 (|00\rangle+\tfrac12|11\rangle)
 (\langle00|+\tfrac12\langle11|)_{ij}.
\]
The ground state $\psi$ is unique up to scalar. Its multiaffine
polynomial $f_\psi(z)=\sum_b\psi_bz^b$ has a zero satisfying
\[
 |z_i|<\frac{499}{500}\sqrt2\qquad(0\leq i\leq5).
\]
Thus the proposed radius $s^{-1/2}$ fails at $s=1/2$ on an
unweighted six-vertex tree.
\end{theorem}

\begin{proof}
We give rational spectral and zero certificates. All inequalities
below are verified by the finite procedure at the end of this
appendix, using only integer arithmetic and exact fractions.

Write $K=-H$. Index the computational basis by integers
$b\in\{0,\ldots,63\}$, with bit $b_i$ at vertex $i$.
For each edge $ij$ and each input with $b_i=b_j=a$, changing both
bits to $c\in\{0,1\}$ contributes $2^{-(a+c)}$ to the corresponding
entry of $K$. All other entries from that edge are zero.
This specifies the symmetric rational matrix $K$ completely.

Consider $M=37I-8K$. The charge
\[
 Q(b)=b_0+b_4+b_5-b_1-b_2-b_3
\]
is conserved. In charge order $-3,-2,\ldots,3$, the block dimensions are
\[
 1,\ 6,\ 15,\ 20,\ 15,\ 6,\ 1.
\]
Order each block by increasing integer $b$.
Successive rational symmetric elimination has pivot
signs as follows:
\[
\begin{array}{c|rrrrrrr}
 Q&-3&-2&-1&0&1&2&3\\ \hline
 \text{negative pivots}&0&0&0&1&0&0&0\\
 \text{positive pivots}&1&6&15&19&15&6&1.
\end{array}
\]
No pivot is zero; the negative pivot is the first in the charge-zero
block. For clarity, the recurrence at step $i$ is
\[
 p_i=M^{(i)}_{ii},\qquad
 M^{(i+1)}_{jk}
 =M^{(i)}_{jk}-M^{(i)}_{ji}M^{(i)}_{ik}/p_i
 \quad(j,k>i).
\]
It proves the stated inertia by exact rational comparison, rather
than taking the pivot signs as external data.
Sylvester's law of inertia shows that $K$ has exactly one
eigenvalue above $37/8$. Since $K_{00}=5$, its largest eigenvalue
is at least $5$ and is simple. Every other eigenvalue is below
$37/8$.

Define a rational vector $v$ as follows. For a bit string use the key
\[
 (a,b,i,j)=(b_0,b_1,b_2+b_3,b_4+b_5).
\]
The amplitude for each individual bit string is the following
numerator divided by $10^{15}$; a missing key has amplitude zero.
\begin{center}
\begin{tabular}{cr}
\toprule
Key & Numerator\\
\midrule
$(0,0,0,0)$ & $1000000000000000$\\
$(0,0,1,1)$ & $8942179087130$\\
$(0,0,2,2)$ & $136748346620$\\
$(0,1,0,1)$ & $171232824560490$\\
$(0,1,1,2)$ & $2452435031309$\\
$(1,0,1,0)$ & $171232824560490$\\
$(1,0,2,1)$ & $2452435031309$\\
$(1,1,0,0)$ & $97255575326164$\\
$(1,1,1,1)$ & $37858301388951$\\
$(1,1,2,2)$ & $1200949534669$\\
\bottomrule
\end{tabular}
\end{center}
These are computational-basis amplitudes, not normalized orbit-basis
coordinates. Put
\[
 \widehat\lambda=\frac{539109343678406}{10^{14}}.
\]
Direct rational summation gives
\[
 \|(K-\widehat\lambda I)v\|_2^2<10^{-24},
 \qquad \widehat\lambda-37/8>1/2.
\]
Let $v_g$ be the orthogonal projection onto the largest eigenspace
of $K$. Expanding in an orthonormal eigenbasis and using the
preceding separation gives
\[
 \|v-v_g\|_2
 \leq2\|(K-\widehat\lambda I)v\|_2<2\cdot10^{-12}.
\]
Since $v_{000000}=1$, $v_g$ is nonzero, hence is a scalar multiple
of the true ground state. Both $v$ and $v_g$ have charge zero;
in particular they have even parity.

For an even-parity vector $w$ put
\[
 \widetilde f_w(x)=f_w(\sqrt2\,x)
        =\sum_b2^{|b|/2}w_bx^b.
\]
The coefficients of $\widetilde f_v$ are rational. Set
\[
 c=\frac{175105+981180\,\mathrm i}{10^6},\qquad
 \ell=\frac{-379834+921467\,\mathrm i}{10^6}.
\]
Exact squaring gives $|c|,|\ell|<499/500$. Fix
$x_0=x_1=c$ and $x_2=x_3=x_4=\ell$. Multiaffinity yields
\[
 \widetilde f_v(c,c,\ell,\ell,\ell,x_5)=A_v+B_vx_5,
 \qquad
 |A_v|^2<(85/1000)^2,\quad |B_v|^2>(852/10000)^2 .
\]
The displayed inequalities are finite sums over the amplitude table.
Each coefficient $A$ or $B$ contains at most $32$ terms. A fixed-variable
monomial has modulus at most one, and its even tilt factor is at most
eight. Therefore
\[
 |A_{v_g}-A_v|,\ |B_{v_g}-B_v|
 <32\cdot8\cdot2\cdot10^{-12}<10^{-8}.
\]
In particular $B_{v_g}\neq0$, and
\[
 x_5=-A_{v_g}/B_{v_g},\qquad
 |x_5|<
 \frac{85/1000+10^{-8}}{852/10000-10^{-8}}<499/500.
\]
The last inequality is rational. This gives a zero of
$\widetilde f_{v_g}$ with all coordinates of modulus below $499/500$.
Rescaling by $\sqrt2$ proves the theorem.
\end{proof}

The certificate locates a zero of the exact ground-state polynomial.
It does not treat a rounded eigenvector as an exact eigenvector.
The existence of some radius greater than one, proved above by
positive-field strictification, is compatible with this example.

\subsection*{Complete finite verification procedure}
The following standard-library Python procedure reproduces all
rational assertions used in the proof. The operations implementing
a Gaussian rational number are written explicitly. Floating-point
arithmetic and numerical eigenvalue computations are not used.
The same source is supplied as
\path{certificates/epr_six_vertex_certificate.py}; a second,
independent rational certificate is included with the source package.
\begin{lstlisting}[basicstyle=\ttfamily\scriptsize,breaklines=true,
 columns=fullflexible,keepspaces=true,showstringspaces=false,
 language=Python,frame=single]
"""Exact standard-library certificate: a zero strictly inside radius sqrt(2).

No floating point or numerical eigensolver is used in any assertion.
K is the negative EPR Hamiltonian at s=1/2 on edges 01,02,03,14,15.
"""
from fractions import Fraction as F

EDGES = [(0, 1), (0, 2), (0, 3), (1, 4), (1, 5)]
N = 64
K = [[F(0) for _ in range(N)] for _ in range(N)]
for b in range(N):
    for i, j in EDGES:
        bi, bj = (b >> i) & 1, (b >> j) & 1
        if bi == bj:
            for out in (0, 1):
                c = (b & ~(1 << i) & ~(1 << j)) | (out << i) | (out << j)
                K[c][b] += F(1, 2) ** (bi + out)
assert all(K[i][j] == K[j][i] for i in range(N) for j in range(N))

# Exact inertia proves that all but the largest eigenvalue are below 37/8.
def charge(b):
    return sum((b >> j) & 1 for j in (0, 4, 5)) - sum((b >> j) & 1 for j in (1, 2, 3))

assert all(not K[i][j] or charge(i) == charge(j) for i in range(N) for j in range(N))

negative = 0
pivot_signs = []
for q in range(-3, 4):
    sector = [b for b in range(N) if charge(b) == q]
    A = [[F(37 if i == j else 0) - 8 * K[i][j] for j in sector] for i in sector]
    signs = []
    for i in range(len(sector)):
        pivot = A[i][i]
        assert pivot != 0
        negative += pivot < 0
        signs.append('+' if pivot > 0 else '-')
        for j in range(i + 1, len(sector)):
            for k in range(j, len(sector)):
                A[k][j] -= A[j][i] * A[k][i] / pivot
                A[j][k] = A[k][j]
    pivot_signs.append((q, ''.join(signs)))
assert negative == 1
assert K[0][0] == 5

states = [(0,0,0,0),(0,0,1,1),(0,0,2,2),(0,1,0,1),(0,1,1,2),
          (1,0,1,0),(1,0,2,1),(1,1,0,0),(1,1,1,1),(1,1,2,2)]
numerators = [1000000000000000,8942179087130,136748346620,171232824560490,
              2452435031309,171232824560490,2452435031309,97255575326164,
              37858301388951,1200949534669]
amps = {state:F(num,10**15) for state,num in zip(states,numerators)}
v = []
for b in range(N):
    state = ((b >> 0)&1,(b >> 1)&1,((b >> 2)&1)+((b >> 3)&1),
             ((b >> 4)&1)+((b >> 5)&1))
    v.append(amps.get(state,F(0)))
assert v[0] == 1
lam = F(539109343678406,10**14)
r = [sum(K[i][j] * v[j] for j in range(N)) - lam*v[i] for i in range(N)]
assert sum(x*x for x in r) < F(1,10**24)
assert lam - F(37,8) > F(1,2)

# Therefore the orthogonal projection v_g onto the true top eigenspace
# obeys ||v-v_g||_2 < 2*10^-12. In particular v_g is nonzero.

def add(x,y): return x[0]+y[0],x[1]+y[1]
def mul(x,y): return x[0]*y[0]-x[1]*y[1],x[0]*y[1]+x[1]*y[0]
def scale(c,x): return c*x[0],c*x[1]
def norm2(x): return x[0]*x[0]+x[1]*x[1]

zc=(F(175105,10**6),F(981180,10**6))
zl=(F(-379834,10**6),F(921467,10**6))
certified_radius=F(499,500)
assert norm2(zc)<certified_radius**2 and norm2(zl)<certified_radius**2
fixed=[zc,zc,zl,zl,zl]
A=(F(0),F(0));B=(F(0),F(0))
for b in range(N):
    if not v[b]: continue
    degree=b.bit_count()
    assert degree%2==0
    term=(v[b]*2**(degree//2),F(0))
    for j in range(5):
        if (b>>j)&1: term=mul(term,fixed[j])
    if (b>>5)&1:B=add(B,term)
    else:A=add(A,term)

# These wide rational bounds leave much more room than the eigenvector error.
assert norm2(A) < F(85,1000)**2
assert norm2(B) > F(852,10000)**2

# At the five fixed coordinates each monomial has modulus<1, and tilting
# by sqrt(2) multiplies each even coefficient by at most8. Each linear
# coefficient A or B involves at most32 terms. Thus replacing v by v_g
# changes A and B by less than 32*8*2e-12 < 1e-8.
err=F(1,10**8)
assert 32*8*F(2,10**12) < err
assert F(85,1000)+err < F(852,10000)-err
assert F(85,1000)+err < certified_radius*(F(852,10000)-err)

print('PASS: exact inertia has one negative pivot:',pivot_signs)
print('PASS: squared eigenvector residual < 10^-24')
print('PASS: |A_v|<85/1000 and |B_v|>852/10000')
print('PASS: true ground polynomial has all |z_i|<499/500 after sqrt(2) scaling')
\end{lstlisting}

% END SECTION counterexamples

\section*{Disclosure of artificial intelligence use}
Most of the mathematical content, proofs, computational certificates,
and text in this article were generated with substantial assistance from
OpenAI GPT-6 Pro in ChatGPT chat mode and GPT-6 Astra in Codex, both
using the ultra reasoning setting. These were the only AI tools used.
The author organized and directed the research through repeated
interactions, including questions, choices of research directions,
requests for extensions and critical review, and selection and
organization of the material. AI assistance was also used for
literature searches, proof audits, exact computational checks, and
manuscript revision. These checks do not constitute complete
proof-assistant verification, which is not claimed.
The author takes full responsibility for the accuracy, originality,
attribution, and conclusions of the article.
\bibliographystyle{amsplain}
\providecommand{\bysame}{\leavevmode\hbox to3em{\hrulefill}\thinspace}
\providecommand{\MR}{\relax\ifhmode\unskip\space\fi MR }
% \MRhref is called by the amsart/book/proc definition of \MR.
\providecommand{\MRhref}[2]{%
  \href{http://www.ams.org/mathscinet-getitem?mr=#1}{#2}
}
\providecommand{\href}[2]{#2}
\begin{thebibliography}{10}

\bibitem{AlNuwairan14}
Muneerah Al~Nuwairan, \emph{The extreme points of {$SU(2)$}-irreducibly
  covariant channels}, Internat. J. Math. \textbf{25} (2014), 1450048.

\bibitem{Aniello10}
P.~Aniello, A.~Kossakowski, G.~Marmo, and F.~Ventriglia, \emph{Brownian motion
  on {Lie} groups and open quantum systems}, J. Phys. A: Math. Theor.
  \textbf{43} (2010), no.~26, 265301.

\bibitem{ALMPSS26}
Anuj Apte, Eunou Lee, Kunal Marwaha, Ojas Parekh, Lennart Sinjorgo, and James
  Sud, \emph{A 0.8395-approximation algorithm for the {EPR} problem}, 34th
  Annual European Symposium on Algorithms ({ESA} 2026), Leibniz International
  Proceedings in Informatics (LIPIcs), vol. 388, Schloss
  Dagstuhl--Leibniz-Zentrum f{\"u}r Informatik, 2026, pp.~144:1--144:13.

\bibitem{APS26}
Anuj Apte, Ojas Parekh, and James Sud, \emph{Conjectured bounds for 2-local
  {Hamiltonians} via token graphs}, arXiv:2506.03441v2, 2026, Version 2, 27 May
  2026; first version 3 June 2025.

\bibitem{AsanoPRL70}
Taro Asano, \emph{{Lee--Yang} theorem and the {Griffiths} inequality for the
  anisotropic {Heisenberg} ferromagnet}, Phys. Rev. Lett. \textbf{24} (1970),
  1409--1411.

\bibitem{AsanoJapanese70}
\bysame, \emph{Rigorous theorems concerning {Heisenberg} ferromagnets}, Bussei
  Kenkyu \textbf{14} (1970), no.~2, 72--96, In Japanese; title translated from
  the original.

\bibitem{ARS24}
Tommaso Aschieri, B{\l}a{\.{z}}ej Ruba, and Jan~Philip Solovej,
  \emph{{SU(2)}-equivariant quantum channels: Semiclassical analysis}, Comm.
  Math. Phys. \textbf{405} (2024), 298.

\bibitem{BBKL26}
Ainesh Bakshi, Arpon Basu, Pravesh Kothari, and Anqi Li, \emph{Sharp bounds on
  the eigenvalues of {Kikuchi} graphs and applications to {Quantum Max Cut}},
  arXiv:2605.14994v1, 2026, Version 1, 14 May 2026.

\bibitem{Boas54}
Ralph~Philip Boas, Jr., \emph{Entire functions}, Pure and Applied Mathematics,
  vol.~5, Academic Press, New York, 1954.

\bibitem{BB09}
Julius Borcea and Petter Br{\"a}nd{\'e}n, \emph{The {Lee--Yang} and
  {P{\'o}lya--Schur} programs. {I}. {Linear} operators preserving stability},
  Invent. Math. \textbf{177} (2009), no.~3, 541--569.

\bibitem{BB10II}
\bysame, \emph{The {Lee--Yang} and {P{\'o}lya--Schur} programs. {II}. {Theory}
  of stable polynomials and applications}, Comm. Pure Appl. Math. \textbf{62}
  (2009), no.~12, 1595--1631.

\bibitem{BB09Circular}
\bysame, \emph{{P{\'o}lya--Schur} master theorems for circular domains and
  their boundaries}, Ann. of Math. (2) \textbf{170} (2009), no.~1, 465--492.

\bibitem{BB10}
\bysame, \emph{Multivariate {P{\'o}lya--Schur} classification problems in the
  {Weyl} algebra}, Proc. Lond. Math. Soc. (3) \textbf{101} (2010), no.~1,
  73--104.

\bibitem{BBL09}
Julius Borcea, Petter Br{\"a}nd{\'e}n, and Thomas~M. Liggett, \emph{Negative
  dependence and the geometry of polynomials}, J. Amer. Math. Soc. \textbf{22}
  (2009), no.~2, 521--567.

\bibitem{Branden10Discrete}
Petter Br{\"a}nd{\'e}n, \emph{Discrete concavity and the half-plane property},
  SIAM Journal on Discrete Mathematics \textbf{24} (2010), no.~3, 921--933.

\bibitem{BDL08}
Sergey Bravyi, David~P. DiVincenzo, and Daniel Loss, \emph{Polynomial-time
  algorithm for simulation of weakly interacting quantum spin systems}, Comm.
  Math. Phys. \textbf{284} (2008), 481--507.

\bibitem{BG17}
Sergey Bravyi and David Gosset, \emph{Polynomial-time classical simulation of
  quantum ferromagnets}, Phys. Rev. Lett. \textbf{119} (2017), 100503.

\bibitem{BGLW26}
Sergey Bravyi, David Gosset, Yinchen Liu, and Benjamin Wong, \emph{Efficient
  quantum algorithm for {Heisenberg} spin systems}, arXiv:2607.14401, 2026,
  Version 1, July 15, 2026.

\bibitem{Burgarth19}
Daniel Burgarth, Paolo Facchi, Hiromichi Nakazato, Saverio Pascazio, and Kazuya
  Yuasa, \emph{Generalized adiabatic theorem and strong-coupling limits},
  Quantum \textbf{3} (2019), 152.

\bibitem{BurkeEaton12}
James~V. Burke and Julia Eaton, \emph{On the subdifferential regularity of max
  root functions for polynomials}, Nonlinear Anal. \textbf{75} (2012),
  1168--1187.

\bibitem{BurkeOverton01}
James~V. Burke and Michael~L. Overton, \emph{Variational analysis of the
  abscissa mapping for polynomials}, SIAM J. Control Optim. \textbf{39} (2001),
  no.~6, 1651--1676.

\bibitem{Chung67}
Kai~Lai Chung, \emph{Markov chains with stationary transition probabilities},
  second ed., Grundlehren der mathematischen Wissenschaften, vol. 104,
  Springer-Verlag, Berlin--Heidelberg, 1967.

\bibitem{DenisMartin22}
J\'{e}r\^{o}me Denis and John Martin, \emph{Extreme depolarization for any
  spin}, Phys. Rev. Research \textbf{4} (2022), no.~1, 013178.

\bibitem{Dubois09}
Lo{\"i}c Dubois, \emph{Projective metrics and contraction principles for
  complex cones}, J. Lond. Math. Soc. (2) \textbf{79} (2009), no.~3, 719--737.

\bibitem{Dunlop79}
Fran{\c c}ois Dunlop, \emph{Analyticity of the pressure for {Heisenberg} and
  plane rotator models}, Comm. Math. Phys. \textbf{69} (1979), 81--88.

\bibitem{FilippovKuzhamuratova19}
Sergey~N. Filippov and Ksenia~V. Kuzhamuratova, \emph{Quantum informational
  properties of the {Landau--Streater} channel}, J. Math. Phys. \textbf{60}
  (2019), no.~4, 042202.

\bibitem{GaubertQu15}
St{\'e}phane Gaubert and Zheng Qu, \emph{{Dobrushin}'s ergodicity coefficient
  for {Markov} operators on cones}, Integral Equations Operator Theory
  \textbf{81} (2015), no.~1, 127--150.

\bibitem{GhorbalSogokonPlatzer17}
Khalil Ghorbal, Andrew Sogokon, and Andr{\'e} Platzer, \emph{A hierarchy of
  proof rules for checking positive invariance of algebraic and semi-algebraic
  sets}, Comput. Lang. Syst. Struct. \textbf{47} (2017), no.~1, 19--43.

\bibitem{GKS76}
Vittorio Gorini, Andrzej Kossakowski, and E.~C.~G. Sudarshan, \emph{Completely
  positive dynamical semigroups of {$N$}-level systems}, J. Math. Phys.
  \textbf{17} (1976), no.~5, 821--825.

\bibitem{HMS20}
Aram~W. Harrow, Saeed Mehraban, and Mehdi Soleimanifar, \emph{Classical
  algorithms, correlation decay, and complex zeros of partition functions of
  quantum many-body systems}, Proceedings of the 52nd Annual {ACM SIGACT}
  Symposium on Theory of Computing, ACM, 2020, pp.~378--386.

\bibitem{HarveyLawson09}
F.~Reese Harvey and H.~Blaine Lawson, Jr., \emph{Hyperbolic polynomials and the
  {Dirichlet} problem}, arXiv:0912.5220, 2009.

\bibitem{Hunt56}
G.~A. Hunt, \emph{Semi-groups of measures on {Lie} groups}, Trans. Amer. Math.
  Soc. \textbf{81} (1956), no.~2, 264--293.

\bibitem{JiangSturmfels21}
Yuhan Jiang and Bernd Sturmfels, \emph{Bad projections of the {PSD} cone},
  Collect. Math. \textbf{72} (2021), 261--280.

\bibitem{KaltonRoberts83}
Nigel~J. Kalton and James~W. Roberts, \emph{Uniformly exhaustive submeasures
  and nearly additive set functions}, Trans. Amer. Math. Soc. \textbf{278}
  (1983), no.~2, 803--816.

\bibitem{Kendall55}
David~G. Kendall, \emph{Some analytical properties of continuous stationary
  {Markov} transition functions}, Trans. Amer. Math. Soc. \textbf{78} (1955),
  no.~2, 529--540.

\bibitem{Knese24}
Greg Knese, \emph{Rational inner functions on the polydisk---a survey},
  arXiv:2409.14604, 2024.

\bibitem{KnillLaflamme97}
Emanuel Knill and Raymond Laflamme, \emph{Theory of quantum error-correcting
  codes}, Phys. Rev. A \textbf{55} (1997), no.~2, 900--911.

\bibitem{LewisParriloRamana05}
Adrian~S. Lewis, Pablo~A. Parrilo, and Motakuri~V. Ramana, \emph{The {Lax}
  conjecture is true}, Proc. Amer. Math. Soc. \textbf{133} (2005), no.~9,
  2495--2499.

\bibitem{LVR10}
Thomas~M. Liggett and Alexander Vandenberg-Rodes, \emph{Stability on
  {$\{0,1,2,\ldots\}^{S}$}: Birth-death chains and particle systems}, Notions
  of Positivity and the Geometry of Polynomials (Petter Br{\"a}nd{\'e}n, Mikael
  Passare, and Mihai Putinar, eds.), Trends in Mathematics,
  Birkh{\"a}user/Springer Basel, Basel, 2011, pp.~311--329.

\bibitem{Lindblad76}
G{\"o}ran Lindblad, \emph{On the generators of quantum dynamical semigroups},
  Comm. Math. Phys. \textbf{48} (1976), 119--130.

\bibitem{MarwahaSud26}
Kunal Marwaha and James Sud, \emph{A complexity phase transition at the {EPR
  Hamiltonian}}, arXiv:2604.13026, 2026.

\bibitem{Nagumo42}
Mitio Nagumo, \emph{{\"U}ber die {Lage} der {Integralkurven} gew{\"o}hnlicher
  {Differentialgleichungen}}, Proc. Phys.-Math. Soc. Japan (3) \textbf{24}
  (1942), 551--559.

\bibitem{Nathanson21}
Melvyn~B. Nathanson, \emph{The {Hermite--Sylvester} criterion for real-rooted
  polynomials}, Math. Gaz. \textbf{105} (2021), no.~562, 122--125.

\bibitem{Neeb92}
Karl-Hermann Neeb, \emph{On the foundations of {Lie} semigroups}, J. Reine
  Angew. Math. \textbf{431} (1992), 165--189.

\bibitem{PatelRegts17}
Viresh Patel and Guus Regts, \emph{Deterministic polynomial-time approximation
  algorithms for partition functions and graph polynomials}, SIAM J. Comput.
  \textbf{46} (2017), no.~6, 1893--1919.

\bibitem{PRSS04}
Victoria Powers, Bruce Reznick, Claus Scheiderer, and Frank Sottile, \emph{A
  new approach to {Hilbert}'s theorem on ternary quartics}, C. R. Math. Acad.
  Sci. Paris \textbf{339} (2004), no.~9, 617--620.

\bibitem{Purbhoo18}
Kevin Purbhoo, \emph{Total nonnegativity and stable polynomials}, Canad. Math.
  Bull. \textbf{61} (2018), no.~4, 836--847.

\bibitem{RT26}
Chaithanya Rayudu and Jun Takahashi, \emph{Spectral gap of {Lee--Yang
  Hamiltonians}}, arXiv:2607.10765, 2026, Version 1, July 12, 2026.

\bibitem{Ruelle10}
David Ruelle, \emph{Characterization of {Lee--Yang} polynomials}, Ann. of Math.
  (2) \textbf{171} (2010), no.~1, 589--603.

\bibitem{Rugh10}
Hans~Henrik Rugh, \emph{Cones and gauges in complex spaces: Spectral gaps and
  complex {Perron--Frobenius} theory}, Ann. of Math. (2) \textbf{171} (2010),
  no.~3, 1707--1752.

\bibitem{SanMartinSantana02}
Luiz A.~B. San~Martin and Alexandre~J. Santana, \emph{The homotopy type of
  {Lie} semigroups in semi-simple {Lie} groups}, Monatsh. Math. \textbf{136}
  (2002), 151--173.

\bibitem{Shkarin11}
Stanislav Shkarin, \emph{Hypercyclic and mixing operator semigroups}, Proc.
  Edinb. Math. Soc. (2) \textbf{54} (2011), no.~3, 761--782.

\bibitem{SuzukiFisher71}
Masuo Suzuki and Michael~E. Fisher, \emph{Zeros of the partition function for
  the {Heisenberg}, ferroelectric, and general {Ising} models}, J. Math. Phys.
  \textbf{12} (1971), no.~2, 235--246.

\bibitem{Turbiner92}
Alexander Turbiner, \emph{{Lie}-algebraic approach to the theory of polynomial
  solutions. {I}. {Ordinary} differential equations and finite-difference
  equations in one variable}, arXiv:hep-th/9209079, 1992.

\bibitem{Turbiner94}
\bysame, \emph{{Lie}-algebras and linear operators with invariant subspaces},
  Lie Algebras, Cohomologies and New Findings in Quantum Mechanics (Niky Kamran
  and Peter~J. Olver, eds.), Contemporary Mathematics, vol. 160, American
  Mathematical Society, Providence, RI, 1994, Revised author version:
  arXiv:funct-an/9301001v2 (1998), pp.~263--310.

\bibitem{Voit03}
Michael Voit, \emph{A {L\'{e}vy}-characterization for {Gaussian} processes on
  matrix groups}, Publ. Math. Debrecen \textbf{63} (2003), no.~1--2, 115--132.

\bibitem{Werner98}
Reinhard~F. Werner, \emph{Optimal cloning of pure states}, Phys. Rev. A
  \textbf{58} (1998), no.~3, 1827--1832.

\bibitem{WildAlhambra23}
Dominik~S. Wild and {\'A}lvaro~M. Alhambra, \emph{Classical simulation of
  short-time quantum dynamics}, PRX Quantum \textbf{4} (2023), 020340.

\bibitem{WongTalk23}
Benjamin Wong, \emph{Quantum {Max-Cut} on bipartite graphs}, URA Seminar,
  University of Waterloo, August 14, 2023, Public talk abstract.

\bibitem{WongThesis26}
\bysame, \emph{{Lee--Yang} tensors in quantum information and algorithms for
  local {Hamiltonian} problems}, Master's thesis, University of Waterloo, 2026.

\bibitem{WBG26}
Benjamin Wong, Sergey Bravyi, David Gosset, and Yinchen Liu, \emph{{Lee--Yang}
  tensors and {Hamiltonian} complexity}, arXiv:2602.03605, 2026, Version 1,
  February 3, 2026.

\bibitem{YYZ22}
Penghui Yao, Yitong Yin, and Xinyuan Zhang, \emph{Polynomial-time approximation
  of zero-free partition functions}, 49th International Colloquium on Automata,
  Languages, and Programming ({ICALP} 2022) (Dagstuhl, Germany) (Miko{\l}aj
  Boja\'{n}czyk, Emanuela Merelli, and David~P. Woodruff, eds.), Leibniz
  International Proceedings in Informatics, vol. 229, Schloss Dagstuhl --
  Leibniz-Zentrum f{\"u}r Informatik, 2022, pp.~108:1--108:20.

\end{thebibliography}

\end{document}
